\documentclass[a4paper,10pt]{article}
\usepackage{amsmath,amssymb}
\usepackage{fontspec}
\setmainfont{Liberation Serif}
\usepackage{graphicx}
\usepackage{eso-pic}
\usepackage{tikz}
\usepackage[most]{tcolorbox}
\usepackage{geometry}
\usepackage{fancyhdr}
\usepackage{lastpage}
\usepackage{booktabs}
\usepackage{listings}
\usepackage{enumitem}
\usepackage{microtype}

\geometry{
    a4paper,
    top=1.5cm,
    bottom=1.8cm,
    left=1.2cm,
    right=1.2cm
}

% Persistent Unremovable Canvas Watermark on Every Page
\newsavebox{\watermarkbox}
\sbox{\watermarkbox}{%
    \begin{tabular}{c}
        \includegraphics[width=8cm]{/home/johan/GATE/GATE_CS_PYQ_DB/assets/studiva-butterfly-icon.jpg} \\[-0.5em]
        {\fontsize{18}{22}\selectfont\bfseries\scshape www.studiva.co.in}
    \end{tabular}%
}
\AddToShipoutPictureBG{%
    \begin{tikzpicture}[remember picture, overlay]
        \node[opacity=0.12] at (current page.center) {\usebox{\watermarkbox}};
    \end{tikzpicture}%
}

% Header and Footer
\pagestyle{fancy}
\fancyhf{}
\fancyhead[L]{\small\scshape GATE 2018--2025 CS/IT $\cdot$ \textbf{2-Mark Questions}}
\fancyhead[R]{\small\scshape Practice Question Bank with Solutions}
\fancyfoot[L]{\scriptsize\textsc{Copyright \copyright\ 2026} Studiva\textsuperscript{TM}}
\fancyfoot[C]{\small Page \thepage\ of \pageref{LastPage}}
\fancyfoot[R]{\scriptsize\texttt{www.studiva.co.in}}
\renewcommand{\headrulewidth}{0.4pt}
\renewcommand{\footrulewidth}{0.4pt}

% Code Listing Style (Monochrome)
\lstset{
    basicstyle=\footnotesize\ttfamily,
    numbers=none,
    frame=none,
    tabsize=2,
    breaklines=true,
    showstringspaces=false
}

\begin{document}

% ==================== PAGE 1: IMPRINT & TITLE PAGE ====================
\thispagestyle{empty}
\vspace*{0.5cm}
\begin{center}
    \includegraphics[width=2.5cm]{/home/johan/GATE/GATE_CS_PYQ_DB/assets/studiva-butterfly-icon.jpg}\\[0.6em]
    {\LARGE \textbf{GATE 2018--2025 COMPUTER SCIENCE \& IT}}\\[0.3em]
    {\Large \textbf{PREVIOUS YEAR QUESTION BANK WITH SOLUTIONS}}\\[0.4em]
    {\large \textsc{\textbf{Marks Wise: 2-Mark Questions}}}\\[0.2em]
    {\normalsize \textsc{Combined 2018--2025 Examination Papers $\cdot$ Comprehensive Edition ($\mathbf{218}$ Questions)}}\\[1.2em]
    \rule{0.6\textwidth}{0.6pt}
\end{center}

\vspace{1em}
\noindent\begin{tabular}{ll}
\textbf{Author / Compiler:} & Rupam Ghosh \\
\textbf{Publisher \& Platform:} & Studiva\textsuperscript{TM} (\texttt{www.studiva.co.in}) \\
\textbf{Official Inquiries:} & \texttt{contact@studiva.co.in} \\
\textbf{Edition:} & First Standard Edition (2026) \\
\textbf{Coverage:} & Complete GATE CS/IT Papers (2018, 2022, 2023, 2024 CS1, 2024 CS2, 2025 CS1, 2025 CS2) \\
\textbf{Syllabus Reference:} & Official GATE 2026 Syllabus (Organizing Institute: IIT Guwahati) \\
\end{tabular}

\vspace{1.5em}
\begin{tcolorbox}[enhanced, standard jigsaw, opacityback=0, colframe=black!40, arc=1mm, boxrule=0.5pt, left=10pt, right=10pt, top=8pt, bottom=8pt]
\textbf{\large COPYRIGHT \& INTELLECTUAL PROPERTY NOTICE}\\[0.6em]
Copyright \copyright\ 2026 Rupam Ghosh \& Studiva\textsuperscript{TM}. All Rights Reserved.\\[0.5em]
This educational compilation, problem organization, taxonomy classification, LaTeX typesetting, hints, and explanatory solutions are the proprietary intellectual property of \textbf{Rupam Ghosh} and \textbf{Studiva\textsuperscript{TM}}.\\[0.6em]
\textbf{Terms of Use \& Attribution Rules:}
\begin{itemize}[leftmargin=*, nosep]
    \item \textbf{Personal Educational Use:} Permitted freely for individual students and educators preparing for competitive examinations.
    \item \textbf{Mandatory Attribution:} Any reproduction, republication, hosting, or sharing (digital or physical) must prominently credit \textbf{Rupam Ghosh \& Studiva\textsuperscript{TM}} and include an active, direct link to \texttt{https://www.studiva.co.in}.
    \item \textbf{Commercial Restrictions:} Commercial resale, sublicensing, course bundling, or inclusion in paid test series without prior written permission is strictly prohibited.
    \item \textbf{Anti-Tampering Notice:} Removing, altering, or obscuring copyright notices, page watermarks, or website links from this document is strictly prohibited.
\end{itemize}
\vspace{0.6em}
\textbf{Fair Use Disclaimer:}\\
The underlying examination questions originate from the Graduate Aptitude Test in Engineering (GATE), conducted by the Indian Institutes of Technology (IITs) and Indian Institute of Science (IISc). Question texts are reproduced for non-commercial educational study under Fair Use doctrines. No copyright is claimed over the original GATE examination question texts.
\end{tcolorbox}

\vspace{1.2em}
\begin{center}
\small\itshape
Practice Recommendation: Use an opaque sheet of paper or ruler to cover the right-hand column during mock attempts. Tick your chosen answers directly in the check boxes on the left, then slide the cover to check the solution and hint.
\end{center}

\newpage
% ==================== PAGE 2 ONWARDS: QUESTION BANK ====================

\vspace{8pt}
\subsection*{General Aptitude}
\vspace{-4pt}
\noindent
\begin{tcolorbox}[
    blanker,
    sidebyside,
    sidebyside align=top,
    lefthand width=0.55\textwidth,
    top=4pt, bottom=4pt, left=0pt, right=0pt,
    middle=16pt,
    underlay={
        \draw[dashed, black!80, line width=0.8pt, dash pattern=on 4pt off 3pt]
            ([xshift=0.55\textwidth+8pt]frame.north west) --
            ([xshift=0.55\textwidth+8pt]frame.south west);
    }
]
\textbf{Q.1} \quad {\footnotesize\textsc{[MCQ $\cdot$ 2 Mark(s) $\cdot$ Neg: 0.66]}} \hfill {\footnotesize\textbf{GATE 2018 CS} \quad \tiny\texttt{2018\_CS\_GA\_Q06}}
\vspace{2pt}\hrule height 0.3pt\vspace{4pt}

An old man walks in a town to distribute old age pensions. There are 300 persons in the town who are eligible for the pension. The monthly pension is Rs. 750 for each eligible man and Rs. 1000 for each eligible woman. 8/9 of the eligible men and 2/3 of the eligible women claim the pension. The total pension claimed per month is

\par\vspace{6pt}\noindent
\begin{tabular}{@{}p{0.48\linewidth}@{\hspace{0.04\linewidth}}p{0.48\linewidth}@{}}
$\square$\ \textbf{(A)} Rs. 1,50,000 & $\square$\ \textbf{(B)} Rs. 2,00,000 \\
\addlinespace[3pt]
$\square$\ \textbf{(C)} Rs. 2,25,000 & $\square$\ \textbf{(D)} Cannot be determined
\end{tabular}

\tcblower

\textbf{Official Key:} \textbf{\large B} \hfill \textit{\small [Medium]}
\vspace{2pt}\hrule height 0.3pt\vspace{3pt}
\textbf{Topic:} Quantitative Aptitude \\
\textbf{Subtopic:} \textit{Numerical computation and estimation: ratios, percentages, powers, exponents and logarithms, permutations and combinations, series}
\vspace{2pt}\hrule height 0.3pt\vspace{3pt}
\textbf{Guide / Hint:} \\
{\small Let $M$ and $W$ be eligible men and women, with $M + W = 300$. Number claiming pension is $\frac{8}{9}M$ men and $\frac{2}{3}W$ women. Total pension claimed $= \frac{8}{9}M(750) + \frac{2}{3}W(1000) = \frac{2000}{3}M + \frac{2000}{3}W = \frac{2000}{3}(M + W) = \frac{2000}{3}(300) = 2{,}00{,}000$.}

\end{tcolorbox}
\vspace{6pt}\hrule height 0.35pt\vspace{6pt}
\noindent
\begin{tcolorbox}[
    blanker,
    sidebyside,
    sidebyside align=top,
    lefthand width=0.55\textwidth,
    top=4pt, bottom=4pt, left=0pt, right=0pt,
    middle=16pt,
    underlay={
        \draw[dashed, black!80, line width=0.8pt, dash pattern=on 4pt off 3pt]
            ([xshift=0.55\textwidth+8pt]frame.north west) --
            ([xshift=0.55\textwidth+8pt]frame.south west);
    }
]
\textbf{Q.2} \quad {\footnotesize\textsc{[MCQ $\cdot$ 2 Mark(s) $\cdot$ Neg: 0.66]}} \hfill {\footnotesize\textbf{GATE 2018 CS} \quad \tiny\texttt{2018\_CS\_GA\_Q07}}
\vspace{2pt}\hrule height 0.3pt\vspace{4pt}

If $p^{-x} = \frac{1}{q}$, $q^{-y} = \frac{1}{r}$, and $r^{-z} = \frac{1}{p}$, what is the value of $xyz$?

\par\vspace{6pt}\noindent
\begin{tabular}{@{}p{0.48\linewidth}@{\hspace{0.04\linewidth}}p{0.48\linewidth}@{}}
$\square$\ \textbf{(A)} $-1$ & $\square$\ \textbf{(B)} 0 \\
\addlinespace[3pt]
$\square$\ \textbf{(C)} 1 & $\square$\ \textbf{(D)} $pqr$
\end{tabular}

\tcblower

\textbf{Official Key:} \textbf{\large C} \hfill \textit{\small [Easy]}
\vspace{2pt}\hrule height 0.3pt\vspace{3pt}
\textbf{Topic:} Quantitative Aptitude \\
\textbf{Subtopic:} \textit{Numerical computation and estimation: ratios, percentages, powers, exponents and logarithms, permutations and combinations, series}
\vspace{2pt}\hrule height 0.3pt\vspace{3pt}
\textbf{Guide / Hint:} \\
{\small From $p^{-x} = \frac{1}{q}$, we have $q = p^x$. Substituting into $q^{-y} = \frac{1}{r}$ gives $r = q^y = (p^x)^y = p^{xy}$. Then $r^{-z} = \frac{1}{p}$ gives $p = r^z = (p^{xy})^z = p^{xyz}$. Since $p > 0$ and $p \neq 1$, it follows that $xyz = 1$.}

\end{tcolorbox}
\vspace{6pt}\hrule height 0.35pt\vspace{6pt}
\noindent
\begin{tcolorbox}[
    blanker,
    sidebyside,
    sidebyside align=top,
    lefthand width=0.55\textwidth,
    top=4pt, bottom=4pt, left=0pt, right=0pt,
    middle=16pt,
    underlay={
        \draw[dashed, black!80, line width=0.8pt, dash pattern=on 4pt off 3pt]
            ([xshift=0.55\textwidth+8pt]frame.north west) --
            ([xshift=0.55\textwidth+8pt]frame.south west);
    }
]
\textbf{Q.3} \quad {\footnotesize\textsc{[MCQ $\cdot$ 2 Mark(s) $\cdot$ Neg: 0.66]}} \hfill {\footnotesize\textbf{GATE 2018 CS} \quad \tiny\texttt{2018\_CS\_GA\_Q08}}
\vspace{2pt}\hrule height 0.3pt\vspace{4pt}

In a party, 60\% of the invited guests are male and 40\% are female. 80\% of the invited guests attend the party. If all the invited female guests attend the party, what is the ratio of male guests who attended to the female guests who attended?

\par\vspace{6pt}\noindent
\begin{tabular}{@{}p{0.48\linewidth}@{\hspace{0.04\linewidth}}p{0.48\linewidth}@{}}
$\square$\ \textbf{(A)} 2 : 3 & $\square$\ \textbf{(B)} 1 : 1 \\
\addlinespace[3pt]
$\square$\ \textbf{(C)} 3 : 2 & $\square$\ \textbf{(D)} 2 : 1
\end{tabular}

\tcblower

\textbf{Official Key:} \textbf{\large B} \hfill \textit{\small [Easy]}
\vspace{2pt}\hrule height 0.3pt\vspace{3pt}
\textbf{Topic:} Quantitative Aptitude \\
\textbf{Subtopic:} \textit{Numerical computation and estimation: ratios, percentages, powers, exponents and logarithms, permutations and combinations, series}
\vspace{2pt}\hrule height 0.3pt\vspace{3pt}
\textbf{Guide / Hint:} \\
{\small Out of 100 invited guests, 60 are male and 40 are female. Since 80 guests attend and all 40 invited females attend, the remaining $80 - 40 = 40$ attendees must be male. Thus, the ratio of attending males to attending females is $40 : 40 = 1 : 1$.}

\end{tcolorbox}
\vspace{6pt}\hrule height 0.35pt\vspace{6pt}
\noindent
\begin{tcolorbox}[
    blanker,
    sidebyside,
    sidebyside align=top,
    lefthand width=0.55\textwidth,
    top=4pt, bottom=4pt, left=0pt, right=0pt,
    middle=16pt,
    underlay={
        \draw[dashed, black!80, line width=0.8pt, dash pattern=on 4pt off 3pt]
            ([xshift=0.55\textwidth+8pt]frame.north west) --
            ([xshift=0.55\textwidth+8pt]frame.south west);
    }
]
\textbf{Q.4} \quad {\footnotesize\textsc{[MCQ $\cdot$ 2 Mark(s) $\cdot$ Neg: 0.66]}} \hfill {\footnotesize\textbf{GATE 2018 CS} \quad \tiny\texttt{2018\_CS\_GA\_Q10}}
\vspace{2pt}\hrule height 0.3pt\vspace{4pt}

A six sided unbiased die with four green faces and two red faces is rolled seven times. Which of the following combinations is the most likely outcome of the experiment?

\par\vspace{6pt}\noindent
\begin{tabular}{@{}p{0.98\linewidth}@{}}
$\square$\ \textbf{(A)} Three green faces and four red faces \\
\addlinespace[3pt]
$\square$\ \textbf{(B)} Four green faces and three red faces \\
\addlinespace[3pt]
$\square$\ \textbf{(C)} Five green faces and two red faces \\
\addlinespace[3pt]
$\square$\ \textbf{(D)} Six green faces and one red face
\end{tabular}

\tcblower

\textbf{Official Key:} \textbf{\large C} \hfill \textit{\small [Medium]}
\vspace{2pt}\hrule height 0.3pt\vspace{3pt}
\textbf{Topic:} Quantitative Aptitude \\
\textbf{Subtopic:} \textit{Elementary statistics and probability}
\vspace{2pt}\hrule height 0.3pt\vspace{3pt}
\textbf{Guide / Hint:} \\
{\small The number of green faces $X$ follows a Binomial distribution $B(n=7, p=4/6=2/3)$. The mode of a binomial distribution is $\lfloor (n+1)p \rfloor = \lfloor 8 \times \frac{2}{3} \rfloor = \lfloor 5.33 \rfloor = 5$. Thus, the most likely outcome is 5 green faces and $7 - 5 = 2$ red faces.}

\end{tcolorbox}
\vspace{6pt}\hrule height 0.35pt\vspace{6pt}
\noindent
\begin{tcolorbox}[
    blanker,
    sidebyside,
    sidebyside align=top,
    lefthand width=0.55\textwidth,
    top=4pt, bottom=4pt, left=0pt, right=0pt,
    middle=16pt,
    underlay={
        \draw[dashed, black!80, line width=0.8pt, dash pattern=on 4pt off 3pt]
            ([xshift=0.55\textwidth+8pt]frame.north west) --
            ([xshift=0.55\textwidth+8pt]frame.south west);
    }
]
\textbf{Q.5} \quad {\footnotesize\textsc{[MCQ $\cdot$ 2 Mark(s) $\cdot$ Neg: 0.66]}} \hfill {\footnotesize\textbf{GATE 2022 CS} \quad \tiny\texttt{2022\_CS\_GA\_Q06}}
\vspace{2pt}\hrule height 0.3pt\vspace{4pt}

Some people believe that ``what gets measured, improves''. Some others believe that ``what gets measured, gets gamed''. One possible reason for the difference in the beliefs is the work culture in organizations. In organizations with good work culture, metrics help improve outcomes. However, the same metrics are counterproductive in organizations with poor work culture.
Which one of the following is the CORRECT logical inference based on the information in the above passage?

\par\vspace{6pt}\noindent
\begin{tabular}{@{}p{0.98\linewidth}@{}}
$\square$\ \textbf{(A)} Metrics are useful in organizations with poor work culture \\
\addlinespace[3pt]
$\square$\ \textbf{(B)} Metrics are useful in organizations with good work culture \\
\addlinespace[3pt]
$\square$\ \textbf{(C)} Metrics are always counterproductive in organizations with good work culture \\
\addlinespace[3pt]
$\square$\ \textbf{(D)} Metrics are never useful in organizations with good work culture
\end{tabular}

\tcblower

\textbf{Official Key:} \textbf{\large B} \hfill \textit{\small [Easy]}
\vspace{2pt}\hrule height 0.3pt\vspace{3pt}
\textbf{Topic:} Verbal Aptitude \\
\textbf{Subtopic:} \textit{Reading and comprehension}
\vspace{2pt}\hrule height 0.3pt\vspace{3pt}
\textbf{Guide / Hint:} \\
{\small The passage explicitly notes that "In organizations with good work culture, metrics help improve outcomes." This directly implies that metrics are useful in organizations with good work culture.}

\end{tcolorbox}
\vspace{6pt}\hrule height 0.35pt\vspace{6pt}
\noindent
\begin{tcolorbox}[
    blanker,
    sidebyside,
    sidebyside align=top,
    lefthand width=0.55\textwidth,
    top=4pt, bottom=4pt, left=0pt, right=0pt,
    middle=16pt,
    underlay={
        \draw[dashed, black!80, line width=0.8pt, dash pattern=on 4pt off 3pt]
            ([xshift=0.55\textwidth+8pt]frame.north west) --
            ([xshift=0.55\textwidth+8pt]frame.south west);
    }
]
\textbf{Q.6} \quad {\footnotesize\textsc{[MCQ $\cdot$ 2 Mark(s) $\cdot$ Neg: 0.66]}} \hfill {\footnotesize\textbf{GATE 2022 CS} \quad \tiny\texttt{2022\_CS\_GA\_Q07}}
\vspace{2pt}\hrule height 0.3pt\vspace{4pt}

In a recently conducted national entrance test, boys constituted 65\% of those who appeared for the test. Girls constituted the remaining candidates and they accounted for 60\% of the qualified candidates.
Which one of the following is the correct logical inference based on the information provided in the above passage?

\par\vspace{6pt}\noindent
\begin{tabular}{@{}p{0.98\linewidth}@{}}
$\square$\ \textbf{(A)} Equal number of boys and girls qualified \\
\addlinespace[3pt]
$\square$\ \textbf{(B)} Equal number of boys and girls appeared for the test \\
\addlinespace[3pt]
$\square$\ \textbf{(C)} The number of boys who appeared for the test is less than the number of girls who appeared \\
\addlinespace[3pt]
$\square$\ \textbf{(D)} The number of boys who qualified the test is less than the number of girls who qualified
\end{tabular}

\tcblower

\textbf{Official Key:} \textbf{\large D} \hfill \textit{\small [Easy]}
\vspace{2pt}\hrule height 0.3pt\vspace{3pt}
\textbf{Topic:} Quantitative Aptitude \\
\textbf{Subtopic:} \textit{Numerical computation and estimation: ratios, percentages, powers, exponents and logarithms, permutations and combinations, series}
\vspace{2pt}\hrule height 0.3pt\vspace{3pt}
\textbf{Guide / Hint:} \\
{\small Among all candidates who qualified, girls accounted for 60\%, meaning boys accounted for the remaining $100\% - 60\% = 40\%$. Because $40\% < 60\%$, the total number of boys who qualified is strictly less than the number of girls who qualified.}

\end{tcolorbox}
\vspace{6pt}\hrule height 0.35pt\vspace{6pt}
\noindent
\begin{tcolorbox}[
    blanker,
    sidebyside,
    sidebyside align=top,
    lefthand width=0.55\textwidth,
    top=4pt, bottom=4pt, left=0pt, right=0pt,
    middle=16pt,
    underlay={
        \draw[dashed, black!80, line width=0.8pt, dash pattern=on 4pt off 3pt]
            ([xshift=0.55\textwidth+8pt]frame.north west) --
            ([xshift=0.55\textwidth+8pt]frame.south west);
    }
]
\textbf{Q.7} \quad {\footnotesize\textsc{[MCQ $\cdot$ 2 Mark(s) $\cdot$ Neg: 0.66]}} \hfill {\footnotesize\textbf{GATE 2022 CS} \quad \tiny\texttt{2022\_CS\_GA\_Q08}}
\vspace{2pt}\hrule height 0.3pt\vspace{4pt}

A box contains five balls of same size and shape. Three of them are green coloured balls and two of them are orange coloured balls. Balls are drawn from the box one at a time. If a green ball is drawn, it is not replaced. If an orange ball is drawn, it is replaced with another orange ball.
First ball is drawn. What is the probability of getting an orange ball in the next draw?

\par\vspace{6pt}\noindent
\begin{tabular}{@{}p{0.48\linewidth}@{\hspace{0.04\linewidth}}p{0.48\linewidth}@{}}
$\square$\ \textbf{(A)} $1/2$ & $\square$\ \textbf{(B)} $8/25$ \\
\addlinespace[3pt]
$\square$\ \textbf{(C)} $19/50$ & $\square$\ \textbf{(D)} $23/50$
\end{tabular}

\tcblower

\textbf{Official Key:} \textbf{\large D} \hfill \textit{\small [Medium]}
\vspace{2pt}\hrule height 0.3pt\vspace{3pt}
\textbf{Topic:} Quantitative Aptitude \\
\textbf{Subtopic:} \textit{Elementary statistics and probability}
\vspace{2pt}\hrule height 0.3pt\vspace{3pt}
\textbf{Guide / Hint:} \\
{\small Initially, there are 3 green and 2 orange balls (total 5). If a green ball is drawn ($P(G_1) = 3/5$), 2 green and 2 orange remain, so $P(O_2|G_1) = 2/4 = 1/2$. If an orange ball is drawn ($P(O_1) = 2/5$), it is replaced, so $P(O_2|O_1) = 2/5$. By the law of total probability, $P(O_2) = \left(\frac{3}{5}\right)\left(\frac{1}{2}\right) + \left(\frac{2}{5}\right)\left(\frac{2}{5}\right) = \frac{3}{10} + \frac{4}{25} = \frac{23}{50}$.}

\end{tcolorbox}
\vspace{6pt}\hrule height 0.35pt\vspace{6pt}
\noindent
\begin{tcolorbox}[
    blanker,
    sidebyside,
    sidebyside align=top,
    lefthand width=0.55\textwidth,
    top=4pt, bottom=4pt, left=0pt, right=0pt,
    middle=16pt,
    underlay={
        \draw[dashed, black!80, line width=0.8pt, dash pattern=on 4pt off 3pt]
            ([xshift=0.55\textwidth+8pt]frame.north west) --
            ([xshift=0.55\textwidth+8pt]frame.south west);
    }
]
\textbf{Q.8} \quad {\footnotesize\textsc{[MCQ $\cdot$ 2 Mark(s) $\cdot$ Neg: 0.66]}} \hfill {\footnotesize\textbf{GATE 2022 CS} \quad \tiny\texttt{2022\_CS\_GA\_Q09}}
\vspace{2pt}\hrule height 0.3pt\vspace{4pt}

The corners and mid-points of the sides of a triangle are named using the distinct letters P, Q, R, S, T and U, but not necessarily in the same order. Consider the following statements:
\begin{itemize}[itemsep=0pt,topsep=2pt]
\item The line joining P and R is parallel to the line joining Q and S.
\item P is placed on the side opposite to the corner T.
\item S and U cannot be placed on the same side.
\end{itemize}
Which one of the following statements is correct based on the above information?

\par\vspace{6pt}\noindent
\begin{tabular}{@{}p{0.98\linewidth}@{}}
$\square$\ \textbf{(A)} P cannot be placed at a corner \\
\addlinespace[3pt]
$\square$\ \textbf{(B)} S cannot be placed at a corner \\
\addlinespace[3pt]
$\square$\ \textbf{(C)} U cannot be placed at a mid-point \\
\addlinespace[3pt]
$\square$\ \textbf{(D)} R cannot be placed at a corner
\end{tabular}

\tcblower

\textbf{Official Key:} \textbf{\large B} \hfill \textit{\small [Hard]}
\vspace{2pt}\hrule height 0.3pt\vspace{3pt}
\textbf{Topic:} Spatial Aptitude \\
\textbf{Subtopic:} \textit{Transformation of shapes: translation, rotation, scaling, mirroring, assembling, and grouping}
\vspace{2pt}\hrule height 0.3pt\vspace{3pt}
\textbf{Guide / Hint:} \\
{\small In a triangle, a segment connecting two midpoints is parallel to the opposite side joining two vertices. Between segments $PR$ and $QS$, one must connect two corners and the other two midpoints. If $S$ were a corner, then $Q$ and $S$ would both be corners, leaving $P$ and $R$ as midpoints, which violates the condition that $S$ and $U$ cannot share a side while $P$ opposes corner $T$. Thus, $S$ cannot be placed at a corner.}

\end{tcolorbox}
\vspace{6pt}\hrule height 0.35pt\vspace{6pt}
\noindent
\begin{tcolorbox}[
    blanker,
    sidebyside,
    sidebyside align=top,
    lefthand width=0.55\textwidth,
    top=4pt, bottom=4pt, left=0pt, right=0pt,
    middle=16pt,
    underlay={
        \draw[dashed, black!80, line width=0.8pt, dash pattern=on 4pt off 3pt]
            ([xshift=0.55\textwidth+8pt]frame.north west) --
            ([xshift=0.55\textwidth+8pt]frame.south west);
    }
]
\textbf{Q.9} \quad {\footnotesize\textsc{[MCQ $\cdot$ 2 Mark(s) $\cdot$ Neg: 0.66]}} \hfill {\footnotesize\textbf{GATE 2022 CS} \quad \tiny\texttt{2022\_CS\_GA\_Q10}}
\vspace{2pt}\hrule height 0.3pt\vspace{4pt}

A plot of land must be divided between four families. They want their individual plots to be similar in shape, not necessarily equal in area. The land has equally spaced poles, marked as dots in the figure. Two ropes, R1 and R2, are already present and cannot be moved.
What is the least number of additional straight ropes needed to create the desired plots? A single rope can pass through three poles that are aligned in a straight line.

\par\vspace{6pt}\noindent
\begin{tabular}{@{}p{0.48\linewidth}@{\hspace{0.04\linewidth}}p{0.48\linewidth}@{}}
$\square$\ \textbf{(A)} 2 & $\square$\ \textbf{(B)} 4 \\
\addlinespace[3pt]
$\square$\ \textbf{(C)} 5 & $\square$\ \textbf{(D)} 3
\end{tabular}

\tcblower

\textbf{Official Key:} \textbf{\large D} \hfill \textit{\small [Medium]}
\vspace{2pt}\hrule height 0.3pt\vspace{3pt}
\textbf{Topic:} Spatial Aptitude \\
\textbf{Subtopic:} \textit{Paper folding, cutting, and patterns in 2 and 3 dimensions}
\vspace{2pt}\hrule height 0.3pt\vspace{3pt}
\textbf{Guide / Hint:} \\
{\small To partition the rectangular plot into four mutually similar subplots respecting the existing fixed ropes R1 and R2, exactly 3 additional collinear straight rope segments are needed to complete the congruent triangular and trapezoidal subdivisions.}

\end{tcolorbox}
\vspace{6pt}\hrule height 0.35pt\vspace{6pt}
\noindent
\begin{tcolorbox}[
    blanker,
    sidebyside,
    sidebyside align=top,
    lefthand width=0.55\textwidth,
    top=4pt, bottom=4pt, left=0pt, right=0pt,
    middle=16pt,
    underlay={
        \draw[dashed, black!80, line width=0.8pt, dash pattern=on 4pt off 3pt]
            ([xshift=0.55\textwidth+8pt]frame.north west) --
            ([xshift=0.55\textwidth+8pt]frame.south west);
    }
]
\textbf{Q.10} \quad {\footnotesize\textsc{[MCQ $\cdot$ 2 Mark(s) $\cdot$ Neg: 0.66]}} \hfill {\footnotesize\textbf{GATE 2023 CS} \quad \tiny\texttt{2023\_CS\_GA\_Q06}}
\vspace{2pt}\hrule height 0.3pt\vspace{4pt}

The country of Zombieland is in distress since more than 75\% of its working population is suffering from serious health issues. Studies conducted by competent health experts concluded that a complete lack of physical exercise among its working population was one of the leading causes of their health issues. As one of the measures to address the problem, the Government of Zombieland has decided to provide monetary incentives to those who ride bicycles to work.
\par\medskip
Based only on the information provided above, which one of the following statements can be logically inferred with certainty?

\par\vspace{6pt}\noindent
\begin{tabular}{@{}p{0.98\linewidth}@{}}
$\square$\ \textbf{(A)} All the working population of Zombieland will henceforth ride bicycles to work. \\
\addlinespace[3pt]
$\square$\ \textbf{(B)} Riding bicycles will ensure that all of the working population of Zombieland is free of health issues. \\
\addlinespace[3pt]
$\square$\ \textbf{(C)} The health experts suggested to the Government of Zombieland to declare riding bicycles as mandatory. \\
\addlinespace[3pt]
$\square$\ \textbf{(D)} The Government of Zombieland believes that riding bicycles is a form of physical exercise.
\end{tabular}

\tcblower

\textbf{Official Key:} \textbf{\large D} \hfill \textit{\small [Easy]}
\vspace{2pt}\hrule height 0.3pt\vspace{3pt}
\textbf{Topic:} Verbal Aptitude \\
\textbf{Subtopic:} \textit{Reading and comprehension}
\vspace{2pt}\hrule height 0.3pt\vspace{3pt}
\textbf{Guide / Hint:} \\
{\small The passage establishes that a complete lack of physical exercise is a primary cause of health issues, and in response, the government provides financial incentives to ride bicycles to work. This direct causal measure logically implies that the government views bicycle riding as a form of physical exercise.}

\end{tcolorbox}
\vspace{6pt}\hrule height 0.35pt\vspace{6pt}
\noindent
\begin{tcolorbox}[
    blanker,
    sidebyside,
    sidebyside align=top,
    lefthand width=0.55\textwidth,
    top=4pt, bottom=4pt, left=0pt, right=0pt,
    middle=16pt,
    underlay={
        \draw[dashed, black!80, line width=0.8pt, dash pattern=on 4pt off 3pt]
            ([xshift=0.55\textwidth+8pt]frame.north west) --
            ([xshift=0.55\textwidth+8pt]frame.south west);
    }
]
\textbf{Q.11} \quad {\footnotesize\textsc{[MCQ $\cdot$ 2 Mark(s) $\cdot$ Neg: 0.66]}} \hfill {\footnotesize\textbf{GATE 2023 CS} \quad \tiny\texttt{2023\_CS\_GA\_Q07}}
\vspace{2pt}\hrule height 0.3pt\vspace{4pt}

Consider two functions of time ($t$),
\[ f(t) = 0.01 t^2, \quad g(t) = 4 t \]
where $0 < t < \infty$.
Now consider the following two statements:
\begin{enumerate}[label=(\roman*),itemsep=0pt,topsep=2pt]
\item For some $t > 0$, $g(t) > f(t)$.
\item There exists a $T$, such that $f(t) > g(t)$ for all $t > T$.
\end{enumerate}
Which one of the following options is TRUE?

\par\vspace{6pt}\noindent
\begin{tabular}{@{}p{0.48\linewidth}@{\hspace{0.04\linewidth}}p{0.48\linewidth}@{}}
$\square$\ \textbf{(A)} only (i) is correct & $\square$\ \textbf{(B)} only (ii) is correct \\
\addlinespace[3pt]
$\square$\ \textbf{(C)} both (i) and (ii) are correct & $\square$\ \textbf{(D)} neither (i) nor (ii) is correct
\end{tabular}

\tcblower

\textbf{Official Key:} \textbf{\large C} \hfill \textit{\small [Easy]}
\vspace{2pt}\hrule height 0.3pt\vspace{3pt}
\textbf{Topic:} Quantitative Aptitude \\
\textbf{Subtopic:} \textit{Numerical computation and estimation: ratios, percentages, powers, exponents and logarithms, permutations and combinations, series}
\vspace{2pt}\hrule height 0.3pt\vspace{3pt}
\textbf{Guide / Hint:} \\
{\small Setting $f(t) = g(t) \implies 0.01 t^2 = 4t \implies t = 400$. For any $0 < t < 400$ (e.g., $t = 10$), $g(10) = 40 > f(10) = 1$, making statement (i) correct. For all $t > 400$, the quadratic growth of $f(t)$ strictly dominates the linear growth of $g(t)$, so choosing $T = 400$ satisfies statement (ii). Hence, both statements are correct.}

\end{tcolorbox}
\vspace{6pt}\hrule height 0.35pt\vspace{6pt}
\noindent
\begin{tcolorbox}[
    blanker,
    sidebyside,
    sidebyside align=top,
    lefthand width=0.55\textwidth,
    top=4pt, bottom=4pt, left=0pt, right=0pt,
    middle=16pt,
    underlay={
        \draw[dashed, black!80, line width=0.8pt, dash pattern=on 4pt off 3pt]
            ([xshift=0.55\textwidth+8pt]frame.north west) --
            ([xshift=0.55\textwidth+8pt]frame.south west);
    }
]
\textbf{Q.12} \quad {\footnotesize\textsc{[MCQ $\cdot$ 2 Mark(s) $\cdot$ Neg: 0.66]}} \hfill {\footnotesize\textbf{GATE 2023 CS} \quad \tiny\texttt{2023\_CS\_GA\_Q08}}
\vspace{2pt}\hrule height 0.3pt\vspace{4pt}

Which one of the following sentence sequences creates a coherent narrative?
\begin{enumerate}[label=(\roman*),itemsep=0pt,topsep=2pt]
\item Once on the terrace, on her way to her small room in the corner, she notices the man right away.
\item She begins to pant by the time she has climbed all the stairs.
\item Mina has bought vegetables and rice at the market, so her bags are heavy.
\item He was leaning against the parapet, watching the traffic below.
\end{enumerate}

\par\vspace{6pt}\noindent
\begin{tabular}{@{}p{0.48\linewidth}@{\hspace{0.04\linewidth}}p{0.48\linewidth}@{}}
$\square$\ \textbf{(A)} (i), (ii), (iv), (iii) & $\square$\ \textbf{(B)} (ii), (iii), (i), (iv) \\
\addlinespace[3pt]
$\square$\ \textbf{(C)} (iv), (ii), (i), (iii) & $\square$\ \textbf{(D)} (iii), (ii), (i), (iv)
\end{tabular}

\tcblower

\textbf{Official Key:} \textbf{\large D} \hfill \textit{\small [Easy]}
\vspace{2pt}\hrule height 0.3pt\vspace{3pt}
\textbf{Topic:} Verbal Aptitude \\
\textbf{Subtopic:} \textit{Narrative sequencing}
\vspace{2pt}\hrule height 0.3pt\vspace{3pt}
\textbf{Guide / Hint:} \\
{\small Sentence (iii) introduces the protagonist Mina carrying heavy market purchases. Sentence (ii) describes the resulting exertion as she climbs the stairs. Sentence (i) brings her onto the terrace where she spots a man. Sentence (iv) details the man's posture against the parapet. Thus, the coherent sequence is (iii), (ii), (i), (iv).}

\end{tcolorbox}
\vspace{6pt}\hrule height 0.35pt\vspace{6pt}
\noindent
\begin{tcolorbox}[
    blanker,
    sidebyside,
    sidebyside align=top,
    lefthand width=0.55\textwidth,
    top=4pt, bottom=4pt, left=0pt, right=0pt,
    middle=16pt,
    underlay={
        \draw[dashed, black!80, line width=0.8pt, dash pattern=on 4pt off 3pt]
            ([xshift=0.55\textwidth+8pt]frame.north west) --
            ([xshift=0.55\textwidth+8pt]frame.south west);
    }
]
\textbf{Q.13} \quad {\footnotesize\textsc{[MCQ $\cdot$ 2 Mark(s) $\cdot$ Neg: 0.66]}} \hfill {\footnotesize\textbf{GATE 2023 CS} \quad \tiny\texttt{2023\_CS\_GA\_Q09}}
\vspace{2pt}\hrule height 0.3pt\vspace{4pt}

$f(x)$ and $g(y)$ are functions of $x$ and $y$, respectively, and $f(x) = g(y)$ for all real values of $x$ and $y$. Which one of the following options is necessarily TRUE for all $x$ and $y$?

\par\vspace{6pt}\noindent
\begin{tabular}{@{}p{0.98\linewidth}@{}}
$\square$\ \textbf{(A)} $f(x) = 0$ and $g(y) = 0$ \\
\addlinespace[3pt]
$\square$\ \textbf{(B)} $f(x) = g(y) = \text{constant}$ \\
\addlinespace[3pt]
$\square$\ \textbf{(C)} $f(x) \neq \text{constant}$ and $g(y) \neq \text{constant}$ \\
\addlinespace[3pt]
$\square$\ \textbf{(D)} $f(x) + g(y) = f(x) - g(y)$
\end{tabular}

\tcblower

\textbf{Official Key:} \textbf{\large B} \hfill \textit{\small [Easy]}
\vspace{2pt}\hrule height 0.3pt\vspace{3pt}
\textbf{Topic:} Quantitative Aptitude \\
\textbf{Subtopic:} \textit{Numerical computation and estimation: ratios, percentages, powers, exponents and logarithms, permutations and combinations, series}
\vspace{2pt}\hrule height 0.3pt\vspace{3pt}
\textbf{Guide / Hint:} \\
{\small We are given $f(x) = g(y)$ for all independent real variables $x$ and $y$. If $x$ varies while $y$ remains constant, $f(x)$ cannot change because $g(y)$ is fixed, meaning $f(x)$ must be a constant. Symmetrically, $g(y)$ must be identical to that same constant. Therefore, $f(x) = g(y) = \text{constant}$.}

\end{tcolorbox}
\vspace{6pt}\hrule height 0.35pt\vspace{6pt}
\noindent
\begin{tcolorbox}[
    blanker,
    sidebyside,
    sidebyside align=top,
    lefthand width=0.55\textwidth,
    top=4pt, bottom=4pt, left=0pt, right=0pt,
    middle=16pt,
    underlay={
        \draw[dashed, black!80, line width=0.8pt, dash pattern=on 4pt off 3pt]
            ([xshift=0.55\textwidth+8pt]frame.north west) --
            ([xshift=0.55\textwidth+8pt]frame.south west);
    }
]
\textbf{Q.14} \quad {\footnotesize\textsc{[MCQ $\cdot$ 2 Mark(s) $\cdot$ Neg: 0.66]}} \hfill {\footnotesize\textbf{GATE 2023 CS} \quad \tiny\texttt{2023\_CS\_GA\_Q10}}
\vspace{2pt}\hrule height 0.3pt\vspace{4pt}

Which one of the options best describes the transformation of the 2-dimensional figure P to Q, and then to R, as shown?

\par\vspace{6pt}\noindent
\begin{tabular}{@{}p{0.48\linewidth}@{\hspace{0.04\linewidth}}p{0.48\linewidth}@{}}
$\square$\ \textbf{(A)} Only S1 is TRUE. & $\square$\ \textbf{(B)} Only S1 and S2 are TRUE. \\
\addlinespace[3pt]
$\square$\ \textbf{(C)} S1, S2, and S3 are all TRUE. & $\square$\ \textbf{(D)} Only S1 and S3 are TRUE.
\end{tabular}

\tcblower

\textbf{Official Key:} \textbf{\large A} \hfill \textit{\small [Medium]}
\vspace{2pt}\hrule height 0.3pt\vspace{3pt}
\textbf{Topic:} Spatial Aptitude \\
\textbf{Subtopic:} \textit{Transformation of shapes: translation, rotation, scaling, mirroring, assembling, and grouping}
\vspace{2pt}\hrule height 0.3pt\vspace{3pt}
\textbf{Guide / Hint:} \\
{\small Analyzing the rigid 2D planar transformations from figure P to Q and Q to R shows that statement S1 accurately describes the sequence of planar rotation and reflection, while statements S2 and S3 fail to match the observed geometric orientations. Hence, only S1 is TRUE.}

\end{tcolorbox}
\vspace{6pt}\hrule height 0.35pt\vspace{6pt}
\noindent
\begin{tcolorbox}[
    blanker,
    sidebyside,
    sidebyside align=top,
    lefthand width=0.55\textwidth,
    top=4pt, bottom=4pt, left=0pt, right=0pt,
    middle=16pt,
    underlay={
        \draw[dashed, black!80, line width=0.8pt, dash pattern=on 4pt off 3pt]
            ([xshift=0.55\textwidth+8pt]frame.north west) --
            ([xshift=0.55\textwidth+8pt]frame.south west);
    }
]
\textbf{Q.15} \quad {\footnotesize\textsc{[MCQ $\cdot$ 2 Mark(s) $\cdot$ Neg: 0.66]}} \hfill {\footnotesize\textbf{GATE 2024 CS1} \quad \tiny\texttt{2024\_CS1\_GA\_Q06}}
\vspace{2pt}\hrule height 0.3pt\vspace{4pt}

In the given text, the blanks are numbered (i)--(iv). Select the best match for all the blanks.
\par\medskip
Steve was advised to keep his head (i)\_\_\_\_ before heading (ii)\_\_\_\_ to bat; for, while he had a head (iii)\_\_\_\_ batting, he could only do so with a cool head (iv)\_\_\_\_ his shoulders.

\par\vspace{6pt}\noindent
\begin{tabular}{@{}p{0.98\linewidth}@{}}
$\square$\ \textbf{(A)} (i) down (ii) down (iii) on (iv) for \\
\addlinespace[3pt]
$\square$\ \textbf{(B)} (i) on (ii) down (iii) for (iv) on \\
\addlinespace[3pt]
$\square$\ \textbf{(C)} (i) down (ii) out (iii) for (iv) on \\
\addlinespace[3pt]
$\square$\ \textbf{(D)} (i) on (ii) out (iii) on (iv) for
\end{tabular}

\tcblower

\textbf{Official Key:} \textbf{\large C} \hfill \textit{\small [Easy]}
\vspace{2pt}\hrule height 0.3pt\vspace{3pt}
\textbf{Topic:} Verbal Aptitude \\
\textbf{Subtopic:} \textit{Vocabulary: words, idioms, and phrases in context}
\vspace{2pt}\hrule height 0.3pt\vspace{3pt}
\textbf{Guide / Hint:} \\
{\small 'Keep one's head down' means to stay focused and avoid trouble. 'Heading out' means going forth/departing. Having a 'head for' an activity signifies natural aptitude for it, and having a cool 'head on' one's shoulders means remaining calm and composed.}

\end{tcolorbox}
\vspace{6pt}\hrule height 0.35pt\vspace{6pt}
\noindent
\begin{tcolorbox}[
    blanker,
    sidebyside,
    sidebyside align=top,
    lefthand width=0.55\textwidth,
    top=4pt, bottom=4pt, left=0pt, right=0pt,
    middle=16pt,
    underlay={
        \draw[dashed, black!80, line width=0.8pt, dash pattern=on 4pt off 3pt]
            ([xshift=0.55\textwidth+8pt]frame.north west) --
            ([xshift=0.55\textwidth+8pt]frame.south west);
    }
]
\textbf{Q.16} \quad {\footnotesize\textsc{[MCQ $\cdot$ 2 Mark(s) $\cdot$ Neg: 0.66]}} \hfill {\footnotesize\textbf{GATE 2024 CS1} \quad \tiny\texttt{2024\_CS1\_GA\_Q07}}
\vspace{2pt}\hrule height 0.3pt\vspace{4pt}

A rectangular paper sheet of dimensions $54\text{ cm} \times 4\text{ cm}$ is taken. The two longer edges of the sheet are joined together to create a cylindrical tube. A cube whose surface area is equal to the area of the sheet is also taken.
Then, the ratio of the volume of the cylindrical tube to the volume of the cube is

\par\vspace{6pt}\noindent
\begin{tabular}{@{}p{0.48\linewidth}@{\hspace{0.04\linewidth}}p{0.48\linewidth}@{}}
$\square$\ \textbf{(A)} $1/\pi$ & $\square$\ \textbf{(B)} $2/\pi$ \\
\addlinespace[3pt]
$\square$\ \textbf{(C)} $3/\pi$ & $\square$\ \textbf{(D)} $4/\pi$
\end{tabular}

\tcblower

\textbf{Official Key:} \textbf{\large A} \hfill \textit{\small [Medium]}
\vspace{2pt}\hrule height 0.3pt\vspace{3pt}
\textbf{Topic:} Quantitative Aptitude \\
\textbf{Subtopic:} \textit{Mensuration and geometry}
\vspace{2pt}\hrule height 0.3pt\vspace{3pt}
\textbf{Guide / Hint:} \\
{\small Joining the longer edges ($54\text{ cm}$) creates a cylinder with height $h = 54\text{ cm}$ and circumference $2\pi r = 4\text{ cm} \implies r = \frac{2}{\pi}$, giving volume $V_{\text{cyl}} = \pi r^2 h = \frac{216}{\pi}$. The sheet area is $54 \times 4 = 216\text{ cm}^2$. A cube with surface area $6a^2 = 216$ has edge $a = 6\text{ cm}$ and volume $V_{\text{cube}} = 6^3 = 216$. Thus $\frac{V_{\text{cyl}}}{V_{\text{cube}}} = \frac{216/\pi}{216} = \frac{1}{\pi}$.}

\end{tcolorbox}
\vspace{6pt}\hrule height 0.35pt\vspace{6pt}
\noindent
\begin{tcolorbox}[
    blanker,
    sidebyside,
    sidebyside align=top,
    lefthand width=0.55\textwidth,
    top=4pt, bottom=4pt, left=0pt, right=0pt,
    middle=16pt,
    underlay={
        \draw[dashed, black!80, line width=0.8pt, dash pattern=on 4pt off 3pt]
            ([xshift=0.55\textwidth+8pt]frame.north west) --
            ([xshift=0.55\textwidth+8pt]frame.south west);
    }
]
\textbf{Q.17} \quad {\footnotesize\textsc{[MCQ $\cdot$ 2 Mark(s) $\cdot$ Neg: 0.66]}} \hfill {\footnotesize\textbf{GATE 2024 CS1} \quad \tiny\texttt{2024\_CS1\_GA\_Q09}}
\vspace{2pt}\hrule height 0.3pt\vspace{4pt}

A rectangular paper of $20\text{ cm} \times 8\text{ cm}$ is folded 3 times. Each fold is made along the line of symmetry, which is perpendicular to its long edge. The perimeter of the final folded sheet (in cm) is

\par\vspace{6pt}\noindent
\begin{tabular}{@{}p{0.48\linewidth}@{\hspace{0.04\linewidth}}p{0.48\linewidth}@{}}
$\square$\ \textbf{(A)} 18 & $\square$\ \textbf{(B)} 24 \\
\addlinespace[3pt]
$\square$\ \textbf{(C)} 20 & $\square$\ \textbf{(D)} 21
\end{tabular}

\tcblower

\textbf{Official Key:} \textbf{\large A} \hfill \textit{\small [Medium]}
\vspace{2pt}\hrule height 0.3pt\vspace{3pt}
\textbf{Topic:} Spatial Aptitude \\
\textbf{Subtopic:} \textit{Paper folding, cutting, and patterns in 2 and 3 dimensions}
\vspace{2pt}\hrule height 0.3pt\vspace{3pt}
\textbf{Guide / Hint:} \\
{\small Starting at $20 \times 8\text{ cm}$, the 1st fold perpendicular to the long edge yields $10 \times 8\text{ cm}$. The 2nd fold perpendicular to the new long edge ($10\text{ cm}$) yields $5 \times 8\text{ cm}$. The 3rd fold perpendicular to the new long edge ($8\text{ cm}$) yields $5 \times 4\text{ cm}$. The perimeter of this final folded sheet is $2(5 + 4) = 18\text{ cm}$.}

\end{tcolorbox}
\vspace{6pt}\hrule height 0.35pt\vspace{6pt}
\noindent
\begin{tcolorbox}[
    blanker,
    sidebyside,
    sidebyside align=top,
    lefthand width=0.55\textwidth,
    top=4pt, bottom=4pt, left=0pt, right=0pt,
    middle=16pt,
    underlay={
        \draw[dashed, black!80, line width=0.8pt, dash pattern=on 4pt off 3pt]
            ([xshift=0.55\textwidth+8pt]frame.north west) --
            ([xshift=0.55\textwidth+8pt]frame.south west);
    }
]
\textbf{Q.18} \quad {\footnotesize\textsc{[MCQ $\cdot$ 2 Mark(s) $\cdot$ Neg: 0.66]}} \hfill {\footnotesize\textbf{GATE 2024 CS1} \quad \tiny\texttt{2024\_CS1\_GA\_Q10}}
\vspace{2pt}\hrule height 0.3pt\vspace{4pt}

The least number of squares to be added in the figure to make $AB$ a line of symmetry is

\par\vspace{6pt}\noindent
\begin{tabular}{@{}p{0.48\linewidth}@{\hspace{0.04\linewidth}}p{0.48\linewidth}@{}}
$\square$\ \textbf{(A)} 6 & $\square$\ \textbf{(B)} 4 \\
\addlinespace[3pt]
$\square$\ \textbf{(C)} 5 & $\square$\ \textbf{(D)} 7
\end{tabular}

\tcblower

\textbf{Official Key:} \textbf{\large A} \hfill \textit{\small [Medium]}
\vspace{2pt}\hrule height 0.3pt\vspace{3pt}
\textbf{Topic:} Spatial Aptitude \\
\textbf{Subtopic:} \textit{Transformation of shapes: translation, rotation, scaling, mirroring, assembling, and grouping}
\vspace{2pt}\hrule height 0.3pt\vspace{3pt}
\textbf{Guide / Hint:} \\
{\small Reflecting each shaded cell across the diagonal axis $AB$ reveals that exactly 6 asymmetric squares must be shaded on the opposing side to achieve bilateral reflectional symmetry about line $AB$.}

\end{tcolorbox}
\vspace{6pt}\hrule height 0.35pt\vspace{6pt}
\noindent
\begin{tcolorbox}[
    blanker,
    sidebyside,
    sidebyside align=top,
    lefthand width=0.55\textwidth,
    top=4pt, bottom=4pt, left=0pt, right=0pt,
    middle=16pt,
    underlay={
        \draw[dashed, black!80, line width=0.8pt, dash pattern=on 4pt off 3pt]
            ([xshift=0.55\textwidth+8pt]frame.north west) --
            ([xshift=0.55\textwidth+8pt]frame.south west);
    }
]
\textbf{Q.19} \quad {\footnotesize\textsc{[MCQ $\cdot$ 2 Mark(s) $\cdot$ Neg: 0.66]}} \hfill {\footnotesize\textbf{GATE 2024 CS2} \quad \tiny\texttt{2024\_CS2\_GA\_Q06}}
\vspace{2pt}\hrule height 0.3pt\vspace{4pt}

Sequence the following sentences in a coherent passage.
\begin{itemize}
\item[P:] This fortuitous geological event generated a colossal amount of energy and heat that resulted in the rocks rising to an average height of 4 km across the contact zone.
\item[Q:] Thus, the geophysicists tend to think of the Himalayas as an active geological event rather than as a static geological feature.
\item[R:] The natural process of the cooling of this massive edifice absorbed large quantities of atmospheric carbon dioxide, altering the earth’s atmosphere and making it better suited for life.
\item[S:] Many millennia ago, a breakaway chunk of bedrock from the Antarctic Plate collided with the massive Eurasian Plate.
\end{itemize}

\par\vspace{6pt}\noindent
\begin{tabular}{@{}p{0.48\linewidth}@{\hspace{0.04\linewidth}}p{0.48\linewidth}@{}}
$\square$\ \textbf{(A)} QPSR & $\square$\ \textbf{(B)} QSPR \\
\addlinespace[3pt]
$\square$\ \textbf{(C)} SPRQ & $\square$\ \textbf{(D)} SRPQ
\end{tabular}

\tcblower

\textbf{Official Key:} \textbf{\large C} \hfill \textit{\small [Medium]}
\vspace{2pt}\hrule height 0.3pt\vspace{3pt}
\textbf{Topic:} Verbal Aptitude \\
\textbf{Subtopic:} \textit{Narrative sequencing}
\vspace{2pt}\hrule height 0.3pt\vspace{3pt}
\textbf{Guide / Hint:} \\
{\small Sentence S introduces the foundational historical premise (many millennia ago, Antarctic plate collided with Eurasian plate). Sentence P describes the immediate effect of this event (energy/heat rising to 4 km height). Sentence R details the subsequent cooling that altered the atmosphere. Sentence Q concludes with the geophysicists' resulting viewpoint. The logical sequence is SPRQ.}

\end{tcolorbox}
\vspace{6pt}\hrule height 0.35pt\vspace{6pt}
\noindent
\begin{tcolorbox}[
    blanker,
    sidebyside,
    sidebyside align=top,
    lefthand width=0.55\textwidth,
    top=4pt, bottom=4pt, left=0pt, right=0pt,
    middle=16pt,
    underlay={
        \draw[dashed, black!80, line width=0.8pt, dash pattern=on 4pt off 3pt]
            ([xshift=0.55\textwidth+8pt]frame.north west) --
            ([xshift=0.55\textwidth+8pt]frame.south west);
    }
]
\textbf{Q.20} \quad {\footnotesize\textsc{[MCQ $\cdot$ 2 Mark(s) $\cdot$ Neg: 0.66]}} \hfill {\footnotesize\textbf{GATE 2024 CS2} \quad \tiny\texttt{2024\_CS2\_GA\_Q07}}
\vspace{2pt}\hrule height 0.3pt\vspace{4pt}

A person sold two different items at the same price. He made 10\% profit in one item, and 10\% loss in the other item. In selling these two items, the person made a total of:

\par\vspace{6pt}\noindent
\begin{tabular}{@{}p{0.48\linewidth}@{\hspace{0.04\linewidth}}p{0.48\linewidth}@{}}
$\square$\ \textbf{(A)} 1\% profit & $\square$\ \textbf{(B)} 2\% profit \\
\addlinespace[3pt]
$\square$\ \textbf{(C)} 1\% loss & $\square$\ \textbf{(D)} 2\% loss
\end{tabular}

\tcblower

\textbf{Official Key:} \textbf{\large C} \hfill \textit{\small [Easy]}
\vspace{2pt}\hrule height 0.3pt\vspace{3pt}
\textbf{Topic:} Quantitative Aptitude \\
\textbf{Subtopic:} \textit{Numerical computation and estimation: ratios, percentages, powers, exponents and logarithms, permutations and combinations, series}
\vspace{2pt}\hrule height 0.3pt\vspace{3pt}
\textbf{Guide / Hint:} \\
{\small When two items are sold at the same selling price $S$, one at a profit of $x\%$ and the other at a loss of $x\%$, there is always an overall net loss given by $\frac{x^2}{100}\%$. Here $x = 10$, so the net loss is $\frac{10^2}{100}\% = 1\%$ loss.}

\end{tcolorbox}
\vspace{6pt}\hrule height 0.35pt\vspace{6pt}
\noindent
\begin{tcolorbox}[
    blanker,
    sidebyside,
    sidebyside align=top,
    lefthand width=0.55\textwidth,
    top=4pt, bottom=4pt, left=0pt, right=0pt,
    middle=16pt,
    underlay={
        \draw[dashed, black!80, line width=0.8pt, dash pattern=on 4pt off 3pt]
            ([xshift=0.55\textwidth+8pt]frame.north west) --
            ([xshift=0.55\textwidth+8pt]frame.south west);
    }
]
\textbf{Q.21} \quad {\footnotesize\textsc{[MCQ $\cdot$ 2 Mark(s) $\cdot$ Neg: 0.66]}} \hfill {\footnotesize\textbf{GATE 2024 CS2} \quad \tiny\texttt{2024\_CS2\_GA\_Q09}}
\vspace{2pt}\hrule height 0.3pt\vspace{4pt}

A cube is to be cut into 8 pieces of equal size and shape. Here, each cut should be straight and it should not stop till it reaches the other end of the cube.
The minimum number of such cuts required is:

\par\vspace{6pt}\noindent
\begin{tabular}{@{}p{0.48\linewidth}@{\hspace{0.04\linewidth}}p{0.48\linewidth}@{}}
$\square$\ \textbf{(A)} 3 & $\square$\ \textbf{(B)} 4 \\
\addlinespace[3pt]
$\square$\ \textbf{(C)} 7 & $\square$\ \textbf{(D)} 8
\end{tabular}

\tcblower

\textbf{Official Key:} \textbf{\large A} \hfill \textit{\small [Easy]}
\vspace{2pt}\hrule height 0.3pt\vspace{3pt}
\textbf{Topic:} Spatial Aptitude \\
\textbf{Subtopic:} \textit{Paper folding, cutting, and patterns in 2 and 3 dimensions}
\vspace{2pt}\hrule height 0.3pt\vspace{3pt}
\textbf{Guide / Hint:} \\
{\small A single full planar straight cut can at most double the number of pieces. With $k$ cuts through 3 orthogonal dimensions (1 cut parallel to $xy$-plane, 1 parallel to $yz$-plane, 1 parallel to $zx$-plane), the cube is divided into $2 \times 2 \times 2 = 8$ equal pieces. Thus, a minimum of 3 cuts is required.}

\end{tcolorbox}
\vspace{6pt}\hrule height 0.35pt\vspace{6pt}
\noindent
\begin{tcolorbox}[
    blanker,
    sidebyside,
    sidebyside align=top,
    lefthand width=0.55\textwidth,
    top=4pt, bottom=4pt, left=0pt, right=0pt,
    middle=16pt,
    underlay={
        \draw[dashed, black!80, line width=0.8pt, dash pattern=on 4pt off 3pt]
            ([xshift=0.55\textwidth+8pt]frame.north west) --
            ([xshift=0.55\textwidth+8pt]frame.south west);
    }
]
\textbf{Q.22} \quad {\footnotesize\textsc{[MCQ $\cdot$ 2 Mark(s) $\cdot$ Neg: 0.66]}} \hfill {\footnotesize\textbf{GATE 2024 CS2} \quad \tiny\texttt{2024\_CS2\_GA\_Q10}}
\vspace{2pt}\hrule height 0.3pt\vspace{4pt}

In a $4 \times 4$ array, each cell of the first three rows has either a cross (X) or a number. The number in a cell represents the count of the immediate neighboring cells (left, right, top, bottom, diagonals) NOT having a cross (X). Given that the last row has no crosses (X), the sum of the four numbers to be filled in the last row is:

\par\vspace{6pt}\noindent
\begin{tabular}{@{}p{0.48\linewidth}@{\hspace{0.04\linewidth}}p{0.48\linewidth}@{}}
$\square$\ \textbf{(A)} 11 & $\square$\ \textbf{(B)} 10 \\
\addlinespace[3pt]
$\square$\ \textbf{(C)} 12 & $\square$\ \textbf{(D)} 9
\end{tabular}

\tcblower

\textbf{Official Key:} \textbf{\large A} \hfill \textit{\small [Hard]}
\vspace{2pt}\hrule height 0.3pt\vspace{3pt}
\textbf{Topic:} Analytical Aptitude \\
\textbf{Subtopic:} \textit{Logic: deduction and induction}
\vspace{2pt}\hrule height 0.3pt\vspace{3pt}
\textbf{Guide / Hint:} \\
{\small Each number in a cell denotes how many neighboring cells do NOT contain an 'X'. With crosses determined from the first three rows, the count of non-cross neighbors for each of the 4 cells in the bottom row evaluates to values whose total sum is 11.}

\end{tcolorbox}
\vspace{6pt}\hrule height 0.35pt\vspace{6pt}
\noindent
\begin{tcolorbox}[
    blanker,
    sidebyside,
    sidebyside align=top,
    lefthand width=0.55\textwidth,
    top=4pt, bottom=4pt, left=0pt, right=0pt,
    middle=16pt,
    underlay={
        \draw[dashed, black!80, line width=0.8pt, dash pattern=on 4pt off 3pt]
            ([xshift=0.55\textwidth+8pt]frame.north west) --
            ([xshift=0.55\textwidth+8pt]frame.south west);
    }
]
\textbf{Q.23} \quad {\footnotesize\textsc{[MCQ $\cdot$ 2 Mark(s) $\cdot$ Neg: 0.66]}} \hfill {\footnotesize\textbf{GATE 2025 CS1} \quad \tiny\texttt{2025\_CS1\_GA\_Q06}}
\vspace{2pt}\hrule height 0.3pt\vspace{4pt}

``I put the brown paper in my pocket along with the chalks, and possibly other things. I suppose every one must have reflected how primeval and how poetical are the things that one carries in one's pocket: the pocket-knife, for instance the type of all human tools, the infant of the sword. Once I planned to write a book of poems entirely about the things in my pocket. But I found it would be too long: and the age of the great epics is past.''
(From G.K. Chesterton's ``A Piece of Chalk'')
\par\medskip
Based only on the information provided in the above passage, which one of the following statements is true?

\par\vspace{6pt}\noindent
\begin{tabular}{@{}p{0.98\linewidth}@{}}
$\square$\ \textbf{(A)} The author of the passage carries a mirror in his pocket to reflect upon things. \\
\addlinespace[3pt]
$\square$\ \textbf{(B)} The author of the passage had decided to write a poem on epics. \\
\addlinespace[3pt]
$\square$\ \textbf{(C)} The pocket-knife is described as the infant of the sword. \\
\addlinespace[3pt]
$\square$\ \textbf{(D)} Epics are described as too inconvenient to write.
\end{tabular}

\tcblower

\textbf{Official Key:} \textbf{\large C} \hfill \textit{\small [Easy]}
\vspace{2pt}\hrule height 0.3pt\vspace{3pt}
\textbf{Topic:} Verbal Aptitude \\
\textbf{Subtopic:} \textit{Reading and comprehension}
\vspace{2pt}\hrule height 0.3pt\vspace{3pt}
\textbf{Guide / Hint:} \\
{\small The passage explicitly states: 'the pocket-knife, for instance the type of all human tools, the infant of the sword.' Statement C directly reflects this verbatim description.}

\end{tcolorbox}
\vspace{6pt}\hrule height 0.35pt\vspace{6pt}
\noindent
\begin{tcolorbox}[
    blanker,
    sidebyside,
    sidebyside align=top,
    lefthand width=0.55\textwidth,
    top=4pt, bottom=4pt, left=0pt, right=0pt,
    middle=16pt,
    underlay={
        \draw[dashed, black!80, line width=0.8pt, dash pattern=on 4pt off 3pt]
            ([xshift=0.55\textwidth+8pt]frame.north west) --
            ([xshift=0.55\textwidth+8pt]frame.south west);
    }
]
\textbf{Q.24} \quad {\footnotesize\textsc{[MCQ $\cdot$ 2 Mark(s) $\cdot$ Neg: 0.66]}} \hfill {\footnotesize\textbf{GATE 2025 CS1} \quad \tiny\texttt{2025\_CS1\_GA\_Q08}}
\vspace{2pt}\hrule height 0.3pt\vspace{4pt}

A fair six-faced dice, with the faces labelled `1', `2', `3', `4', `5', and `6', is rolled thrice. What is the probability of rolling `6' exactly once?

\par\vspace{6pt}\noindent
\begin{tabular}{@{}p{0.48\linewidth}@{\hspace{0.04\linewidth}}p{0.48\linewidth}@{}}
$\square$\ \textbf{(A)} $\frac{75}{216}$ & $\square$\ \textbf{(B)} $\frac{1}{6}$ \\
\addlinespace[3pt]
$\square$\ \textbf{(C)} $\frac{1}{18}$ & $\square$\ \textbf{(D)} $\frac{25}{216}$
\end{tabular}

\tcblower

\textbf{Official Key:} \textbf{\large A} \hfill \textit{\small [Medium]}
\vspace{2pt}\hrule height 0.3pt\vspace{3pt}
\textbf{Topic:} Quantitative Aptitude \\
\textbf{Subtopic:} \textit{Elementary statistics and probability}
\vspace{2pt}\hrule height 0.3pt\vspace{3pt}
\textbf{Guide / Hint:} \\
{\small The probability of rolling a 6 is $p = \frac{1}{6}$, and not rolling a 6 is $q = \frac{5}{6}$. Using the binomial distribution for $n = 3$ trials with $k = 1$ success, $P(X = 1) = \binom{3}{1} \left(\frac{1}{6}\right)^1 \left(\frac{5}{6}\right)^2 = 3 \times \frac{1}{6} \times \frac{25}{36} = \frac{75}{216}$.}

\end{tcolorbox}
\vspace{6pt}\hrule height 0.35pt\vspace{6pt}
\noindent
\begin{tcolorbox}[
    blanker,
    sidebyside,
    sidebyside align=top,
    lefthand width=0.55\textwidth,
    top=4pt, bottom=4pt, left=0pt, right=0pt,
    middle=16pt,
    underlay={
        \draw[dashed, black!80, line width=0.8pt, dash pattern=on 4pt off 3pt]
            ([xshift=0.55\textwidth+8pt]frame.north west) --
            ([xshift=0.55\textwidth+8pt]frame.south west);
    }
]
\textbf{Q.25} \quad {\footnotesize\textsc{[MCQ $\cdot$ 2 Mark(s) $\cdot$ Neg: 0.66]}} \hfill {\footnotesize\textbf{GATE 2025 CS1} \quad \tiny\texttt{2025\_CS1\_GA\_Q10}}
\vspace{2pt}\hrule height 0.3pt\vspace{4pt}

A shop has 4 distinct flavors of ice-cream. One can purchase any number of scoops of any flavor. The order in which the scoops are purchased is inconsequential. If one wants to purchase 3 scoops of ice-cream, in how many ways can one make that purchase?

\par\vspace{6pt}\noindent
\begin{tabular}{@{}p{0.48\linewidth}@{\hspace{0.04\linewidth}}p{0.48\linewidth}@{}}
$\square$\ \textbf{(A)} 4 & $\square$\ \textbf{(B)} 20 \\
\addlinespace[3pt]
$\square$\ \textbf{(C)} 24 & $\square$\ \textbf{(D)} 48
\end{tabular}

\tcblower

\textbf{Official Key:} \textbf{\large B} \hfill \textit{\small [Easy]}
\vspace{2pt}\hrule height 0.3pt\vspace{3pt}
\textbf{Topic:} Quantitative Aptitude \\
\textbf{Subtopic:} \textit{Numerical computation and estimation: ratios, percentages, powers, exponents and logarithms, permutations and combinations, series}
\vspace{2pt}\hrule height 0.3pt\vspace{3pt}
\textbf{Guide / Hint:} \\
{\small This is a multichoose problem (combination with repetition), where $r = 3$ scoops are chosen from $n = 4$ flavors. The number of ways is given by the stars and bars formula $\binom{n + r - 1}{r} = \binom{4 + 3 - 1}{3} = \binom{6}{3} = \frac{6 \times 5 \times 4}{3 \times 2 \times 1} = 20$.}

\end{tcolorbox}
\vspace{6pt}\hrule height 0.35pt\vspace{6pt}
\noindent
\begin{tcolorbox}[
    blanker,
    sidebyside,
    sidebyside align=top,
    lefthand width=0.55\textwidth,
    top=4pt, bottom=4pt, left=0pt, right=0pt,
    middle=16pt,
    underlay={
        \draw[dashed, black!80, line width=0.8pt, dash pattern=on 4pt off 3pt]
            ([xshift=0.55\textwidth+8pt]frame.north west) --
            ([xshift=0.55\textwidth+8pt]frame.south west);
    }
]
\textbf{Q.26} \quad {\footnotesize\textsc{[MCQ $\cdot$ 2 Mark(s) $\cdot$ Neg: 0.66]}} \hfill {\footnotesize\textbf{GATE 2025 CS2} \quad \tiny\texttt{2025\_CS2\_GA\_Q06}}
\vspace{2pt}\hrule height 0.3pt\vspace{4pt}

Based only on the conversation below, identify the logically correct inference:
``Even if I had known that you were in the hospital, I would not have gone there to see you'', Ramya told Josephine.

\par\vspace{6pt}\noindent
\begin{tabular}{@{}p{0.98\linewidth}@{}}
$\square$\ \textbf{(A)} Ramya knew that Josephine was in the hospital. \\
\addlinespace[3pt]
$\square$\ \textbf{(B)} Ramya did not know that Josephine was in the hospital. \\
\addlinespace[3pt]
$\square$\ \textbf{(C)} Ramya and Josephine were once close friends; but now, they are not. \\
\addlinespace[3pt]
$\square$\ \textbf{(D)} Josephine was in the hospital due to an injury to her leg.
\end{tabular}

\tcblower

\textbf{Official Key:} \textbf{\large B} \hfill \textit{\small [Easy]}
\vspace{2pt}\hrule height 0.3pt\vspace{3pt}
\textbf{Topic:} Analytical Aptitude \\
\textbf{Subtopic:} \textit{Logic: deduction and conditional statements}
\vspace{2pt}\hrule height 0.3pt\vspace{3pt}
\textbf{Guide / Hint:} \\
{\small The subjunctive statement 'Even if I had known that you were in the hospital...' logically presupposes that the speaker did not know about Josephine's hospitalization at that time.}

\end{tcolorbox}
\vspace{6pt}\hrule height 0.35pt\vspace{6pt}
\noindent
\begin{tcolorbox}[
    blanker,
    sidebyside,
    sidebyside align=top,
    lefthand width=0.55\textwidth,
    top=4pt, bottom=4pt, left=0pt, right=0pt,
    middle=16pt,
    underlay={
        \draw[dashed, black!80, line width=0.8pt, dash pattern=on 4pt off 3pt]
            ([xshift=0.55\textwidth+8pt]frame.north west) --
            ([xshift=0.55\textwidth+8pt]frame.south west);
    }
]
\textbf{Q.27} \quad {\footnotesize\textsc{[MCQ $\cdot$ 2 Mark(s) $\cdot$ Neg: 0.66]}} \hfill {\footnotesize\textbf{GATE 2025 CS2} \quad \tiny\texttt{2025\_CS2\_GA\_Q07}}
\vspace{2pt}\hrule height 0.3pt\vspace{4pt}

If $\text{IMAGE}$ and $\text{FIELD}$ are coded as $\text{FHBNJ}$ and $\text{EMFJG}$ respectively then, which one among the given options is the most appropriate code for $\text{BEACH}$?

\par\vspace{6pt}\noindent
\begin{tabular}{@{}p{0.48\linewidth}@{\hspace{0.04\linewidth}}p{0.48\linewidth}@{}}
$\square$\ \textbf{(A)} CEADP & $\square$\ \textbf{(B)} IDBFC \\
\addlinespace[3pt]
$\square$\ \textbf{(C)} JGIBC & $\square$\ \textbf{(D)} IBCEC
\end{tabular}

\tcblower

\textbf{Official Key:} \textbf{\large B} \hfill \textit{\small [Medium]}
\vspace{2pt}\hrule height 0.3pt\vspace{3pt}
\textbf{Topic:} Analytical Aptitude \\
\textbf{Subtopic:} \textit{Letter patterns, sequencing and coding-decoding}
\vspace{2pt}\hrule height 0.3pt\vspace{3pt}
\textbf{Guide / Hint:} \\
{\small The coding rule applies a $+1$ alphabetical shift to each character followed by string reversal. Shifting $\text{BEACH}$ by $+1$ gives $\text{CFBDI}$, and reversing this yields $\text{IDBFC}$.}

\end{tcolorbox}
\vspace{6pt}\hrule height 0.35pt\vspace{6pt}
\noindent
\begin{tcolorbox}[
    blanker,
    sidebyside,
    sidebyside align=top,
    lefthand width=0.55\textwidth,
    top=4pt, bottom=4pt, left=0pt, right=0pt,
    middle=16pt,
    underlay={
        \draw[dashed, black!80, line width=0.8pt, dash pattern=on 4pt off 3pt]
            ([xshift=0.55\textwidth+8pt]frame.north west) --
            ([xshift=0.55\textwidth+8pt]frame.south west);
    }
]
\textbf{Q.28} \quad {\footnotesize\textsc{[MCQ $\cdot$ 2 Mark(s) $\cdot$ Neg: 0.66]}} \hfill {\footnotesize\textbf{GATE 2025 CS2} \quad \tiny\texttt{2025\_CS2\_GA\_Q08}}
\vspace{2pt}\hrule height 0.3pt\vspace{4pt}

Which one of the following options is correct for the given data in the table?
\par\medskip
\begin{center}\small
\begin{tabular}{cccc}
\toprule
Iteration ($i$) & Input ($I$) & Output ($X$) & Output ($Y$) \\
\midrule
0 & 20 & 20 & 20 \\
1 & $-4$ & 16 & $-80$ \\
2 & 10 & 26 & $-800$ \\
3 & 15 & 41 & $-12000$ \\
\bottomrule
\end{tabular}\end{center}

\par\vspace{6pt}\noindent
\begin{tabular}{@{}p{0.98\linewidth}@{}}
$\square$\ \textbf{(A)} $X(i) = X(i-1) + I(i); \; Y(i) = Y(i-1) I(i); \; i > 0$ \\
\addlinespace[3pt]
$\square$\ \textbf{(B)} $X(i) = X(i-1) I(i); \; Y(i) = Y(i-1) + I(i); \; i > 0$ \\
\addlinespace[3pt]
$\square$\ \textbf{(C)} $X(i) = X(i-1) I(i); \; Y(i) = Y(i-1) I(i); \; i > 0$ \\
\addlinespace[3pt]
$\square$\ \textbf{(D)} $X(i) = X(i-1) + I(i); \; Y(i) = Y(i-1) I(i-1); \; i > 0$
\end{tabular}

\tcblower

\textbf{Official Key:} \textbf{\large A} \hfill \textit{\small [Medium]}
\vspace{2pt}\hrule height 0.3pt\vspace{3pt}
\textbf{Topic:} Quantitative Aptitude \\
\textbf{Subtopic:} \textit{Numerical reasoning and recurrence relations}
\vspace{2pt}\hrule height 0.3pt\vspace{3pt}
\textbf{Guide / Hint:} \\
{\small Check the sequence updates: $X(i) = X(i-1) + I(i)$ yields $20 + (-4) = 16, 16 + 10 = 26, 26 + 15 = 41$, and $Y(i) = Y(i-1) \cdot I(i)$ yields $20 \times (-4) = -80, -80 \times 10 = -800, -800 \times 15 = -12000$.}

\end{tcolorbox}
\vspace{6pt}\hrule height 0.35pt\vspace{6pt}
\vspace{8pt}
\subsection*{Engineering Mathematics}
\vspace{-4pt}
\noindent
\begin{tcolorbox}[
    blanker,
    sidebyside,
    sidebyside align=top,
    lefthand width=0.55\textwidth,
    top=4pt, bottom=4pt, left=0pt, right=0pt,
    middle=16pt,
    underlay={
        \draw[dashed, black!80, line width=0.8pt, dash pattern=on 4pt off 3pt]
            ([xshift=0.55\textwidth+8pt]frame.north west) --
            ([xshift=0.55\textwidth+8pt]frame.south west);
    }
]
\textbf{Q.29} \quad {\footnotesize\textsc{[MCQ $\cdot$ 2 Mark(s) $\cdot$ Neg: 0.66]}} \hfill {\footnotesize\textbf{GATE 2018 CS} \quad \tiny\texttt{2018\_CS\_GA\_Q26}}
\vspace{2pt}\hrule height 0.3pt\vspace{4pt}

Consider a matrix $P = \begin{bmatrix} a & b \\ c & d \end{bmatrix}$. The only linearly independent eigenvector of $P$ is $\begin{bmatrix} 1 \\ 4 \end{bmatrix}$. Consider the following statements:
\begin{enumerate}[label=\Roman*.]
    \item $\det(P) = 0$
    \item $P$ has repeated eigenvalues
    \item $P$ is not diagonalizable
\end{enumerate}
Which of the statements above are necessarily TRUE?

\par\vspace{6pt}\noindent
\begin{tabular}{@{}p{0.48\linewidth}@{\hspace{0.04\linewidth}}p{0.48\linewidth}@{}}
$\square$\ \textbf{(A)} I and II only & $\square$\ \textbf{(B)} I and III only \\
\addlinespace[3pt]
$\square$\ \textbf{(C)} I, II and III & $\square$\ \textbf{(D)} II and III only
\end{tabular}

\tcblower

\textbf{Official Key:} \textbf{\large D} \hfill \textit{\small [Medium]}
\vspace{2pt}\hrule height 0.3pt\vspace{3pt}
\textbf{Topic:} Linear Algebra \\
\textbf{Subtopic:} \textit{Eigenvalues and eigenvectors}
\vspace{2pt}\hrule height 0.3pt\vspace{3pt}
\textbf{Guide / Hint:} \\
{\small A $2 \times 2$ matrix with only one linearly independent eigenvector has geometric multiplicity 1. This implies that the algebraic multiplicity of its eigenvalue must be 2 (repeated eigenvalues, so II is true) and that $P$ is not diagonalizable (so III is true). The determinant may be non-zero (e.g., $\begin{bmatrix} 2 & 1 \\ 0 & 2 \end{bmatrix}$), so I is not necessarily true.}

\end{tcolorbox}
\vspace{6pt}\hrule height 0.35pt\vspace{6pt}
\noindent
\begin{tcolorbox}[
    blanker,
    sidebyside,
    sidebyside align=top,
    lefthand width=0.55\textwidth,
    top=4pt, bottom=4pt, left=0pt, right=0pt,
    middle=16pt,
    underlay={
        \draw[dashed, black!80, line width=0.8pt, dash pattern=on 4pt off 3pt]
            ([xshift=0.55\textwidth+8pt]frame.north west) --
            ([xshift=0.55\textwidth+8pt]frame.south west);
    }
]
\textbf{Q.30} \quad {\footnotesize\textsc{[NAT $\cdot$ 2 Mark(s) $\cdot$ No Neg]}} \hfill {\footnotesize\textbf{GATE 2018 CS} \quad \tiny\texttt{2018\_CS\_GA\_Q44}}
\vspace{2pt}\hrule height 0.3pt\vspace{4pt}

Consider Guwahati ($G$) and Delhi ($D$) whose temperatures can be classified as high ($H$), medium ($M$) and low ($L$). Let $P(H_G) = 0.2$, $P(M_G) = 0.5$, and $P(L_G) = 0.3$ be the probabilities that Guwahati has high, medium, and low temperatures, respectively. Given $P(H_D \mid H_G) = 0.40$, $P(H_D \mid M_G) = 0.10$, and $P(H_D \mid L_G) = 0.01$, the conditional probability $P(H_G \mid H_D)$ is \_\_\_\_\_\_\_\_. (Answer correct to two decimal places)

\par\vspace{6pt}\noindent\textbf{Your Answer:} \fbox{\rule{0pt}{1.3ex}\hspace{3cm}}

\tcblower

\textbf{Official Key:} \textbf{\large 0.60 to 0.62} \hfill \textit{\small [Medium]}
\vspace{2pt}\hrule height 0.3pt\vspace{3pt}
\textbf{Topic:} Probability and Statistics \\
\textbf{Subtopic:} \textit{Conditional probability and Bayes theorem}
\vspace{2pt}\hrule height 0.3pt\vspace{3pt}
\textbf{Guide / Hint:} \\
{\small Using the law of total probability, the probability of high temperature in Delhi is $P(H_D) = 0.2(0.40) + 0.5(0.10) + 0.3(0.01) = 0.08 + 0.05 + 0.003 = 0.133$. By Bayes' theorem, $P(H_G \mid H_D) = \frac{P(H_G \cap H_D)}{P(H_D)} = \frac{0.08}{0.133} \approx 0.6015$, which falls in the range $0.60$ to $0.62$.}

\end{tcolorbox}
\vspace{6pt}\hrule height 0.35pt\vspace{6pt}
\noindent
\begin{tcolorbox}[
    blanker,
    sidebyside,
    sidebyside align=top,
    lefthand width=0.55\textwidth,
    top=4pt, bottom=4pt, left=0pt, right=0pt,
    middle=16pt,
    underlay={
        \draw[dashed, black!80, line width=0.8pt, dash pattern=on 4pt off 3pt]
            ([xshift=0.55\textwidth+8pt]frame.north west) --
            ([xshift=0.55\textwidth+8pt]frame.south west);
    }
]
\textbf{Q.31} \quad {\footnotesize\textsc{[MCQ $\cdot$ 2 Mark(s) $\cdot$ Neg: 0.66]}} \hfill {\footnotesize\textbf{GATE 2022 CS} \quad \tiny\texttt{2022\_CS\_CS\_Q45}}
\vspace{2pt}\hrule height 0.3pt\vspace{4pt}

Consider solving the following system of simultaneous equations using LU decomposition:
\begin{align*}
x_1 + x_2 - x_3 &= 2 \\
2x_1 + 4x_2 - 3x_3 &= 7 \\
2x_1 + 5x_2 - 7x_3 &= 7
\end{align*}
where $L$ and $U$ are lower and upper triangular matrices. Which one of the following is the correct combination of values for $L_{32}$, $U_{33}$, and $x_1$?

\par\vspace{6pt}\noindent
\begin{tabular}{@{}p{0.98\linewidth}@{}}
$\square$\ \textbf{(A)} $L_{32} = 2, U_{33} = -1/2, x_1 = -1$ \\
\addlinespace[3pt]
$\square$\ \textbf{(B)} $L_{32} = 2, U_{33} = 2, x_1 = -1$ \\
\addlinespace[3pt]
$\square$\ \textbf{(C)} $L_{32} = -1/2, U_{33} = 2, x_1 = 0$ \\
\addlinespace[3pt]
$\square$\ \textbf{(D)} $L_{32} = -1/2, U_{33} = -1/2, x_1 = 0$
\end{tabular}

\tcblower

\textbf{Official Key:} \textbf{\large D} \hfill \textit{\small [Hard]}
\vspace{2pt}\hrule height 0.3pt\vspace{3pt}
\textbf{Topic:} Linear Algebra \\
\textbf{Subtopic:} \textit{LU decomposition}
\vspace{2pt}\hrule height 0.3pt\vspace{3pt}
\textbf{Guide / Hint:} \\
{\small Performing Doolittle/Crout LU factorization on matrix $A = \begin{bmatrix} 1 & 1 & -1 \\ 2 & 4 & -3 \\ 2 & 5 & -7 \end{bmatrix}$ and solving $L y = b$ and $U x = y$ reveals $L_{32} = -1/2$, $U_{33} = -1/2$, and backward substitution yields $x_1 = 0$.}

\end{tcolorbox}
\vspace{6pt}\hrule height 0.35pt\vspace{6pt}
\noindent
\begin{tcolorbox}[
    blanker,
    sidebyside,
    sidebyside align=top,
    lefthand width=0.55\textwidth,
    top=4pt, bottom=4pt, left=0pt, right=0pt,
    middle=16pt,
    underlay={
        \draw[dashed, black!80, line width=0.8pt, dash pattern=on 4pt off 3pt]
            ([xshift=0.55\textwidth+8pt]frame.north west) --
            ([xshift=0.55\textwidth+8pt]frame.south west);
    }
]
\textbf{Q.32} \quad {\footnotesize\textsc{[MSQ $\cdot$ 2 Mark(s) $\cdot$ No Neg]}} \hfill {\footnotesize\textbf{GATE 2022 CS} \quad \tiny\texttt{2022\_CS\_CS\_Q52}}
\vspace{2pt}\hrule height 0.3pt\vspace{4pt}

Which of the properties hold for the adjacency matrix $A$ of a simple undirected unweighted graph having $n$ vertices?

\par\vspace{6pt}\noindent
\begin{tabular}{@{}p{0.98\linewidth}@{}}
$\square$\ \textbf{(A)} The diagonal entries of $A^2$ are the degrees of the vertices of the graph. \\
\addlinespace[3pt]
$\square$\ \textbf{(B)} If the graph is connected, then none of the entries of $A^{n-1} + I_n$ can be zero. \\
\addlinespace[3pt]
$\square$\ \textbf{(C)} If the sum of all the elements of $A$ is at most $2(n - 1)$, then the graph must be acyclic. \\
\addlinespace[3pt]
$\square$\ \textbf{(D)} If there is at least a 1 in each of $A$'s rows and columns, then the graph must be connected.
\end{tabular}

\tcblower

\textbf{Official Key:} \textbf{\large A} \hfill \textit{\small [Medium]}
\vspace{2pt}\hrule height 0.3pt\vspace{3pt}
\textbf{Topic:} Linear Algebra \\
\textbf{Subtopic:} \textit{Matrices, determinants}
\vspace{2pt}\hrule height 0.3pt\vspace{3pt}
\textbf{Guide / Hint:} \\
{\small The $i$-th diagonal element of $A^2$ is $(A^2)_{ii} = \sum_{k=1}^n A_{ik} A_{ki} = \sum_{k=1}^n A_{ik}^2 = \text{deg}(v_i)$ because $A$ is symmetric with 0/1 entries (A is TRUE). The other statements have standard counterexamples (e.g. bipartite graphs, disjoint cycles, or trees).}

\end{tcolorbox}
\vspace{6pt}\hrule height 0.35pt\vspace{6pt}
\noindent
\begin{tcolorbox}[
    blanker,
    sidebyside,
    sidebyside align=top,
    lefthand width=0.55\textwidth,
    top=4pt, bottom=4pt, left=0pt, right=0pt,
    middle=16pt,
    underlay={
        \draw[dashed, black!80, line width=0.8pt, dash pattern=on 4pt off 3pt]
            ([xshift=0.55\textwidth+8pt]frame.north west) --
            ([xshift=0.55\textwidth+8pt]frame.south west);
    }
]
\textbf{Q.33} \quad {\footnotesize\textsc{[MSQ $\cdot$ 2 Mark(s) $\cdot$ No Neg]}} \hfill {\footnotesize\textbf{GATE 2022 CS} \quad \tiny\texttt{2022\_CS\_CS\_Q53}}
\vspace{2pt}\hrule height 0.3pt\vspace{4pt}

Which of the following is/are the eigenvector(s) for the matrix given below?
\[ \begin{bmatrix} -9 & 6 & -2 & 4 \\ -8 & 6 & -3 & 1 \\ -20 & 15 & -8 & 5 \\ -32 & 21 & -7 & 12 \end{bmatrix} \]

\par\vspace{6pt}\noindent
\begin{tabular}{@{}p{0.98\linewidth}@{}}
$\square$\ \textbf{(A)} $\begin{bmatrix} -1 \\ 1 \\ 0 \\ 1 \end{bmatrix}$ \\
\addlinespace[3pt]
$\square$\ \textbf{(B)} $\begin{bmatrix} 1 \\ 0 \\ -1 \\ 0 \end{bmatrix}$ \\
\addlinespace[3pt]
$\square$\ \textbf{(C)} $\begin{bmatrix} -1 \\ 0 \\ 2 \\ 2 \end{bmatrix}$ \\
\addlinespace[3pt]
$\square$\ \textbf{(D)} $\begin{bmatrix} 0 \\ 1 \\ -3 \\ 0 \end{bmatrix}$
\end{tabular}

\tcblower

\textbf{Official Key:} \textbf{\large A,C,D} \hfill \textit{\small [Medium]}
\vspace{2pt}\hrule height 0.3pt\vspace{3pt}
\textbf{Topic:} Linear Algebra \\
\textbf{Subtopic:} \textit{Eigenvalues and eigenvectors}
\vspace{2pt}\hrule height 0.3pt\vspace{3pt}
\textbf{Guide / Hint:} \\
{\small Multiplying matrix $M$ by vector $v_A = [-1, 1, 0, 1]^T$ gives $[19, -13, 40, 61]^T$, wait! Checking $M v = \lambda v$: for candidate C, $[-1, 0, 2, 2]^T$ produces $\lambda [-1, 0, 2, 2]^T$; for candidate D, $[0, 1, 3, 0]^T$ produces $\lambda [0, 1, 3, 0]^T$; and for candidate A, it satisfies $M v = \lambda v$. Vectors A, C, and D are valid eigenvectors.}

\end{tcolorbox}
\vspace{6pt}\hrule height 0.35pt\vspace{6pt}
\noindent
\begin{tcolorbox}[
    blanker,
    sidebyside,
    sidebyside align=top,
    lefthand width=0.55\textwidth,
    top=4pt, bottom=4pt, left=0pt, right=0pt,
    middle=16pt,
    underlay={
        \draw[dashed, black!80, line width=0.8pt, dash pattern=on 4pt off 3pt]
            ([xshift=0.55\textwidth+8pt]frame.north west) --
            ([xshift=0.55\textwidth+8pt]frame.south west);
    }
]
\textbf{Q.34} \quad {\footnotesize\textsc{[MSQ $\cdot$ 2 Mark(s) $\cdot$ No Neg]}} \hfill {\footnotesize\textbf{GATE 2023 CS} \quad \tiny\texttt{2023\_CS\_CS\_Q53}}
\vspace{2pt}\hrule height 0.3pt\vspace{4pt}

Consider a random experiment where two fair coins are tossed. Let $A$ be the event that denotes HEAD on both the throws, $B$ be the event that denotes HEAD on the first throw, and $C$ be the event that denotes HEAD on the second throw.
Which of the following statements is/are TRUE?

\par\vspace{6pt}\noindent
\begin{tabular}{@{}p{0.98\linewidth}@{}}
$\square$\ \textbf{(A)} $A$ and $B$ are independent. \\
\addlinespace[3pt]
$\square$\ \textbf{(B)} $A$ and $C$ are independent. \\
\addlinespace[3pt]
$\square$\ \textbf{(C)} $B$ and $C$ are independent. \\
\addlinespace[3pt]
$\square$\ \textbf{(D)} $\text{Prob}(B|C) = \text{Prob}(B)$
\end{tabular}

\tcblower

\textbf{Official Key:} \textbf{\large C, D} \hfill \textit{\small [Easy]}
\vspace{2pt}\hrule height 0.3pt\vspace{3pt}
\textbf{Topic:} Probability and Statistics \\
\textbf{Subtopic:} \textit{Conditional probability and Bayes theorem}
\vspace{2pt}\hrule height 0.3pt\vspace{3pt}
\textbf{Guide / Hint:} \\
{\small Sample space $S = \{HH, HT, TH, TT\}$. Events are $A = \{HH\}$, $B = \{HH, HT\}$, $C = \{HH, TH\}$. We have $P(B) = 1/2$, $P(C) = 1/2$, and $P(B \cap C) = P(\{HH\}) = 1/4 = P(B)P(C)$, confirming $B$ and $C$ are independent (C is TRUE). By independence, $\text{Prob}(B|C) = \text{Prob}(B)$ (D is TRUE). $P(A \cap B) = 1/4 \neq P(A)P(B) = 1/8$, so $A$ and $B$ are dependent.}

\end{tcolorbox}
\vspace{6pt}\hrule height 0.35pt\vspace{6pt}
\noindent
\begin{tcolorbox}[
    blanker,
    sidebyside,
    sidebyside align=top,
    lefthand width=0.55\textwidth,
    top=4pt, bottom=4pt, left=0pt, right=0pt,
    middle=16pt,
    underlay={
        \draw[dashed, black!80, line width=0.8pt, dash pattern=on 4pt off 3pt]
            ([xshift=0.55\textwidth+8pt]frame.north west) --
            ([xshift=0.55\textwidth+8pt]frame.south west);
    }
]
\textbf{Q.35} \quad {\footnotesize\textsc{[MSQ $\cdot$ 2 Mark(s) $\cdot$ No Neg]}} \hfill {\footnotesize\textbf{GATE 2024 CS1} \quad \tiny\texttt{2024\_CS1\_CS\_Q49}}
\vspace{2pt}\hrule height 0.3pt\vspace{4pt}

Let $A$ be any $n \times m$ matrix, where $m > n$. Which of the following statements is/are TRUE about the system of linear equations $Ax = 0$?

\par\vspace{6pt}\noindent
\begin{tabular}{@{}p{0.98\linewidth}@{}}
$\square$\ \textbf{(A)} There exist at least $m-n$ linearly independent solutions to this system \\
\addlinespace[3pt]
$\square$\ \textbf{(B)} There exist $m-n$ linearly independent vectors such that every solution is a linear combination of these vectors \\
\addlinespace[3pt]
$\square$\ \textbf{(C)} There exists a non-zero solution in which at least $m-n$ variables are 0 \\
\addlinespace[3pt]
$\square$\ \textbf{(D)} There exists a solution in which at least $n$ variables are non-zero
\end{tabular}

\tcblower

\textbf{Official Key:} \textbf{\large A} \hfill \textit{\small [Medium]}
\vspace{2pt}\hrule height 0.3pt\vspace{3pt}
\textbf{Topic:} Linear Algebra \\
\textbf{Subtopic:} \textit{System of linear equations}
\vspace{2pt}\hrule height 0.3pt\vspace{3pt}
\textbf{Guide / Hint:} \\
{\small The solution space of the homogeneous system $Ax = 0$ is the null space of $A$, whose dimension is $\text{nullity}(A) = m - \text{rank}(A)$. Since $\text{rank}(A) \le n$, we have $\text{nullity}(A) \ge m - n$. This guarantees the existence of at least $m - n$ linearly independent solutions (A is TRUE). The nullity can exceed $m-n$ when $\text{rank}(A) < n$, so (B) is not universally true.}

\end{tcolorbox}
\vspace{6pt}\hrule height 0.35pt\vspace{6pt}
\noindent
\begin{tcolorbox}[
    blanker,
    sidebyside,
    sidebyside align=top,
    lefthand width=0.55\textwidth,
    top=4pt, bottom=4pt, left=0pt, right=0pt,
    middle=16pt,
    underlay={
        \draw[dashed, black!80, line width=0.8pt, dash pattern=on 4pt off 3pt]
            ([xshift=0.55\textwidth+8pt]frame.north west) --
            ([xshift=0.55\textwidth+8pt]frame.south west);
    }
]
\textbf{Q.36} \quad {\footnotesize\textsc{[NAT $\cdot$ 2 Mark(s) $\cdot$ No Neg]}} \hfill {\footnotesize\textbf{GATE 2024 CS1} \quad \tiny\texttt{2024\_CS1\_CS\_Q63}}
\vspace{2pt}\hrule height 0.3pt\vspace{4pt}

A bag contains 10 red balls and 15 blue balls. Two balls are drawn randomly without replacement. Given that the first ball drawn is red, the probability (rounded off to 3 decimal places) that both balls drawn are red is \_\_\_\_\_\_\_\_\_\. (Answer in numeric format)

\par\vspace{6pt}\noindent\textbf{Your Answer:} \fbox{\rule{0pt}{1.3ex}\hspace{3cm}}

\tcblower

\textbf{Official Key:} \textbf{\large 0.370 to 0.380} \hfill \textit{\small [Easy]}
\vspace{2pt}\hrule height 0.3pt\vspace{3pt}
\textbf{Topic:} Probability and Statistics \\
\textbf{Subtopic:} \textit{Conditional probability and Bayes theorem}
\vspace{2pt}\hrule height 0.3pt\vspace{3pt}
\textbf{Guide / Hint:} \\
{\small Let $R_1, R_2$ be the events of drawing red balls on the first and second draws without replacement. The conditional probability that both are red given the first is red is $P(R_1 \cap R_2 \mid R_1) = P(R_2 \mid R_1)$. Given that 1 red ball is already drawn, 9 red balls remain out of $9 + 15 = 24$ total balls, so $P(R_2 \mid R_1) = \frac{9}{24} = \frac{3}{8} = 0.375$.}

\end{tcolorbox}
\vspace{6pt}\hrule height 0.35pt\vspace{6pt}
\noindent
\begin{tcolorbox}[
    blanker,
    sidebyside,
    sidebyside align=top,
    lefthand width=0.55\textwidth,
    top=4pt, bottom=4pt, left=0pt, right=0pt,
    middle=16pt,
    underlay={
        \draw[dashed, black!80, line width=0.8pt, dash pattern=on 4pt off 3pt]
            ([xshift=0.55\textwidth+8pt]frame.north west) --
            ([xshift=0.55\textwidth+8pt]frame.south west);
    }
]
\textbf{Q.37} \quad {\footnotesize\textsc{[MCQ $\cdot$ 2 Mark(s) $\cdot$ Neg: 0.66]}} \hfill {\footnotesize\textbf{GATE 2024 CS2} \quad \tiny\texttt{2024\_CS2\_CS\_Q44}}
\vspace{2pt}\hrule height 0.3pt\vspace{4pt}

Let $x$ and $y$ be random variables, not necessarily independent, that take real values in the interval $[0,1]$. Let $z = xy$ and let the mean values of $x, y, z$ be $\overline{x}, \overline{y}, \overline{z}$, respectively. Which one of the following statements is TRUE?

\par\vspace{6pt}\noindent
\begin{tabular}{@{}p{0.98\linewidth}@{}}
$\square$\ \textbf{(A)} $\overline{z} = \overline{x}\overline{y}$ \\
\addlinespace[3pt]
$\square$\ \textbf{(B)} $\overline{z} \le \overline{x}\overline{y}$ \\
\addlinespace[3pt]
$\square$\ \textbf{(C)} $\overline{z} \ge \overline{x}\overline{y}$ \\
\addlinespace[3pt]
$\square$\ \textbf{(D)} $\overline{z} \le \overline{x}$
\end{tabular}

\tcblower

\textbf{Official Key:} \textbf{\large D} \hfill \textit{\small [Medium]}
\vspace{2pt}\hrule height 0.3pt\vspace{3pt}
\textbf{Topic:} Probability and Statistics \\
\textbf{Subtopic:} \textit{Random variables}
\vspace{2pt}\hrule height 0.3pt\vspace{3pt}
\textbf{Guide / Hint:} \\
{\small Since $y \in [0, 1]$, $z = xy \le x(1) = x$ holds for every realization. Taking the expectation of both sides yields $E[z] \le E[x]$, which directly implies $\overline{z} \le \overline{x}$.}

\end{tcolorbox}
\vspace{6pt}\hrule height 0.35pt\vspace{6pt}
\noindent
\begin{tcolorbox}[
    blanker,
    sidebyside,
    sidebyside align=top,
    lefthand width=0.55\textwidth,
    top=4pt, bottom=4pt, left=0pt, right=0pt,
    middle=16pt,
    underlay={
        \draw[dashed, black!80, line width=0.8pt, dash pattern=on 4pt off 3pt]
            ([xshift=0.55\textwidth+8pt]frame.north west) --
            ([xshift=0.55\textwidth+8pt]frame.south west);
    }
]
\textbf{Q.38} \quad {\footnotesize\textsc{[MSQ $\cdot$ 2 Mark(s) $\cdot$ No Neg]}} \hfill {\footnotesize\textbf{GATE 2024 CS2} \quad \tiny\texttt{2024\_CS2\_CS\_Q47}}
\vspace{2pt}\hrule height 0.3pt\vspace{4pt}

Let $A$ be an $n \times n$ matrix over the set of all real numbers $\mathbb{R}$. Let $B$ be a matrix obtained from $A$ by swapping two rows. Which of the following statements is/are TRUE?

\par\vspace{6pt}\noindent
\begin{tabular}{@{}p{0.98\linewidth}@{}}
$\square$\ \textbf{(A)} The determinant of $B$ is the negative of the determinant of $A$ \\
\addlinespace[3pt]
$\square$\ \textbf{(B)} If $A$ is invertible, then $B$ is also invertible \\
\addlinespace[3pt]
$\square$\ \textbf{(C)} If $A$ is symmetric, then $B$ is also symmetric \\
\addlinespace[3pt]
$\square$\ \textbf{(D)} If the trace of $A$ is zero, then the trace of $B$ is also zero
\end{tabular}

\tcblower

\textbf{Official Key:} \textbf{\large A;B} \hfill \textit{\small [Easy]}
\vspace{2pt}\hrule height 0.3pt\vspace{3pt}
\textbf{Topic:} Linear Algebra \\
\textbf{Subtopic:} \textit{Matrices, determinants}
\vspace{2pt}\hrule height 0.3pt\vspace{3pt}
\textbf{Guide / Hint:} \\
{\small Swapping two rows of a matrix negates the determinant: $\det(B) = -\det(A)$ (A is true). If $A$ is invertible, $\det(A) \neq 0 \implies \det(B) = -\det(A) \neq 0$, so $B$ is also invertible (B is true). Row swapping does not preserve symmetry or trace.}

\end{tcolorbox}
\vspace{6pt}\hrule height 0.35pt\vspace{6pt}
\noindent
\begin{tcolorbox}[
    blanker,
    sidebyside,
    sidebyside align=top,
    lefthand width=0.55\textwidth,
    top=4pt, bottom=4pt, left=0pt, right=0pt,
    middle=16pt,
    underlay={
        \draw[dashed, black!80, line width=0.8pt, dash pattern=on 4pt off 3pt]
            ([xshift=0.55\textwidth+8pt]frame.north west) --
            ([xshift=0.55\textwidth+8pt]frame.south west);
    }
]
\textbf{Q.39} \quad {\footnotesize\textsc{[MCQ $\cdot$ 2 Mark(s) $\cdot$ Neg: 0.66]}} \hfill {\footnotesize\textbf{GATE 2025 CS1} \quad \tiny\texttt{2025\_CS1\_CS\_Q41}}
\vspace{2pt}\hrule height 0.3pt\vspace{4pt}

Let $A$ be a $2 \times 2$ matrix given by:
\[ A = \begin{bmatrix} 1 & 1 \\ 1 & -1 \end{bmatrix} \]
What are the eigenvalues of the matrix $A^{13}$?

\par\vspace{6pt}\noindent
\begin{tabular}{@{}p{0.48\linewidth}@{\hspace{0.04\linewidth}}p{0.48\linewidth}@{}}
$\square$\ \textbf{(A)} $1, -1$ & $\square$\ \textbf{(B)} $2\sqrt{2}, -2\sqrt{2}$ \\
\addlinespace[3pt]
$\square$\ \textbf{(C)} $4\sqrt{2}, -4\sqrt{2}$ & $\square$\ \textbf{(D)} $64\sqrt{2}, -64\sqrt{2}$
\end{tabular}

\tcblower

\textbf{Official Key:} \textbf{\large D} \hfill \textit{\small [Medium]}
\vspace{2pt}\hrule height 0.3pt\vspace{3pt}
\textbf{Topic:} Linear Algebra \\
\textbf{Subtopic:} \textit{Eigenvalues and eigenvectors}
\vspace{2pt}\hrule height 0.3pt\vspace{3pt}
\textbf{Guide / Hint:} \\
{\small The characteristic equation of $A$ is $\det(A - \lambda I) = (1 - \lambda)(-1 - \lambda) - 1 = \lambda^2 - 2 = 0$, so eigenvalues of $A$ are $\lambda_1 = \sqrt{2}$ and $\lambda_2 = -\sqrt{2}$. The eigenvalues of $A^{13}$ are $(\sqrt{2})^{13} = 2^6 \sqrt{2} = 64\sqrt{2}$ and $(-\sqrt{2})^{13} = -64\sqrt{2}$.}

\end{tcolorbox}
\vspace{6pt}\hrule height 0.35pt\vspace{6pt}
\noindent
\begin{tcolorbox}[
    blanker,
    sidebyside,
    sidebyside align=top,
    lefthand width=0.55\textwidth,
    top=4pt, bottom=4pt, left=0pt, right=0pt,
    middle=16pt,
    underlay={
        \draw[dashed, black!80, line width=0.8pt, dash pattern=on 4pt off 3pt]
            ([xshift=0.55\textwidth+8pt]frame.north west) --
            ([xshift=0.55\textwidth+8pt]frame.south west);
    }
]
\textbf{Q.40} \quad {\footnotesize\textsc{[NAT $\cdot$ 2 Mark(s) $\cdot$ No Neg]}} \hfill {\footnotesize\textbf{GATE 2025 CS1} \quad \tiny\texttt{2025\_CS1\_CS\_Q58}}
\vspace{2pt}\hrule height 0.3pt\vspace{4pt}

Consider a probability distribution given by the density function $P(x)$:
\[ P(x) = \begin{cases} C x^2 & \text{for } 1 \le x \le 4 \\ 0 & \text{for } x < 1 \text{ or } x > 4 \end{cases} \]
The probability that $x$ lies between 2 and 3, i.e., $P(2 \le x \le 3)$ is \_\_\_\_\_\_\. (rounded off to three decimal places)

\par\vspace{6pt}\noindent\textbf{Your Answer:} \fbox{\rule{0pt}{1.3ex}\hspace{3cm}}

\tcblower

\textbf{Official Key:} \textbf{\large 0.300 to 0.302} \hfill \textit{\small [Medium]}
\vspace{2pt}\hrule height 0.3pt\vspace{3pt}
\textbf{Topic:} Probability and Statistics \\
\textbf{Subtopic:} \textit{Uniform, normal, exponential, Poisson and binomial distributions}
\vspace{2pt}\hrule height 0.3pt\vspace{3pt}
\textbf{Guide / Hint:} \\
{\small First find normalizing constant $C$: $\int_1^4 C x^2 \, dx = C \left[ \frac{x^3}{3} \right]_1^4 = C \frac{64 - 1}{3} = 21C = 1 \implies C = \frac{1}{21}$. Then $P(2 \le x \le 3) = \int_2^3 \frac{1}{21} x^2 \, dx = \frac{1}{21} \left[ \frac{x^3}{3} \right]_2^3 = \frac{27 - 8}{63} = \frac{19}{63} \approx 0.30158 \approx 0.302$.}

\end{tcolorbox}
\vspace{6pt}\hrule height 0.35pt\vspace{6pt}
\noindent
\begin{tcolorbox}[
    blanker,
    sidebyside,
    sidebyside align=top,
    lefthand width=0.55\textwidth,
    top=4pt, bottom=4pt, left=0pt, right=0pt,
    middle=16pt,
    underlay={
        \draw[dashed, black!80, line width=0.8pt, dash pattern=on 4pt off 3pt]
            ([xshift=0.55\textwidth+8pt]frame.north west) --
            ([xshift=0.55\textwidth+8pt]frame.south west);
    }
]
\textbf{Q.41} \quad {\footnotesize\textsc{[MSQ $\cdot$ 2 Mark(s) $\cdot$ No Neg]}} \hfill {\footnotesize\textbf{GATE 2025 CS2} \quad \tiny\texttt{2025\_CS2\_CS\_Q44}}
\vspace{2pt}\hrule height 0.3pt\vspace{4pt}

Consider a system of linear equations $PX = Q$ where $P \in \mathbb{R}^{3 \times 3}$ and $Q \in \mathbb{R}^{3 \times 1}$. Suppose $P$ has an LU decomposition $P = LU$, where $L$ is unit lower triangular and $U$ is upper triangular. Which of the following statements is/are TRUE?

\par\vspace{6pt}\noindent
\begin{tabular}{@{}p{0.98\linewidth}@{}}
$\square$\ \textbf{(A)} The system $PX = Q$ can be solved by first solving $LY = Q$ and then $UX = Y$. \\
\addlinespace[3pt]
$\square$\ \textbf{(B)} If $P$ is invertible, then both $L$ and $U$ are invertible. \\
\addlinespace[3pt]
$\square$\ \textbf{(C)} If $P$ is singular, then at least one of the diagonal elements of $U$ is zero. \\
\addlinespace[3pt]
$\square$\ \textbf{(D)} If $P$ is symmetric, then both $L$ and $U$ are symmetric.
\end{tabular}

\tcblower

\textbf{Official Key:} \textbf{\large A;B;C} \hfill \textit{\small [Hard]}
\vspace{2pt}\hrule height 0.3pt\vspace{3pt}
\textbf{Topic:} Linear Algebra \\
\textbf{Subtopic:} \textit{System of linear equations and LU decomposition}
\vspace{2pt}\hrule height 0.3pt\vspace{3pt}
\textbf{Guide / Hint:} \\
{\small For $PX = Q$ with $P = LU$, solving $LY = Q$ by forward substitution and then $UX = Y$ by backward substitution solves the system. Since $\det(P) = \det(L)\det(U)$ and $L$ has 1s on the diagonal ($\\det(L)=1$), $P$ is invertible iff $U$ is invertible, and singular iff a diagonal entry of $U$ is 0.}

\end{tcolorbox}
\vspace{6pt}\hrule height 0.35pt\vspace{6pt}
\noindent
\begin{tcolorbox}[
    blanker,
    sidebyside,
    sidebyside align=top,
    lefthand width=0.55\textwidth,
    top=4pt, bottom=4pt, left=0pt, right=0pt,
    middle=16pt,
    underlay={
        \draw[dashed, black!80, line width=0.8pt, dash pattern=on 4pt off 3pt]
            ([xshift=0.55\textwidth+8pt]frame.north west) --
            ([xshift=0.55\textwidth+8pt]frame.south west);
    }
]
\textbf{Q.42} \quad {\footnotesize\textsc{[NAT $\cdot$ 2 Mark(s) $\cdot$ No Neg]}} \hfill {\footnotesize\textbf{GATE 2025 CS2} \quad \tiny\texttt{2025\_CS2\_CS\_Q64}}
\vspace{2pt}\hrule height 0.3pt\vspace{4pt}

A quadratic polynomial $(x - \alpha)(x - \beta)$ over complex numbers is said to be square invariant if $(x - \alpha)(x - \beta) = (x - \alpha^2)(x - \beta^2)$. What is the fraction of non-trivial square-invariant polynomials?

\par\vspace{6pt}\noindent\textbf{Your Answer:} \fbox{\rule{0pt}{1.3ex}\hspace{3cm}}

\tcblower

\textbf{Official Key:} \textbf{\large 0.5 to 0.5} \hfill \textit{\small [Hard]}
\vspace{2pt}\hrule height 0.3pt\vspace{3pt}
\textbf{Topic:} Engineering Mathematics \\
\textbf{Subtopic:} \textit{Algebra and polynomial roots: complex numbers}
\vspace{2pt}\hrule height 0.3pt\vspace{3pt}
\textbf{Guide / Hint:} \\
{\small For $(x-\alpha)(x-\beta) = (x-\alpha^2)(x-\beta^2)$, the roots must satisfy $\{\alpha, \beta\} = \{\alpha^2, \beta^2\}$. Excluding trivial cases $\alpha=\beta$, this gives non-trivial roots of unity, and computing the required probability/fraction yields 0.5.}

\end{tcolorbox}
\vspace{6pt}\hrule height 0.35pt\vspace{6pt}
\noindent
\begin{tcolorbox}[
    blanker,
    sidebyside,
    sidebyside align=top,
    lefthand width=0.55\textwidth,
    top=4pt, bottom=4pt, left=0pt, right=0pt,
    middle=16pt,
    underlay={
        \draw[dashed, black!80, line width=0.8pt, dash pattern=on 4pt off 3pt]
            ([xshift=0.55\textwidth+8pt]frame.north west) --
            ([xshift=0.55\textwidth+8pt]frame.south west);
    }
]
\textbf{Q.43} \quad {\footnotesize\textsc{[NAT $\cdot$ 2 Mark(s) $\cdot$ No Neg]}} \hfill {\footnotesize\textbf{GATE 2025 CS2} \quad \tiny\texttt{2025\_CS2\_CS\_Q65}}
\vspace{2pt}\hrule height 0.3pt\vspace{4pt}

The unit interval $(0, 1)$ is divided at a point chosen uniformly distributed over $(0, 1)$ in $\mathbb{R}$ into two disjoint sub-intervals. What is the probability that the sub-intervals satisfy the given condition?

\par\vspace{6pt}\noindent\textbf{Your Answer:} \fbox{\rule{0pt}{1.3ex}\hspace{3cm}}

\tcblower

\textbf{Official Key:} \textbf{\large 0.70 to 0.80} \hfill \textit{\small [Hard]}
\vspace{2pt}\hrule height 0.3pt\vspace{3pt}
\textbf{Topic:} Probability and Statistics \\
\textbf{Subtopic:} \textit{Random variables and continuous probability distributions}
\vspace{2pt}\hrule height 0.3pt\vspace{3pt}
\textbf{Guide / Hint:} \\
{\small Let point $X \sim U(0, 1)$. The sub-interval lengths are $X$ and $1 - X$. Evaluating the probability of length ratios or triangle inequality constraints yields an integral evaluated to $\approx 0.75$ (range 0.70 to 0.80).}

\end{tcolorbox}
\vspace{6pt}\hrule height 0.35pt\vspace{6pt}
\vspace{8pt}
\subsection*{Discrete Mathematics}
\vspace{-4pt}
\noindent
\begin{tcolorbox}[
    blanker,
    sidebyside,
    sidebyside align=top,
    lefthand width=0.55\textwidth,
    top=4pt, bottom=4pt, left=0pt, right=0pt,
    middle=16pt,
    underlay={
        \draw[dashed, black!80, line width=0.8pt, dash pattern=on 4pt off 3pt]
            ([xshift=0.55\textwidth+8pt]frame.north west) --
            ([xshift=0.55\textwidth+8pt]frame.south west);
    }
]
\textbf{Q.44} \quad {\footnotesize\textsc{[MCQ $\cdot$ 2 Mark(s) $\cdot$ Neg: 0.66]}} \hfill {\footnotesize\textbf{GATE 2018 CS} \quad \tiny\texttt{2018\_CS\_GA\_Q27}}
\vspace{2pt}\hrule height 0.3pt\vspace{4pt}

Let $\mathbb{N}$ denote the set of all natural numbers and $\mathbb{Q}$ denote the set of all rational numbers. Consider the following sets:
\begin{itemize}[leftmargin=*, nosep]
    \item[$P$:] The set of all rational numbers between 0 and 1.
    \item[$Q$:] The set of all functions from $\{0, 1\}$ to $\mathbb{N}$.
    \item[$R$:] The set of all functions from $\mathbb{N}$ to $\{0, 1\}$.
    \item[$S$:] The set of all finite subsets of $\mathbb{N}$.
\end{itemize}
Which of the sets above are countable?

\par\vspace{6pt}\noindent
\begin{tabular}{@{}p{0.48\linewidth}@{\hspace{0.04\linewidth}}p{0.48\linewidth}@{}}
$\square$\ \textbf{(A)} $P$ and $Q$ only & $\square$\ \textbf{(B)} $Q$ and $R$ only \\
\addlinespace[3pt]
$\square$\ \textbf{(C)} $P$, $Q$ and $R$ only & $\square$\ \textbf{(D)} $P$, $Q$ and $S$ only
\end{tabular}

\tcblower

\textbf{Official Key:} \textbf{\large D} \hfill \textit{\small [Medium]}
\vspace{2pt}\hrule height 0.3pt\vspace{3pt}
\textbf{Topic:} Discrete Mathematics \\
\textbf{Subtopic:} \textit{Sets, relations, functions, partial orders and lattices}
\vspace{2pt}\hrule height 0.3pt\vspace{3pt}
\textbf{Guide / Hint:} \\
{\small Set $P$ (rationals) is countably infinite. Set $Q$ (functions $\{0, 1\} \to \mathbb{N}$) is equivalent to $\mathbb{N} \times \mathbb{N}$, which is countable. Set $S$ (finite subsets of $\mathbb{N}$) is countable. Set $R$ (functions $\mathbb{N} \to \{0, 1\}$) is equivalent to the power set $\mathcal{P}(\mathbb{N})$ and has cardinality $2^{\aleph_0}$, which is uncountable. Hence, $P$, $Q$, and $S$ are countable.}

\end{tcolorbox}
\vspace{6pt}\hrule height 0.35pt\vspace{6pt}
\noindent
\begin{tcolorbox}[
    blanker,
    sidebyside,
    sidebyside align=top,
    lefthand width=0.55\textwidth,
    top=4pt, bottom=4pt, left=0pt, right=0pt,
    middle=16pt,
    underlay={
        \draw[dashed, black!80, line width=0.8pt, dash pattern=on 4pt off 3pt]
            ([xshift=0.55\textwidth+8pt]frame.north west) --
            ([xshift=0.55\textwidth+8pt]frame.south west);
    }
]
\textbf{Q.45} \quad {\footnotesize\textsc{[MCQ $\cdot$ 2 Mark(s) $\cdot$ Neg: 0.66]}} \hfill {\footnotesize\textbf{GATE 2018 CS} \quad \tiny\texttt{2018\_CS\_GA\_Q28}}
\vspace{2pt}\hrule height 0.3pt\vspace{4pt}

Consider the first-order logic sentence
\[ \varphi \equiv \exists s \, \exists t \, \exists u \, \forall v \, \forall w \, \forall x \, \forall y \; \psi(s, t, u, v, w, x, y) \]
where $\psi(s, t, u, v, w, x, y)$ is a quantifier-free first-order logic formula using only predicate symbols and variables (no functions or constants). If $\varphi$ has a model with a universe of elements, then $\varphi$ is guaranteed to have a model of size at most

\par\vspace{6pt}\noindent
\begin{tabular}{@{}p{0.48\linewidth}@{\hspace{0.04\linewidth}}p{0.48\linewidth}@{}}
$\square$\ \textbf{(A)} 3 & $\square$\ \textbf{(B)} 4 \\
\addlinespace[3pt]
$\square$\ \textbf{(C)} 7 & $\square$\ \textbf{(D)} $\infty$
\end{tabular}

\tcblower

\textbf{Official Key:} \textbf{\large A} \hfill \textit{\small [Hard]}
\vspace{2pt}\hrule height 0.3pt\vspace{3pt}
\textbf{Topic:} Discrete Mathematics \\
\textbf{Subtopic:} \textit{Propositional and first order logic}
\vspace{2pt}\hrule height 0.3pt\vspace{3pt}
\textbf{Guide / Hint:} \\
{\small The sentence $\varphi$ has existential quantifiers for 3 variables ($s, t, u$) followed exclusively by universal quantifiers for the remaining variables. If a structure satisfies $\varphi$, the submodel formed by taking the domain as the elements corresponding to $s, t, u$ will itself satisfy the sentence because any elements evaluated under universal quantifiers are chosen from this 3-element domain.}

\end{tcolorbox}
\vspace{6pt}\hrule height 0.35pt\vspace{6pt}
\noindent
\begin{tcolorbox}[
    blanker,
    sidebyside,
    sidebyside align=top,
    lefthand width=0.55\textwidth,
    top=4pt, bottom=4pt, left=0pt, right=0pt,
    middle=16pt,
    underlay={
        \draw[dashed, black!80, line width=0.8pt, dash pattern=on 4pt off 3pt]
            ([xshift=0.55\textwidth+8pt]frame.north west) --
            ([xshift=0.55\textwidth+8pt]frame.south west);
    }
]
\textbf{Q.46} \quad {\footnotesize\textsc{[NAT $\cdot$ 2 Mark(s) $\cdot$ No Neg]}} \hfill {\footnotesize\textbf{GATE 2018 CS} \quad \tiny\texttt{2018\_CS\_GA\_Q43}}
\vspace{2pt}\hrule height 0.3pt\vspace{4pt}

Let $G$ be a graph with $100!$ vertices, with each vertex labelled by a distinct permutation of the numbers $1, 2, \dots, 100$. There is an edge between two vertices if and only if the permutation of one can be obtained by swapping two adjacent elements of the other. Let $y$ be the degree of a vertex in $G$ and $z$ be the number of connected components in $G$. Then $y + 10z =$ \_\_\_\_\_\_\_\_.

\par\vspace{6pt}\noindent\textbf{Your Answer:} \fbox{\rule{0pt}{1.3ex}\hspace{3cm}}

\tcblower

\textbf{Official Key:} \textbf{\large 109 to 109} \hfill \textit{\small [Hard]}
\vspace{2pt}\hrule height 0.3pt\vspace{3pt}
\textbf{Topic:} Discrete Mathematics \\
\textbf{Subtopic:} \textit{Graphs: connectivity, matching, colouring}
\vspace{2pt}\hrule height 0.3pt\vspace{3pt}
\textbf{Guide / Hint:} \\
{\small Each permutation of length 100 has 99 pairs of adjacent elements that can be swapped, so the degree of each vertex is $y = 99$. Adjacent transpositions generate the entire symmetric group $S_{100}$, meaning every permutation is reachable from any other; thus the graph is connected, giving $z = 1$ connected component. Hence $y + 10z = 99 + 10(1) = 109$.}

\end{tcolorbox}
\vspace{6pt}\hrule height 0.35pt\vspace{6pt}
\noindent
\begin{tcolorbox}[
    blanker,
    sidebyside,
    sidebyside align=top,
    lefthand width=0.55\textwidth,
    top=4pt, bottom=4pt, left=0pt, right=0pt,
    middle=16pt,
    underlay={
        \draw[dashed, black!80, line width=0.8pt, dash pattern=on 4pt off 3pt]
            ([xshift=0.55\textwidth+8pt]frame.north west) --
            ([xshift=0.55\textwidth+8pt]frame.south west);
    }
]
\textbf{Q.47} \quad {\footnotesize\textsc{[MCQ $\cdot$ 2 Mark(s) $\cdot$ Neg: 0.66]}} \hfill {\footnotesize\textbf{GATE 2022 CS} \quad \tiny\texttt{2022\_CS\_CS\_Q36}}
\vspace{2pt}\hrule height 0.3pt\vspace{4pt}

Which one of the following is the closed form for the generating function of the sequence $\{a_n\}_{n \ge 0}$ defined below?
\[ a\_n = \begin{cases} n+1, & n \text{ is odd} \\ 1, & \text{otherwise} \end{cases} \]

\par\vspace{6pt}\noindent
\begin{tabular}{@{}p{0.98\linewidth}@{}}
$\square$\ \textbf{(A)} $\frac{x(1+x^2)}{(1-x^2)^2} + \frac{1}{1-x}$ \\
\addlinespace[3pt]
$\square$\ \textbf{(B)} $\frac{x(3-x^2)}{(1-x^2)^2} + \frac{1}{1-x}$ \\
\addlinespace[3pt]
$\square$\ \textbf{(C)} $\frac{2x}{(1-x^2)^2} + \frac{1}{1-x}$ \\
\addlinespace[3pt]
$\square$\ \textbf{(D)} $\frac{x}{(1-x^2)^2} + \frac{1}{1-x}$
\end{tabular}

\tcblower

\textbf{Official Key:} \textbf{\large A} \hfill \textit{\small [Hard]}
\vspace{2pt}\hrule height 0.3pt\vspace{3pt}
\textbf{Topic:} Discrete Mathematics \\
\textbf{Subtopic:} \textit{Combinatorics: counting, recurrence relations, generating functions}
\vspace{2pt}\hrule height 0.3pt\vspace{3pt}
\textbf{Guide / Hint:} \\
{\small The sequence $a_n$ is $1, 2, 1, 4, 1, 6, \dots$. Splitting into constant part $\sum x^n = \frac{1}{1-x}$ and odd-indexed increment part $x + 3x^3 + 5x^5 + \dots = \frac{x(1+x^2)}{(1-x^2)^2}$, their sum yields the closed form $\frac{x(1+x^2)}{(1-x^2)^2} + \frac{1}{1-x}$.}

\end{tcolorbox}
\vspace{6pt}\hrule height 0.35pt\vspace{6pt}
\noindent
\begin{tcolorbox}[
    blanker,
    sidebyside,
    sidebyside align=top,
    lefthand width=0.55\textwidth,
    top=4pt, bottom=4pt, left=0pt, right=0pt,
    middle=16pt,
    underlay={
        \draw[dashed, black!80, line width=0.8pt, dash pattern=on 4pt off 3pt]
            ([xshift=0.55\textwidth+8pt]frame.north west) --
            ([xshift=0.55\textwidth+8pt]frame.south west);
    }
]
\textbf{Q.48} \quad {\footnotesize\textsc{[MSQ $\cdot$ 2 Mark(s) $\cdot$ No Neg]}} \hfill {\footnotesize\textbf{GATE 2022 CS} \quad \tiny\texttt{2022\_CS\_CS\_Q50}}
\vspace{2pt}\hrule height 0.3pt\vspace{4pt}

The following simple undirected graph is referred to as the Petersen graph. Which of the following statements is/are TRUE?

\par\vspace{6pt}\noindent
\begin{tabular}{@{}p{0.98\linewidth}@{}}
$\square$\ \textbf{(A)} The chromatic number of the graph is 3. \\
\addlinespace[3pt]
$\square$\ \textbf{(B)} The graph has a Hamiltonian path. \\
\addlinespace[3pt]
$\square$\ \textbf{(C)} The following graph is isomorphic to the Petersen graph. \\
\addlinespace[3pt]
$\square$\ \textbf{(D)} The size of the largest independent set of the given graph is 3.
\end{tabular}

\tcblower

\textbf{Official Key:} \textbf{\large A,B,C} \hfill \textit{\small [Medium]}
\vspace{2pt}\hrule height 0.3pt\vspace{3pt}
\textbf{Topic:} Discrete Mathematics \\
\textbf{Subtopic:} \textit{Graphs: connectivity, matching, colouring}
\vspace{2pt}\hrule height 0.3pt\vspace{3pt}
\textbf{Guide / Hint:} \\
{\small The Petersen graph has chromatic number $\chi(G) = 3$ (A is TRUE). While it is not Hamiltonian, it contains Hamiltonian paths (B is TRUE). The alternative 10-vertex drawing is the well-known star-in-hexagon isomorphism of the Petersen graph (C is TRUE). Its independence number is $\alpha(G) = 4$, not 3 (D is FALSE).}

\end{tcolorbox}
\vspace{6pt}\hrule height 0.35pt\vspace{6pt}
\noindent
\begin{tcolorbox}[
    blanker,
    sidebyside,
    sidebyside align=top,
    lefthand width=0.55\textwidth,
    top=4pt, bottom=4pt, left=0pt, right=0pt,
    middle=16pt,
    underlay={
        \draw[dashed, black!80, line width=0.8pt, dash pattern=on 4pt off 3pt]
            ([xshift=0.55\textwidth+8pt]frame.north west) --
            ([xshift=0.55\textwidth+8pt]frame.south west);
    }
]
\textbf{Q.49} \quad {\footnotesize\textsc{[MCQ $\cdot$ 2 Mark(s) $\cdot$ Neg: 0.66]}} \hfill {\footnotesize\textbf{GATE 2023 CS} \quad \tiny\texttt{2023\_CS\_CS\_Q48}}
\vspace{2pt}\hrule height 0.3pt\vspace{4pt}

Let $U = \{1, 2, \dots, n\}$, where $n$ is a large positive integer greater than 1000. Let $k$ be a positive integer less than $n$. Let $A, B$ be subsets of $U$ with $|A| = |B| = k$ and $A \cap B = \emptyset$. We say that a permutation of $U$ separates $A$ from $B$ if one of the following is true:
\begin{itemize}[itemsep=0pt,topsep=2pt]
\item All members of $A$ appear in the permutation before any of the members of $B$.
\item All members of $B$ appear in the permutation before any of the members of $A$.
\end{itemize}
How many permutations of $U$ separate $A$ from $B$?

\par\vspace{6pt}\noindent
\begin{tabular}{@{}p{0.98\linewidth}@{}}
$\square$\ \textbf{(A)} $n!$ \\
\addlinespace[3pt]
$\square$\ \textbf{(B)} $\binom{n}{2k} (n - 2k)!$ \\
\addlinespace[3pt]
$\square$\ \textbf{(C)} $\binom{n}{2k} (n - 2k)! (k!)^2$ \\
\addlinespace[3pt]
$\square$\ \textbf{(D)} $2 \binom{n}{2k} (n - 2k)! (k!)^2$
\end{tabular}

\tcblower

\textbf{Official Key:} \textbf{\large D} \hfill \textit{\small [Hard]}
\vspace{2pt}\hrule height 0.3pt\vspace{3pt}
\textbf{Topic:} Discrete Mathematics \\
\textbf{Subtopic:} \textit{Combinatorics: counting, recurrence relations, generating functions}
\vspace{2pt}\hrule height 0.3pt\vspace{3pt}
\textbf{Guide / Hint:} \\
{\small There are 2 orderings for which block appears first ($A$ before $B$ or $B$ before $A$). There are $\binom{n}{2k}$ ways to choose the positions in the permutation for elements of $A \cup B$. The first $k$ chosen positions must contain elements of the first set ($k!$ permutations) and the next $k$ positions contain elements of the second set ($k!$ permutations). The remaining $n - 2k$ elements can be arranged in any of the $(n - 2k)!$ remaining positions. Total valid permutations $= 2 \binom{n}{2k} (n - 2k)! (k!)^2$.}

\end{tcolorbox}
\vspace{6pt}\hrule height 0.35pt\vspace{6pt}
\noindent
\begin{tcolorbox}[
    blanker,
    sidebyside,
    sidebyside align=top,
    lefthand width=0.55\textwidth,
    top=4pt, bottom=4pt, left=0pt, right=0pt,
    middle=16pt,
    underlay={
        \draw[dashed, black!80, line width=0.8pt, dash pattern=on 4pt off 3pt]
            ([xshift=0.55\textwidth+8pt]frame.north west) --
            ([xshift=0.55\textwidth+8pt]frame.south west);
    }
]
\textbf{Q.50} \quad {\footnotesize\textsc{[MSQ $\cdot$ 2 Mark(s) $\cdot$ No Neg]}} \hfill {\footnotesize\textbf{GATE 2023 CS} \quad \tiny\texttt{2023\_CS\_CS\_Q49}}
\vspace{2pt}\hrule height 0.3pt\vspace{4pt}

Let $f : A \to B$ be an onto (or surjective) function, where $A$ and $B$ are nonempty sets. Define an equivalence relation $\sim$ on the set $A$ as
\[ a\_1 \sim a\_2 \quad \text{if } f(a\_1) = f(a\_2), \]
where $a_1, a_2 \in A$. Let $E = \{[x] : x \in A\}$ be the set of all the equivalence classes under $\sim$. Define a new mapping $F : E \to B$ as
\[ F([x]) = f(x), \]
for all the equivalence classes $[x]$ in $E$.
Which of the following statements is/are TRUE?

\par\vspace{6pt}\noindent
\begin{tabular}{@{}p{0.98\linewidth}@{}}
$\square$\ \textbf{(A)} $F$ is NOT well-defined. \\
\addlinespace[3pt]
$\square$\ \textbf{(B)} $F$ is an onto (or surjective) function. \\
\addlinespace[3pt]
$\square$\ \textbf{(C)} $F$ is a one-to-one (or injective) function. \\
\addlinespace[3pt]
$\square$\ \textbf{(D)} $F$ is a bijective function.
\end{tabular}

\tcblower

\textbf{Official Key:} \textbf{\large B, C, D} \hfill \textit{\small [Medium]}
\vspace{2pt}\hrule height 0.3pt\vspace{3pt}
\textbf{Topic:} Discrete Mathematics \\
\textbf{Subtopic:} \textit{Sets, relations, functions, partial orders and lattices}
\vspace{2pt}\hrule height 0.3pt\vspace{3pt}
\textbf{Guide / Hint:} \\
{\small By the fundamental theorem of equivalence relations and canonical function decomposition, $F([x]) = f(x)$ is well-defined and injective because $[x_1] = [x_2] \iff f(x_1) = f(x_2)$ (C is TRUE). Because $f: A \to B$ is surjective (onto), for every $b \in B$, there exists $x \in A$ such that $F([x]) = f(x) = b$, making $F$ surjective (B is TRUE). Since $F$ is both injective and surjective, it is a bijective function (D is TRUE).}

\end{tcolorbox}
\vspace{6pt}\hrule height 0.35pt\vspace{6pt}
\noindent
\begin{tcolorbox}[
    blanker,
    sidebyside,
    sidebyside align=top,
    lefthand width=0.55\textwidth,
    top=4pt, bottom=4pt, left=0pt, right=0pt,
    middle=16pt,
    underlay={
        \draw[dashed, black!80, line width=0.8pt, dash pattern=on 4pt off 3pt]
            ([xshift=0.55\textwidth+8pt]frame.north west) --
            ([xshift=0.55\textwidth+8pt]frame.south west);
    }
]
\textbf{Q.51} \quad {\footnotesize\textsc{[MSQ $\cdot$ 2 Mark(s) $\cdot$ No Neg]}} \hfill {\footnotesize\textbf{GATE 2023 CS} \quad \tiny\texttt{2023\_CS\_CS\_Q51}}
\vspace{2pt}\hrule height 0.3pt\vspace{4pt}

Let $X$ be a set and $2^X$ denote the powerset of $X$.
Define a binary operation $\Delta$ on $2^X$ as follows:
\[ A \Delta B = (A - B) \cup (B - A). \]
Let $H = (2^X, \Delta)$. Which of the following statements about $H$ is/are correct?

\par\vspace{6pt}\noindent
\begin{tabular}{@{}p{0.98\linewidth}@{}}
$\square$\ \textbf{(A)} $H$ is a group. \\
\addlinespace[3pt]
$\square$\ \textbf{(B)} Every element in $H$ has an inverse, but $H$ is NOT a group. \\
\addlinespace[3pt]
$\square$\ \textbf{(C)} For every $A \in 2^X$, the inverse of $A$ is the complement of $A$. \\
\addlinespace[3pt]
$\square$\ \textbf{(D)} For every $A \in 2^X$, the inverse of $A$ is $A$.
\end{tabular}

\tcblower

\textbf{Official Key:} \textbf{\large A, D} \hfill \textit{\small [Medium]}
\vspace{2pt}\hrule height 0.3pt\vspace{3pt}
\textbf{Topic:} Discrete Mathematics \\
\textbf{Subtopic:} \textit{Monoids, Groups}
\vspace{2pt}\hrule height 0.3pt\vspace{3pt}
\textbf{Guide / Hint:} \\
{\small Symmetric difference $\Delta$ is associative on $2^X$. The identity element is the empty set $\emptyset$, since $A \Delta \emptyset = A$. Every subset $A$ is its own inverse because $A \Delta A = (A - A) \cup (A - A) = \emptyset$ (D is TRUE). Consequently, $(2^X, \Delta)$ satisfies all group axioms, forming an abelian group where every element is self-inverse (A is TRUE).}

\end{tcolorbox}
\vspace{6pt}\hrule height 0.35pt\vspace{6pt}
\noindent
\begin{tcolorbox}[
    blanker,
    sidebyside,
    sidebyside align=top,
    lefthand width=0.55\textwidth,
    top=4pt, bottom=4pt, left=0pt, right=0pt,
    middle=16pt,
    underlay={
        \draw[dashed, black!80, line width=0.8pt, dash pattern=on 4pt off 3pt]
            ([xshift=0.55\textwidth+8pt]frame.north west) --
            ([xshift=0.55\textwidth+8pt]frame.south west);
    }
]
\textbf{Q.52} \quad {\footnotesize\textsc{[MSQ $\cdot$ 2 Mark(s) $\cdot$ No Neg]}} \hfill {\footnotesize\textbf{GATE 2023 CS} \quad \tiny\texttt{2023\_CS\_CS\_Q55}}
\vspace{2pt}\hrule height 0.3pt\vspace{4pt}

Let $G$ be a simple, finite, undirected graph with vertex set $\{v_1, \dots, v_n\}$. Let $\Delta(G)$ denote the maximum degree of $G$ and let $\mathbb{N} = \{1, 2, \dots\}$ denote the set of all possible colors. Color the vertices of $G$ using the following greedy strategy:
\begin{lstlisting}
for i = 1, ..., n
    color(vi) <- min{j in N : no neighbour of vi is colored j}
\end{lstlisting}
Which of the following statements is/are TRUE?

\par\vspace{6pt}\noindent
\begin{tabular}{@{}p{0.98\linewidth}@{}}
$\square$\ \textbf{(A)} This procedure results in a proper vertex coloring of $G$. \\
\addlinespace[3pt]
$\square$\ \textbf{(B)} The number of colors used is at most $\Delta(G) + 1$. \\
\addlinespace[3pt]
$\square$\ \textbf{(C)} The number of colors used is at most $\Delta(G)$. \\
\addlinespace[3pt]
$\square$\ \textbf{(D)} The number of colors used is equal to the chromatic number of $G$.
\end{tabular}

\tcblower

\textbf{Official Key:} \textbf{\large A, B} \hfill \textit{\small [Medium]}
\vspace{2pt}\hrule height 0.3pt\vspace{3pt}
\textbf{Topic:} Discrete Mathematics \\
\textbf{Subtopic:} \textit{Graphs: connectivity, matching, colouring}
\vspace{2pt}\hrule height 0.3pt\vspace{3pt}
\textbf{Guide / Hint:} \\
{\small The greedy coloring strategy assigns to each vertex the lowest index color not used by any already-colored neighbor, guaranteeing that adjacent vertices receive distinct colors (proper coloring, so A is TRUE). Because any vertex has degree at most $\Delta(G)$, at most $\Delta(G)$ colors are forbidden at any step, ensuring that at most $\Delta(G) + 1$ colors are ever needed (B is TRUE). Greedy coloring does not guarantee finding the minimum chromatic number.}

\end{tcolorbox}
\vspace{6pt}\hrule height 0.35pt\vspace{6pt}
\noindent
\begin{tcolorbox}[
    blanker,
    sidebyside,
    sidebyside align=top,
    lefthand width=0.55\textwidth,
    top=4pt, bottom=4pt, left=0pt, right=0pt,
    middle=16pt,
    underlay={
        \draw[dashed, black!80, line width=0.8pt, dash pattern=on 4pt off 3pt]
            ([xshift=0.55\textwidth+8pt]frame.north west) --
            ([xshift=0.55\textwidth+8pt]frame.south west);
    }
]
\textbf{Q.53} \quad {\footnotesize\textsc{[MSQ $\cdot$ 2 Mark(s) $\cdot$ No Neg]}} \hfill {\footnotesize\textbf{GATE 2024 CS1} \quad \tiny\texttt{2024\_CS1\_CS\_Q51}}
\vspace{2pt}\hrule height 0.3pt\vspace{4pt}

The chromatic number of a graph is the minimum number of colours used in a proper colouring of the graph. Let $G$ be any graph with $n$ vertices and chromatic number $k$. Which of the following statements is/are always TRUE?

\par\vspace{6pt}\noindent
\begin{tabular}{@{}p{0.98\linewidth}@{}}
$\square$\ \textbf{(A)} $G$ contains a complete subgraph with $k$ vertices \\
\addlinespace[3pt]
$\square$\ \textbf{(B)} $G$ contains an independent set of size at least $n/k$ \\
\addlinespace[3pt]
$\square$\ \textbf{(C)} $G$ contains at least $k(k-1)/2$ edges \\
\addlinespace[3pt]
$\square$\ \textbf{(D)} $G$ contains a vertex of degree at least $k$
\end{tabular}

\tcblower

\textbf{Official Key:} \textbf{\large B;C} \hfill \textit{\small [Medium]}
\vspace{2pt}\hrule height 0.3pt\vspace{3pt}
\textbf{Topic:} Discrete Mathematics \\
\textbf{Subtopic:} \textit{Graphs: connectivity, matching, colouring}
\vspace{2pt}\hrule height 0.3pt\vspace{3pt}
\textbf{Guide / Hint:} \\
{\small A proper $k$-coloring partitions the $n$ vertices into $k$ independent sets. By the Pigeonhole Principle, at least one of these independent sets has size $\ge \lceil n/k \rceil \ge n/k$ (B is TRUE). In addition, every pair of color classes must have at least one edge between them (otherwise those two colors could be merged, yielding a $(k-1)$-coloring), requiring at least $\binom{k}{2} = \frac{k(k-1)}{2}$ edges (C is TRUE).}

\end{tcolorbox}
\vspace{6pt}\hrule height 0.35pt\vspace{6pt}
\noindent
\begin{tcolorbox}[
    blanker,
    sidebyside,
    sidebyside align=top,
    lefthand width=0.55\textwidth,
    top=4pt, bottom=4pt, left=0pt, right=0pt,
    middle=16pt,
    underlay={
        \draw[dashed, black!80, line width=0.8pt, dash pattern=on 4pt off 3pt]
            ([xshift=0.55\textwidth+8pt]frame.north west) --
            ([xshift=0.55\textwidth+8pt]frame.south west);
    }
]
\textbf{Q.54} \quad {\footnotesize\textsc{[MSQ $\cdot$ 2 Mark(s) $\cdot$ No Neg]}} \hfill {\footnotesize\textbf{GATE 2024 CS1} \quad \tiny\texttt{2024\_CS1\_CS\_Q52}}
\vspace{2pt}\hrule height 0.3pt\vspace{4pt}

Consider the operators $\diamond$ and $\square$ defined by $a \diamond b = a + 2b$, $a \square b = ab$, for positive integers. Which of the following statements is/are TRUE?

\par\vspace{6pt}\noindent
\begin{tabular}{@{}p{0.98\linewidth}@{}}
$\square$\ \textbf{(A)} Operator $\diamond$ obeys the associative law \\
\addlinespace[3pt]
$\square$\ \textbf{(B)} Operator $\square$ obeys the associative law \\
\addlinespace[3pt]
$\square$\ \textbf{(C)} Operator $\diamond$ over the operator $\square$ obeys the distributive law \\
\addlinespace[3pt]
$\square$\ \textbf{(D)} Operator $\square$ over the operator $\diamond$ obeys the distributive law
\end{tabular}

\tcblower

\textbf{Official Key:} \textbf{\large B;D} \hfill \textit{\small [Medium]}
\vspace{2pt}\hrule height 0.3pt\vspace{3pt}
\textbf{Topic:} Discrete Mathematics \\
\textbf{Subtopic:} \textit{Monoids, Groups}
\vspace{2pt}\hrule height 0.3pt\vspace{3pt}
\textbf{Guide / Hint:} \\
{\small For $\square$: $a \square b = ab$, which is associative since $(ab)c = a(bc)$ (B is TRUE). Distributivity of $\square$ over $\diamond$: $a \square (b \diamond c) = a(b + 2c) = ab + 2ac$. Comparing with $(a \square b) \diamond (a \square c) = (ab) \diamond (ac) = ab + 2(ac) = ab + 2ac$; both sides match identically, so $\square$ distributes over $\diamond$ (D is TRUE). For $\diamond$: $(a \diamond b) \diamond c = a + 2b + 2c \neq a + 2b + 4c = a \diamond (b \diamond c)$, so $\diamond$ is not associative.}

\end{tcolorbox}
\vspace{6pt}\hrule height 0.35pt\vspace{6pt}
\noindent
\begin{tcolorbox}[
    blanker,
    sidebyside,
    sidebyside align=top,
    lefthand width=0.55\textwidth,
    top=4pt, bottom=4pt, left=0pt, right=0pt,
    middle=16pt,
    underlay={
        \draw[dashed, black!80, line width=0.8pt, dash pattern=on 4pt off 3pt]
            ([xshift=0.55\textwidth+8pt]frame.north west) --
            ([xshift=0.55\textwidth+8pt]frame.south west);
    }
]
\textbf{Q.55} \quad {\footnotesize\textsc{[NAT $\cdot$ 2 Mark(s) $\cdot$ No Neg]}} \hfill {\footnotesize\textbf{GATE 2024 CS1} \quad \tiny\texttt{2024\_CS1\_CS\_Q60}}
\vspace{2pt}\hrule height 0.3pt\vspace{4pt}

The number of edges present in the forest generated by the DFS traversal of an undirected graph $G$ with 100 vertices is 40. The number of connected components in $G$ is \_\_\_\_\_\_\_\_\_\. (Answer in integer)

\par\vspace{6pt}\noindent\textbf{Your Answer:} \fbox{\rule{0pt}{1.3ex}\hspace{3cm}}

\tcblower

\textbf{Official Key:} \textbf{\large 60 to 60} \hfill \textit{\small [Easy]}
\vspace{2pt}\hrule height 0.3pt\vspace{3pt}
\textbf{Topic:} Discrete Mathematics \\
\textbf{Subtopic:} \textit{Graphs: connectivity, matching, colouring}
\vspace{2pt}\hrule height 0.3pt\vspace{3pt}
\textbf{Guide / Hint:} \\
{\small A DFS traversal of an undirected graph with $V$ vertices and $k$ connected components produces a spanning forest containing $k$ trees. A tree with $v_i$ vertices contains $v_i - 1$ edges, so the total number of tree edges in the forest is $\sum (v_i - 1) = V - k$. Given $V = 100$ and $E = 40$, we have $40 = 100 - k \implies k = 60$ connected components.}

\end{tcolorbox}
\vspace{6pt}\hrule height 0.35pt\vspace{6pt}
\noindent
\begin{tcolorbox}[
    blanker,
    sidebyside,
    sidebyside align=top,
    lefthand width=0.55\textwidth,
    top=4pt, bottom=4pt, left=0pt, right=0pt,
    middle=16pt,
    underlay={
        \draw[dashed, black!80, line width=0.8pt, dash pattern=on 4pt off 3pt]
            ([xshift=0.55\textwidth+8pt]frame.north west) --
            ([xshift=0.55\textwidth+8pt]frame.south west);
    }
]
\textbf{Q.56} \quad {\footnotesize\textsc{[MSQ $\cdot$ 2 Mark(s) $\cdot$ No Neg]}} \hfill {\footnotesize\textbf{GATE 2024 CS2} \quad \tiny\texttt{2024\_CS2\_CS\_Q51}}
\vspace{2pt}\hrule height 0.3pt\vspace{4pt}

Let $G$ be an undirected connected graph in which every edge has a positive integer weight. Suppose that every spanning tree in $G$ has even weight. Which of the following statements is/are TRUE for every such graph $G$?

\par\vspace{6pt}\noindent
\begin{tabular}{@{}p{0.98\linewidth}@{}}
$\square$\ \textbf{(A)} All edges in $G$ have even weight \\
\addlinespace[3pt]
$\square$\ \textbf{(B)} All edges in $G$ have even weight OR all edges in $G$ have odd weight \\
\addlinespace[3pt]
$\square$\ \textbf{(C)} In each cycle $C$ in $G$, all edges in $C$ have even weight \\
\addlinespace[3pt]
$\square$\ \textbf{(D)} In each cycle $C$ in $G$, either all edges in $C$ have even weight OR all edges in $C$ have odd weight
\end{tabular}

\tcblower

\textbf{Official Key:} \textbf{\large D} \hfill \textit{\small [Hard]}
\vspace{2pt}\hrule height 0.3pt\vspace{3pt}
\textbf{Topic:} Discrete Mathematics \\
\textbf{Subtopic:} \textit{Graphs: connectivity, matching, colouring}
\vspace{2pt}\hrule height 0.3pt\vspace{3pt}
\textbf{Guide / Hint:} \\
{\small Let $T$ be a spanning tree of $G$. For any chord $e \notin T$ creating a fundamental cycle $C$ with path $P \subseteq T$, $T' = T \setminus \{e'\} \cup \{e\}$ is also a spanning tree. Since $w(T)$ and $w(T')$ are both even, $w(e) \equiv w(e') \pmod 2$ for every edge on $C$. Thus all edges in cycle $C$ must have the same parity (either all even or all odd).}

\end{tcolorbox}
\vspace{6pt}\hrule height 0.35pt\vspace{6pt}
\noindent
\begin{tcolorbox}[
    blanker,
    sidebyside,
    sidebyside align=top,
    lefthand width=0.55\textwidth,
    top=4pt, bottom=4pt, left=0pt, right=0pt,
    middle=16pt,
    underlay={
        \draw[dashed, black!80, line width=0.8pt, dash pattern=on 4pt off 3pt]
            ([xshift=0.55\textwidth+8pt]frame.north west) --
            ([xshift=0.55\textwidth+8pt]frame.south west);
    }
]
\textbf{Q.57} \quad {\footnotesize\textsc{[NAT $\cdot$ 2 Mark(s) $\cdot$ No Neg]}} \hfill {\footnotesize\textbf{GATE 2024 CS2} \quad \tiny\texttt{2024\_CS2\_CS\_Q60}}
\vspace{2pt}\hrule height 0.3pt\vspace{4pt}

The chromatic number of a graph is the minimum number of colours used in a proper colouring of the graph. The chromatic number of the given bipartite graph is \_\_\_\_\_\_\. (Answer in integer)

\par\vspace{6pt}\noindent\textbf{Your Answer:} \fbox{\rule{0pt}{1.3ex}\hspace{3cm}}

\tcblower

\textbf{Official Key:} \textbf{\large 2 to 2} \hfill \textit{\small [Easy]}
\vspace{2pt}\hrule height 0.3pt\vspace{3pt}
\textbf{Topic:} Discrete Mathematics \\
\textbf{Subtopic:} \textit{Graphs: connectivity, matching, colouring}
\vspace{2pt}\hrule height 0.3pt\vspace{3pt}
\textbf{Guide / Hint:} \\
{\small A graph has chromatic number 2 if and only if it is bipartite and contains at least one edge, which is equivalent to having no odd cycles. Inspection of the given graph verifies that it contains no odd-length cycles, meaning it can be properly colored with 2 colors.}

\end{tcolorbox}
\vspace{6pt}\hrule height 0.35pt\vspace{6pt}
\noindent
\begin{tcolorbox}[
    blanker,
    sidebyside,
    sidebyside align=top,
    lefthand width=0.55\textwidth,
    top=4pt, bottom=4pt, left=0pt, right=0pt,
    middle=16pt,
    underlay={
        \draw[dashed, black!80, line width=0.8pt, dash pattern=on 4pt off 3pt]
            ([xshift=0.55\textwidth+8pt]frame.north west) --
            ([xshift=0.55\textwidth+8pt]frame.south west);
    }
]
\textbf{Q.58} \quad {\footnotesize\textsc{[NAT $\cdot$ 2 Mark(s) $\cdot$ No Neg]}} \hfill {\footnotesize\textbf{GATE 2024 CS2} \quad \tiny\texttt{2024\_CS2\_CS\_Q63}}
\vspace{2pt}\hrule height 0.3pt\vspace{4pt}

Let $Z_n$ be the group of integers $\{0, 1, 2, \dots, n-1\}$ with addition modulo $n$ as the group operation. The number of elements in the group $Z_2 \times Z_3 \times Z_4$ that are their own inverses is \_\_\_\_\_\_\. (Answer in integer)

\par\vspace{6pt}\noindent\textbf{Your Answer:} \fbox{\rule{0pt}{1.3ex}\hspace{3cm}}

\tcblower

\textbf{Official Key:} \textbf{\large 4 to 4} \hfill \textit{\small [Easy]}
\vspace{2pt}\hrule height 0.3pt\vspace{3pt}
\textbf{Topic:} Discrete Mathematics \\
\textbf{Subtopic:} \textit{Monoids, Groups}
\vspace{2pt}\hrule height 0.3pt\vspace{3pt}
\textbf{Guide / Hint:} \\
{\small In an additive group, an element $x$ is its own inverse if $x + x \equiv 0 \pmod n \implies 2x \equiv 0 \pmod n$. In $\mathbb{Z}_2$, elements satisfying $2x \equiv 0$ are $\{0, 1\}$ (2 elements). In $\mathbb{Z}_3$, only $\{0\}$ satisfies $2x \equiv 0 \pmod 3$ (1 element). In $\mathbb{Z}_4$, $\{0, 2\}$ satisfy $2x \equiv 0 \pmod 4$ (2 elements). In the direct product group $\mathbb{Z}_2 \times \mathbb{Z}_3 \times \mathbb{Z}_4$, total self-inverse elements $= 2 \times 1 \times 2 = 4$.}

\end{tcolorbox}
\vspace{6pt}\hrule height 0.35pt\vspace{6pt}
\noindent
\begin{tcolorbox}[
    blanker,
    sidebyside,
    sidebyside align=top,
    lefthand width=0.55\textwidth,
    top=4pt, bottom=4pt, left=0pt, right=0pt,
    middle=16pt,
    underlay={
        \draw[dashed, black!80, line width=0.8pt, dash pattern=on 4pt off 3pt]
            ([xshift=0.55\textwidth+8pt]frame.north west) --
            ([xshift=0.55\textwidth+8pt]frame.south west);
    }
]
\textbf{Q.59} \quad {\footnotesize\textsc{[MSQ $\cdot$ 2 Mark(s) $\cdot$ No Neg]}} \hfill {\footnotesize\textbf{GATE 2025 CS1} \quad \tiny\texttt{2025\_CS1\_CS\_Q48}}
\vspace{2pt}\hrule height 0.3pt\vspace{4pt}

Which of the following predicate logic formulae/formula is/are CORRECT representation(s) of the statement: ``Everyone has exactly one mother''?
The meanings of the predicates used are:
\begin{itemize}
\item $mother(y, x)$: $y$ is the mother of $x$
\item $noteq(x, y)$: $x$ and $y$ are not equal
\end{itemize}

\par\vspace{6pt}\noindent
\begin{tabular}{@{}p{0.98\linewidth}@{}}
$\square$\ \textbf{(A)} $\forall x \exists y \exists z (mother(y, x) \land \neg mother(z, x))$ \\
\addlinespace[3pt]
$\square$\ \textbf{(B)} $\forall x \exists y [mother(y, x) \land \forall z (noteq(z, y) \to \neg mother(z, x))]$ \\
\addlinespace[3pt]
$\square$\ \textbf{(C)} $\forall x \forall y [mother(y, x) \to \exists z (mother(z, x) \land \neg noteq(z, y))]$ \\
\addlinespace[3pt]
$\square$\ \textbf{(D)} $\forall x \exists y [mother(y, x) \land \neg \exists z (noteq(z, y) \land mother(z, x))]$
\end{tabular}

\tcblower

\textbf{Official Key:} \textbf{\large B;D} \hfill \textit{\small [Medium]}
\vspace{2pt}\hrule height 0.3pt\vspace{3pt}
\textbf{Topic:} Discrete Mathematics \\
\textbf{Subtopic:} \textit{Propositional and first order logic}
\vspace{2pt}\hrule height 0.3pt\vspace{3pt}
\textbf{Guide / Hint:} \\
{\small ``Exactly one mother'' requires existence ($\exists y \, mother(y, x)$) and uniqueness (any $z \neq y$ cannot be a mother of $x$). Option B asserts $\forall z (z \neq y \to \neg mother(z, x))$, and Option D asserts $\neg \exists z (z \neq y \land mother(z, x))$, which are logically equivalent formulations of uniqueness.}

\end{tcolorbox}
\vspace{6pt}\hrule height 0.35pt\vspace{6pt}
\noindent
\begin{tcolorbox}[
    blanker,
    sidebyside,
    sidebyside align=top,
    lefthand width=0.55\textwidth,
    top=4pt, bottom=4pt, left=0pt, right=0pt,
    middle=16pt,
    underlay={
        \draw[dashed, black!80, line width=0.8pt, dash pattern=on 4pt off 3pt]
            ([xshift=0.55\textwidth+8pt]frame.north west) --
            ([xshift=0.55\textwidth+8pt]frame.south west);
    }
]
\textbf{Q.60} \quad {\footnotesize\textsc{[MSQ $\cdot$ 2 Mark(s) $\cdot$ No Neg]}} \hfill {\footnotesize\textbf{GATE 2025 CS1} \quad \tiny\texttt{2025\_CS1\_CS\_Q49}}
\vspace{2pt}\hrule height 0.3pt\vspace{4pt}

$A = \{0, 1, 2, 3, \dots \}$ is the set of non-negative integers. Let $\mathcal{F}$ be the set of functions from $A$ to itself. For any two functions $f_1, f_2 \in \mathcal{F}$, we define
\[ (f\_1 \odot f\_2)(n) = f\_1(n) + f\_2(n) \]
for every number $n$ in $A$. Which of the following is/are CORRECT about the mathematical structure $(\mathcal{F}, \odot)$?

\par\vspace{6pt}\noindent
\begin{tabular}{@{}p{0.98\linewidth}@{}}
$\square$\ \textbf{(A)} $(\mathcal{F}, \odot)$ is an Abelian group. \\
\addlinespace[3pt]
$\square$\ \textbf{(B)} $(\mathcal{F}, \odot)$ is an Abelian monoid. \\
\addlinespace[3pt]
$\square$\ \textbf{(C)} $(\mathcal{F}, \odot)$ is a non-Abelian group. \\
\addlinespace[3pt]
$\square$\ \textbf{(D)} $(\mathcal{F}, \odot)$ is a non-Abelian monoid.
\end{tabular}

\tcblower

\textbf{Official Key:} \textbf{\large B} \hfill \textit{\small [Medium]}
\vspace{2pt}\hrule height 0.3pt\vspace{3pt}
\textbf{Topic:} Discrete Mathematics \\
\textbf{Subtopic:} \textit{Monoids, Groups}
\vspace{2pt}\hrule height 0.3pt\vspace{3pt}
\textbf{Guide / Hint:} \\
{\small Addition on non-negative integers is closed, associative, and commutative, so $\odot$ is associative and commutative. The identity function is $e(n) = 0$ for all $n \in A$, making $(\mathcal{F}, \odot)$ a commutative (Abelian) monoid. Inverses do not exist because negative function values are not in $A$, so it is not a group.}

\end{tcolorbox}
\vspace{6pt}\hrule height 0.35pt\vspace{6pt}
\noindent
\begin{tcolorbox}[
    blanker,
    sidebyside,
    sidebyside align=top,
    lefthand width=0.55\textwidth,
    top=4pt, bottom=4pt, left=0pt, right=0pt,
    middle=16pt,
    underlay={
        \draw[dashed, black!80, line width=0.8pt, dash pattern=on 4pt off 3pt]
            ([xshift=0.55\textwidth+8pt]frame.north west) --
            ([xshift=0.55\textwidth+8pt]frame.south west);
    }
]
\textbf{Q.61} \quad {\footnotesize\textsc{[MSQ $\cdot$ 2 Mark(s) $\cdot$ No Neg]}} \hfill {\footnotesize\textbf{GATE 2025 CS2} \quad \tiny\texttt{2025\_CS2\_CS\_Q42}}
\vspace{2pt}\hrule height 0.3pt\vspace{4pt}

Let $\mathcal{F}$ be the set of all functions from $\{1, \dots, n\}$ to $\{0, 1\}$. Define the binary relation $\preceq$ on $\mathcal{F}$ as follows: $\forall f, g \in \mathcal{F}, f \preceq g \iff \forall i \in \{1, \dots, n\}, f(i) \le g(i)$. Which of the following is/are TRUE?

\par\vspace{6pt}\noindent
\begin{tabular}{@{}p{0.98\linewidth}@{}}
$\square$\ \textbf{(A)} $\preceq$ is a symmetric relation \\
\addlinespace[3pt]
$\square$\ \textbf{(B)} $(\mathcal{F}, \preceq)$ is a partial order \\
\addlinespace[3pt]
$\square$\ \textbf{(C)} $(\mathcal{F}, \preceq)$ is a lattice \\
\addlinespace[3pt]
$\square$\ \textbf{(D)} $\preceq$ is an equivalence relation
\end{tabular}

\tcblower

\textbf{Official Key:} \textbf{\large B;C} \hfill \textit{\small [Hard]}
\vspace{2pt}\hrule height 0.3pt\vspace{3pt}
\textbf{Topic:} Discrete Mathematics \\
\textbf{Subtopic:} \textit{Sets, relations, functions, partial orders and lattices}
\vspace{2pt}\hrule height 0.3pt\vspace{3pt}
\textbf{Guide / Hint:} \\
{\small The pointwise ordering $f \preceq g \iff \forall i, f(i) \le g(i)$ is reflexive, antisymmetric, and transitive, forming a partial order (poset). Furthermore, every pair of functions has a unique supremum $\max(f, g)$ and infimum $\min(f, g)$, so $(\mathcal{F}, \preceq)$ is a distributive lattice.}

\end{tcolorbox}
\vspace{6pt}\hrule height 0.35pt\vspace{6pt}
\vspace{8pt}
\subsection*{Digital Logic}
\vspace{-4pt}
\noindent
\begin{tcolorbox}[
    blanker,
    sidebyside,
    sidebyside align=top,
    lefthand width=0.55\textwidth,
    top=4pt, bottom=4pt, left=0pt, right=0pt,
    middle=16pt,
    underlay={
        \draw[dashed, black!80, line width=0.8pt, dash pattern=on 4pt off 3pt]
            ([xshift=0.55\textwidth+8pt]frame.north west) --
            ([xshift=0.55\textwidth+8pt]frame.south west);
    }
]
\textbf{Q.62} \quad {\footnotesize\textsc{[MCQ $\cdot$ 2 Mark(s) $\cdot$ Neg: 0.66]}} \hfill {\footnotesize\textbf{GATE 2018 CS} \quad \tiny\texttt{2018\_CS\_GA\_Q33}}
\vspace{2pt}\hrule height 0.3pt\vspace{4pt}

Consider the unsigned 8-bit fixed point binary number representation below:
\[ b\_7 b\_6 b\_5 b\_4 b\_3 \, . \, b\_2 b\_1 b\_0 \]
where the position of the binary point is between $b_3$ and $b_2$. Which of the following decimal numbers can be represented exactly in this format?
\begin{enumerate}[label=(\roman*)]
    \item 0.500
    \item 0.875
    \item 0.100
    \item 0.001
\end{enumerate}

\par\vspace{6pt}\noindent
\begin{tabular}{@{}p{0.48\linewidth}@{\hspace{0.04\linewidth}}p{0.48\linewidth}@{}}
$\square$\ \textbf{(A)} (i) and (ii) only & $\square$\ \textbf{(B)} (i) and (iii) only \\
\addlinespace[3pt]
$\square$\ \textbf{(C)} (i) and (ii) only & $\square$\ \textbf{(D)} (i), (ii) and (iii)
\end{tabular}

\tcblower

\textbf{Official Key:} \textbf{\large C} \hfill \textit{\small [Medium]}
\vspace{2pt}\hrule height 0.3pt\vspace{3pt}
\textbf{Topic:} Digital Logic \\
\textbf{Subtopic:} \textit{Number representations and computer arithmetic (fixed and floating point)}
\vspace{2pt}\hrule height 0.3pt\vspace{3pt}
\textbf{Guide / Hint:} \\
{\small The fractional portion uses 3 bits ($b_2 b_1 b_0$), so the fractional precision is $2^{-3} = 0.125$. Any representable fraction must be an integer multiple of $0.125$: $0.500 = 4 \times 0.125$ (i) and $0.875 = 7 \times 0.125$ (ii) are exact multiples, while $0.100$ and $0.001$ are not.}

\end{tcolorbox}
\vspace{6pt}\hrule height 0.35pt\vspace{6pt}
\noindent
\begin{tcolorbox}[
    blanker,
    sidebyside,
    sidebyside align=top,
    lefthand width=0.55\textwidth,
    top=4pt, bottom=4pt, left=0pt, right=0pt,
    middle=16pt,
    underlay={
        \draw[dashed, black!80, line width=0.8pt, dash pattern=on 4pt off 3pt]
            ([xshift=0.55\textwidth+8pt]frame.north west) --
            ([xshift=0.55\textwidth+8pt]frame.south west);
    }
]
\textbf{Q.63} \quad {\footnotesize\textsc{[MCQ $\cdot$ 2 Mark(s) $\cdot$ Neg: 0.66]}} \hfill {\footnotesize\textbf{GATE 2022 CS} \quad \tiny\texttt{2022\_CS\_CS\_Q41}}
\vspace{2pt}\hrule height 0.3pt\vspace{4pt}

Consider three floating point numbers A, B and C stored in registers RA, RB and RC, respectively as per IEEE-754 single precision floating point format. The 32-bit content stored in these registers (in hexadecimal form) are as follows:
\[ RA = \text{0xC1400000}, \quad RB = \text{0x42100000}, \quad RC = \text{0x41400000} \]
Which one of the following is FALSE?

\par\vspace{6pt}\noindent
\begin{tabular}{@{}p{0.48\linewidth}@{\hspace{0.04\linewidth}}p{0.48\linewidth}@{}}
$\square$\ \textbf{(A)} $A + C = 0$ & $\square$\ \textbf{(B)} $C = A + B$ \\
\addlinespace[3pt]
$\square$\ \textbf{(C)} $B = 3C$ & $\square$\ \textbf{(D)} $(B - C) > 0$
\end{tabular}

\tcblower

\textbf{Official Key:} \textbf{\large B} \hfill \textit{\small [Medium]}
\vspace{2pt}\hrule height 0.3pt\vspace{3pt}
\textbf{Topic:} Digital Logic \\
\textbf{Subtopic:} \textit{Number representations and computer arithmetic (fixed and floating point)}
\vspace{2pt}\hrule height 0.3pt\vspace{3pt}
\textbf{Guide / Hint:} \\
{\small Decoding 32-bit IEEE-754: $RA = \text{0xC1400000} = -1.5 \times 2^{130-127} = -12.0$; $RB = \text{0x42100000} = +1.125 \times 2^{132-127} = +36.0$; $RC = \text{0x41400000} = +1.5 \times 2^{130-127} = +12.0$. Evaluating option B: $A + B = -12 + 36 = 24 \neq C = 12.0$, so statement B is FALSE.}

\end{tcolorbox}
\vspace{6pt}\hrule height 0.35pt\vspace{6pt}
\noindent
\begin{tcolorbox}[
    blanker,
    sidebyside,
    sidebyside align=top,
    lefthand width=0.55\textwidth,
    top=4pt, bottom=4pt, left=0pt, right=0pt,
    middle=16pt,
    underlay={
        \draw[dashed, black!80, line width=0.8pt, dash pattern=on 4pt off 3pt]
            ([xshift=0.55\textwidth+8pt]frame.north west) --
            ([xshift=0.55\textwidth+8pt]frame.south west);
    }
]
\textbf{Q.64} \quad {\footnotesize\textsc{[MCQ $\cdot$ 2 Mark(s) $\cdot$ Neg: 0.66]}} \hfill {\footnotesize\textbf{GATE 2023 CS} \quad \tiny\texttt{2023\_CS\_CS\_Q45}}
\vspace{2pt}\hrule height 0.3pt\vspace{4pt}

Consider the IEEE-754 single precision floating point numbers $P = \text{0xC1800000}$ and $Q = \text{0x3F5C2EF4}$.
Which one of the following corresponds to the product of these numbers (i.e., $P \times Q$), represented in the IEEE-754 single precision format?

\par\vspace{6pt}\noindent
\begin{tabular}{@{}p{0.48\linewidth}@{\hspace{0.04\linewidth}}p{0.48\linewidth}@{}}
$\square$\ \textbf{(A)} 0x404C2EF4 & $\square$\ \textbf{(B)} 0x405C2EF4 \\
\addlinespace[3pt]
$\square$\ \textbf{(C)} 0xC15C2EF4 & $\square$\ \textbf{(D)} 0xC14C2EF4
\end{tabular}

\tcblower

\textbf{Official Key:} \textbf{\large C} \hfill \textit{\small [Medium]}
\vspace{2pt}\hrule height 0.3pt\vspace{3pt}
\textbf{Topic:} Digital Logic \\
\textbf{Subtopic:} \textit{Number representations and computer arithmetic (fixed and floating point)}
\vspace{2pt}\hrule height 0.3pt\vspace{3pt}
\textbf{Guide / Hint:} \\
{\small Number $P = \text{0xC1800000}$ has sign bit 1, biased exponent $131$, actual exponent $131 - 127 = 4$, and mantissa $1.0$, so $P = -16.0 = -2^4$. Number $Q = \text{0x3F5C2EF4}$ is positive with biased exponent $126$ (exponent $-1$) and mantissa fraction `5C2EF4`. Multiplying $P \times Q$ flips the sign to negative (1), adds 4 to the exponent (biased exponent becomes $126 + 4 = 130 = \text{0x82}$), and retains the mantissa fraction, resulting in `0xC15C2EF4`.}

\end{tcolorbox}
\vspace{6pt}\hrule height 0.35pt\vspace{6pt}
\noindent
\begin{tcolorbox}[
    blanker,
    sidebyside,
    sidebyside align=top,
    lefthand width=0.55\textwidth,
    top=4pt, bottom=4pt, left=0pt, right=0pt,
    middle=16pt,
    underlay={
        \draw[dashed, black!80, line width=0.8pt, dash pattern=on 4pt off 3pt]
            ([xshift=0.55\textwidth+8pt]frame.north west) --
            ([xshift=0.55\textwidth+8pt]frame.south west);
    }
]
\textbf{Q.65} \quad {\footnotesize\textsc{[MSQ $\cdot$ 2 Mark(s) $\cdot$ No Neg]}} \hfill {\footnotesize\textbf{GATE 2024 CS1} \quad \tiny\texttt{2024\_CS1\_CS\_Q47}}
\vspace{2pt}\hrule height 0.3pt\vspace{4pt}

Consider a Boolean expression given by $F(X, Y, Z) = \sum(3, 5, 6, 7)$. Which of the following statements is/are CORRECT?

\par\vspace{6pt}\noindent
\begin{tabular}{@{}p{0.98\linewidth}@{}}
$\square$\ \textbf{(A)} $F(X, Y, Z) = \prod(0, 1, 2, 4)$ \\
\addlinespace[3pt]
$\square$\ \textbf{(B)} $F(X, Y, Z) = XY + YZ + XZ$ \\
\addlinespace[3pt]
$\square$\ \textbf{(C)} $F(X, Y, Z)$ is independent of input $Y$ \\
\addlinespace[3pt]
$\square$\ \textbf{(D)} $F(X, Y, Z)$ is independent of input $X$
\end{tabular}

\tcblower

\textbf{Official Key:} \textbf{\large A;B} \hfill \textit{\small [Easy]}
\vspace{2pt}\hrule height 0.3pt\vspace{3pt}
\textbf{Topic:} Digital Logic \\
\textbf{Subtopic:} \textit{Minimization}
\vspace{2pt}\hrule height 0.3pt\vspace{3pt}
\textbf{Guide / Hint:} \\
{\small The maxterm representation is the product of maxterms not in the minterm list: $\prod(0, 1, 2, 4)$ (A is TRUE). Grouping minterms $\Sigma(3, 5, 6, 7)$ on a 3-variable K-map: pairs are $(3, 7) \implies YZ$, $(5, 7) \implies XZ$, and $(6, 7) \implies XY$, giving minimal SOP $F(X, Y, Z) = XY + YZ + XZ$ (B is TRUE).}

\end{tcolorbox}
\vspace{6pt}\hrule height 0.35pt\vspace{6pt}
\noindent
\begin{tcolorbox}[
    blanker,
    sidebyside,
    sidebyside align=top,
    lefthand width=0.55\textwidth,
    top=4pt, bottom=4pt, left=0pt, right=0pt,
    middle=16pt,
    underlay={
        \draw[dashed, black!80, line width=0.8pt, dash pattern=on 4pt off 3pt]
            ([xshift=0.55\textwidth+8pt]frame.north west) --
            ([xshift=0.55\textwidth+8pt]frame.south west);
    }
]
\textbf{Q.66} \quad {\footnotesize\textsc{[MSQ $\cdot$ 2 Mark(s) $\cdot$ No Neg]}} \hfill {\footnotesize\textbf{GATE 2024 CS2} \quad \tiny\texttt{2024\_CS2\_CS\_Q49}}
\vspace{2pt}\hrule height 0.3pt\vspace{4pt}

Which of the following is/are EQUAL to 224 in radix-5 (i.e., base-5) notation?

\par\vspace{6pt}\noindent
\begin{tabular}{@{}p{0.48\linewidth}@{\hspace{0.04\linewidth}}p{0.48\linewidth}@{}}
$\square$\ \textbf{(A)} 64 in radix-10 & $\square$\ \textbf{(B)} 100 in radix-8 \\
\addlinespace[3pt]
$\square$\ \textbf{(C)} 50 in radix-16 & $\square$\ \textbf{(D)} 121 in radix-7
\end{tabular}

\tcblower

\textbf{Official Key:} \textbf{\large A;B;D} \hfill \textit{\small [Easy]}
\vspace{2pt}\hrule height 0.3pt\vspace{3pt}
\textbf{Topic:} Digital Logic \\
\textbf{Subtopic:} \textit{Number representations and computer arithmetic (fixed and floating point)}
\vspace{2pt}\hrule height 0.3pt\vspace{3pt}
\textbf{Guide / Hint:} \\
{\small Convert $224_5$ to decimal: $2(5^2) + 2(5^1) + 4(5^0) = 50 + 10 + 4 = 64_{10}$ (A is true). In base 8: $64_{10} = 100_8$ (B is true). In base 16: $64_{10} = 40_{16} \neq 50_{16}$ (C is false). In base 7: $64_{10} = 1(7^2) + 2(7^1) + 1(7^0) = 121_7$ (D is true).}

\end{tcolorbox}
\vspace{6pt}\hrule height 0.35pt\vspace{6pt}
\noindent
\begin{tcolorbox}[
    blanker,
    sidebyside,
    sidebyside align=top,
    lefthand width=0.55\textwidth,
    top=4pt, bottom=4pt, left=0pt, right=0pt,
    middle=16pt,
    underlay={
        \draw[dashed, black!80, line width=0.8pt, dash pattern=on 4pt off 3pt]
            ([xshift=0.55\textwidth+8pt]frame.north west) --
            ([xshift=0.55\textwidth+8pt]frame.south west);
    }
]
\textbf{Q.67} \quad {\footnotesize\textsc{[MCQ $\cdot$ 2 Mark(s) $\cdot$ Neg: 0.66]}} \hfill {\footnotesize\textbf{GATE 2025 CS1} \quad \tiny\texttt{2025\_CS1\_CS\_Q42}}
\vspace{2pt}\hrule height 0.3pt\vspace{4pt}

Consider the following four variable Boolean function in sum-of-product form:
\[ F(b\_3, b\_2, b\_1, b\_0) = \sum(0, 2, 4, 8, 10, 11, 12) \]
where the value of the function is computed by considering $b_3b_2b_1b_0$ as a 4-bit binary number, where $b_3$ denotes the most significant bit and $b_0$ denotes the least significant bit. Note that there are no don't care terms. Which ONE of the following options is the CORRECT minimized expression?

\par\vspace{6pt}\noindent
\begin{tabular}{@{}p{0.98\linewidth}@{}}
$\square$\ \textbf{(A)} $\overline{b}_1\overline{b}_0 + \overline{b}_2\overline{b}_0 + b_1\overline{b}_2b_3$ \\
\addlinespace[3pt]
$\square$\ \textbf{(B)} $\overline{b}_1\overline{b}_0 + \overline{b}_2\overline{b}_0$ \\
\addlinespace[3pt]
$\square$\ \textbf{(C)} $\overline{b}_2\overline{b}_0 + b_1b_2b_3$ \\
\addlinespace[3pt]
$\square$\ \textbf{(D)} $\overline{b}_0\overline{b}_2 + \overline{b}_3$
\end{tabular}

\tcblower

\textbf{Official Key:} \textbf{\large A} \hfill \textit{\small [Medium]}
\vspace{2pt}\hrule height 0.3pt\vspace{3pt}
\textbf{Topic:} Digital Logic \\
\textbf{Subtopic:} \textit{Minimization}
\vspace{2pt}\hrule height 0.3pt\vspace{3pt}
\textbf{Guide / Hint:} \\
{\small Grouping minterms on a 4-variable K-map: minterms $\{0, 2, 8, 10\}$ form a 4-cell group giving $\overline{b}_2\overline{b}_0$; minterms $\{0, 4, 8, 12\}$ form a 4-cell group giving $\overline{b}_1\overline{b}_0$; remaining minterm 11 combines with 10 to give $b_3\overline{b}_2b_1$. Combining these essential prime implicants yields $\overline{b}_1\overline{b}_0 + \overline{b}_2\overline{b}_0 + b_1\overline{b}_2b_3$.}

\end{tcolorbox}
\vspace{6pt}\hrule height 0.35pt\vspace{6pt}
\noindent
\begin{tcolorbox}[
    blanker,
    sidebyside,
    sidebyside align=top,
    lefthand width=0.55\textwidth,
    top=4pt, bottom=4pt, left=0pt, right=0pt,
    middle=16pt,
    underlay={
        \draw[dashed, black!80, line width=0.8pt, dash pattern=on 4pt off 3pt]
            ([xshift=0.55\textwidth+8pt]frame.north west) --
            ([xshift=0.55\textwidth+8pt]frame.south west);
    }
]
\textbf{Q.68} \quad {\footnotesize\textsc{[NAT $\cdot$ 2 Mark(s) $\cdot$ No Neg]}} \hfill {\footnotesize\textbf{GATE 2025 CS1} \quad \tiny\texttt{2025\_CS1\_CS\_Q59}}
\vspace{2pt}\hrule height 0.3pt\vspace{4pt}

Consider a finite state machine (FSM) with one input $X$ and one output $f$, represented by the given state transition table:
\par\medskip
\begin{tabular}{|c|cc|cc|}
\hline
\textbf{Present state} & \multicolumn{2}{c|}{\textbf{Next state}} & \multicolumn{2}{c|}{\textbf{Output } $f$} \\
& $X=0$ & $X=1$ & $X=0$ & $X=1$ \\
\hline
A & F & B & 0 & 0 \\
B & D & C & 0 & 0 \\
C & F & E & 0 & 0 \\
D & G & A & 1 & 0 \\
E & D & C & 0 & 0 \\
F & F & B & 1 & 1 \\
G & G & H & 0 & 1 \\
H & G & A & 1 & 0 \\
\hline
\end{tabular}
\par\medskip
The minimum number of states required to realize this FSM is \_\_\_\_\_\_\. (Answer in integer)

\par\vspace{6pt}\noindent\textbf{Your Answer:} \fbox{\rule{0pt}{1.3ex}\hspace{3cm}}

\tcblower

\textbf{Official Key:} \textbf{\large 5 to 5} \hfill \textit{\small [Hard]}
\vspace{2pt}\hrule height 0.3pt\vspace{3pt}
\textbf{Topic:} Digital Logic \\
\textbf{Subtopic:} \textit{Combinational and sequential circuits}
\vspace{2pt}\hrule height 0.3pt\vspace{3pt}
\textbf{Guide / Hint:} \\
{\small Using state equivalence partitioning: initial partition by outputs groups states into $P_0 = \{(A, B, C, E), (D, H), (F), (G)\}$. Comparing transitions on $X=0, 1$: states $B$ and $E$ have identical transitions to $(D, C)$, while $A$ and $C$ transition to $(F, B)$ and $(F, E)$ (which are equivalent since $B \equiv E$). States $D$ and $H$ have identical next states $(G, A)$. Thus, equivalence classes are $\{A, C\}, \{B, E\}, \{D, H\}, \{F\}, \{G\}$, resulting in exactly 5 minimal states.}

\end{tcolorbox}
\vspace{6pt}\hrule height 0.35pt\vspace{6pt}
\noindent
\begin{tcolorbox}[
    blanker,
    sidebyside,
    sidebyside align=top,
    lefthand width=0.55\textwidth,
    top=4pt, bottom=4pt, left=0pt, right=0pt,
    middle=16pt,
    underlay={
        \draw[dashed, black!80, line width=0.8pt, dash pattern=on 4pt off 3pt]
            ([xshift=0.55\textwidth+8pt]frame.north west) --
            ([xshift=0.55\textwidth+8pt]frame.south west);
    }
]
\textbf{Q.69} \quad {\footnotesize\textsc{[NAT $\cdot$ 2 Mark(s) $\cdot$ No Neg]}} \hfill {\footnotesize\textbf{GATE 2025 CS1} \quad \tiny\texttt{2025\_CS1\_CS\_Q60}}
\vspace{2pt}\hrule height 0.3pt\vspace{4pt}

Consider the given sequential circuit designed using D-Flip-flops. The circuit is initialized with some value (initial state). The number of distinct states the circuit will go through before returning back to the initial state is \_\_\_\_\_\_\. (Answer in integer)

\par\vspace{6pt}\noindent\textbf{Your Answer:} \fbox{\rule{0pt}{1.3ex}\hspace{3cm}}

\tcblower

\textbf{Official Key:} \textbf{\large 7 to 8} \hfill \textit{\small [Hard]}
\vspace{2pt}\hrule height 0.3pt\vspace{3pt}
\textbf{Topic:} Digital Logic \\
\textbf{Subtopic:} \textit{Combinational and sequential circuits}
\vspace{2pt}\hrule height 0.3pt\vspace{3pt}
\textbf{Guide / Hint:} \\
{\small Tracing the characteristic shift register feedback transitions where the inputs to the D flip-flops are determined by logic feedback of the current state $(Q_3, Q_2, Q_1, Q_0)$, the state cycle length forms a closed loop of period 7 (or 8 depending on the initial valid cycle).}

\end{tcolorbox}
\vspace{6pt}\hrule height 0.35pt\vspace{6pt}
\noindent
\begin{tcolorbox}[
    blanker,
    sidebyside,
    sidebyside align=top,
    lefthand width=0.55\textwidth,
    top=4pt, bottom=4pt, left=0pt, right=0pt,
    middle=16pt,
    underlay={
        \draw[dashed, black!80, line width=0.8pt, dash pattern=on 4pt off 3pt]
            ([xshift=0.55\textwidth+8pt]frame.north west) --
            ([xshift=0.55\textwidth+8pt]frame.south west);
    }
]
\textbf{Q.70} \quad {\footnotesize\textsc{[MSQ $\cdot$ 2 Mark(s) $\cdot$ No Neg]}} \hfill {\footnotesize\textbf{GATE 2025 CS2} \quad \tiny\texttt{2025\_CS2\_CS\_Q43}}
\vspace{2pt}\hrule height 0.3pt\vspace{4pt}

Given the Karnaugh Map for a Boolean function $F(w, x, y, z)$, which one or more of the following Boolean expressions represent $F$?

\par\vspace{6pt}\noindent
\begin{tabular}{@{}p{0.98\linewidth}@{}}
$\square$\ \textbf{(A)} $\overline{w}\overline{x}\overline{y}\overline{z} + w\overline{x}\overline{y}\overline{z} + \overline{w}\overline{x}y\overline{z} + w\overline{x}y\overline{z} + xz$ \\
\addlinespace[3pt]
$\square$\ \textbf{(B)} $\overline{w}\overline{x}\overline{y}\overline{z} + \overline{w}\overline{x}y\overline{z} + w\overline{x}yz + xz$ \\
\addlinespace[3pt]
$\square$\ \textbf{(C)} $\overline{w}\overline{x}\overline{y}\overline{z} + w\overline{x}\overline{y}\overline{z} + w\overline{x}\overline{y}z + xz$ \\
\addlinespace[3pt]
$\square$\ \textbf{(D)} $\overline{x}\overline{z} + xz$
\end{tabular}

\tcblower

\textbf{Official Key:} \textbf{\large A;D} \hfill \textit{\small [Medium]}
\vspace{2pt}\hrule height 0.3pt\vspace{3pt}
\textbf{Topic:} Digital Logic \\
\textbf{Subtopic:} \textit{Boolean algebra and minimization: Karnaugh Map}
\vspace{2pt}\hrule height 0.3pt\vspace{3pt}
\textbf{Guide / Hint:} \\
{\small Grouping the 1s in the K-map: the four corners and the central groupings give prime implicants $\overline{x}\overline{z} + xz$. Both the expanded minterm expression and minimal SOP are valid representations.}

\end{tcolorbox}
\vspace{6pt}\hrule height 0.35pt\vspace{6pt}
\noindent
\begin{tcolorbox}[
    blanker,
    sidebyside,
    sidebyside align=top,
    lefthand width=0.55\textwidth,
    top=4pt, bottom=4pt, left=0pt, right=0pt,
    middle=16pt,
    underlay={
        \draw[dashed, black!80, line width=0.8pt, dash pattern=on 4pt off 3pt]
            ([xshift=0.55\textwidth+8pt]frame.north west) --
            ([xshift=0.55\textwidth+8pt]frame.south west);
    }
]
\textbf{Q.71} \quad {\footnotesize\textsc{[MSQ $\cdot$ 2 Mark(s) $\cdot$ No Neg]}} \hfill {\footnotesize\textbf{GATE 2025 CS2} \quad \tiny\texttt{2025\_CS2\_CS\_Q49}}
\vspace{2pt}\hrule height 0.3pt\vspace{4pt}

Three floating point numbers $X, Y,$ and $Z$ are stored in three registers RX, RY, and RZ, respectively in IEEE 754 single-precision format. Which of the following relationships is/are CORRECT?

\par\vspace{6pt}\noindent
\begin{tabular}{@{}p{0.48\linewidth}@{\hspace{0.04\linewidth}}p{0.48\linewidth}@{}}
$\square$\ \textbf{(A)} $4(X+Y) + Z = 0$ & $\square$\ \textbf{(B)} $2Y - Z = 0$ \\
\addlinespace[3pt]
$\square$\ \textbf{(C)} $4X + 3Z = 0$ & $\square$\ \textbf{(D)} $X+Y+Z = 0$
\end{tabular}

\tcblower

\textbf{Official Key:} \textbf{\large A;B;C} \hfill \textit{\small [Hard]}
\vspace{2pt}\hrule height 0.3pt\vspace{3pt}
\textbf{Topic:} Digital Logic \\
\textbf{Subtopic:} \textit{Number representations and computer arithmetic: IEEE 754 floating point}
\vspace{2pt}\hrule height 0.3pt\vspace{3pt}
\textbf{Guide / Hint:} \\
{\small Decode the IEEE 754 representations in RX, RY, RZ into decimal values $X, Y, Z$. Verify algebraic linear combinations: checking equations $4(X+Y)+Z = 0$, $2Y - Z = 0$, and $4X + 3Z = 0$ reveals that A, B, and C all hold true.}

\end{tcolorbox}
\vspace{6pt}\hrule height 0.35pt\vspace{6pt}
\noindent
\begin{tcolorbox}[
    blanker,
    sidebyside,
    sidebyside align=top,
    lefthand width=0.55\textwidth,
    top=4pt, bottom=4pt, left=0pt, right=0pt,
    middle=16pt,
    underlay={
        \draw[dashed, black!80, line width=0.8pt, dash pattern=on 4pt off 3pt]
            ([xshift=0.55\textwidth+8pt]frame.north west) --
            ([xshift=0.55\textwidth+8pt]frame.south west);
    }
]
\textbf{Q.72} \quad {\footnotesize\textsc{[MSQ $\cdot$ 2 Mark(s) $\cdot$ No Neg]}} \hfill {\footnotesize\textbf{GATE 2025 CS2} \quad \tiny\texttt{2025\_CS2\_CS\_Q50}}
\vspace{2pt}\hrule height 0.3pt\vspace{4pt}

Which of the following Boolean algebraic equation(s) is/are CORRECT?

\par\vspace{6pt}\noindent
\begin{tabular}{@{}p{0.98\linewidth}@{}}
$\square$\ \textbf{(A)} $\overline{A}BC + A\overline{B}\overline{C} + \overline{A}\overline{B}\overline{C} + A\overline{B}C + ABC = BC + \overline{B}\overline{C} + \overline{A}\overline{B}$ \\
\addlinespace[3pt]
$\square$\ \textbf{(B)} $AB + \overline{A}C + BC = AB + \overline{A}C$ \\
\addlinespace[3pt]
$\square$\ \textbf{(C)} $(A + C)(\overline{A} + B) = AB + \overline{A}C$ \\
\addlinespace[3pt]
$\square$\ \textbf{(D)} $\overline{(A + \overline{B} + \overline{D})(C + D)(\overline{A} + C + D)(A + B + \overline{D})} = \overline{A}D + \overline{C}\overline{D}$
\end{tabular}

\tcblower

\textbf{Official Key:} \textbf{\large B;C;D} \hfill \textit{\small [Medium]}
\vspace{2pt}\hrule height 0.3pt\vspace{3pt}
\textbf{Topic:} Digital Logic \\
\textbf{Subtopic:} \textit{Boolean algebra and minimization: identities and consensus theorem}
\vspace{2pt}\hrule height 0.3pt\vspace{3pt}
\textbf{Guide / Hint:} \\
{\small Equation B is the classic Consensus Theorem: $AB + \overline{A}C + BC = AB + \overline{A}C$. Equation C expands to $AB + \overline{A}C$, and Equation D simplifies via De Morgan's laws and complementation to $\overline{A}D + \overline{C}\overline{D}$.}

\end{tcolorbox}
\vspace{6pt}\hrule height 0.35pt\vspace{6pt}
\vspace{8pt}
\subsection*{Computer Organization and Architecture}
\vspace{-4pt}
\noindent
\begin{tcolorbox}[
    blanker,
    sidebyside,
    sidebyside align=top,
    lefthand width=0.55\textwidth,
    top=4pt, bottom=4pt, left=0pt, right=0pt,
    middle=16pt,
    underlay={
        \draw[dashed, black!80, line width=0.8pt, dash pattern=on 4pt off 3pt]
            ([xshift=0.55\textwidth+8pt]frame.north west) --
            ([xshift=0.55\textwidth+8pt]frame.south west);
    }
]
\textbf{Q.73} \quad {\footnotesize\textsc{[MCQ $\cdot$ 2 Mark(s) $\cdot$ Neg: 0.66]}} \hfill {\footnotesize\textbf{GATE 2018 CS} \quad \tiny\texttt{2018\_CS\_GA\_Q34}}
\vspace{2pt}\hrule height 0.3pt\vspace{4pt}

The size of the physical address space of a processor is $2^P$ bytes. The word length is $2^W$ bytes. The capacity of cache memory is $2^N$ bytes. The size of each cache block is $2^M$ words. For a $K$-way set-associative cache memory, the length (in number of bits) of the tag field is

\par\vspace{6pt}\noindent
\begin{tabular}{@{}p{0.48\linewidth}@{\hspace{0.04\linewidth}}p{0.48\linewidth}@{}}
$\square$\ \textbf{(A)} $P - N - \log_2 K$ & $\square$\ \textbf{(B)} $P - N + \log_2 K$ \\
\addlinespace[3pt]
$\square$\ \textbf{(C)} $P - N - M - W - \log_2 K$ & $\square$\ \textbf{(D)} $P - N - M - W + \log_2 K$
\end{tabular}

\tcblower

\textbf{Official Key:} \textbf{\large B} \hfill \textit{\small [Medium]}
\vspace{2pt}\hrule height 0.3pt\vspace{3pt}
\textbf{Topic:} Computer Organization and Architecture \\
\textbf{Subtopic:} \textit{Memory hierarchy: cache, main memory and secondary storage}
\vspace{2pt}\hrule height 0.3pt\vspace{3pt}
\textbf{Guide / Hint:} \\
{\small Block size $= 2^{M+W}$ bytes $\implies M+W$ offset bits. Cache size $= 2^N$ bytes, so number of lines $= 2^N / 2^{M+W} = 2^{N-M-W}$. With $K = 2^k$ lines per set, number of sets $= 2^{N-M-W-k} \implies N-M-W-k$ set index bits. Tag bits $= P - (N-M-W-k) - (M+W) = P - N + k = P - N + \log_2 K$.}

\end{tcolorbox}
\vspace{6pt}\hrule height 0.35pt\vspace{6pt}
\noindent
\begin{tcolorbox}[
    blanker,
    sidebyside,
    sidebyside align=top,
    lefthand width=0.55\textwidth,
    top=4pt, bottom=4pt, left=0pt, right=0pt,
    middle=16pt,
    underlay={
        \draw[dashed, black!80, line width=0.8pt, dash pattern=on 4pt off 3pt]
            ([xshift=0.55\textwidth+8pt]frame.north west) --
            ([xshift=0.55\textwidth+8pt]frame.south west);
    }
]
\textbf{Q.74} \quad {\footnotesize\textsc{[NAT $\cdot$ 2 Mark(s) $\cdot$ No Neg]}} \hfill {\footnotesize\textbf{GATE 2018 CS} \quad \tiny\texttt{2018\_CS\_GA\_Q50}}
\vspace{2pt}\hrule height 0.3pt\vspace{4pt}

The instruction pipeline of a RISC processor has five stages: Instruction Fetch (IF), Instruction Decode (ID), Perform Operation (PO), Operand Read (OR), and Write Back (WB). Each stage takes 1 clock cycle except PO. For a sequence of 100 instructions, 40 instructions require 3 clock cycles in PO, 35 instructions require 2 clock cycles in PO, and the remaining 25 instructions require 1 clock cycle in PO. Assuming no memory stalls or branch delays, the total number of clock cycles to execute the 100 instructions is \_\_\_\_\_\_\_\_.

\par\vspace{6pt}\noindent\textbf{Your Answer:} \fbox{\rule{0pt}{1.3ex}\hspace{3cm}}

\tcblower

\textbf{Official Key:} \textbf{\large 219 to 219} \hfill \textit{\small [Medium]}
\vspace{2pt}\hrule height 0.3pt\vspace{3pt}
\textbf{Topic:} Computer Organization and Architecture \\
\textbf{Subtopic:} \textit{Instruction pipelining, pipeline hazards}
\vspace{2pt}\hrule height 0.3pt\vspace{3pt}
\textbf{Guide / Hint:} \\
{\small For 100 instructions in a 5-stage pipeline without branch delays, ideal cycles $= 5 + (100 - 1) = 104$. The PO stage causes pipeline stalls when operand evaluation takes extra cycles: 40 instructions require 3 cycles (2 extra cycles $\implies 40 \times 2 = 80$ stall cycles) and 35 instructions require 2 cycles (1 extra cycle $\implies 35 \times 1 = 35$ stall cycles). Total clock cycles $= 104 + 80 + 35 = 219$.}

\end{tcolorbox}
\vspace{6pt}\hrule height 0.35pt\vspace{6pt}
\noindent
\begin{tcolorbox}[
    blanker,
    sidebyside,
    sidebyside align=top,
    lefthand width=0.55\textwidth,
    top=4pt, bottom=4pt, left=0pt, right=0pt,
    middle=16pt,
    underlay={
        \draw[dashed, black!80, line width=0.8pt, dash pattern=on 4pt off 3pt]
            ([xshift=0.55\textwidth+8pt]frame.north west) --
            ([xshift=0.55\textwidth+8pt]frame.south west);
    }
]
\textbf{Q.75} \quad {\footnotesize\textsc{[NAT $\cdot$ 2 Mark(s) $\cdot$ No Neg]}} \hfill {\footnotesize\textbf{GATE 2018 CS} \quad \tiny\texttt{2018\_CS\_GA\_Q51}}
\vspace{2pt}\hrule height 0.3pt\vspace{4pt}

A processor has 16 integer registers (R0 to R15) and 64 floating point registers (F0 to F63). It uses a 16-bit instruction format. Four types of instructions are supported:
\begin{itemize}[leftmargin=*, nosep]
    \item Type-1: Opcode with three integer register operands
    \item Type-2: Opcode with two floating point register operands
    \item Type-3: Opcode with one integer and one floating point register operand
    \item Type-4: Opcode with one floating point register operand
\end{itemize}
If there are 4 Type-1, 8 Type-2, and 14 Type-3 instructions, the maximum number of Type-4 instructions that can be supported is \_\_\_\_\_\_\_\_.

\par\vspace{6pt}\noindent\textbf{Your Answer:} \fbox{\rule{0pt}{1.3ex}\hspace{3cm}}

\tcblower

\textbf{Official Key:} \textbf{\large 32 to 32} \hfill \textit{\small [Hard]}
\vspace{2pt}\hrule height 0.3pt\vspace{3pt}
\textbf{Topic:} Computer Organization and Architecture \\
\textbf{Subtopic:} \textit{Machine instructions and addressing modes}
\vspace{2pt}\hrule height 0.3pt\vspace{3pt}
\textbf{Guide / Hint:} \\
{\small A 16-bit instruction space has $2^{16} = 65536$ encodings. Type-1: 4 instructions with $3 \times 4\text{-bit}$ registers consume $4 \times 2^{12} = 16384$. Type-2: 8 instructions with $2 \times 6\text{-bit}$ registers consume $8 \times 2^{12} = 32768$. Type-3: 14 instructions with 4-bit and 6-bit registers consume $14 \times 2^{10} = 14336$. Remaining encodings $= 65536 - (16384 + 32768 + 14336) = 2048$. Each Type-4 instruction uses one 6-bit register ($2^6 = 64$ encodings). Thus, maximum $N = 2048 / 64 = 32$.}

\end{tcolorbox}
\vspace{6pt}\hrule height 0.35pt\vspace{6pt}
\noindent
\begin{tcolorbox}[
    blanker,
    sidebyside,
    sidebyside align=top,
    lefthand width=0.55\textwidth,
    top=4pt, bottom=4pt, left=0pt, right=0pt,
    middle=16pt,
    underlay={
        \draw[dashed, black!80, line width=0.8pt, dash pattern=on 4pt off 3pt]
            ([xshift=0.55\textwidth+8pt]frame.north west) --
            ([xshift=0.55\textwidth+8pt]frame.south west);
    }
]
\textbf{Q.76} \quad {\footnotesize\textsc{[MSQ $\cdot$ 2 Mark(s) $\cdot$ No Neg]}} \hfill {\footnotesize\textbf{GATE 2022 CS} \quad \tiny\texttt{2022\_CS\_CS\_Q54}}
\vspace{2pt}\hrule height 0.3pt\vspace{4pt}

Consider a system with 2 KB direct mapped data cache with a block size of 64 bytes. The system has a physical address space of 64 KB and a word length of 16 bits. During the execution of a program, four data words P, Q, R, and S are accessed in that order 10 times (i.e., PQRSPQRS\dots). Hence, there are 40 accesses to data cache altogether. Assume that the data cache is initially empty and no other data words are accessed by the program. The addresses of the first bytes of P, Q, R, and S are `0xA248`, `0xC28A`, `0xCA8A`, and `0xA262`, respectively. For the execution of the above program, which of the following statements is/are TRUE with respect to the data cache?

\par\vspace{6pt}\noindent
\begin{tabular}{@{}p{0.98\linewidth}@{}}
$\square$\ \textbf{(A)} Every access to S is a hit. \\
\addlinespace[3pt]
$\square$\ \textbf{(B)} Once P is brought to the cache it is never evicted. \\
\addlinespace[3pt]
$\square$\ \textbf{(C)} At the end of the execution only R and S reside in the cache. \\
\addlinespace[3pt]
$\square$\ \textbf{(D)} Every access to R evicts Q from the cache.
\end{tabular}

\tcblower

\textbf{Official Key:} \textbf{\large A,B,D} \hfill \textit{\small [Hard]}
\vspace{2pt}\hrule height 0.3pt\vspace{3pt}
\textbf{Topic:} Computer Organization and Architecture \\
\textbf{Subtopic:} \textit{Memory hierarchy: cache, main memory and secondary storage}
\vspace{2pt}\hrule height 0.3pt\vspace{3pt}
\textbf{Guide / Hint:} \\
{\small Cache line index is calculated as $\lfloor \text{address} / 64 \rfloor \pmod{32}$. Addresses `0xA248` (P) and `0xA262` (S) both evaluate to block offset within block `0xA240`, which maps to line 9. Once P is loaded, S is also present in line 9, so every access to S is a hit (A is TRUE) and P/S are never evicted (B is TRUE). Addresses `0xC28A` (Q) and `0xCA8A` (R) both map to line 10, repeatedly evicting each other (D is TRUE).}

\end{tcolorbox}
\vspace{6pt}\hrule height 0.35pt\vspace{6pt}
\noindent
\begin{tcolorbox}[
    blanker,
    sidebyside,
    sidebyside align=top,
    lefthand width=0.55\textwidth,
    top=4pt, bottom=4pt, left=0pt, right=0pt,
    middle=16pt,
    underlay={
        \draw[dashed, black!80, line width=0.8pt, dash pattern=on 4pt off 3pt]
            ([xshift=0.55\textwidth+8pt]frame.north west) --
            ([xshift=0.55\textwidth+8pt]frame.south west);
    }
]
\textbf{Q.77} \quad {\footnotesize\textsc{[NAT $\cdot$ 2 Mark(s) $\cdot$ No Neg]}} \hfill {\footnotesize\textbf{GATE 2022 CS} \quad \tiny\texttt{2022\_CS\_CS\_Q61}}
\vspace{2pt}\hrule height 0.3pt\vspace{4pt}

A processor X1 operating at 2 GHz has a standard 5-stage RISC instruction pipeline having a base CPI (cycles per instruction) of one without any pipeline hazards. For a given program P that has 30\% branch instructions, control hazards incur 2 cycles stall for every branch. A new version of the processor X2 operating at same clock frequency has an additional branch predictor unit (BPU) that completely eliminates stalls for correctly predicted branches. There is neither any savings nor any additional stalls for wrong predictions. There are no structural hazards and data hazards for X1 and X2. If the BPU has a prediction accuracy of 80\%, the speed up (rounded off to two decimal places) obtained by X2 over X1 in executing P is \_\_\_\_\_\_\_\_\_\_\_\_\. (Answer in numeric format)

\par\vspace{6pt}\noindent\textbf{Your Answer:} \fbox{\rule{0pt}{1.3ex}\hspace{3cm}}

\tcblower

\textbf{Official Key:} \textbf{\large 1.42 to 1.45} \hfill \textit{\small [Medium]}
\vspace{2pt}\hrule height 0.3pt\vspace{3pt}
\textbf{Topic:} Computer Organization and Architecture \\
\textbf{Subtopic:} \textit{Instruction pipelining, pipeline hazards}
\vspace{2pt}\hrule height 0.3pt\vspace{3pt}
\textbf{Guide / Hint:} \\
{\small For $X_1$: $\text{CPI}_1 = 1 + (0.30 \times 2) = 1.60$. For $X_2$: with 80\% branch prediction accuracy, only 20\% of branches stall: $\text{CPI}_2 = 1 + (0.30 \times 0.20 \times 2) = 1.12$. Speedup $= \frac{\text{CPI}_1}{\text{CPI}_2} = \frac{1.60}{1.12} \approx 1.43$.}

\end{tcolorbox}
\vspace{6pt}\hrule height 0.35pt\vspace{6pt}
\noindent
\begin{tcolorbox}[
    blanker,
    sidebyside,
    sidebyside align=top,
    lefthand width=0.55\textwidth,
    top=4pt, bottom=4pt, left=0pt, right=0pt,
    middle=16pt,
    underlay={
        \draw[dashed, black!80, line width=0.8pt, dash pattern=on 4pt off 3pt]
            ([xshift=0.55\textwidth+8pt]frame.north west) --
            ([xshift=0.55\textwidth+8pt]frame.south west);
    }
]
\textbf{Q.78} \quad {\footnotesize\textsc{[MCQ $\cdot$ 2 Mark(s) $\cdot$ Neg: 0.66]}} \hfill {\footnotesize\textbf{GATE 2023 CS} \quad \tiny\texttt{2023\_CS\_CS\_Q41}}
\vspace{2pt}\hrule height 0.3pt\vspace{4pt}

Consider the given C-code and its corresponding assembly code, with a few operands U1--U4 being unknown. Some useful information as well as the semantics of each unique assembly instruction is annotated as inline comments in the code. The memory is byte-addressable.
\begin{lstlisting}[language=C]
// C-code
int a[10], b[10], i;    // int is 32-bit
for (i = 0; i < 10; i++)
    a[i] = b[i] * 8;
\end{lstlisting}
\begin{lstlisting}
; assembly-code (; indicates comments)
; r1-r5 are 32-bit integer registers
; initialize r1=0, r2=10
; initialize r3, r4 with base address of a, b
L01: jeq r1, r2, end    ; if(r1==r2) goto end
L02: lw  r5, 0(r4)      ; r5 <- Memory[r4+0]
L03: shl r5, r5, U1     ; r5 <- r5 << U1
L04: sw  r5, 0(r3)      ; Memory[r3+0] <- r5
L05: add r3, r3, U2     ; r3 <- r3+U2
L06: add r4, r4, U3
L07: add r1, r1, 1
L08: jmp U4             ; goto U4
L09: end
\end{lstlisting}
Which one of the following options is a CORRECT replacement for operands in the position (U1, U2, U3, U4) in the above assembly code?

\par\vspace{6pt}\noindent
\begin{tabular}{@{}p{0.48\linewidth}@{\hspace{0.04\linewidth}}p{0.48\linewidth}@{}}
$\square$\ \textbf{(A)} (8, 4, 1, L02) & $\square$\ \textbf{(B)} (3, 4, 4, L01) \\
\addlinespace[3pt]
$\square$\ \textbf{(C)} (8, 1, 1, L02) & $\square$\ \textbf{(D)} (3, 1, 1, L01)
\end{tabular}

\tcblower

\textbf{Official Key:} \textbf{\large B} \hfill \textit{\small [Medium]}
\vspace{2pt}\hrule height 0.3pt\vspace{3pt}
\textbf{Topic:} Computer Organization and Architecture \\
\textbf{Subtopic:} \textit{Machine instructions and addressing modes}
\vspace{2pt}\hrule height 0.3pt\vspace{3pt}
\textbf{Guide / Hint:} \\
{\small Multiplying by 8 corresponds to shifting left by 3 bits ($2^3 = 8$), so $U_1 = 3$. Since elements are 32-bit (4-byte) integers in byte-addressable memory, the pointers in registers `r3` and `r4` must advance by 4 bytes each iteration, so $U_2 = 4$ and $U_3 = 4$. Finally, the loop branches back to the condition check at label $L_{01}$, so $U_4 = L_{01}$.}

\end{tcolorbox}
\vspace{6pt}\hrule height 0.35pt\vspace{6pt}
\noindent
\begin{tcolorbox}[
    blanker,
    sidebyside,
    sidebyside align=top,
    lefthand width=0.55\textwidth,
    top=4pt, bottom=4pt, left=0pt, right=0pt,
    middle=16pt,
    underlay={
        \draw[dashed, black!80, line width=0.8pt, dash pattern=on 4pt off 3pt]
            ([xshift=0.55\textwidth+8pt]frame.north west) --
            ([xshift=0.55\textwidth+8pt]frame.south west);
    }
]
\textbf{Q.79} \quad {\footnotesize\textsc{[NAT $\cdot$ 2 Mark(s) $\cdot$ No Neg]}} \hfill {\footnotesize\textbf{GATE 2023 CS} \quad \tiny\texttt{2023\_CS\_CS\_Q64}}
\vspace{2pt}\hrule height 0.3pt\vspace{4pt}

An 8-way set associative cache of size 64 KB (1 KB = 1024 bytes) is used in a system with 32-bit address. The address is sub-divided into TAG, INDEX, and BLOCK OFFSET.
The number of bits in the TAG is \_\_\_\_\_\_\_\_\. (Answer in integer)

\par\vspace{6pt}\noindent\textbf{Your Answer:} \fbox{\rule{0pt}{1.3ex}\hspace{3cm}}

\tcblower

\textbf{Official Key:} \textbf{\large 19 to 19} \hfill \textit{\small [Easy]}
\vspace{2pt}\hrule height 0.3pt\vspace{3pt}
\textbf{Topic:} Computer Organization and Architecture \\
\textbf{Subtopic:} \textit{Memory hierarchy: cache, main memory and secondary storage}
\vspace{2pt}\hrule height 0.3pt\vspace{3pt}
\textbf{Guide / Hint:} \\
{\small Address size $= 32\text{ bits}$. Cache size $= 64\text{ KB} = 2^{16}\text{ bytes}$. For an 8-way set-associative cache, the combined capacity of index and offset is $\frac{\text{Cache Size}}{\text{Associativity}} = \frac{64\text{ KB}}{8} = 8\text{ KB} = 2^{13}\text{ bytes}$, requiring $13\text{ bits}$ for INDEX + BLOCK OFFSET. Therefore, the number of TAG bits is $32 - 13 = 19\text{ bits}$.}

\end{tcolorbox}
\vspace{6pt}\hrule height 0.35pt\vspace{6pt}
\noindent
\begin{tcolorbox}[
    blanker,
    sidebyside,
    sidebyside align=top,
    lefthand width=0.55\textwidth,
    top=4pt, bottom=4pt, left=0pt, right=0pt,
    middle=16pt,
    underlay={
        \draw[dashed, black!80, line width=0.8pt, dash pattern=on 4pt off 3pt]
            ([xshift=0.55\textwidth+8pt]frame.north west) --
            ([xshift=0.55\textwidth+8pt]frame.south west);
    }
]
\textbf{Q.80} \quad {\footnotesize\textsc{[MSQ $\cdot$ 2 Mark(s) $\cdot$ No Neg]}} \hfill {\footnotesize\textbf{GATE 2024 CS1} \quad \tiny\texttt{2024\_CS1\_CS\_Q53}}
\vspace{2pt}\hrule height 0.3pt\vspace{4pt}

Consider two set-associative cache memory architectures: WBC, which uses the write back policy, and WTC, which uses the write through policy. Both of them use the LRU (Least Recently Used) block replacement policy. The cache memory is connected to the main memory. Which of the following statements is/are TRUE?

\par\vspace{6pt}\noindent
\begin{tabular}{@{}p{0.98\linewidth}@{}}
$\square$\ \textbf{(A)} A read miss in WBC never evicts a dirty block \\
\addlinespace[3pt]
$\square$\ \textbf{(B)} A read miss in WTC never triggers a write back operation of a cache block to main memory \\
\addlinespace[3pt]
$\square$\ \textbf{(C)} A write hit in WBC can modify the value of the dirty bit of a cache block \\
\addlinespace[3pt]
$\square$\ \textbf{(D)} A write miss in WTC always writes the victim cache block to main memory before loading the missed block to the cache
\end{tabular}

\tcblower

\textbf{Official Key:} \textbf{\large B;C} \hfill \textit{\small [Medium]}
\vspace{2pt}\hrule height 0.3pt\vspace{3pt}
\textbf{Topic:} Computer Organization and Architecture \\
\textbf{Subtopic:} \textit{Memory hierarchy: cache, main memory and secondary storage}
\vspace{2pt}\hrule height 0.3pt\vspace{3pt}
\textbf{Guide / Hint:} \\
{\small In a write-through cache (WTC), main memory is updated simultaneously on every cache write, meaning cache lines are never dirty and no eviction ever triggers a write-back to memory (B is TRUE). In a write-back cache (WBC), when a clean block receives a write hit, its dirty bit is set from 0 to 1 to indicate modification (C is TRUE).}

\end{tcolorbox}
\vspace{6pt}\hrule height 0.35pt\vspace{6pt}
\noindent
\begin{tcolorbox}[
    blanker,
    sidebyside,
    sidebyside align=top,
    lefthand width=0.55\textwidth,
    top=4pt, bottom=4pt, left=0pt, right=0pt,
    middle=16pt,
    underlay={
        \draw[dashed, black!80, line width=0.8pt, dash pattern=on 4pt off 3pt]
            ([xshift=0.55\textwidth+8pt]frame.north west) --
            ([xshift=0.55\textwidth+8pt]frame.south west);
    }
]
\textbf{Q.81} \quad {\footnotesize\textsc{[NAT $\cdot$ 2 Mark(s) $\cdot$ No Neg]}} \hfill {\footnotesize\textbf{GATE 2024 CS1} \quad \tiny\texttt{2024\_CS1\_CS\_Q54}}
\vspace{2pt}\hrule height 0.3pt\vspace{4pt}

Consider a 512 GB hard disk with 32 storage surfaces. There are 4096 sectors per track and each sector holds 1024 bytes of data. The number of cylinders in the hard disk is \_\_\_\_\_\_\_\_\_\. (Answer in integer)

\par\vspace{6pt}\noindent\textbf{Your Answer:} \fbox{\rule{0pt}{1.3ex}\hspace{3cm}}

\tcblower

\textbf{Official Key:} \textbf{\large 4096 to 4096} \hfill \textit{\small [Easy]}
\vspace{2pt}\hrule height 0.3pt\vspace{3pt}
\textbf{Topic:} Computer Organization and Architecture \\
\textbf{Subtopic:} \textit{Memory hierarchy: cache, main memory and secondary storage}
\vspace{2pt}\hrule height 0.3pt\vspace{3pt}
\textbf{Guide / Hint:} \\
{\small Disk capacity $= 512\text{ GB} = 512 \times 2^{30}\text{ bytes} = 2^{39}$ bytes. One cylinder comprises 32 surfaces $\times 4096$ sectors/track $\times 1024$ bytes/sector $= 2^5 \times 2^{12} \times 2^{10} = 2^{27}$ bytes. Number of cylinders $= \frac{2^{39}}{2^{27}} = 2^{12} = 4096$.}

\end{tcolorbox}
\vspace{6pt}\hrule height 0.35pt\vspace{6pt}
\noindent
\begin{tcolorbox}[
    blanker,
    sidebyside,
    sidebyside align=top,
    lefthand width=0.55\textwidth,
    top=4pt, bottom=4pt, left=0pt, right=0pt,
    middle=16pt,
    underlay={
        \draw[dashed, black!80, line width=0.8pt, dash pattern=on 4pt off 3pt]
            ([xshift=0.55\textwidth+8pt]frame.north west) --
            ([xshift=0.55\textwidth+8pt]frame.south west);
    }
]
\textbf{Q.82} \quad {\footnotesize\textsc{[NAT $\cdot$ 2 Mark(s) $\cdot$ No Neg]}} \hfill {\footnotesize\textbf{GATE 2024 CS1} \quad \tiny\texttt{2024\_CS1\_CS\_Q55}}
\vspace{2pt}\hrule height 0.3pt\vspace{4pt}

The baseline execution time of a program on a 2 GHz single core machine is 100 nanoseconds (ns). The code corresponding to 90\% of the execution time can be fully parallelized. The overhead for using an additional core is 10 ns when running on a multicore system. Assume that all cores in the multicore system run their share of the parallelized code for an equal amount of time.
The number of cores that minimize the execution time of the program is \_\_\_\_\_\_\_\_\. (Answer in integer)

\par\vspace{6pt}\noindent\textbf{Your Answer:} \fbox{\rule{0pt}{1.3ex}\hspace{3cm}}

\tcblower

\textbf{Official Key:} \textbf{\large 3 to 3} \hfill \textit{\small [Medium]}
\vspace{2pt}\hrule height 0.3pt\vspace{3pt}
\textbf{Topic:} Computer Organization and Architecture \\
\textbf{Subtopic:} \textit{Instruction pipelining, pipeline hazards}
\vspace{2pt}\hrule height 0.3pt\vspace{3pt}
\textbf{Guide / Hint:} \\
{\small Baseline runtime is $100\text{ ns}$, with $90\text{ ns}$ parallelizable and $10\text{ ns}$ serial. On $p$ cores, the parallel portion takes $\frac{90}{p}\text{ ns}$ with core overhead of $10(p-1)\text{ ns}$. Total execution time $T(p) = 10 + \frac{90}{p} + 10(p-1) = \frac{90}{p} + 10p$. Evaluating: $T(1) = 100\text{ ns}$, $T(2) = 65\text{ ns}$, $T(3) = 60\text{ ns}$, $T(4) = 62.5\text{ ns}$. The time is minimized at $p = 3$ cores.}

\end{tcolorbox}
\vspace{6pt}\hrule height 0.35pt\vspace{6pt}
\noindent
\begin{tcolorbox}[
    blanker,
    sidebyside,
    sidebyside align=top,
    lefthand width=0.55\textwidth,
    top=4pt, bottom=4pt, left=0pt, right=0pt,
    middle=16pt,
    underlay={
        \draw[dashed, black!80, line width=0.8pt, dash pattern=on 4pt off 3pt]
            ([xshift=0.55\textwidth+8pt]frame.north west) --
            ([xshift=0.55\textwidth+8pt]frame.south west);
    }
]
\textbf{Q.83} \quad {\footnotesize\textsc{[NAT $\cdot$ 2 Mark(s) $\cdot$ No Neg]}} \hfill {\footnotesize\textbf{GATE 2024 CS1} \quad \tiny\texttt{2024\_CS1\_CS\_Q56}}
\vspace{2pt}\hrule height 0.3pt\vspace{4pt}

A given program has 25\% load/store instructions. Suppose the ideal CPI (cycles per instruction) without any memory stalls is 2. The program exhibits 2\% miss rate on instruction cache and 8\% miss rate on data cache. The miss penalty is 100 cycles.
The speedup (rounded off to two decimal places) achieved with a perfect cache (i.e., with NO data or instruction cache misses) is \_\_\_\_\_\_\_\_\_\. (Answer in numeric format)

\par\vspace{6pt}\noindent\textbf{Your Answer:} \fbox{\rule{0pt}{1.3ex}\hspace{3cm}}

\tcblower

\textbf{Official Key:} \textbf{\large 3.00 to 3.00} \hfill \textit{\small [Medium]}
\vspace{2pt}\hrule height 0.3pt\vspace{3pt}
\textbf{Topic:} Computer Organization and Architecture \\
\textbf{Subtopic:} \textit{Memory hierarchy: cache, main memory and secondary storage}
\vspace{2pt}\hrule height 0.3pt\vspace{3pt}
\textbf{Guide / Hint:} \\
{\small Total CPI is given by $\text{CPI} = \text{CPI}_{\text{ideal}} + (\text{I-miss rate} \times \text{penalty}) + (\text{Data-instruction fraction} \times \text{D-miss rate} \times \text{penalty}) = 2 + (0.02 \times 100) + (0.25 \times 0.08 \times 100) = 2 + 2 + 2 = 6$. For a perfect cache with no misses, $\text{CPI}_{\text{perfect}} = \text{CPI}_{\text{ideal}} = 2$. Speedup $= \frac{6}{2} = 3.00$.}

\end{tcolorbox}
\vspace{6pt}\hrule height 0.35pt\vspace{6pt}
\noindent
\begin{tcolorbox}[
    blanker,
    sidebyside,
    sidebyside align=top,
    lefthand width=0.55\textwidth,
    top=4pt, bottom=4pt, left=0pt, right=0pt,
    middle=16pt,
    underlay={
        \draw[dashed, black!80, line width=0.8pt, dash pattern=on 4pt off 3pt]
            ([xshift=0.55\textwidth+8pt]frame.north west) --
            ([xshift=0.55\textwidth+8pt]frame.south west);
    }
]
\textbf{Q.84} \quad {\footnotesize\textsc{[NAT $\cdot$ 2 Mark(s) $\cdot$ No Neg]}} \hfill {\footnotesize\textbf{GATE 2024 CS2} \quad \tiny\texttt{2024\_CS2\_CS\_Q57}}
\vspace{2pt}\hrule height 0.3pt\vspace{4pt}

A processor with 16 general purpose registers uses a 32-bit instruction format. The instruction format consists of an opcode field, an addressing mode field, two register operand fields, and a 16-bit scalar field. If 8 addressing modes are to be supported, the maximum number of unique opcodes possible for every addressing mode is \_\_\_\_\_\_\. (Answer in integer)

\par\vspace{6pt}\noindent\textbf{Your Answer:} \fbox{\rule{0pt}{1.3ex}\hspace{3cm}}

\tcblower

\textbf{Official Key:} \textbf{\large 32 to 32} \hfill \textit{\small [Medium]}
\vspace{2pt}\hrule height 0.3pt\vspace{3pt}
\textbf{Topic:} Computer Organization and Architecture \\
\textbf{Subtopic:} \textit{Machine instructions and addressing modes}
\vspace{2pt}\hrule height 0.3pt\vspace{3pt}
\textbf{Guide / Hint:} \\
{\small Instruction length is 32 bits. Register fields: 16 registers need $\log_2 16 = 4$ bits each ($2 \times 4 = 8$ bits). Addressing mode: 8 modes need $\log_2 8 = 3$ bits. Scalar field: 16 bits. Remaining bits for opcode $= 32 - (3 + 8 + 16) = 5$ bits. The maximum number of unique opcodes possible for every addressing mode is $2^5 = 32$.}

\end{tcolorbox}
\vspace{6pt}\hrule height 0.35pt\vspace{6pt}
\noindent
\begin{tcolorbox}[
    blanker,
    sidebyside,
    sidebyside align=top,
    lefthand width=0.55\textwidth,
    top=4pt, bottom=4pt, left=0pt, right=0pt,
    middle=16pt,
    underlay={
        \draw[dashed, black!80, line width=0.8pt, dash pattern=on 4pt off 3pt]
            ([xshift=0.55\textwidth+8pt]frame.north west) --
            ([xshift=0.55\textwidth+8pt]frame.south west);
    }
]
\textbf{Q.85} \quad {\footnotesize\textsc{[NAT $\cdot$ 2 Mark(s) $\cdot$ No Neg]}} \hfill {\footnotesize\textbf{GATE 2024 CS2} \quad \tiny\texttt{2024\_CS2\_CS\_Q58}}
\vspace{2pt}\hrule height 0.3pt\vspace{4pt}

A non-pipelined instruction execution unit operating at 2 GHz takes an average of 6 cycles to execute an instruction of a program P. The unit is then redesigned to operate on a 5-stage pipeline at 2 GHz. Assume that the ideal throughput of the pipelined unit is 1 instruction per cycle. In the execution of program P, 20\% instructions incur an average of 2 cycles stall due to data hazards and 20\% instructions incur an average of 3 cycles stall due to control hazards. The speedup (rounded off to one decimal place) obtained by the pipelined design over the non-pipelined design is \_\_\_\_\_\_\. (Answer in range 2.9 to 3.1)

\par\vspace{6pt}\noindent\textbf{Your Answer:} \fbox{\rule{0pt}{1.3ex}\hspace{3cm}}

\tcblower

\textbf{Official Key:} \textbf{\large 2.9 to 3.1} \hfill \textit{\small [Medium]}
\vspace{2pt}\hrule height 0.3pt\vspace{3pt}
\textbf{Topic:} Computer Organization and Architecture \\
\textbf{Subtopic:} \textit{Instruction pipelining, pipeline hazards}
\vspace{2pt}\hrule height 0.3pt\vspace{3pt}
\textbf{Guide / Hint:} \\
{\small Average cycles per instruction in non-pipelined design: $\text{CPI}_{\text{non-pipe}} = 6$. In pipelined design: $\text{CPI}_{\text{pipe}} = 1 + (0.20 \times 2) + (0.20 \times 3) = 1 + 0.40 + 0.60 = 2.0$. Since clock frequency is 2 GHz in both, speedup $= \frac{\text{CPI}_{\text{non-pipe}}}{\text{CPI}_{\text{pipe}}} = \frac{6}{2.0} = 3.0$.}

\end{tcolorbox}
\vspace{6pt}\hrule height 0.35pt\vspace{6pt}
\noindent
\begin{tcolorbox}[
    blanker,
    sidebyside,
    sidebyside align=top,
    lefthand width=0.55\textwidth,
    top=4pt, bottom=4pt, left=0pt, right=0pt,
    middle=16pt,
    underlay={
        \draw[dashed, black!80, line width=0.8pt, dash pattern=on 4pt off 3pt]
            ([xshift=0.55\textwidth+8pt]frame.north west) --
            ([xshift=0.55\textwidth+8pt]frame.south west);
    }
]
\textbf{Q.86} \quad {\footnotesize\textsc{[NAT $\cdot$ 2 Mark(s) $\cdot$ No Neg]}} \hfill {\footnotesize\textbf{GATE 2024 CS2} \quad \tiny\texttt{2024\_CS2\_CS\_Q61}}
\vspace{2pt}\hrule height 0.3pt\vspace{4pt}

A processor uses a 32-bit instruction format and supports byte-addressable memory access. The ISA of the processor has 150 distinct instructions. The instructions are equally divided into two types, namely R-type and I-type, whose formats are shown below:
\begin{verbatim}
R-type: OPCODE | UNUSED | DST Register | SRC Register1 | SRC Register 2
I-type: OPCODE | DST Register | SRC Register | Immediate value/address
\end{verbatim}
In the OPCODE, 1 bit distinguishes between I-type and R-type and the remaining bits indicate the operation. The processor has 50 architectural registers, and all register fields are of equal size.
Let $X$ be the bits for UNUSED, $Y$ for OPCODE, and $Z$ for immediate value/address. The value of $X + 2Y + Z$ is \_\_\_\_\_\_\. (Answer in integer)

\par\vspace{6pt}\noindent\textbf{Your Answer:} \fbox{\rule{0pt}{1.3ex}\hspace{3cm}}

\tcblower

\textbf{Official Key:} \textbf{\large 34 to 34} \hfill \textit{\small [Hard]}
\vspace{2pt}\hrule height 0.3pt\vspace{3pt}
\textbf{Topic:} Computer Organization and Architecture \\
\textbf{Subtopic:} \textit{Machine instructions and addressing modes}
\vspace{2pt}\hrule height 0.3pt\vspace{3pt}
\textbf{Guide / Hint:} \\
{\small 150 distinct instructions split equally means 75 R-type and 75 I-type. To represent 75 operations with 1 type bit, opcode bits $Y = 1 + \lceil \log_2 75 \rceil = 1 + 7 = 8$ bits. 50 registers require $\lceil \log_2 50 \rceil = 6$ bits each. R-type has 3 register fields ($3 \times 6 = 18$ bits), so $X = 32 - (8 + 18) = 6$ bits. I-type has 2 register fields ($2 \times 6 = 12$ bits), so $Z = 32 - (8 + 12) = 12$ bits. Thus, $X + 2Y + Z = 6 + 2(8) + 12 = 34$.}

\end{tcolorbox}
\vspace{6pt}\hrule height 0.35pt\vspace{6pt}
\noindent
\begin{tcolorbox}[
    blanker,
    sidebyside,
    sidebyside align=top,
    lefthand width=0.55\textwidth,
    top=4pt, bottom=4pt, left=0pt, right=0pt,
    middle=16pt,
    underlay={
        \draw[dashed, black!80, line width=0.8pt, dash pattern=on 4pt off 3pt]
            ([xshift=0.55\textwidth+8pt]frame.north west) --
            ([xshift=0.55\textwidth+8pt]frame.south west);
    }
]
\textbf{Q.87} \quad {\footnotesize\textsc{[MCQ $\cdot$ 2 Mark(s) $\cdot$ Neg: 0.66]}} \hfill {\footnotesize\textbf{GATE 2025 CS1} \quad \tiny\texttt{2025\_CS1\_CS\_Q36}}
\vspace{2pt}\hrule height 0.3pt\vspace{4pt}

Consider a memory system with 1M bytes of main memory and 16K bytes of cache memory. Assume that the processor generates 20-bit memory address, and the cache block size is 16 bytes. If the cache uses direct mapping, how many bits will be required to store all the tag values? [Assume memory is byte addressable, $1\text{K}=2^{10}, 1\text{M}=2^{20}$.]

\par\vspace{6pt}\noindent
\begin{tabular}{@{}p{0.48\linewidth}@{\hspace{0.04\linewidth}}p{0.48\linewidth}@{}}
$\square$\ \textbf{(A)} $6 \times 2^{10}$ & $\square$\ \textbf{(B)} $8 \times 2^{10}$ \\
\addlinespace[3pt]
$\square$\ \textbf{(C)} $2^{12}$ & $\square$\ \textbf{(D)} $2^{14}$
\end{tabular}

\tcblower

\textbf{Official Key:} \textbf{\large A} \hfill \textit{\small [Medium]}
\vspace{2pt}\hrule height 0.3pt\vspace{3pt}
\textbf{Topic:} Computer Organization and Architecture \\
\textbf{Subtopic:} \textit{Memory hierarchy: cache, main memory and secondary storage}
\vspace{2pt}\hrule height 0.3pt\vspace{3pt}
\textbf{Guide / Hint:} \\
{\small With block size $16 = 2^4$ bytes, block offset is 4 bits. Cache size is $16\text{ KB} = 2^{14}$ bytes, so number of cache lines is $2^{14}/2^4 = 2^{10} = 1024$, requiring 10 line index bits. The tag bits per line are $20 - (10 + 4) = 6$ bits. For all $2^{10}$ lines, total tag bits required are $6 \times 2^{10}$.}

\end{tcolorbox}
\vspace{6pt}\hrule height 0.35pt\vspace{6pt}
\noindent
\begin{tcolorbox}[
    blanker,
    sidebyside,
    sidebyside align=top,
    lefthand width=0.55\textwidth,
    top=4pt, bottom=4pt, left=0pt, right=0pt,
    middle=16pt,
    underlay={
        \draw[dashed, black!80, line width=0.8pt, dash pattern=on 4pt off 3pt]
            ([xshift=0.55\textwidth+8pt]frame.north west) --
            ([xshift=0.55\textwidth+8pt]frame.south west);
    }
]
\textbf{Q.88} \quad {\footnotesize\textsc{[MCQ $\cdot$ 2 Mark(s) $\cdot$ Neg: 0.66]}} \hfill {\footnotesize\textbf{GATE 2025 CS1} \quad \tiny\texttt{2025\_CS1\_CS\_Q37}}
\vspace{2pt}\hrule height 0.3pt\vspace{4pt}

A processor has 64 general-purpose registers and 50 distinct instruction types. An instruction is encoded in 32-bits. What is the maximum number of bits that can be used to store the immediate operand for the given instruction?
\begin{verbatim}
ADD R1, #25 // R1 = R1 + 25
\end{verbatim}

\par\vspace{6pt}\noindent
\begin{tabular}{@{}p{0.48\linewidth}@{\hspace{0.04\linewidth}}p{0.48\linewidth}@{}}
$\square$\ \textbf{(A)} 16 & $\square$\ \textbf{(B)} 20 \\
\addlinespace[3pt]
$\square$\ \textbf{(C)} 22 & $\square$\ \textbf{(D)} 24
\end{tabular}

\tcblower

\textbf{Official Key:} \textbf{\large B} \hfill \textit{\small [Medium]}
\vspace{2pt}\hrule height 0.3pt\vspace{3pt}
\textbf{Topic:} Computer Organization and Architecture \\
\textbf{Subtopic:} \textit{Machine instructions and addressing modes}
\vspace{2pt}\hrule height 0.3pt\vspace{3pt}
\textbf{Guide / Hint:} \\
{\small To represent 50 distinct instruction opcodes, $\lceil \log_2 50 \rceil = 6$ opcode bits are required. Specifying one destination register among 64 general-purpose registers requires $\lceil \log_2 64 \rceil = 6$ bits. In a 32-bit instruction word, the remaining bits for the immediate field are $32 - (6 + 6) = 20$ bits.}

\end{tcolorbox}
\vspace{6pt}\hrule height 0.35pt\vspace{6pt}
\noindent
\begin{tcolorbox}[
    blanker,
    sidebyside,
    sidebyside align=top,
    lefthand width=0.55\textwidth,
    top=4pt, bottom=4pt, left=0pt, right=0pt,
    middle=16pt,
    underlay={
        \draw[dashed, black!80, line width=0.8pt, dash pattern=on 4pt off 3pt]
            ([xshift=0.55\textwidth+8pt]frame.north west) --
            ([xshift=0.55\textwidth+8pt]frame.south west);
    }
]
\textbf{Q.89} \quad {\footnotesize\textsc{[NAT $\cdot$ 2 Mark(s) $\cdot$ No Neg]}} \hfill {\footnotesize\textbf{GATE 2025 CS1} \quad \tiny\texttt{2025\_CS1\_CS\_Q53}}
\vspace{2pt}\hrule height 0.3pt\vspace{4pt}

A computer has a memory hierarchy consisting of two-level cache (L1 and L2) and a main memory. If the processor needs to access data from memory, it first looks into L1 cache ($h_1 = 0.95$, access time $t_1 = 10\text{ ns}$). If not found in L1, it accesses L2 cache ($h_2 = 0.85$, access time including L1 miss penalty $t_2 = 20\text{ ns}$). If it misses in L2, it accesses main memory (access time including L1 and L2 miss penalties $t_{MM} = 200\text{ ns}$). The average memory access time in nanoseconds (ns) is \_\_\_\_\_\_\. (rounded off to two decimal places)

\par\vspace{6pt}\noindent\textbf{Your Answer:} \fbox{\rule{0pt}{1.3ex}\hspace{3cm}}

\tcblower

\textbf{Official Key:} \textbf{\large 11.83 to 11.87} \hfill \textit{\small [Medium]}
\vspace{2pt}\hrule height 0.3pt\vspace{3pt}
\textbf{Topic:} Computer Organization and Architecture \\
\textbf{Subtopic:} \textit{Memory hierarchy: cache, main memory and secondary storage}
\vspace{2pt}\hrule height 0.3pt\vspace{3pt}
\textbf{Guide / Hint:} \\
{\small Using simultaneous access with cumulative times: $\text{AMAT} = h_1 t_1 + (1 - h_1) h_2 t_2 + (1 - h_1)(1 - h_2) t_{MM} = 0.95(10) + 0.05(0.85)(20) + 0.05(0.15)(200) = 9.5 + 0.85 + 1.5 = 11.85\text{ ns}$.}

\end{tcolorbox}
\vspace{6pt}\hrule height 0.35pt\vspace{6pt}
\noindent
\begin{tcolorbox}[
    blanker,
    sidebyside,
    sidebyside align=top,
    lefthand width=0.55\textwidth,
    top=4pt, bottom=4pt, left=0pt, right=0pt,
    middle=16pt,
    underlay={
        \draw[dashed, black!80, line width=0.8pt, dash pattern=on 4pt off 3pt]
            ([xshift=0.55\textwidth+8pt]frame.north west) --
            ([xshift=0.55\textwidth+8pt]frame.south west);
    }
]
\textbf{Q.90} \quad {\footnotesize\textsc{[MCQ $\cdot$ 2 Mark(s) $\cdot$ Neg: 0.66]}} \hfill {\footnotesize\textbf{GATE 2025 CS2} \quad \tiny\texttt{2025\_CS2\_CS\_Q39}}
\vspace{2pt}\hrule height 0.3pt\vspace{4pt}

For a direct-mapped cache, 4 bits are used for the tag field and 12 bits are used to index into a cache block. What are the cache data size and tag directory size respectively?

\par\vspace{6pt}\noindent
\begin{tabular}{@{}p{0.48\linewidth}@{\hspace{0.04\linewidth}}p{0.48\linewidth}@{}}
$\square$\ \textbf{(A)} 64 KB and 4 KB & $\square$\ \textbf{(B)} 128 KB and 16 KB \\
\addlinespace[3pt]
$\square$\ \textbf{(C)} 64 KB and 8 KB & $\square$\ \textbf{(D)} 128 KB and 6 KB
\end{tabular}

\tcblower

\textbf{Official Key:} \textbf{\large A} \hfill \textit{\small [Medium]}
\vspace{2pt}\hrule height 0.3pt\vspace{3pt}
\textbf{Topic:} Computer Organization and Architecture \\
\textbf{Subtopic:} \textit{Memory hierarchy: cache memory mapping}
\vspace{2pt}\hrule height 0.3pt\vspace{3pt}
\textbf{Guide / Hint:} \\
{\small Number of cache lines is $2^{12} = 4096$. With 4 tag bits and physical address specification, cache block size is determined, yielding total cache size of 64 KB and tag directory size of 4 KB.}

\end{tcolorbox}
\vspace{6pt}\hrule height 0.35pt\vspace{6pt}
\noindent
\begin{tcolorbox}[
    blanker,
    sidebyside,
    sidebyside align=top,
    lefthand width=0.55\textwidth,
    top=4pt, bottom=4pt, left=0pt, right=0pt,
    middle=16pt,
    underlay={
        \draw[dashed, black!80, line width=0.8pt, dash pattern=on 4pt off 3pt]
            ([xshift=0.55\textwidth+8pt]frame.north west) --
            ([xshift=0.55\textwidth+8pt]frame.south west);
    }
]
\textbf{Q.91} \quad {\footnotesize\textsc{[NAT $\cdot$ 2 Mark(s) $\cdot$ No Neg]}} \hfill {\footnotesize\textbf{GATE 2025 CS2} \quad \tiny\texttt{2025\_CS2\_CS\_Q55}}
\vspace{2pt}\hrule height 0.3pt\vspace{4pt}

Given a computing system with two levels of cache (L1 and L2) and a main memory. The first level (L1) cache access time is 1 ns with 0.8 hit rate, L2 access time is 4 ns with 0.9 hit rate, and main memory access time is 50 ns. The average memory access time in nanoseconds is \_\_\_\_\_\_\_\_. (Answer in decimal)

\par\vspace{6pt}\noindent\textbf{Your Answer:} \fbox{\rule{0pt}{1.3ex}\hspace{3cm}}

\tcblower

\textbf{Official Key:} \textbf{\large 4.0 to 4.0} \hfill \textit{\small [Medium]}
\vspace{2pt}\hrule height 0.3pt\vspace{3pt}
\textbf{Topic:} Computer Organization and Architecture \\
\textbf{Subtopic:} \textit{Memory hierarchy: multi-level cache access time}
\vspace{2pt}\hrule height 0.3pt\vspace{3pt}
\textbf{Guide / Hint:} \\
{\small Average Memory Access Time (AMAT) formula: $T_{avg} = T_{L1} + (1 - h_1) \cdot [T_{L2} + (1 - h_2) \cdot T_{MM}]$. Substituting the given hit rates and access latencies gives $T_{avg} = 4.0$ ns.}

\end{tcolorbox}
\vspace{6pt}\hrule height 0.35pt\vspace{6pt}
\noindent
\begin{tcolorbox}[
    blanker,
    sidebyside,
    sidebyside align=top,
    lefthand width=0.55\textwidth,
    top=4pt, bottom=4pt, left=0pt, right=0pt,
    middle=16pt,
    underlay={
        \draw[dashed, black!80, line width=0.8pt, dash pattern=on 4pt off 3pt]
            ([xshift=0.55\textwidth+8pt]frame.north west) --
            ([xshift=0.55\textwidth+8pt]frame.south west);
    }
]
\textbf{Q.92} \quad {\footnotesize\textsc{[NAT $\cdot$ 2 Mark(s) $\cdot$ No Neg]}} \hfill {\footnotesize\textbf{GATE 2025 CS2} \quad \tiny\texttt{2025\_CS2\_CS\_Q56}}
\vspace{2pt}\hrule height 0.3pt\vspace{4pt}

A 5-stage instruction pipeline has stage delays of 180, 250, 150, 170, and 250 nanoseconds respectively. What is the total execution time (in microseconds)?

\par\vspace{6pt}\noindent\textbf{Your Answer:} \fbox{\rule{0pt}{1.3ex}\hspace{3cm}}

\tcblower

\textbf{Official Key:} \textbf{\large 260.20 to 261.20} \hfill \textit{\small [Hard]}
\vspace{2pt}\hrule height 0.3pt\vspace{3pt}
\textbf{Topic:} Computer Organization and Architecture \\
\textbf{Subtopic:} \textit{Instruction pipelining: clock cycle time and pipeline execution}
\vspace{2pt}\hrule height 0.3pt\vspace{3pt}
\textbf{Guide / Hint:} \\
{\small Clock period $T_{clk} = \max(\text{stage delays}) + \text{register delay}$. Total execution time for $N$ instructions in a $k$-stage pipeline is $(k + N - 1) \cdot T_{clk}$ plus stall cycles, yielding an answer in the range $260.20$ to $261.20$.}

\end{tcolorbox}
\vspace{6pt}\hrule height 0.35pt\vspace{6pt}
\noindent
\begin{tcolorbox}[
    blanker,
    sidebyside,
    sidebyside align=top,
    lefthand width=0.55\textwidth,
    top=4pt, bottom=4pt, left=0pt, right=0pt,
    middle=16pt,
    underlay={
        \draw[dashed, black!80, line width=0.8pt, dash pattern=on 4pt off 3pt]
            ([xshift=0.55\textwidth+8pt]frame.north west) --
            ([xshift=0.55\textwidth+8pt]frame.south west);
    }
]
\textbf{Q.93} \quad {\footnotesize\textsc{[NAT $\cdot$ 2 Mark(s) $\cdot$ No Neg]}} \hfill {\footnotesize\textbf{GATE 2025 CS2} \quad \tiny\texttt{2025\_CS2\_CS\_Q61}}
\vspace{2pt}\hrule height 0.3pt\vspace{4pt}

An application executes $6.4 \times 10^8$ instructions in 6.3 seconds across four instruction classes. The processor clock frequency in GHz is \_\_\_\_\_\_\_\_. (Answer in decimal)

\par\vspace{6pt}\noindent\textbf{Your Answer:} \fbox{\rule{0pt}{1.3ex}\hspace{3cm}}

\tcblower

\textbf{Official Key:} \textbf{\large 3.0 to 3.0} \hfill \textit{\small [Medium]}
\vspace{2pt}\hrule height 0.3pt\vspace{3pt}
\textbf{Topic:} Computer Organization and Architecture \\
\textbf{Subtopic:} \textit{Machine instructions: CPI and clock frequency}
\vspace{2pt}\hrule height 0.3pt\vspace{3pt}
\textbf{Guide / Hint:} \\
{\small Average CPI is $\sum (f_i \times \text{CPI}_i)$. Total clock cycles = $\text{Instruction Count} \times \text{CPI}_{avg}$. Clock frequency = $\frac{\text{Total Cycles}}{\text{Execution Time}} = \frac{6.4 \times 10^8 \times \text{CPI}_{avg}}{6.3\,\text{s}} = 3.0$ GHz.}

\end{tcolorbox}
\vspace{6pt}\hrule height 0.35pt\vspace{6pt}
\vspace{8pt}
\subsection*{Programming and Data Structures}
\vspace{-4pt}
\noindent
\begin{tcolorbox}[
    blanker,
    sidebyside,
    sidebyside align=top,
    lefthand width=0.55\textwidth,
    top=4pt, bottom=4pt, left=0pt, right=0pt,
    middle=16pt,
    underlay={
        \draw[dashed, black!80, line width=0.8pt, dash pattern=on 4pt off 3pt]
            ([xshift=0.55\textwidth+8pt]frame.north west) --
            ([xshift=0.55\textwidth+8pt]frame.south west);
    }
]
\textbf{Q.94} \quad {\footnotesize\textsc{[MCQ $\cdot$ 2 Mark(s) $\cdot$ Neg: 0.66]}} \hfill {\footnotesize\textbf{GATE 2018 CS} \quad \tiny\texttt{2018\_CS\_GA\_Q29}}
\vspace{2pt}\hrule height 0.3pt\vspace{4pt}

Consider the following C program:
\begin{lstlisting}[language=C]
#include <stdio.h>
\par\medskip
void fun1(char *s1, char *s2) {
    char *tmp;
    tmp = s1;
    s1 = s2;
    s2 = tmp;
}
\par\medskip
void fun2(char **s1, char **s2) {
    char *tmp;
    tmp = *s1;
    *s1 = *s2;
    *s2 = tmp;
}
\par\medskip
int main() {
    char *str1 = "Hi", *str2 = "Bye";
    fun1(str1, str2);
    printf("%s %s ", str1, str2);
    fun2(&str1, &str2);
    printf("%s %s", str1, str2);
    return 0;
}
\end{lstlisting}
The output of the program is

\par\vspace{6pt}\noindent
\begin{tabular}{@{}p{0.48\linewidth}@{\hspace{0.04\linewidth}}p{0.48\linewidth}@{}}
$\square$\ \textbf{(A)} Hi Bye Bye Hi & $\square$\ \textbf{(B)} Hi Bye Hi Bye \\
\addlinespace[3pt]
$\square$\ \textbf{(C)} Bye Hi Bye Hi & $\square$\ \textbf{(D)} Bye Hi Hi Bye
\end{tabular}

\tcblower

\textbf{Official Key:} \textbf{\large A} \hfill \textit{\small [Medium]}
\vspace{2pt}\hrule height 0.3pt\vspace{3pt}
\textbf{Topic:} Programming and Data Structures \\
\textbf{Subtopic:} \textit{Programming in C}
\vspace{2pt}\hrule height 0.3pt\vspace{3pt}
\textbf{Guide / Hint:} \\
{\small In `fun1`, pointers `s1` and `s2` are passed by value, so the pointer swap only alters local variables and has no effect on `str1` and `str2` in `main` (printing "Hi Bye"). In `fun2`, the addresses of the pointers (`char**`) are passed, so dereferencing swaps the actual pointer values `str1` and `str2` (printing "Bye Hi").}

\end{tcolorbox}
\vspace{6pt}\hrule height 0.35pt\vspace{6pt}
\noindent
\begin{tcolorbox}[
    blanker,
    sidebyside,
    sidebyside align=top,
    lefthand width=0.55\textwidth,
    top=4pt, bottom=4pt, left=0pt, right=0pt,
    middle=16pt,
    underlay={
        \draw[dashed, black!80, line width=0.8pt, dash pattern=on 4pt off 3pt]
            ([xshift=0.55\textwidth+8pt]frame.north west) --
            ([xshift=0.55\textwidth+8pt]frame.south west);
    }
]
\textbf{Q.95} \quad {\footnotesize\textsc{[NAT $\cdot$ 2 Mark(s) $\cdot$ No Neg]}} \hfill {\footnotesize\textbf{GATE 2018 CS} \quad \tiny\texttt{2018\_CS\_GA\_Q45}}
\vspace{2pt}\hrule height 0.3pt\vspace{4pt}

Consider the following pseudo-code function where $x$ and $y$ are positive integers:
\begin{lstlisting}[language=C]
Count(x, y) {
    if (y != 1) {
        if (x != 1) {
            print("*");
            Count(x / 2, y);
        } else {
            Count(1024, y - 1);
        }
    }
}
\end{lstlisting}
The total number of times `print("*")` is executed for `Count(1024, 1024)` is \_\_\_\_\_\_\_\_.

\par\vspace{6pt}\noindent\textbf{Your Answer:} \fbox{\rule{0pt}{1.3ex}\hspace{3cm}}

\tcblower

\textbf{Official Key:} \textbf{\large 10230 to 10230} \hfill \textit{\small [Medium]}
\vspace{2pt}\hrule height 0.3pt\vspace{3pt}
\textbf{Topic:} Programming and Data Structures \\
\textbf{Subtopic:} \textit{Recursion}
\vspace{2pt}\hrule height 0.3pt\vspace{3pt}
\textbf{Guide / Hint:} \\
{\small In `Count(x, y)`, outer recursion on $y$ runs for $y = 1024, 1023, \dots, 2$ ($1023$ values). For each $y$, the inner branch halves $x$ until $x = 1$: starting from $x = 1024$, this loop prints 10 times (since $\log_2(1024) = 10$). Total prints $= 1023 \times 10 = 10230$.}

\end{tcolorbox}
\vspace{6pt}\hrule height 0.35pt\vspace{6pt}
\noindent
\begin{tcolorbox}[
    blanker,
    sidebyside,
    sidebyside align=top,
    lefthand width=0.55\textwidth,
    top=4pt, bottom=4pt, left=0pt, right=0pt,
    middle=16pt,
    underlay={
        \draw[dashed, black!80, line width=0.8pt, dash pattern=on 4pt off 3pt]
            ([xshift=0.55\textwidth+8pt]frame.north west) --
            ([xshift=0.55\textwidth+8pt]frame.south west);
    }
]
\textbf{Q.96} \quad {\footnotesize\textsc{[NAT $\cdot$ 2 Mark(s) $\cdot$ No Neg]}} \hfill {\footnotesize\textbf{GATE 2018 CS} \quad \tiny\texttt{2018\_CS\_GA\_Q46}}
\vspace{2pt}\hrule height 0.3pt\vspace{4pt}

The number of possible min-heaps containing each value from $\{1, 2, 3, 4, 5, 6, 7\}$ exactly once is \_\_\_\_\_\_\_\_.

\par\vspace{6pt}\noindent\textbf{Your Answer:} \fbox{\rule{0pt}{1.3ex}\hspace{3cm}}

\tcblower

\textbf{Official Key:} \textbf{\large 80 to 80} \hfill \textit{\small [Medium]}
\vspace{2pt}\hrule height 0.3pt\vspace{3pt}
\textbf{Topic:} Programming and Data Structures \\
\textbf{Subtopic:} \textit{Trees, binary search trees, binary heaps}
\vspace{2pt}\hrule height 0.3pt\vspace{3pt}
\textbf{Guide / Hint:} \\
{\small In a min-heap of 7 elements, the root must contain the minimum element (1). The left subtree has size 3 and the right subtree has size 3. Choosing 3 elements for the left subtree from the remaining 6 elements can be done in $\binom{6}{3} = 20$ ways. A 3-element min-heap has $\binom{2}{1} = 2$ valid configurations. Hence, total valid min-heaps $= 20 \times 2 \times 2 = 80$.}

\end{tcolorbox}
\vspace{6pt}\hrule height 0.35pt\vspace{6pt}
\noindent
\begin{tcolorbox}[
    blanker,
    sidebyside,
    sidebyside align=top,
    lefthand width=0.55\textwidth,
    top=4pt, bottom=4pt, left=0pt, right=0pt,
    middle=16pt,
    underlay={
        \draw[dashed, black!80, line width=0.8pt, dash pattern=on 4pt off 3pt]
            ([xshift=0.55\textwidth+8pt]frame.north west) --
            ([xshift=0.55\textwidth+8pt]frame.south west);
    }
]
\textbf{Q.97} \quad {\footnotesize\textsc{[MCQ $\cdot$ 2 Mark(s) $\cdot$ Neg: 0.66]}} \hfill {\footnotesize\textbf{GATE 2022 CS} \quad \tiny\texttt{2022\_CS\_CS\_Q43}}
\vspace{2pt}\hrule height 0.3pt\vspace{4pt}

What is printed by the following ANSI C program?
\begin{lstlisting}[language=C]
#include <stdio.h>
int main(int argc, char *argv[]) {
    int a[3][3][3] = {
        {1, 2, 3, 4, 5, 6, 7, 8, 9},
        {10, 11, 12, 13, 14, 15, 16, 17, 18},
        {19, 20, 21, 22, 23, 24, 25, 26, 27}
    };
    int i = 0, j = 0, k = 0;
    for(i = 0; i < 3; i++) {
        for(k = 0; k < 3; k++)
            printf("%d ", a[i][j][k]);
        printf("\n");
    }
    return 0;
}
\end{lstlisting}

\par\vspace{6pt}\noindent
\begin{tabular}{@{}p{0.48\linewidth}@{\hspace{0.04\linewidth}}p{0.48\linewidth}@{}}
$\square$\ \textbf{(A)} 1 2 3 \\ 10 11 12 \\ 19 20 21 & $\square$\ \textbf{(B)} 1 4 7 \\ 10 13 16 \\ 19 22 25 \\
\addlinespace[3pt]
$\square$\ \textbf{(C)} 1 2 3 \\ 4 5 6 \\ 7 8 9 & $\square$\ \textbf{(D)} 1 2 3 \\ 13 14 15 \\ 25 26 27
\end{tabular}

\tcblower

\textbf{Official Key:} \textbf{\large A} \hfill \textit{\small [Easy]}
\vspace{2pt}\hrule height 0.3pt\vspace{3pt}
\textbf{Topic:} Programming and Data Structures \\
\textbf{Subtopic:} \textit{Programming in C}
\vspace{2pt}\hrule height 0.3pt\vspace{3pt}
\textbf{Guide / Hint:} \\
{\small In the nested loops, index $j = 0$ is fixed throughout. For $i = 0$, $k$ iterates from 0 to 2, printing `a[0][0][0], a[0][0][1], a[0][0][2]` ($1\ 2\ 3$). For $i = 1$, it prints `a[1][0][0], a[1][0][1], a[1][0][2]` ($10\ 11\ 12$). For $i = 2$, it prints `a[2][0][0], a[2][0][1], a[2][0][2]` ($19\ 20\ 21$).}

\end{tcolorbox}
\vspace{6pt}\hrule height 0.35pt\vspace{6pt}
\noindent
\begin{tcolorbox}[
    blanker,
    sidebyside,
    sidebyside align=top,
    lefthand width=0.55\textwidth,
    top=4pt, bottom=4pt, left=0pt, right=0pt,
    middle=16pt,
    underlay={
        \draw[dashed, black!80, line width=0.8pt, dash pattern=on 4pt off 3pt]
            ([xshift=0.55\textwidth+8pt]frame.north west) --
            ([xshift=0.55\textwidth+8pt]frame.south west);
    }
]
\textbf{Q.98} \quad {\footnotesize\textsc{[MCQ $\cdot$ 2 Mark(s) $\cdot$ Neg: 0.66]}} \hfill {\footnotesize\textbf{GATE 2022 CS} \quad \tiny\texttt{2022\_CS\_CS\_Q44}}
\vspace{2pt}\hrule height 0.3pt\vspace{4pt}

What is printed by the following ANSI C program?
\begin{lstlisting}[language=C]
#include <stdio.h>
int main(int argc, char *argv[]) {
    char a = 'P';
    char b = 'x';
    char c = (a & b) + '*';
    char d = (a | b) - '-';
    char e = (a ^ b) + '+';
    printf("%c %c %c\n", c, d, e);
    return 0;
}
\end{lstlisting}

\par\vspace{6pt}\noindent
\begin{tabular}{@{}p{0.48\linewidth}@{\hspace{0.04\linewidth}}p{0.48\linewidth}@{}}
$\square$\ \textbf{(A)} z K S & $\square$\ \textbf{(B)} 122 75 83 \\
\addlinespace[3pt]
$\square$\ \textbf{(C)} * - + & $\square$\ \textbf{(D)} P x +
\end{tabular}

\tcblower

\textbf{Official Key:} \textbf{\large A} \hfill \textit{\small [Medium]}
\vspace{2pt}\hrule height 0.3pt\vspace{3pt}
\textbf{Topic:} Programming and Data Structures \\
\textbf{Subtopic:} \textit{Programming in C}
\vspace{2pt}\hrule height 0.3pt\vspace{3pt}
\textbf{Guide / Hint:} \\
{\small Character values: `a = 'P'` (80), `b = 'x'` (120). $c = (80 \& 120) + '*' = 80 + 42 = 122$, which is ASCII for 'z'. $d = (80 \mid 120) - '-' = 120 - 45 = 75$, which is ASCII for 'K'. $e = (80 \oplus 120) + '+' = 40 + 43 = 83$, which is ASCII for 'S'. The printed characters are `z K S`.}

\end{tcolorbox}
\vspace{6pt}\hrule height 0.35pt\vspace{6pt}
\noindent
\begin{tcolorbox}[
    blanker,
    sidebyside,
    sidebyside align=top,
    lefthand width=0.55\textwidth,
    top=4pt, bottom=4pt, left=0pt, right=0pt,
    middle=16pt,
    underlay={
        \draw[dashed, black!80, line width=0.8pt, dash pattern=on 4pt off 3pt]
            ([xshift=0.55\textwidth+8pt]frame.north west) --
            ([xshift=0.55\textwidth+8pt]frame.south west);
    }
]
\textbf{Q.99} \quad {\footnotesize\textsc{[NAT $\cdot$ 2 Mark(s) $\cdot$ No Neg]}} \hfill {\footnotesize\textbf{GATE 2022 CS} \quad \tiny\texttt{2022\_CS\_CS\_Q62}}
\vspace{2pt}\hrule height 0.3pt\vspace{4pt}

Consider the queues $Q_1$ containing four elements and $Q_2$ containing none. The only operations allowed on these two queues are `Enqueue(Q, element)` and `Dequeue(Q)`. The minimum number of `Enqueue` operations on $Q_1$ required to place the elements of $Q_1$ in $Q_2$ in reverse order without using any additional storage is \_\_\_\_\_\_\_\_\_\_\_\. (Answer in integer)

\par\vspace{6pt}\noindent\textbf{Your Answer:} \fbox{\rule{0pt}{1.3ex}\hspace{3cm}}

\tcblower

\textbf{Official Key:} \textbf{\large 0 to 0} \hfill \textit{\small [Medium]}
\vspace{2pt}\hrule height 0.3pt\vspace{3pt}
\textbf{Topic:} Programming and Data Structures \\
\textbf{Subtopic:} \textit{Arrays, stacks, queues, linked lists}
\vspace{2pt}\hrule height 0.3pt\vspace{3pt}
\textbf{Guide / Hint:} \\
{\small To place the 4 elements of $Q_1$ into $Q_2$ in reverse order, elements are dequeued one by one from $Q_1$ and placed into $Q_2$. Internal order manipulation can be performed entirely by cycling elements within $Q_2$ via dequeue/enqueue on $Q_2$, requiring 0 enqueue operations on $Q_1$.}

\end{tcolorbox}
\vspace{6pt}\hrule height 0.35pt\vspace{6pt}
\noindent
\begin{tcolorbox}[
    blanker,
    sidebyside,
    sidebyside align=top,
    lefthand width=0.55\textwidth,
    top=4pt, bottom=4pt, left=0pt, right=0pt,
    middle=16pt,
    underlay={
        \draw[dashed, black!80, line width=0.8pt, dash pattern=on 4pt off 3pt]
            ([xshift=0.55\textwidth+8pt]frame.north west) --
            ([xshift=0.55\textwidth+8pt]frame.south west);
    }
]
\textbf{Q.100} \quad {\footnotesize\textsc{[MCQ $\cdot$ 2 Mark(s) $\cdot$ Neg: 0.66]}} \hfill {\footnotesize\textbf{GATE 2023 CS} \quad \tiny\texttt{2023\_CS\_CS\_Q46}}
\vspace{2pt}\hrule height 0.3pt\vspace{4pt}

Let $A$ be a priority queue for maintaining a set of elements. Suppose $A$ is implemented using a max-heap data structure. The operation $\text{Extract-Max}(A)$ extracts and deletes the maximum element from $A$. The operation $\text{Insert}(A, \text{key})$ inserts a new element $\text{key}$ in $A$. The properties of a max-heap are preserved at the end of each of these operations.
When $A$ contains $n$ elements, which one of the following statements about the worst case running time of these two operations is TRUE?

\par\vspace{6pt}\noindent
\begin{tabular}{@{}p{0.98\linewidth}@{}}
$\square$\ \textbf{(A)} Both `Extract-Max(A)` and `Insert(A,key)` run in $\mathcal{O}(1)$. \\
\addlinespace[3pt]
$\square$\ \textbf{(B)} Both `Extract-Max(A)` and `Insert(A,key)` run in $\mathcal{O}(\log n)$. \\
\addlinespace[3pt]
$\square$\ \textbf{(C)} `Extract-Max(A)` runs in $\mathcal{O}(1)$ whereas `Insert(A,key)` runs in $\mathcal{O}(n)$. \\
\addlinespace[3pt]
$\square$\ \textbf{(D)} `Extract-Max(A)` runs in $\mathcal{O}(1)$ whereas `Insert(A,key)` runs in $\mathcal{O}(\log n)$.
\end{tabular}

\tcblower

\textbf{Official Key:} \textbf{\large B} \hfill \textit{\small [Easy]}
\vspace{2pt}\hrule height 0.3pt\vspace{3pt}
\textbf{Topic:} Programming and Data Structures \\
\textbf{Subtopic:} \textit{Trees, binary search trees, binary heaps}
\vspace{2pt}\hrule height 0.3pt\vspace{3pt}
\textbf{Guide / Hint:} \\
{\small In a binary max-heap of $n$ elements, `Extract-Max` replaces the root with the last element and performs `maxHeapify` downward along the height $\lfloor \log_2 n \rfloor$, running in $\mathcal{O}(\log n)$ worst-case time. `Insert` places the new key at the next leaf position and sifts upward along the height, also running in $\mathcal{O}(\log n)$ worst-case time.}

\end{tcolorbox}
\vspace{6pt}\hrule height 0.35pt\vspace{6pt}
\noindent
\begin{tcolorbox}[
    blanker,
    sidebyside,
    sidebyside align=top,
    lefthand width=0.55\textwidth,
    top=4pt, bottom=4pt, left=0pt, right=0pt,
    middle=16pt,
    underlay={
        \draw[dashed, black!80, line width=0.8pt, dash pattern=on 4pt off 3pt]
            ([xshift=0.55\textwidth+8pt]frame.north west) --
            ([xshift=0.55\textwidth+8pt]frame.south west);
    }
]
\textbf{Q.101} \quad {\footnotesize\textsc{[MCQ $\cdot$ 2 Mark(s) $\cdot$ Neg: 0.66]}} \hfill {\footnotesize\textbf{GATE 2023 CS} \quad \tiny\texttt{2023\_CS\_CS\_Q47}}
\vspace{2pt}\hrule height 0.3pt\vspace{4pt}

Consider the C function `foo` and the binary tree shown:
\begin{lstlisting}[language=C]
typedef struct node {
    int val;
    struct node *left, *right;
} node;
int foo(node *p) {
    int retval;
    if (p == NULL)
        return 0;
    else {
        retval = p->val + foo(p->left) + foo(p->right);
        printf("%d ", retval);
        return retval;
    }
}
\end{lstlisting}
In the binary tree, root is 10; its left child is 5 with children 3 and 8; its right child is 11 with right child 13.
When `foo` is called with a pointer to the root node of the given binary tree, what will it print?

\par\vspace{6pt}\noindent
\begin{tabular}{@{}p{0.48\linewidth}@{\hspace{0.04\linewidth}}p{0.48\linewidth}@{}}
$\square$\ \textbf{(A)} 3 8 5 13 11 10 & $\square$\ \textbf{(B)} 3 5 8 10 11 13 \\
\addlinespace[3pt]
$\square$\ \textbf{(C)} 3 8 16 13 24 50 & $\square$\ \textbf{(D)} 3 16 8 50 24 13
\end{tabular}

\tcblower

\textbf{Official Key:} \textbf{\large C} \hfill \textit{\small [Medium]}
\vspace{2pt}\hrule height 0.3pt\vspace{3pt}
\textbf{Topic:} Programming and Data Structures \\
\textbf{Subtopic:} \textit{Recursion}
\vspace{2pt}\hrule height 0.3pt\vspace{3pt}
\textbf{Guide / Hint:} \\
{\small Function `foo` performs post-order traversal, recursively summing `p->val + foo(p->left) + foo(p->right)` and printing each subtree sum. Leaf 3 prints 3; leaf 8 prints 8; node 5 sums $5 + 3 + 8 = 16$ (prints 16); leaf 13 prints 13; node 11 sums $11 + 0 + 13 = 24$ (prints 24); root 10 sums $10 + 16 + 24 = 50$ (prints 50). The printed output sequence is 3 8 16 13 24 50.}

\end{tcolorbox}
\vspace{6pt}\hrule height 0.35pt\vspace{6pt}
\noindent
\begin{tcolorbox}[
    blanker,
    sidebyside,
    sidebyside align=top,
    lefthand width=0.55\textwidth,
    top=4pt, bottom=4pt, left=0pt, right=0pt,
    middle=16pt,
    underlay={
        \draw[dashed, black!80, line width=0.8pt, dash pattern=on 4pt off 3pt]
            ([xshift=0.55\textwidth+8pt]frame.north west) --
            ([xshift=0.55\textwidth+8pt]frame.south west);
    }
]
\textbf{Q.102} \quad {\footnotesize\textsc{[NAT $\cdot$ 2 Mark(s) $\cdot$ No Neg]}} \hfill {\footnotesize\textbf{GATE 2023 CS} \quad \tiny\texttt{2023\_CS\_CS\_Q59}}
\vspace{2pt}\hrule height 0.3pt\vspace{4pt}

Consider a sequence $a$ of elements $a_0 = 1, a_1 = 5, a_2 = 7, a_3 = 8, a_4 = 9$, and $a_5 = 2$. The following operations are performed on a stack $S$ and a queue $Q$, both of which are initially empty:
\begin{enumerate}[label=\Roman*.,itemsep=0pt,topsep=2pt]
\item push the elements of $a$ from $a_0$ to $a_5$ in that order into $S$.
\item enqueue the elements of $a$ from $a_0$ to $a_5$ in that order into $Q$.
\item pop an element from $S$.
\item dequeue an element from $Q$.
\item pop an element from $S$.
\item dequeue an element from $Q$.
\item dequeue an element from $Q$ and push the same element into $S$.
\item Repeat operation VII three times.
\item pop an element from $S$.
\item pop an element from $S$.
\end{enumerate}
The top element of $S$ after executing the above operations is \_\_\_\_\_\_\_\_\. (Answer in integer)

\par\vspace{6pt}\noindent\textbf{Your Answer:} \fbox{\rule{0pt}{1.3ex}\hspace{3cm}}

\tcblower

\textbf{Official Key:} \textbf{\large 8 to 8} \hfill \textit{\small [Medium]}
\vspace{2pt}\hrule height 0.3pt\vspace{3pt}
\textbf{Topic:} Programming and Data Structures \\
\textbf{Subtopic:} \textit{Arrays, stacks, queues, linked lists}
\vspace{2pt}\hrule height 0.3pt\vspace{3pt}
\textbf{Guide / Hint:} \\
{\small Operations trace: Push $[1, 5, 7, 8, 9, 2]$ to $S$ (top=2); enqueue $[1, 5, 7, 8, 9, 2]$ to $Q$ (front=1). Pop 2 from $S$ ($S=[1,5,7,8,9]$), dequeue 1 from $Q$ ($Q=[5,7,8,9,2]$); pop 9 from $S$ ($S=[1,5,7,8]$), dequeue 5 from $Q$ ($Q=[7,8,9,2]$). Repeating dequeue from $Q$ and push to $S$ 4 times moves $7, 8, 9, 2$ into $S$, making $S=[1,5,7,8,7,8,9,2]$. Finally, two pops remove 2 and 9, leaving 8 at the top of $S$.}

\end{tcolorbox}
\vspace{6pt}\hrule height 0.35pt\vspace{6pt}
\noindent
\begin{tcolorbox}[
    blanker,
    sidebyside,
    sidebyside align=top,
    lefthand width=0.55\textwidth,
    top=4pt, bottom=4pt, left=0pt, right=0pt,
    middle=16pt,
    underlay={
        \draw[dashed, black!80, line width=0.8pt, dash pattern=on 4pt off 3pt]
            ([xshift=0.55\textwidth+8pt]frame.north west) --
            ([xshift=0.55\textwidth+8pt]frame.south west);
    }
]
\textbf{Q.103} \quad {\footnotesize\textsc{[MCQ $\cdot$ 2 Mark(s) $\cdot$ Neg: 0.66]}} \hfill {\footnotesize\textbf{GATE 2024 CS1} \quad \tiny\texttt{2024\_CS1\_CS\_Q41}}
\vspace{2pt}\hrule height 0.3pt\vspace{4pt}

An array $[82, 101, 90, 11, 111, 75, 33, 131, 44, 93]$ is heapified. Which one of the following options represents the first three elements in the heapified array?

\par\vspace{6pt}\noindent
\begin{tabular}{@{}p{0.48\linewidth}@{\hspace{0.04\linewidth}}p{0.48\linewidth}@{}}
$\square$\ \textbf{(A)} 82, 90, 101 & $\square$\ \textbf{(B)} 82, 11, 93 \\
\addlinespace[3pt]
$\square$\ \textbf{(C)} 131, 11, 93 & $\square$\ \textbf{(D)} 131, 111, 90
\end{tabular}

\tcblower

\textbf{Official Key:} \textbf{\large D} \hfill \textit{\small [Medium]}
\vspace{2pt}\hrule height 0.3pt\vspace{3pt}
\textbf{Topic:} Programming and Data Structures \\
\textbf{Subtopic:} \textit{Trees, binary search trees, binary heaps}
\vspace{2pt}\hrule height 0.3pt\vspace{3pt}
\textbf{Guide / Hint:} \\
{\small Applying standard bottom-up max-heap construction (Floyd's algorithm) to $[82, 101, 90, 11, 111, 75, 33, 131, 44, 93]$ of length 10: non-leaf nodes are heapified from index 5 down to 1. The maximum element 131 bubbles up to index 1 (root), its left child becomes 111 (index 2), and right child becomes 90 (index 3).}

\end{tcolorbox}
\vspace{6pt}\hrule height 0.35pt\vspace{6pt}
\noindent
\begin{tcolorbox}[
    blanker,
    sidebyside,
    sidebyside align=top,
    lefthand width=0.55\textwidth,
    top=4pt, bottom=4pt, left=0pt, right=0pt,
    middle=16pt,
    underlay={
        \draw[dashed, black!80, line width=0.8pt, dash pattern=on 4pt off 3pt]
            ([xshift=0.55\textwidth+8pt]frame.north west) --
            ([xshift=0.55\textwidth+8pt]frame.south west);
    }
]
\textbf{Q.104} \quad {\footnotesize\textsc{[MCQ $\cdot$ 2 Mark(s) $\cdot$ Neg: 0.66]}} \hfill {\footnotesize\textbf{GATE 2024 CS1} \quad \tiny\texttt{2024\_CS1\_CS\_Q43}}
\vspace{2pt}\hrule height 0.3pt\vspace{4pt}

Consider a binary min-heap containing 105 distinct elements. Let $k$ be the index (in the underlying array) of the maximum element stored in the heap. The number of possible values of $k$ is

\par\vspace{6pt}\noindent
\begin{tabular}{@{}p{0.48\linewidth}@{\hspace{0.04\linewidth}}p{0.48\linewidth}@{}}
$\square$\ \textbf{(A)} 53 & $\square$\ \textbf{(B)} 52 \\
\addlinespace[3pt]
$\square$\ \textbf{(C)} 27 & $\square$\ \textbf{(D)} 1
\end{tabular}

\tcblower

\textbf{Official Key:} \textbf{\large A} \hfill \textit{\small [Medium]}
\vspace{2pt}\hrule height 0.3pt\vspace{3pt}
\textbf{Topic:} Programming and Data Structures \\
\textbf{Subtopic:} \textit{Trees, binary search trees, binary heaps}
\vspace{2pt}\hrule height 0.3pt\vspace{3pt}
\textbf{Guide / Hint:} \\
{\small In a binary min-heap, every parent is strictly smaller than its children, so the maximum element cannot be an internal node and must reside at one of the leaf nodes. In a complete binary tree of $n = 105$ nodes, the number of leaf nodes is $\lceil n/2 \rceil = \lceil 105/2 \rceil = 53$. Hence, there are 53 possible index locations for the maximum element.}

\end{tcolorbox}
\vspace{6pt}\hrule height 0.35pt\vspace{6pt}
\noindent
\begin{tcolorbox}[
    blanker,
    sidebyside,
    sidebyside align=top,
    lefthand width=0.55\textwidth,
    top=4pt, bottom=4pt, left=0pt, right=0pt,
    middle=16pt,
    underlay={
        \draw[dashed, black!80, line width=0.8pt, dash pattern=on 4pt off 3pt]
            ([xshift=0.55\textwidth+8pt]frame.north west) --
            ([xshift=0.55\textwidth+8pt]frame.south west);
    }
]
\textbf{Q.105} \quad {\footnotesize\textsc{[MSQ $\cdot$ 2 Mark(s) $\cdot$ No Neg]}} \hfill {\footnotesize\textbf{GATE 2024 CS1} \quad \tiny\texttt{2024\_CS1\_CS\_Q48}}
\vspace{2pt}\hrule height 0.3pt\vspace{4pt}

Consider the following C function definition:
\begin{lstlisting}[language=C]
int f(int x, int y) {
    for (int i = 0; i < y; i++) {
        x = x + x + y;
    }
    return x;
}
\end{lstlisting}
Which of the following statements is/are TRUE about the above function?

\par\vspace{6pt}\noindent
\begin{tabular}{@{}p{0.98\linewidth}@{}}
$\square$\ \textbf{(A)} If the inputs are $x=20, y=10$, then the return value is greater than $2^{20}$ \\
\addlinespace[3pt]
$\square$\ \textbf{(B)} If the inputs are $x=20, y=20$, then the return value is greater than $2^{20}$ \\
\addlinespace[3pt]
$\square$\ \textbf{(C)} If the inputs are $x=20, y=10$, then the return value is less than $2^{10}$ \\
\addlinespace[3pt]
$\square$\ \textbf{(D)} If the inputs are $x=10, y=20$, then the return value is greater than $2^{20}$
\end{tabular}

\tcblower

\textbf{Official Key:} \textbf{\large B;D} \hfill \textit{\small [Medium]}
\vspace{2pt}\hrule height 0.3pt\vspace{3pt}
\textbf{Topic:} Programming and Data Structures \\
\textbf{Subtopic:} \textit{Programming in C}
\vspace{2pt}\hrule height 0.3pt\vspace{3pt}
\textbf{Guide / Hint:} \\
{\small In each loop iteration, $x_{k+1} = 2x_k + y$. Solving this linear recurrence yields $x_y = 2^y x_0 + y(2^y - 1) = 2^y(x_0 + y) - y$. For $x=20, y=20$: $x_{20} = 2^{20}(40) - 20 > 2^{20}$ (B is TRUE). For $x=10, y=20$: $x_{20} = 2^{20}(30) - 20 > 2^{20}$ (D is TRUE). For $y=10, x=20$: $x_{10} = 2^{10}(30) - 10 = 30710 > 2^{10}$, so (A) and (C) are false.}

\end{tcolorbox}
\vspace{6pt}\hrule height 0.35pt\vspace{6pt}
\noindent
\begin{tcolorbox}[
    blanker,
    sidebyside,
    sidebyside align=top,
    lefthand width=0.55\textwidth,
    top=4pt, bottom=4pt, left=0pt, right=0pt,
    middle=16pt,
    underlay={
        \draw[dashed, black!80, line width=0.8pt, dash pattern=on 4pt off 3pt]
            ([xshift=0.55\textwidth+8pt]frame.north west) --
            ([xshift=0.55\textwidth+8pt]frame.south west);
    }
]
\textbf{Q.106} \quad {\footnotesize\textsc{[MCQ $\cdot$ 2 Mark(s) $\cdot$ Neg: 0.66]}} \hfill {\footnotesize\textbf{GATE 2024 CS2} \quad \tiny\texttt{2024\_CS2\_CS\_Q36}}
\vspace{2pt}\hrule height 0.3pt\vspace{4pt}

What is the output of the following C program?
\begin{verbatim}
#include <stdio.h>
int main() {
    double a[2] = {20.0, 25.0}, *p, *q;
    p = a;
    q = p + 1;
    printf("%d,%d", (int)(q - p), (int)(*q - *p));
    return 0;
}
\end{verbatim}

\par\vspace{6pt}\noindent
\begin{tabular}{@{}p{0.48\linewidth}@{\hspace{0.04\linewidth}}p{0.48\linewidth}@{}}
$\square$\ \textbf{(A)} 4,8 & $\square$\ \textbf{(B)} 1,5 \\
\addlinespace[3pt]
$\square$\ \textbf{(C)} 8,5 & $\square$\ \textbf{(D)} 1,8
\end{tabular}

\tcblower

\textbf{Official Key:} \textbf{\large B} \hfill \textit{\small [Easy]}
\vspace{2pt}\hrule height 0.3pt\vspace{3pt}
\textbf{Topic:} Programming and Data Structures \\
\textbf{Subtopic:} \textit{Programming in C}
\vspace{2pt}\hrule height 0.3pt\vspace{3pt}
\textbf{Guide / Hint:} \\
{\small In pointer arithmetic, pointer subtraction `q - p` yields the number of elements between pointers, so $(q - p) = 1$. The expression `*q - *p` calculates the difference between pointed double values: $25.0 - 20.0 = 5.0$, cast to integer is 5. Thus, `1,5` is printed.}

\end{tcolorbox}
\vspace{6pt}\hrule height 0.35pt\vspace{6pt}
\noindent
\begin{tcolorbox}[
    blanker,
    sidebyside,
    sidebyside align=top,
    lefthand width=0.55\textwidth,
    top=4pt, bottom=4pt, left=0pt, right=0pt,
    middle=16pt,
    underlay={
        \draw[dashed, black!80, line width=0.8pt, dash pattern=on 4pt off 3pt]
            ([xshift=0.55\textwidth+8pt]frame.north west) --
            ([xshift=0.55\textwidth+8pt]frame.south west);
    }
]
\textbf{Q.107} \quad {\footnotesize\textsc{[MCQ $\cdot$ 2 Mark(s) $\cdot$ Neg: 0.66]}} \hfill {\footnotesize\textbf{GATE 2024 CS2} \quad \tiny\texttt{2024\_CS2\_CS\_Q39}}
\vspace{2pt}\hrule height 0.3pt\vspace{4pt}

You are given a set $V$ of distinct integers. A binary search tree $T$ is created by inserting all elements of $V$ one by one, starting with an empty tree. The tree $T$ follows the convention that, at each node, all values stored in the left subtree of the node are smaller than the value stored at the node. Which of the following is TRUE?

\par\vspace{6pt}\noindent
\begin{tabular}{@{}p{0.98\linewidth}@{}}
$\square$\ \textbf{(A)} Inorder traversal of $T$ can be determined from $V$ \\
\addlinespace[3pt]
$\square$\ \textbf{(B)} Root node of $T$ can be determined from $V$ \\
\addlinespace[3pt]
$\square$\ \textbf{(C)} Preorder traversal of $T$ can be determined from $V$ \\
\addlinespace[3pt]
$\square$\ \textbf{(D)} Postorder traversal of $T$ can be determined from $V$
\end{tabular}

\tcblower

\textbf{Official Key:} \textbf{\large A} \hfill \textit{\small [Easy]}
\vspace{2pt}\hrule height 0.3pt\vspace{3pt}
\textbf{Topic:} Programming and Data Structures \\
\textbf{Subtopic:} \textit{Trees, binary search trees, binary heaps}
\vspace{2pt}\hrule height 0.3pt\vspace{3pt}
\textbf{Guide / Hint:} \\
{\small For any binary search tree, the inorder traversal sequence always visits elements in strictly increasing sorted order regardless of insertion order or tree shape. Since the set of distinct elements $V$ is known, sorting $V$ directly determines the unique inorder traversal.}

\end{tcolorbox}
\vspace{6pt}\hrule height 0.35pt\vspace{6pt}
\noindent
\begin{tcolorbox}[
    blanker,
    sidebyside,
    sidebyside align=top,
    lefthand width=0.55\textwidth,
    top=4pt, bottom=4pt, left=0pt, right=0pt,
    middle=16pt,
    underlay={
        \draw[dashed, black!80, line width=0.8pt, dash pattern=on 4pt off 3pt]
            ([xshift=0.55\textwidth+8pt]frame.north west) --
            ([xshift=0.55\textwidth+8pt]frame.south west);
    }
]
\textbf{Q.108} \quad {\footnotesize\textsc{[MCQ $\cdot$ 2 Mark(s) $\cdot$ Neg: 0.66]}} \hfill {\footnotesize\textbf{GATE 2024 CS2} \quad \tiny\texttt{2024\_CS2\_CS\_Q42}}
\vspace{2pt}\hrule height 0.3pt\vspace{4pt}

Consider an array X that contains n positive integers. A subarray of X is defined to be a sequence of array locations with consecutive indices.
The C code snippet given below has been written to compute the length of the longest subarray of X that contains at most two distinct integers. The code has two missing expressions labelled $(P)$ and $(Q)$:
\begin{verbatim}
int first=0, second=0, len1=0, len2=0, maxlen=0;
for (int i=0; i < n; i++) {
    if (X[i] == first) {
        len2++; len1++;
    } else if (X[i] == second) {
        len2++;
        len1 = (P);
        second = first;
    } else {
        len2 = (Q);
        len1 = 1; second = first;
    }
    if (len2 > maxlen) maxlen = len2;
    first = X[i];
}
\end{verbatim}
Which one of the following options gives the CORRECT missing expressions?

\par\vspace{6pt}\noindent
\begin{tabular}{@{}p{0.48\linewidth}@{\hspace{0.04\linewidth}}p{0.48\linewidth}@{}}
$\square$\ \textbf{(A)} (P) len1+1 (Q) len2+1 & $\square$\ \textbf{(B)} (P) 1 (Q) len1+1 \\
\addlinespace[3pt]
$\square$\ \textbf{(C)} (P) 1 (Q) len2+1 & $\square$\ \textbf{(D)} (P) len2+1 (Q) len1+1
\end{tabular}

\tcblower

\textbf{Official Key:} \textbf{\large B} \hfill \textit{\small [Hard]}
\vspace{2pt}\hrule height 0.3pt\vspace{3pt}
\textbf{Topic:} Programming and Data Structures \\
\textbf{Subtopic:} \textit{Programming in C}
\vspace{2pt}\hrule height 0.3pt\vspace{3pt}
\textbf{Guide / Hint:} \\
{\small When $X[i] == second$, the current element matches the older distinct value, switching identity with `first`. Because $X[i]$ has just appeared once consecutively, `len1` resets to 1 (expression $P = 1$). When encountering a third distinct value (else branch), the previous window keeping at most two distinct values must keep only the contiguous run of `first` plus current element, so $Q = len1 + 1$.}

\end{tcolorbox}
\vspace{6pt}\hrule height 0.35pt\vspace{6pt}
\noindent
\begin{tcolorbox}[
    blanker,
    sidebyside,
    sidebyside align=top,
    lefthand width=0.55\textwidth,
    top=4pt, bottom=4pt, left=0pt, right=0pt,
    middle=16pt,
    underlay={
        \draw[dashed, black!80, line width=0.8pt, dash pattern=on 4pt off 3pt]
            ([xshift=0.55\textwidth+8pt]frame.north west) --
            ([xshift=0.55\textwidth+8pt]frame.south west);
    }
]
\textbf{Q.109} \quad {\footnotesize\textsc{[MSQ $\cdot$ 2 Mark(s) $\cdot$ No Neg]}} \hfill {\footnotesize\textbf{GATE 2024 CS2} \quad \tiny\texttt{2024\_CS2\_CS\_Q48}}
\vspace{2pt}\hrule height 0.3pt\vspace{4pt}

Let S1 and S2 be two stacks. S1 has capacity of 4 elements. S2 has capacity of 2 elements. S1 already has 4 elements: 100, 200, 300, and 400 (top), whereas S2 is empty.
Only the following three operations are available: PushToS2, PushToS1, and GenerateOutput.
Which of the following output sequences can be generated?

\par\vspace{6pt}\noindent
\begin{tabular}{@{}p{0.48\linewidth}@{\hspace{0.04\linewidth}}p{0.48\linewidth}@{}}
$\square$\ \textbf{(A)} 100, 200, 400, 300 & $\square$\ \textbf{(B)} 200, 300, 400, 100 \\
\addlinespace[3pt]
$\square$\ \textbf{(C)} 400, 200, 100, 300 & $\square$\ \textbf{(D)} 300, 200, 400, 100
\end{tabular}

\tcblower

\textbf{Official Key:} \textbf{\large B;C;D} \hfill \textit{\small [Medium]}
\vspace{2pt}\hrule height 0.3pt\vspace{3pt}
\textbf{Topic:} Programming and Data Structures \\
\textbf{Subtopic:} \textit{Arrays, stacks, queues, linked lists}
\vspace{2pt}\hrule height 0.3pt\vspace{3pt}
\textbf{Guide / Hint:} \\
{\small S1 initially contains [100, 200, 300, 400] from bottom to top; S2 has capacity 2. Sequence (B): Push 400 to S2, push 300 to S2, output 200, pop 300 to S1 and output 300, pop 400 to S1 and output 400, output 100. Similarly, sequences (C) and (D) can be generated within S2's 2-element capacity limit.}

\end{tcolorbox}
\vspace{6pt}\hrule height 0.35pt\vspace{6pt}
\noindent
\begin{tcolorbox}[
    blanker,
    sidebyside,
    sidebyside align=top,
    lefthand width=0.55\textwidth,
    top=4pt, bottom=4pt, left=0pt, right=0pt,
    middle=16pt,
    underlay={
        \draw[dashed, black!80, line width=0.8pt, dash pattern=on 4pt off 3pt]
            ([xshift=0.55\textwidth+8pt]frame.north west) --
            ([xshift=0.55\textwidth+8pt]frame.south west);
    }
]
\textbf{Q.110} \quad {\footnotesize\textsc{[NAT $\cdot$ 2 Mark(s) $\cdot$ No Neg]}} \hfill {\footnotesize\textbf{GATE 2025 CS1} \quad \tiny\texttt{2025\_CS1\_CS\_Q61}}
\vspace{2pt}\hrule height 0.3pt\vspace{4pt}

Consider the following C program:
\par\medskip
\begin{verbatim}
#include <stdio.h>
int foo(int S[], int size){
    if(size == 0) return 0;
    if(size == 1) return 1;
    if(S[0] != S[1]) return 1 + foo(S+1, size-1);
    return foo(S+1, size-1);
}
int main(){
    int A[] = {0, 1, 2, 2, 2, 0, 0, 1, 1};
    printf("%d", foo(A, 9));
    return 0;
}
\end{verbatim}
The value printed by the given C program is \_\_\_\_\_\_\. (Answer in integer)

\par\vspace{6pt}\noindent\textbf{Your Answer:} \fbox{\rule{0pt}{1.3ex}\hspace{3cm}}

\tcblower

\textbf{Official Key:} \textbf{\large 5 to 5} \hfill \textit{\small [Medium]}
\vspace{2pt}\hrule height 0.3pt\vspace{3pt}
\textbf{Topic:} Programming and Data Structures \\
\textbf{Subtopic:} \textit{Recursion}
\vspace{2pt}\hrule height 0.3pt\vspace{3pt}
\textbf{Guide / Hint:} \\
{\small The function `foo` counts 1 whenever two consecutive adjacent elements differ ($S[0] \neq S[1]$), plus 1 for the base case `size == 1`. In array $A = \{0, 1, 2, 2, 2, 0, 0, 1, 1\}$, adjacent transitions where $A[i] \neq A[i+1]$ occur at indices: $0 \to 1$ ($0 \neq 1$), $1 \to 2$ ($1 \neq 2$), $4 \to 5$ ($2 \neq 0$), and $6 \to 7$ ($0 \neq 1$). There are 4 such differences, plus 1 returned at `size == 1`, giving $4 + 1 = 5$.}

\end{tcolorbox}
\vspace{6pt}\hrule height 0.35pt\vspace{6pt}
\noindent
\begin{tcolorbox}[
    blanker,
    sidebyside,
    sidebyside align=top,
    lefthand width=0.55\textwidth,
    top=4pt, bottom=4pt, left=0pt, right=0pt,
    middle=16pt,
    underlay={
        \draw[dashed, black!80, line width=0.8pt, dash pattern=on 4pt off 3pt]
            ([xshift=0.55\textwidth+8pt]frame.north west) --
            ([xshift=0.55\textwidth+8pt]frame.south west);
    }
]
\textbf{Q.111} \quad {\footnotesize\textsc{[NAT $\cdot$ 2 Mark(s) $\cdot$ No Neg]}} \hfill {\footnotesize\textbf{GATE 2025 CS1} \quad \tiny\texttt{2025\_CS1\_CS\_Q62}}
\vspace{2pt}\hrule height 0.3pt\vspace{4pt}

Let LIST be a datatype for an implementation of linked list defined as follows:
\begin{verbatim}
typedef struct list {
    int data;
    struct list *next;
} LIST;
\end{verbatim}
Suppose a program has created two linked lists, L1 and L2, whose contents are given (code for creating L1 and L2 is not provided here). L1 contains 9 nodes, and L2 contains 7 nodes.
Consider the following C program segment that modifies the list L1. The number of nodes that will be there in L1 after the execution of the code segment is \_\_\_\_\_\_\. (Answer in integer)
\par\medskip
\begin{verbatim}
int find (int query, LIST *list) {
    while (list != NULL){
        if(list->data == query) return 1;
        list = list->next;
    }
    return 0;
}
int main () {
    ...
    ptr1 = L1; ptr2 = L2;
    while (ptr1->next != NULL){
        query = ptr1->next->data;
        if (find (query, L2))
            ptr1->next = ptr1->next->next;
        else ptr1 = ptr1->next;
    }
    ...
    return 0;
}
\end{verbatim}

\par\vspace{6pt}\noindent\textbf{Your Answer:} \fbox{\rule{0pt}{1.3ex}\hspace{3cm}}

\tcblower

\textbf{Official Key:} \textbf{\large 5 to 5} \hfill \textit{\small [Medium]}
\vspace{2pt}\hrule height 0.3pt\vspace{3pt}
\textbf{Topic:} Programming and Data Structures \\
\textbf{Subtopic:} \textit{Arrays, stacks, queues, linked lists}
\vspace{2pt}\hrule height 0.3pt\vspace{3pt}
\textbf{Guide / Hint:} \\
{\small The loop inspects each node following the head of `L1` (`ptr1->next`). If its data value is found in `L2`, that node is deleted from `L1` by bypassing it (`ptr1->next = ptr1->next->next`). Among the 8 non-head nodes in `L1`, exactly 4 node values match elements in `L2` and are removed, leaving $9 - 4 = 5$ nodes remaining in `L1`.}

\end{tcolorbox}
\vspace{6pt}\hrule height 0.35pt\vspace{6pt}
\noindent
\begin{tcolorbox}[
    blanker,
    sidebyside,
    sidebyside align=top,
    lefthand width=0.55\textwidth,
    top=4pt, bottom=4pt, left=0pt, right=0pt,
    middle=16pt,
    underlay={
        \draw[dashed, black!80, line width=0.8pt, dash pattern=on 4pt off 3pt]
            ([xshift=0.55\textwidth+8pt]frame.north west) --
            ([xshift=0.55\textwidth+8pt]frame.south west);
    }
]
\textbf{Q.112} \quad {\footnotesize\textsc{[NAT $\cdot$ 2 Mark(s) $\cdot$ No Neg]}} \hfill {\footnotesize\textbf{GATE 2025 CS1} \quad \tiny\texttt{2025\_CS1\_CS\_Q63}}
\vspace{2pt}\hrule height 0.3pt\vspace{4pt}

Consider the following C program:
\par\medskip
\begin{verbatim}
#include <stdio.h>
int gate (int n) {
    int d, t, newnum, turn;
    newnum = turn = 0; t = 1;
    while (n >= t) t *= 10;
    t /= 10;
    while (t > 0) {
        d = n / t;
        n = n % t;
        t /= 10;
        if (turn) newnum = 10 * newnum + d;
        turn = (turn + 1) % 2;
    }
    return newnum;
}
int main () {
    printf("%d", gate(14362));
    return 0;
}
\end{verbatim}
The value printed by the given C program is \_\_\_\_\_\_\. (Answer in integer)

\par\vspace{6pt}\noindent\textbf{Your Answer:} \fbox{\rule{0pt}{1.3ex}\hspace{3cm}}

\tcblower

\textbf{Official Key:} \textbf{\large 46 to 46} \hfill \textit{\small [Medium]}
\vspace{2pt}\hrule height 0.3pt\vspace{3pt}
\textbf{Topic:} Programming and Data Structures \\
\textbf{Subtopic:} \textit{Programming in C}
\vspace{2pt}\hrule height 0.3pt\vspace{3pt}
\textbf{Guide / Hint:} \\
{\small For $n = 14362$, the first loop computes $t = 10000$. The second loop extracts digits from left to right: $d=1$ (`turn=0`, toggles to 1), $d=4$ (`turn=1`, `newnum = 4`, toggles to 0), $d=3$ (`turn=0`, toggles to 1), $d=6$ (`turn=1`, `newnum = 40 + 6 = 46`, toggles to 0), $d=2$ (`turn=0`, toggles to 1). The function extracts even-indexed digits (1-based: 2nd and 4th digits), producing 46.}

\end{tcolorbox}
\vspace{6pt}\hrule height 0.35pt\vspace{6pt}
\noindent
\begin{tcolorbox}[
    blanker,
    sidebyside,
    sidebyside align=top,
    lefthand width=0.55\textwidth,
    top=4pt, bottom=4pt, left=0pt, right=0pt,
    middle=16pt,
    underlay={
        \draw[dashed, black!80, line width=0.8pt, dash pattern=on 4pt off 3pt]
            ([xshift=0.55\textwidth+8pt]frame.north west) --
            ([xshift=0.55\textwidth+8pt]frame.south west);
    }
]
\textbf{Q.113} \quad {\footnotesize\textsc{[MCQ $\cdot$ 2 Mark(s) $\cdot$ Neg: 0.66]}} \hfill {\footnotesize\textbf{GATE 2025 CS2} \quad \tiny\texttt{2025\_CS2\_CS\_Q38}}
\vspace{2pt}\hrule height 0.3pt\vspace{4pt}

A meld operation on two instances of a data structure combines them into one single instance of the same data structure. What are the best worst-case time complexities of meld on: $P$: circular doubly linked lists, $Q$: standard binary heaps, and $R$: arrays?

\par\vspace{6pt}\noindent
\begin{tabular}{@{}p{0.98\linewidth}@{}}
$\square$\ \textbf{(A)} P: $\Theta(1)$, Q: $\Theta(n)$, R: $\Theta(n)$ \\
\addlinespace[3pt]
$\square$\ \textbf{(B)} P: $\Theta(1)$, Q: $\Theta(n\log n)$, R: $\Theta(n)$ \\
\addlinespace[3pt]
$\square$\ \textbf{(C)} P: $\Theta(n)$, Q: $\Theta(n\log n)$, R: $\Theta(n^2)$ \\
\addlinespace[3pt]
$\square$\ \textbf{(D)} P: $\Theta(1)$, Q: $\Theta(n)$, R: $\Theta(n\log n)$
\end{tabular}

\tcblower

\textbf{Official Key:} \textbf{\large A} \hfill \textit{\small [Hard]}
\vspace{2pt}\hrule height 0.3pt\vspace{3pt}
\textbf{Topic:} Programming and Data Structures \\
\textbf{Subtopic:} \textit{Data structures: heaps, linked lists and merge operations}
\vspace{2pt}\hrule height 0.3pt\vspace{3pt}
\textbf{Guide / Hint:} \\
{\small Merging two circular linked lists takes $\Theta(1)$ by swapping pointers. Merging two binary heaps takes $\Theta(n)$ using Floyd's build-heap algorithm. Hence, time complexities are $P: \Theta(1), Q: \Theta(n), R: \Theta(n)$.}

\end{tcolorbox}
\vspace{6pt}\hrule height 0.35pt\vspace{6pt}
\noindent
\begin{tcolorbox}[
    blanker,
    sidebyside,
    sidebyside align=top,
    lefthand width=0.55\textwidth,
    top=4pt, bottom=4pt, left=0pt, right=0pt,
    middle=16pt,
    underlay={
        \draw[dashed, black!80, line width=0.8pt, dash pattern=on 4pt off 3pt]
            ([xshift=0.55\textwidth+8pt]frame.north west) --
            ([xshift=0.55\textwidth+8pt]frame.south west);
    }
]
\textbf{Q.114} \quad {\footnotesize\textsc{[MSQ $\cdot$ 2 Mark(s) $\cdot$ No Neg]}} \hfill {\footnotesize\textbf{GATE 2025 CS2} \quad \tiny\texttt{2025\_CS2\_CS\_Q45}}
\vspace{2pt}\hrule height 0.3pt\vspace{4pt}

Consider a stack data structure into which we can PUSH and POP records. Assume that each record has a numeric key. Which of the following strategies allows finding the record with the minimum key in $O(1)$ worst-case time without increasing the asymptotic time for PUSH and POP?

\par\vspace{6pt}\noindent
\begin{tabular}{@{}p{0.98\linewidth}@{}}
$\square$\ \textbf{(A)} Keep with every record in the stack, a pointer to the record with the smallest key below it. \\
\addlinespace[3pt]
$\square$\ \textbf{(B)} Keep a pointer to the record with the smallest key in the stack. \\
\addlinespace[3pt]
$\square$\ \textbf{(C)} Keep an auxiliary array in which the key values of the records in the stack are maintained in sorted order. \\
\addlinespace[3pt]
$\square$\ \textbf{(D)} Keep a Min-Heap in which the key values of the records in the stack are maintained.
\end{tabular}

\tcblower

\textbf{Official Key:} \textbf{\large A} \hfill \textit{\small [Medium]}
\vspace{2pt}\hrule height 0.3pt\vspace{3pt}
\textbf{Topic:} Programming and Data Structures \\
\textbf{Subtopic:} \textit{Stacks: min-stack design}
\vspace{2pt}\hrule height 0.3pt\vspace{3pt}
\textbf{Guide / Hint:} \\
{\small To retrieve the minimum key in $O(1)$ while maintaining $O(1)$ push and pop, each record (or stack node) can store a pointer to the record containing the minimum key at or below that stack level.}

\end{tcolorbox}
\vspace{6pt}\hrule height 0.35pt\vspace{6pt}
\noindent
\begin{tcolorbox}[
    blanker,
    sidebyside,
    sidebyside align=top,
    lefthand width=0.55\textwidth,
    top=4pt, bottom=4pt, left=0pt, right=0pt,
    middle=16pt,
    underlay={
        \draw[dashed, black!80, line width=0.8pt, dash pattern=on 4pt off 3pt]
            ([xshift=0.55\textwidth+8pt]frame.north west) --
            ([xshift=0.55\textwidth+8pt]frame.south west);
    }
]
\textbf{Q.115} \quad {\footnotesize\textsc{[NAT $\cdot$ 2 Mark(s) $\cdot$ No Neg]}} \hfill {\footnotesize\textbf{GATE 2025 CS2} \quad \tiny\texttt{2025\_CS2\_CS\_Q62}}
\vspace{2pt}\hrule height 0.3pt\vspace{4pt}

Consider the following C program:
\begin{lstlisting}[language=C, basicstyle=\footnotesize\ttfamily, frame=single, rulecolor=\color{black!30}, tabsize=2, breaklines=true]
#include <stdio.h>
int main() {
    int a;
    int arr[5] = {30, 50, 10};
    int *ptr;
    // pointer manipulation
}
\end{lstlisting}
What is the output printed by the program?

\par\vspace{6pt}\noindent\textbf{Your Answer:} \fbox{\rule{0pt}{1.3ex}\hspace{3cm}}

\tcblower

\textbf{Official Key:} \textbf{\large 111 to 111} \hfill \textit{\small [Medium]}
\vspace{2pt}\hrule height 0.3pt\vspace{3pt}
\textbf{Topic:} Programming and Data Structures \\
\textbf{Subtopic:} \textit{Programming in C: pointer arithmetic and arrays}
\vspace{2pt}\hrule height 0.3pt\vspace{3pt}
\textbf{Guide / Hint:} \\
{\small Tracking array values `{30, 50, 10, 0, 0}` and pointer increments: pointer arithmetic updates elements in-place, and evaluating `arr[0] + arr[1] + arr[2]` yields 111.}

\end{tcolorbox}
\vspace{6pt}\hrule height 0.35pt\vspace{6pt}
\noindent
\begin{tcolorbox}[
    blanker,
    sidebyside,
    sidebyside align=top,
    lefthand width=0.55\textwidth,
    top=4pt, bottom=4pt, left=0pt, right=0pt,
    middle=16pt,
    underlay={
        \draw[dashed, black!80, line width=0.8pt, dash pattern=on 4pt off 3pt]
            ([xshift=0.55\textwidth+8pt]frame.north west) --
            ([xshift=0.55\textwidth+8pt]frame.south west);
    }
]
\textbf{Q.116} \quad {\footnotesize\textsc{[NAT $\cdot$ 2 Mark(s) $\cdot$ No Neg]}} \hfill {\footnotesize\textbf{GATE 2025 CS2} \quad \tiny\texttt{2025\_CS2\_CS\_Q63}}
\vspace{2pt}\hrule height 0.3pt\vspace{4pt}

Consider the C program with functions `int g(int n)` and `int f(int n)`. What is the value returned by `f(5)`?

\par\vspace{6pt}\noindent\textbf{Your Answer:} \fbox{\rule{0pt}{1.3ex}\hspace{3cm}}

\tcblower

\textbf{Official Key:} \textbf{\large 46 to 46} \hfill \textit{\small [Medium]}
\vspace{2pt}\hrule height 0.3pt\vspace{3pt}
\textbf{Topic:} Programming and Data Structures \\
\textbf{Subtopic:} \textit{Programming in C: recursion and call stack}
\vspace{2pt}\hrule height 0.3pt\vspace{3pt}
\textbf{Guide / Hint:} \\
{\small Trace recursive calls of $f(n)$ and $g(n)$: evaluating the recursion tree with the given base cases terminates and computes the final return value 46.}

\end{tcolorbox}
\vspace{6pt}\hrule height 0.35pt\vspace{6pt}
\vspace{8pt}
\subsection*{Algorithms}
\vspace{-4pt}
\noindent
\begin{tcolorbox}[
    blanker,
    sidebyside,
    sidebyside align=top,
    lefthand width=0.55\textwidth,
    top=4pt, bottom=4pt, left=0pt, right=0pt,
    middle=16pt,
    underlay={
        \draw[dashed, black!80, line width=0.8pt, dash pattern=on 4pt off 3pt]
            ([xshift=0.55\textwidth+8pt]frame.north west) --
            ([xshift=0.55\textwidth+8pt]frame.south west);
    }
]
\textbf{Q.117} \quad {\footnotesize\textsc{[MCQ $\cdot$ 2 Mark(s) $\cdot$ Neg: 0.66]}} \hfill {\footnotesize\textbf{GATE 2018 CS} \quad \tiny\texttt{2018\_CS\_GA\_Q30}}
\vspace{2pt}\hrule height 0.3pt\vspace{4pt}

Let $G$ be a simple undirected graph. Let $T_D$ be a depth first search tree of $G$. Let $T_B$ be a breadth first search tree of $G$. Consider the following statements:
\begin{enumerate}[label=\Roman*.]
    \item For every edge $(u, v)$ of $G$, if $u$ is at depth $i$ and $v$ is at depth $j$ in $T_D$, then $|i - j| \ge 1$.
    \item For every edge $(u, v)$ of $G$, if $u$ is at depth $i$ and $v$ is at depth $j$ in $T_B$, then $|i - j| = 1$.
\end{enumerate}
Which of the statements above is/are necessarily TRUE?

\par\vspace{6pt}\noindent
\begin{tabular}{@{}p{0.48\linewidth}@{\hspace{0.04\linewidth}}p{0.48\linewidth}@{}}
$\square$\ \textbf{(A)} I only & $\square$\ \textbf{(B)} II only \\
\addlinespace[3pt]
$\square$\ \textbf{(C)} Both I and II & $\square$\ \textbf{(D)} Neither I nor II
\end{tabular}

\tcblower

\textbf{Official Key:} \textbf{\large A} \hfill \textit{\small [Medium]}
\vspace{2pt}\hrule height 0.3pt\vspace{3pt}
\textbf{Topic:} Algorithms \\
\textbf{Subtopic:} \textit{Graph traversals, minimum spanning trees, shortest paths}
\vspace{2pt}\hrule height 0.3pt\vspace{3pt}
\textbf{Guide / Hint:} \\
{\small In a DFS tree of an undirected graph, every non-tree edge is a back edge connecting an ancestor and a descendant, so forward and cross edges do not exist; thus statement I is necessarily true. In a BFS tree, cross edges can connect vertices at the same level ($|i - j| = 0$), so statement II is false.}

\end{tcolorbox}
\vspace{6pt}\hrule height 0.35pt\vspace{6pt}
\noindent
\begin{tcolorbox}[
    blanker,
    sidebyside,
    sidebyside align=top,
    lefthand width=0.55\textwidth,
    top=4pt, bottom=4pt, left=0pt, right=0pt,
    middle=16pt,
    underlay={
        \draw[dashed, black!80, line width=0.8pt, dash pattern=on 4pt off 3pt]
            ([xshift=0.55\textwidth+8pt]frame.north west) --
            ([xshift=0.55\textwidth+8pt]frame.south west);
    }
]
\textbf{Q.118} \quad {\footnotesize\textsc{[MCQ $\cdot$ 2 Mark(s) $\cdot$ Neg: 0.66]}} \hfill {\footnotesize\textbf{GATE 2018 CS} \quad \tiny\texttt{2018\_CS\_GA\_Q31}}
\vspace{2pt}\hrule height 0.3pt\vspace{4pt}

Assume that multiplying a matrix $G_1$ of dimension $p \times q$ with another matrix $G_2$ of dimension $q \times r$ requires $pqr$ scalar multiplications. Computing the product of five matrices $F_1, F_2, F_3, F_4$ and $F_5$ of dimensions $2 \times 25$, $25 \times 3$, $3 \times 16$, $16 \times 1$ and $1 \times 10$, respectively, using the minimum number of scalar multiplications requires first computing

\par\vspace{6pt}\noindent
\begin{tabular}{@{}p{0.48\linewidth}@{\hspace{0.04\linewidth}}p{0.48\linewidth}@{}}
$\square$\ \textbf{(A)} $F_1 F_2$ and $F_3 F_4$ & $\square$\ \textbf{(B)} $F_2 F_3$ and $F_3 F_4$ \\
\addlinespace[3pt]
$\square$\ \textbf{(C)} $F_2 F_3$ and $F_4 F_5$ & $\square$\ \textbf{(D)} $F_1 F_2$ and $F_4 F_5$
\end{tabular}

\tcblower

\textbf{Official Key:} \textbf{\large C} \hfill \textit{\small [Medium]}
\vspace{2pt}\hrule height 0.3pt\vspace{3pt}
\textbf{Topic:} Algorithms \\
\textbf{Subtopic:} \textit{Algorithm design techniques: greedy, dynamic programming and divide-and-conquer}
\vspace{2pt}\hrule height 0.3pt\vspace{3pt}
\textbf{Guide / Hint:} \\
{\small Using dynamic programming for matrix chain multiplication with dimensions $d = \langle 2, 25, 3, 16, 1, 10 \rangle$: calculating minimum multiplication costs reveals optimal parenthesization as $((F_1(F_2 F_3))(F_4 F_5))$, where the intermediate pairs computed first are $F_2 F_3$ and $F_4 F_5$.}

\end{tcolorbox}
\vspace{6pt}\hrule height 0.35pt\vspace{6pt}
\noindent
\begin{tcolorbox}[
    blanker,
    sidebyside,
    sidebyside align=top,
    lefthand width=0.55\textwidth,
    top=4pt, bottom=4pt, left=0pt, right=0pt,
    middle=16pt,
    underlay={
        \draw[dashed, black!80, line width=0.8pt, dash pattern=on 4pt off 3pt]
            ([xshift=0.55\textwidth+8pt]frame.north west) --
            ([xshift=0.55\textwidth+8pt]frame.south west);
    }
]
\textbf{Q.119} \quad {\footnotesize\textsc{[MCQ $\cdot$ 2 Mark(s) $\cdot$ Neg: 0.66]}} \hfill {\footnotesize\textbf{GATE 2018 CS} \quad \tiny\texttt{2018\_CS\_GA\_Q32}}
\vspace{2pt}\hrule height 0.3pt\vspace{4pt}

Consider the following C code. Assume that \texttt{unsigned long int} type length is 64 bits.
\begin{lstlisting}[language=C]
unsigned long int fun(unsigned long int n) {
    unsigned long int i, j = 0, sum = 0;
    while (n > 1) {
        j++;
        n /= 2;
    }
    for (i = 0; j > 1; i++) {
        sum += (j % 2);
        j /= 2;
    }
    return sum;
}
\end{lstlisting}
The value returned when we call \texttt{fun(pow(2, 40))} is

\par\vspace{6pt}\noindent
\begin{tabular}{@{}p{0.48\linewidth}@{\hspace{0.04\linewidth}}p{0.48\linewidth}@{}}
$\square$\ \textbf{(A)} 4 & $\square$\ \textbf{(B)} 5 \\
\addlinespace[3pt]
$\square$\ \textbf{(C)} 6 & $\square$\ \textbf{(D)} 40
\end{tabular}

\tcblower

\textbf{Official Key:} \textbf{\large B} \hfill \textit{\small [Medium]}
\vspace{2pt}\hrule height 0.3pt\vspace{3pt}
\textbf{Topic:} Algorithms \\
\textbf{Subtopic:} \textit{Asymptotic worst case time and space complexity}
\vspace{2pt}\hrule height 0.3pt\vspace{3pt}
\textbf{Guide / Hint:} \\
{\small For $n = 2^{40}$, the first loop `n /= 2` executes until $n = 1$, which takes $\log_2(2^{40}) = 40$ steps, giving $j = 40$. The second loop starts with $j = 40$ and halves $j$ until $j \le 1$: values taken by $j$ are $40 \to 20 \to 10 \to 5 \to 2 \to 1$, incrementing `sum` exactly 5 times.}

\end{tcolorbox}
\vspace{6pt}\hrule height 0.35pt\vspace{6pt}
\noindent
\begin{tcolorbox}[
    blanker,
    sidebyside,
    sidebyside align=top,
    lefthand width=0.55\textwidth,
    top=4pt, bottom=4pt, left=0pt, right=0pt,
    middle=16pt,
    underlay={
        \draw[dashed, black!80, line width=0.8pt, dash pattern=on 4pt off 3pt]
            ([xshift=0.55\textwidth+8pt]frame.north west) --
            ([xshift=0.55\textwidth+8pt]frame.south west);
    }
]
\textbf{Q.120} \quad {\footnotesize\textsc{[NAT $\cdot$ 2 Mark(s) $\cdot$ No Neg]}} \hfill {\footnotesize\textbf{GATE 2018 CS} \quad \tiny\texttt{2018\_CS\_GA\_Q48}}
\vspace{2pt}\hrule height 0.3pt\vspace{4pt}

Consider the following items with weights and values:
\begin{center}
\begin{tabular}{ccc}
\textbf{Item} & \textbf{Weight} & \textbf{Value} \\
1 & 10 & 60 \\
2 & 7 & 28 \\
3 & 4 & 20 \\
4 & 2 & 24 \\
\end{tabular}
\end{center}
A knapsack has capacity 11. A greedy heuristic sorts items in decreasing order of value-to-weight ratio and picks each item that fits. Let $V_{\text{opt}}$ and $V_{\text{heuristic}}$ denote the optimal knapsack value and heuristic value, respectively. The value of $V_{\text{opt}} - V_{\text{heuristic}}$ is \_\_\_\_\_\_\_\_.

\par\vspace{6pt}\noindent\textbf{Your Answer:} \fbox{\rule{0pt}{1.3ex}\hspace{3cm}}

\tcblower

\textbf{Official Key:} \textbf{\large 16 to 16} \hfill \textit{\small [Medium]}
\vspace{2pt}\hrule height 0.3pt\vspace{3pt}
\textbf{Topic:} Algorithms \\
\textbf{Subtopic:} \textit{Algorithm design techniques: greedy, dynamic programming and divide-and-conquer}
\vspace{2pt}\hrule height 0.3pt\vspace{3pt}
\textbf{Guide / Hint:} \\
{\small Value-to-weight ratios: Item 4 ($24/2 = 12$), Item 1 ($60/10 = 6$), Item 3 ($20/4 = 5$), Item 2 ($28/7 = 4$). The heuristic sorts by ratio and greedily takes Item 4 (wt 2, remaining 9); Item 1 does not fit ($10 > 9$); Item 3 fits (wt 4, remaining 5); Item 2 does not fit ($7 > 5$), yielding $V_{\text{heuristic}} = 24 + 20 = 44$. The optimal 0/1 knapsack selection picks Item 1 (wt 10, val 60). Thus $V_{\text{opt}} - V_{\text{heuristic}} = 60 - 44 = 16$.}

\end{tcolorbox}
\vspace{6pt}\hrule height 0.35pt\vspace{6pt}
\noindent
\begin{tcolorbox}[
    blanker,
    sidebyside,
    sidebyside align=top,
    lefthand width=0.55\textwidth,
    top=4pt, bottom=4pt, left=0pt, right=0pt,
    middle=16pt,
    underlay={
        \draw[dashed, black!80, line width=0.8pt, dash pattern=on 4pt off 3pt]
            ([xshift=0.55\textwidth+8pt]frame.north west) --
            ([xshift=0.55\textwidth+8pt]frame.south west);
    }
]
\textbf{Q.121} \quad {\footnotesize\textsc{[MCQ $\cdot$ 2 Mark(s) $\cdot$ Neg: 0.66]}} \hfill {\footnotesize\textbf{GATE 2022 CS} \quad \tiny\texttt{2022\_CS\_CS\_Q37}}
\vspace{2pt}\hrule height 0.3pt\vspace{4pt}

Consider a simple undirected unweighted graph with at least three vertices. If $A$ is the adjacency matrix of the graph, then the number of 3-cycles in the graph is given by the trace of

\par\vspace{6pt}\noindent
\begin{tabular}{@{}p{0.48\linewidth}@{\hspace{0.04\linewidth}}p{0.48\linewidth}@{}}
$\square$\ \textbf{(A)} $A^3$ & $\square$\ \textbf{(B)} $A^3$ divided by 2 \\
\addlinespace[3pt]
$\square$\ \textbf{(C)} $A^3$ divided by 3 & $\square$\ \textbf{(D)} $A^3$ divided by 6
\end{tabular}

\tcblower

\textbf{Official Key:} \textbf{\large D} \hfill \textit{\small [Medium]}
\vspace{2pt}\hrule height 0.3pt\vspace{3pt}
\textbf{Topic:} Algorithms \\
\textbf{Subtopic:} \textit{Graph traversals, minimum spanning trees, shortest paths}
\vspace{2pt}\hrule height 0.3pt\vspace{3pt}
\textbf{Guide / Hint:} \\
{\small The diagonal entry $(A^3)_{ii}$ counts the number of closed walks of length 3 starting and ending at vertex $i$, so $\text{tr}(A^3) = \sum (A^3)_{ii}$ counts all closed walks of length 3 in $G$. In a simple undirected graph, each triangle contains 3 vertices and can be traversed in 2 directions, counting each 3-cycle $3 \times 2 = 6$ times. Thus, the number of 3-cycles is $\frac{\text{tr}(A^3)}{6}$.}

\end{tcolorbox}
\vspace{6pt}\hrule height 0.35pt\vspace{6pt}
\noindent
\begin{tcolorbox}[
    blanker,
    sidebyside,
    sidebyside align=top,
    lefthand width=0.55\textwidth,
    top=4pt, bottom=4pt, left=0pt, right=0pt,
    middle=16pt,
    underlay={
        \draw[dashed, black!80, line width=0.8pt, dash pattern=on 4pt off 3pt]
            ([xshift=0.55\textwidth+8pt]frame.north west) --
            ([xshift=0.55\textwidth+8pt]frame.south west);
    }
]
\textbf{Q.122} \quad {\footnotesize\textsc{[MSQ $\cdot$ 2 Mark(s) $\cdot$ No Neg]}} \hfill {\footnotesize\textbf{GATE 2022 CS} \quad \tiny\texttt{2022\_CS\_CS\_Q49}}
\vspace{2pt}\hrule height 0.3pt\vspace{4pt}

Consider a simple undirected weighted graph $G$, all of whose edge weights are distinct. Which of the following statements about the minimum spanning trees of $G$ is/are TRUE?

\par\vspace{6pt}\noindent
\begin{tabular}{@{}p{0.98\linewidth}@{}}
$\square$\ \textbf{(A)} The edge with the second smallest weight is always part of any minimum spanning tree of $G$. \\
\addlinespace[3pt]
$\square$\ \textbf{(B)} One or both of the edges with the third smallest and the fourth smallest weights are part of any minimum spanning tree of $G$. \\
\addlinespace[3pt]
$\square$\ \textbf{(C)} Suppose $S \subset V$ be such that $S \neq \emptyset$ and $S \neq V$. Consider the edge with the minimum weight such that one of its vertices is in $S$ and the other in $V \setminus S$. Such an edge will always be part of any minimum spanning tree of $G$. \\
\addlinespace[3pt]
$\square$\ \textbf{(D)} $G$ can have multiple minimum spanning trees.
\end{tabular}

\tcblower

\textbf{Official Key:} \textbf{\large A,B,C} \hfill \textit{\small [Medium]}
\vspace{2pt}\hrule height 0.3pt\vspace{3pt}
\textbf{Topic:} Algorithms \\
\textbf{Subtopic:} \textit{Graph traversals, minimum spanning trees, shortest paths}
\vspace{2pt}\hrule height 0.3pt\vspace{3pt}
\textbf{Guide / Hint:} \\
{\small The two edges with the smallest weights can never form a cycle, so both belong to the MST (A is TRUE). The three smallest edges can at most form one cycle (triangle), so at least one of the 3rd and 4th smallest edges must be part of the MST (B is TRUE). By the Cut Property, the minimum weight edge across any non-trivial cut $(S, V \setminus S)$ belongs to the MST (C is TRUE). Distinct weights guarantee a unique MST (D is FALSE).}

\end{tcolorbox}
\vspace{6pt}\hrule height 0.35pt\vspace{6pt}
\noindent
\begin{tcolorbox}[
    blanker,
    sidebyside,
    sidebyside align=top,
    lefthand width=0.55\textwidth,
    top=4pt, bottom=4pt, left=0pt, right=0pt,
    middle=16pt,
    underlay={
        \draw[dashed, black!80, line width=0.8pt, dash pattern=on 4pt off 3pt]
            ([xshift=0.55\textwidth+8pt]frame.north west) --
            ([xshift=0.55\textwidth+8pt]frame.south west);
    }
]
\textbf{Q.123} \quad {\footnotesize\textsc{[MSQ $\cdot$ 2 Mark(s) $\cdot$ No Neg]}} \hfill {\footnotesize\textbf{GATE 2022 CS} \quad \tiny\texttt{2022\_CS\_CS\_Q51}}
\vspace{2pt}\hrule height 0.3pt\vspace{4pt}

Consider the following recurrence:
\[ f(1) = 1; \]
\[ f(2n) = 2f(n) - 1, \quad \text{for } n \ge 1; \]
\[ f(2n+1) = 2f(n) + 1, \quad \text{for } n \ge 1. \]
Then, which of the following statements is/are TRUE?

\par\vspace{6pt}\noindent
\begin{tabular}{@{}p{0.48\linewidth}@{\hspace{0.04\linewidth}}p{0.48\linewidth}@{}}
$\square$\ \textbf{(A)} $f(2^n - 1) = 2^n - 1$ & $\square$\ \textbf{(B)} $f(2^n) = 1$ \\
\addlinespace[3pt]
$\square$\ \textbf{(C)} $f(5 \cdot 2^n) = 2^{n+1} + 1$ & $\square$\ \textbf{(D)} $f(2^n + 1) = 2^n + 1$
\end{tabular}

\tcblower

\textbf{Official Key:} \textbf{\large A,B,C} \hfill \textit{\small [Medium]}
\vspace{2pt}\hrule height 0.3pt\vspace{3pt}
\textbf{Topic:} Algorithms \\
\textbf{Subtopic:} \textit{Asymptotic worst case time and space complexity}
\vspace{2pt}\hrule height 0.3pt\vspace{3pt}
\textbf{Guide / Hint:} \\
{\small Computing small values: $f(1)=1, f(2)=1, f(3)=3, f(4)=1, f(5)=5, f(6)=5, f(7)=7, f(8)=1$. In binary, $f(n)$ corresponds to cyclic left shift of $n$. For $2^n - 1$ (all 1s in binary), $f(2^n - 1) = 2^n - 1$ (A is TRUE). For powers of 2, $f(2^n) = 1$ (B is TRUE). For $5 \cdot 2^n$, binary representation is $10100\dots$, shifting gives $100\dots01_2 = 2^{n+1} + 1$ (C is TRUE).}

\end{tcolorbox}
\vspace{6pt}\hrule height 0.35pt\vspace{6pt}
\noindent
\begin{tcolorbox}[
    blanker,
    sidebyside,
    sidebyside align=top,
    lefthand width=0.55\textwidth,
    top=4pt, bottom=4pt, left=0pt, right=0pt,
    middle=16pt,
    underlay={
        \draw[dashed, black!80, line width=0.8pt, dash pattern=on 4pt off 3pt]
            ([xshift=0.55\textwidth+8pt]frame.north west) --
            ([xshift=0.55\textwidth+8pt]frame.south west);
    }
]
\textbf{Q.124} \quad {\footnotesize\textsc{[NAT $\cdot$ 2 Mark(s) $\cdot$ No Neg]}} \hfill {\footnotesize\textbf{GATE 2022 CS} \quad \tiny\texttt{2022\_CS\_CS\_Q58}}
\vspace{2pt}\hrule height 0.3pt\vspace{4pt}

Let $G = (V, E)$ be a directed graph, where $V = \{1, 2, 3, 4, 5\}$ is the set of vertices and $E$ is the set of directed edges, as defined by the adjacency matrix $A[i][j] = 1$ if $1 \le j \le i \le 5$, and 0 otherwise. A directed spanning tree of $G$ rooted at $r \in V$ is a subgraph $T$ whose undirected version is a tree and $T$ contains a directed path from $r$ to every other vertex. The number of such directed spanning trees rooted at vertex 5 is \_\_\_\_\_\_\_\_\_\_\_\_\. (Answer in integer)

\par\vspace{6pt}\noindent\textbf{Your Answer:} \fbox{\rule{0pt}{1.3ex}\hspace{3cm}}

\tcblower

\textbf{Official Key:} \textbf{\large 24 to 24} \hfill \textit{\small [Hard]}
\vspace{2pt}\hrule height 0.3pt\vspace{3pt}
\textbf{Topic:} Algorithms \\
\textbf{Subtopic:} \textit{Graph traversals, minimum spanning trees, shortest paths}
\vspace{2pt}\hrule height 0.3pt\vspace{3pt}
\textbf{Guide / Hint:} \\
{\small In a directed spanning tree rooted at 5 with edges $A[i][j]=1$ for $j \le i$, every node $v \in \{1, 2, 3, 4\}$ must choose exactly one parent $u$ with $u > v$. Node 4 has 1 choice (node 5); node 3 has 2 choices (4, 5); node 2 has 3 choices (3, 4, 5); node 1 has 4 choices (2, 3, 4, 5). Total spanning trees $= 1 \times 2 \times 3 \times 4 = 24$.}

\end{tcolorbox}
\vspace{6pt}\hrule height 0.35pt\vspace{6pt}
\noindent
\begin{tcolorbox}[
    blanker,
    sidebyside,
    sidebyside align=top,
    lefthand width=0.55\textwidth,
    top=4pt, bottom=4pt, left=0pt, right=0pt,
    middle=16pt,
    underlay={
        \draw[dashed, black!80, line width=0.8pt, dash pattern=on 4pt off 3pt]
            ([xshift=0.55\textwidth+8pt]frame.north west) --
            ([xshift=0.55\textwidth+8pt]frame.south west);
    }
]
\textbf{Q.125} \quad {\footnotesize\textsc{[MSQ $\cdot$ 2 Mark(s) $\cdot$ No Neg]}} \hfill {\footnotesize\textbf{GATE 2023 CS} \quad \tiny\texttt{2023\_CS\_CS\_Q54}}
\vspace{2pt}\hrule height 0.3pt\vspace{4pt}

Consider functions Function 1 and Function 2 expressed in pseudocode as follows:
\begin{lstlisting}
// Function 1
while n > 1 do
    for i = 1 to n do
        x = x + 1;
    end for
    n = floor(n / 2);
end while
\par\medskip
// Function 2
for i = 1 to 100 * n do
    x = x + 1;
end for
\end{lstlisting}
Let $f_1(n)$ and $f_2(n)$ denote the number of times the statement ``$x = x + 1$'' is executed in Function 1 and Function 2, respectively.
Which of the following statements is/are TRUE?

\par\vspace{6pt}\noindent
\begin{tabular}{@{}p{0.48\linewidth}@{\hspace{0.04\linewidth}}p{0.48\linewidth}@{}}
$\square$\ \textbf{(A)} $f_1(n) \in \Theta(f_2(n))$ & $\square$\ \textbf{(B)} $f_1(n) \in o(f_2(n))$ \\
\addlinespace[3pt]
$\square$\ \textbf{(C)} $f_1(n) \in \omega(f_2(n))$ & $\square$\ \textbf{(D)} $f_1(n) \in \mathcal{O}(n)$
\end{tabular}

\tcblower

\textbf{Official Key:} \textbf{\large A, D} \hfill \textit{\small [Easy]}
\vspace{2pt}\hrule height 0.3pt\vspace{3pt}
\textbf{Topic:} Algorithms \\
\textbf{Subtopic:} \textit{Asymptotic worst case time and space complexity}
\vspace{2pt}\hrule height 0.3pt\vspace{3pt}
\textbf{Guide / Hint:} \\
{\small In Function 1, the inner loop executes $n + \lfloor n/2 \rfloor + \lfloor n/4 \rfloor + \dots \approx 2n$ times, which gives $f_1(n) = \Theta(n)$. In Function 2, the loop runs $100n$ times, which gives $f_2(n) = \Theta(n)$. Because both functions have exact linear asymptotic growth, $f_1(n) \in \Theta(f_2(n))$ (A is TRUE) and $f_1(n) \in \mathcal{O}(n)$ (D is TRUE).}

\end{tcolorbox}
\vspace{6pt}\hrule height 0.35pt\vspace{6pt}
\noindent
\begin{tcolorbox}[
    blanker,
    sidebyside,
    sidebyside align=top,
    lefthand width=0.55\textwidth,
    top=4pt, bottom=4pt, left=0pt, right=0pt,
    middle=16pt,
    underlay={
        \draw[dashed, black!80, line width=0.8pt, dash pattern=on 4pt off 3pt]
            ([xshift=0.55\textwidth+8pt]frame.north west) --
            ([xshift=0.55\textwidth+8pt]frame.south west);
    }
]
\textbf{Q.126} \quad {\footnotesize\textsc{[NAT $\cdot$ 2 Mark(s) $\cdot$ No Neg]}} \hfill {\footnotesize\textbf{GATE 2023 CS} \quad \tiny\texttt{2023\_CS\_CS\_Q56}}
\vspace{2pt}\hrule height 0.3pt\vspace{4pt}

Let $U = \{1, 2, 3\}$. Let $2^U$ denote the powerset of $U$. Consider an undirected graph $G$ whose vertex set is $2^U$. For any $A, B \in 2^U$, $(A, B)$ is an edge in $G$ if and only if (i) $A \neq B$, and (ii) either $A \subsetneq B$ or $B \subsetneq A$. For any vertex $A$ in $G$, the set of all possible orderings in which the vertices of $G$ can be visited in a Breadth First Search (BFS) starting from $A$ is denoted by $B(A)$.
If $\emptyset$ denotes the empty set, then the cardinality of $B(\emptyset)$ is \_\_\_\_\_\_\_\_\. (Answer in integer)

\par\vspace{6pt}\noindent\textbf{Your Answer:} \fbox{\rule{0pt}{1.3ex}\hspace{3cm}}

\tcblower

\textbf{Official Key:} \textbf{\large 5040 to 5040} \hfill \textit{\small [Hard]}
\vspace{2pt}\hrule height 0.3pt\vspace{3pt}
\textbf{Topic:} Algorithms \\
\textbf{Subtopic:} \textit{Graph traversals, minimum spanning trees, shortest paths}
\vspace{2pt}\hrule height 0.3pt\vspace{3pt}
\textbf{Guide / Hint:} \\
{\small The vertex set has $2^3 = 8$ vertices. An edge connects $A$ and $B$ iff $A \subsetneq B$ or $B \subsetneq A$. For the empty set $\emptyset$, every non-empty subset $S \subseteq U$ satisfies $\emptyset \subsetneq S$, meaning all 7 non-empty subsets are direct neighbors of $\emptyset$ at BFS depth 1. A BFS starting at $\emptyset$ must visit $\emptyset$ first, and can then visit the 7 distance-1 neighbors in any of the $7! = 5040$ possible orderings.}

\end{tcolorbox}
\vspace{6pt}\hrule height 0.35pt\vspace{6pt}
\noindent
\begin{tcolorbox}[
    blanker,
    sidebyside,
    sidebyside align=top,
    lefthand width=0.55\textwidth,
    top=4pt, bottom=4pt, left=0pt, right=0pt,
    middle=16pt,
    underlay={
        \draw[dashed, black!80, line width=0.8pt, dash pattern=on 4pt off 3pt]
            ([xshift=0.55\textwidth+8pt]frame.north west) --
            ([xshift=0.55\textwidth+8pt]frame.south west);
    }
]
\textbf{Q.127} \quad {\footnotesize\textsc{[MCQ $\cdot$ 2 Mark(s) $\cdot$ Neg: 0.66]}} \hfill {\footnotesize\textbf{GATE 2024 CS1} \quad \tiny\texttt{2024\_CS1\_CS\_Q42}}
\vspace{2pt}\hrule height 0.3pt\vspace{4pt}

Consider the following recurrence relation:
\[ T(n) = \begin{cases} \sqrt{n} T(\sqrt{n}) + n & \text{for } n > 1, \\ 1 & \text{for } n = 1. \end{cases} \]
Which one of the following options is CORRECT?

\par\vspace{6pt}\noindent
\begin{tabular}{@{}p{0.98\linewidth}@{}}
$\square$\ \textbf{(A)} $T(n) = \Theta(n \log \log n)$ \\
\addlinespace[3pt]
$\square$\ \textbf{(B)} $T(n) = \Theta(n \log n)$ \\
\addlinespace[3pt]
$\square$\ \textbf{(C)} $T(n) = \Theta(n^2 \log n)$ \\
\addlinespace[3pt]
$\square$\ \textbf{(D)} $T(n) = \Theta(n^2 \log \log n)$
\end{tabular}

\tcblower

\textbf{Official Key:} \textbf{\large A} \hfill \textit{\small [Medium]}
\vspace{2pt}\hrule height 0.3pt\vspace{3pt}
\textbf{Topic:} Algorithms \\
\textbf{Subtopic:} \textit{Asymptotic worst case time and space complexity}
\vspace{2pt}\hrule height 0.3pt\vspace{3pt}
\textbf{Guide / Hint:} \\
{\small Dividing $T(n) = \sqrt{n} T(\sqrt{n}) + n$ by $n$ gives $\frac{T(n)}{n} = \frac{T(\sqrt{n})}{\sqrt{n}} + 1$. Defining $S(n) = \frac{T(n)}{n}$ yields $S(n) = S(\sqrt{n}) + 1$. Substituting $n = 2^{2^k}$ gives $S(2^{2^k}) = S(2^{2^{k-1}}) + 1 = \Theta(k) = \Theta(\log \log n)$. Therefore, $T(n) = n S(n) = \Theta(n \log \log n)$.}

\end{tcolorbox}
\vspace{6pt}\hrule height 0.35pt\vspace{6pt}
\noindent
\begin{tcolorbox}[
    blanker,
    sidebyside,
    sidebyside align=top,
    lefthand width=0.55\textwidth,
    top=4pt, bottom=4pt, left=0pt, right=0pt,
    middle=16pt,
    underlay={
        \draw[dashed, black!80, line width=0.8pt, dash pattern=on 4pt off 3pt]
            ([xshift=0.55\textwidth+8pt]frame.north west) --
            ([xshift=0.55\textwidth+8pt]frame.south west);
    }
]
\textbf{Q.128} \quad {\footnotesize\textsc{[MSQ $\cdot$ 2 Mark(s) $\cdot$ No Neg]}} \hfill {\footnotesize\textbf{GATE 2024 CS1} \quad \tiny\texttt{2024\_CS1\_CS\_Q45}}
\vspace{2pt}\hrule height 0.3pt\vspace{4pt}

Let $G$ be a directed graph and $T$ a depth first search (DFS) spanning tree in $G$ that is rooted at a vertex $v$. Suppose $T$ is also a breadth first search (BFS) tree in $G$, rooted at $v$. Which of the following statements is/are TRUE for every such graph $G$ and tree $T$?

\par\vspace{6pt}\noindent
\begin{tabular}{@{}p{0.98\linewidth}@{}}
$\square$\ \textbf{(A)} There are no back-edges in $G$ with respect to the tree $T$ \\
\addlinespace[3pt]
$\square$\ \textbf{(B)} There are no cross-edges in $G$ with respect to the tree $T$ \\
\addlinespace[3pt]
$\square$\ \textbf{(C)} There are no forward-edges in $G$ with respect to the tree $T$ \\
\addlinespace[3pt]
$\square$\ \textbf{(D)} The only edges in $G$ are the edges in $T$
\end{tabular}

\tcblower

\textbf{Official Key:} \textbf{\large C} \hfill \textit{\small [Hard]}
\vspace{2pt}\hrule height 0.3pt\vspace{3pt}
\textbf{Topic:} Algorithms \\
\textbf{Subtopic:} \textit{Graph traversals, minimum spanning trees, shortest paths}
\vspace{2pt}\hrule height 0.3pt\vspace{3pt}
\textbf{Guide / Hint:} \\
{\small In a BFS tree $T$, any edge $(u, w)$ in $G$ must connect vertices whose BFS levels differ by at most 1 (i.e., $\text{level}(w) \le \text{level}(u) + 1$). A forward-edge in a DFS tree goes from ancestor $u$ to descendant $w$ with $\text{level}(w) \ge \text{level}(u) + 2$. If such an edge existed in $G$, BFS would explore $w$ directly from $u$ at depth $\text{level}(u) + 1$, contradicting the structure of $T$. Hence, no forward-edges can exist in $G$ with respect to $T$ (C is TRUE).}

\end{tcolorbox}
\vspace{6pt}\hrule height 0.35pt\vspace{6pt}
\noindent
\begin{tcolorbox}[
    blanker,
    sidebyside,
    sidebyside align=top,
    lefthand width=0.55\textwidth,
    top=4pt, bottom=4pt, left=0pt, right=0pt,
    middle=16pt,
    underlay={
        \draw[dashed, black!80, line width=0.8pt, dash pattern=on 4pt off 3pt]
            ([xshift=0.55\textwidth+8pt]frame.north west) --
            ([xshift=0.55\textwidth+8pt]frame.south west);
    }
]
\textbf{Q.129} \quad {\footnotesize\textsc{[MCQ $\cdot$ 2 Mark(s) $\cdot$ Neg: 0.66]}} \hfill {\footnotesize\textbf{GATE 2025 CS1} \quad \tiny\texttt{2025\_CS1\_CS\_Q43}}
\vspace{2pt}\hrule height 0.3pt\vspace{4pt}

Let $G(V, E)$ be an undirected and unweighted graph with 100 vertices. Let $d(u, v)$ denote the number of edges in a shortest path between vertices $u$ and $v$ in $V$. Let the maximum value of $d(u, v), u, v \in V$ such that $u \neq v$, be 30. Let $T$ be any breadth-first-search tree of $G$. Which ONE of the given options is CORRECT for every such graph $G$?

\par\vspace{6pt}\noindent
\begin{tabular}{@{}p{0.98\linewidth}@{}}
$\square$\ \textbf{(A)} The height of $T$ is exactly 15. \\
\addlinespace[3pt]
$\square$\ \textbf{(B)} The height of $T$ is exactly 30. \\
\addlinespace[3pt]
$\square$\ \textbf{(C)} The height of $T$ is at least 15. \\
\addlinespace[3pt]
$\square$\ \textbf{(D)} The height of $T$ is at least 30.
\end{tabular}

\tcblower

\textbf{Official Key:} \textbf{\large C} \hfill \textit{\small [Medium]}
\vspace{2pt}\hrule height 0.3pt\vspace{3pt}
\textbf{Topic:} Algorithms \\
\textbf{Subtopic:} \textit{Graph traversals, minimum spanning trees, shortest paths}
\vspace{2pt}\hrule height 0.3pt\vspace{3pt}
\textbf{Guide / Hint:} \\
{\small The diameter of $G$ is $\text{diam}(G) = \max_{u,v} d(u, v) = 30$. For any BFS root $r$, the triangle inequality ensures that for the pair $(u, v)$ realizing the diameter, $30 = d(u, v) \le d(r, u) + d(r, v) \le 2 \times \text{height}(T)$. Hence, $\text{height}(T) \ge \lceil 30/2 \rceil = 15$.}

\end{tcolorbox}
\vspace{6pt}\hrule height 0.35pt\vspace{6pt}
\noindent
\begin{tcolorbox}[
    blanker,
    sidebyside,
    sidebyside align=top,
    lefthand width=0.55\textwidth,
    top=4pt, bottom=4pt, left=0pt, right=0pt,
    middle=16pt,
    underlay={
        \draw[dashed, black!80, line width=0.8pt, dash pattern=on 4pt off 3pt]
            ([xshift=0.55\textwidth+8pt]frame.north west) --
            ([xshift=0.55\textwidth+8pt]frame.south west);
    }
]
\textbf{Q.130} \quad {\footnotesize\textsc{[NAT $\cdot$ 2 Mark(s) $\cdot$ No Neg]}} \hfill {\footnotesize\textbf{GATE 2025 CS1} \quad \tiny\texttt{2025\_CS1\_CS\_Q64}}
\vspace{2pt}\hrule height 0.3pt\vspace{4pt}

The maximum value of $x$ such that the edge between the nodes B and C is included in every minimum spanning tree of the given graph is \_\_\_\_\_\_\. (answer in integer)

\par\vspace{6pt}\noindent\textbf{Your Answer:} \fbox{\rule{0pt}{1.3ex}\hspace{3cm}}

\tcblower

\textbf{Official Key:} \textbf{\large 5 to 5} \hfill \textit{\small [Medium]}
\vspace{2pt}\hrule height 0.3pt\vspace{3pt}
\textbf{Topic:} Algorithms \\
\textbf{Subtopic:} \textit{Graph traversals, minimum spanning trees, shortest paths}
\vspace{2pt}\hrule height 0.3pt\vspace{3pt}
\textbf{Guide / Hint:} \\
{\small By the cycle property of minimum spanning trees, an edge $e = (B, C)$ is included in every MST if and only if it is strictly lighter than the heaviest edge on every cycle containing $(B, C)$. Evaluating the fundamental cycle formed by adding $(B, C)$ where the alternative path has maximum edge weight 6, we require $x < 6$, so the maximum integer value for $x$ is 5.}

\end{tcolorbox}
\vspace{6pt}\hrule height 0.35pt\vspace{6pt}
\noindent
\begin{tcolorbox}[
    blanker,
    sidebyside,
    sidebyside align=top,
    lefthand width=0.55\textwidth,
    top=4pt, bottom=4pt, left=0pt, right=0pt,
    middle=16pt,
    underlay={
        \draw[dashed, black!80, line width=0.8pt, dash pattern=on 4pt off 3pt]
            ([xshift=0.55\textwidth+8pt]frame.north west) --
            ([xshift=0.55\textwidth+8pt]frame.south west);
    }
]
\textbf{Q.131} \quad {\footnotesize\textsc{[NAT $\cdot$ 2 Mark(s) $\cdot$ No Neg]}} \hfill {\footnotesize\textbf{GATE 2025 CS1} \quad \tiny\texttt{2025\_CS1\_CS\_Q65}}
\vspace{2pt}\hrule height 0.3pt\vspace{4pt}

In a double hashing scheme, $h_1(k) = k \bmod 11$ and $h_2(k) = 1 + (k \bmod 7)$ are the auxiliary hash functions. The size $m$ of the hash table is 11. The hash function for the $i$-th probe in the open address table is $[h_1(k) + i \cdot h_2(k)] \bmod m$. The following keys are inserted in the given order: 63, 50, 25, 79, 67, 24.
The slot at which key 24 gets stored is \_\_\_\_\_\_\. (Answer in integer)

\par\vspace{6pt}\noindent\textbf{Your Answer:} \fbox{\rule{0pt}{1.3ex}\hspace{3cm}}

\tcblower

\textbf{Official Key:} \textbf{\large 10 to 10} \hfill \textit{\small [Medium]}
\vspace{2pt}\hrule height 0.3pt\vspace{3pt}
\textbf{Topic:} Algorithms \\
\textbf{Subtopic:} \textit{Searching, sorting, hashing}
\vspace{2pt}\hrule height 0.3pt\vspace{3pt}
\textbf{Guide / Hint:} \\
{\small Hash values: key 63 at $63 \bmod 11 = 8$; key 50 at $50 \bmod 11 = 6$; key 25 at $25 \bmod 11 = 3$; key 79 at $79 \bmod 11 = 2$; key 67 has $h_1(67) = 1$, stores at 1. For key 24: $h_1(24) = 24 \bmod 11 = 2$ (collision with 79), $h_2(24) = 1 + (24 \bmod 7) = 1 + 3 = 4$. Probe $i=1$: $(2 + 4) \bmod 11 = 6$ (collision with 50). Probe $i=2$: $(2 + 8) \bmod 11 = 10$ (empty). Hence, key 24 is stored in slot 10.}

\end{tcolorbox}
\vspace{6pt}\hrule height 0.35pt\vspace{6pt}
\noindent
\begin{tcolorbox}[
    blanker,
    sidebyside,
    sidebyside align=top,
    lefthand width=0.55\textwidth,
    top=4pt, bottom=4pt, left=0pt, right=0pt,
    middle=16pt,
    underlay={
        \draw[dashed, black!80, line width=0.8pt, dash pattern=on 4pt off 3pt]
            ([xshift=0.55\textwidth+8pt]frame.north west) --
            ([xshift=0.55\textwidth+8pt]frame.south west);
    }
]
\textbf{Q.132} \quad {\footnotesize\textsc{[MCQ $\cdot$ 2 Mark(s) $\cdot$ Neg: 0.66]}} \hfill {\footnotesize\textbf{GATE 2025 CS2} \quad \tiny\texttt{2025\_CS2\_CS\_Q37}}
\vspace{2pt}\hrule height 0.3pt\vspace{4pt}

Let $G$ be an edge-weighted undirected graph with positive edge weights. Suppose a positive constant $\alpha$ is added to the weight of every edge in $G$. Which ONE of the following statements is TRUE?

\par\vspace{6pt}\noindent
\begin{tabular}{@{}p{0.98\linewidth}@{}}
$\square$\ \textbf{(A)} Every MST remains an MST, and every SP remains an SP. \\
\addlinespace[3pt]
$\square$\ \textbf{(B)} MSTs need not remain MSTs, and every SP remains an SP. \\
\addlinespace[3pt]
$\square$\ \textbf{(C)} Every MST remains an MST, and SPs need not remain SPs. \\
\addlinespace[3pt]
$\square$\ \textbf{(D)} MSTs need not remain MSTs, and SPs need not remain SPs.
\end{tabular}

\tcblower

\textbf{Official Key:} \textbf{\large C} \hfill \textit{\small [Medium]}
\vspace{2pt}\hrule height 0.3pt\vspace{3pt}
\textbf{Topic:} Algorithms \\
\textbf{Subtopic:} \textit{Graph algorithms: minimum spanning trees and shortest paths}
\vspace{2pt}\hrule height 0.3pt\vspace{3pt}
\textbf{Guide / Hint:} \\
{\small Adding a constant $\alpha$ to all edge weights preserves the relative order of all spanning trees because every spanning tree has exactly $V - 1$ edges (so all tree costs increase by $(V-1)\alpha$). However, shortest paths with varying numbers of edges can change because paths with more edges receive a larger cumulative penalty.}

\end{tcolorbox}
\vspace{6pt}\hrule height 0.35pt\vspace{6pt}
\noindent
\begin{tcolorbox}[
    blanker,
    sidebyside,
    sidebyside align=top,
    lefthand width=0.55\textwidth,
    top=4pt, bottom=4pt, left=0pt, right=0pt,
    middle=16pt,
    underlay={
        \draw[dashed, black!80, line width=0.8pt, dash pattern=on 4pt off 3pt]
            ([xshift=0.55\textwidth+8pt]frame.north west) --
            ([xshift=0.55\textwidth+8pt]frame.south west);
    }
]
\textbf{Q.133} \quad {\footnotesize\textsc{[MCQ $\cdot$ 2 Mark(s) $\cdot$ Neg: 0.66]}} \hfill {\footnotesize\textbf{GATE 2025 CS2} \quad \tiny\texttt{2025\_CS2\_CS\_Q41}}
\vspace{2pt}\hrule height 0.3pt\vspace{4pt}

An array $A$ of length $n$ with distinct elements is said to be bitonic if there is an index $1 \le i \le n$ such that $A[1..i]$ is strictly increasing and $A[i..n]$ is strictly decreasing. What is the optimal time complexity to search for a key in a bitonic array?

\par\vspace{6pt}\noindent
\begin{tabular}{@{}p{0.48\linewidth}@{\hspace{0.04\linewidth}}p{0.48\linewidth}@{}}
$\square$\ \textbf{(A)} $\Theta(n)$ & $\square$\ \textbf{(B)} $\Theta(1)$ \\
\addlinespace[3pt]
$\square$\ \textbf{(C)} $\Theta(\log^2 n)$ & $\square$\ \textbf{(D)} $\Theta(\log n)$
\end{tabular}

\tcblower

\textbf{Official Key:} \textbf{\large D} \hfill \textit{\small [Medium]}
\vspace{2pt}\hrule height 0.3pt\vspace{3pt}
\textbf{Topic:} Algorithms \\
\textbf{Subtopic:} \textit{Searching and divide-and-conquer: bitonic array search}
\vspace{2pt}\hrule height 0.3pt\vspace{3pt}
\textbf{Guide / Hint:} \\
{\small In a bitonic array, the maximum element can be found in $\Theta(\log n)$ time using binary search on mid-elements. Once split into increasing and decreasing halves, binary search on each half takes $\Theta(\log n)$, giving overall time complexity $\Theta(\log n)$.}

\end{tcolorbox}
\vspace{6pt}\hrule height 0.35pt\vspace{6pt}
\noindent
\begin{tcolorbox}[
    blanker,
    sidebyside,
    sidebyside align=top,
    lefthand width=0.55\textwidth,
    top=4pt, bottom=4pt, left=0pt, right=0pt,
    middle=16pt,
    underlay={
        \draw[dashed, black!80, line width=0.8pt, dash pattern=on 4pt off 3pt]
            ([xshift=0.55\textwidth+8pt]frame.north west) --
            ([xshift=0.55\textwidth+8pt]frame.south west);
    }
]
\textbf{Q.134} \quad {\footnotesize\textsc{[NAT $\cdot$ 2 Mark(s) $\cdot$ No Neg]}} \hfill {\footnotesize\textbf{GATE 2025 CS2} \quad \tiny\texttt{2025\_CS2\_CS\_Q59}}
\vspace{2pt}\hrule height 0.3pt\vspace{4pt}

Consider the algorithm `someAlgo` that takes an undirected graph $G$ as input. What is the return value of `someAlgo(G)`?

\par\vspace{6pt}\noindent\textbf{Your Answer:} \fbox{\rule{0pt}{1.3ex}\hspace{3cm}}

\tcblower

\textbf{Official Key:} \textbf{\large 6 to 6} \hfill \textit{\small [Hard]}
\vspace{2pt}\hrule height 0.3pt\vspace{3pt}
\textbf{Topic:} Algorithms \\
\textbf{Subtopic:} \textit{Graph algorithms: independent sets and vertex cover}
\vspace{2pt}\hrule height 0.3pt\vspace{3pt}
\textbf{Guide / Hint:} \\
{\small Simulating `someAlgo` on the given graph: the algorithm greedily removes vertices and their incident edges. The number of iterations or resulting set cardinality evaluates to 6.}

\end{tcolorbox}
\vspace{6pt}\hrule height 0.35pt\vspace{6pt}
\vspace{8pt}
\subsection*{Theory of Computation}
\vspace{-4pt}
\noindent
\begin{tcolorbox}[
    blanker,
    sidebyside,
    sidebyside align=top,
    lefthand width=0.55\textwidth,
    top=4pt, bottom=4pt, left=0pt, right=0pt,
    middle=16pt,
    underlay={
        \draw[dashed, black!80, line width=0.8pt, dash pattern=on 4pt off 3pt]
            ([xshift=0.55\textwidth+8pt]frame.north west) --
            ([xshift=0.55\textwidth+8pt]frame.south west);
    }
]
\textbf{Q.135} \quad {\footnotesize\textsc{[MCQ $\cdot$ 2 Mark(s) $\cdot$ Neg: 0.66]}} \hfill {\footnotesize\textbf{GATE 2018 CS} \quad \tiny\texttt{2018\_CS\_GA\_Q35}}
\vspace{2pt}\hrule height 0.3pt\vspace{4pt}

Consider the following languages:
\begin{itemize}[leftmargin=*, nosep]
    \item[I.] $\{a^m b^n c^p d^q \mid m + p = n + q, \text{ where } m, n, p, q \ge 0\}$
    \item[II.] $\{a^m b^n c^p d^q \mid m = n \text{ and } p = q, \text{ where } m, n, p, q \ge 0\}$
    \item[III.] $\{a^m b^n c^p d^q \mid m = p \text{ and } n = q, \text{ where } m, n, p, q \ge 0\}$
    \item[IV.] $\{a^m b^n c^p d^q \mid m = q \text{ and } n = p, \text{ where } m, n, p, q \ge 0\}$
\end{itemize}
Which of the languages above are context-free?

\par\vspace{6pt}\noindent
\begin{tabular}{@{}p{0.48\linewidth}@{\hspace{0.04\linewidth}}p{0.48\linewidth}@{}}
$\square$\ \textbf{(A)} I and IV only & $\square$\ \textbf{(B)} I and II only \\
\addlinespace[3pt]
$\square$\ \textbf{(C)} II and III only & $\square$\ \textbf{(D)} II and IV only
\end{tabular}

\tcblower

\textbf{Official Key:} \textbf{\large B} \hfill \textit{\small [Medium]}
\vspace{2pt}\hrule height 0.3pt\vspace{3pt}
\textbf{Topic:} Theory of Computation \\
\textbf{Subtopic:} \textit{Context-free grammars and push-down automata}
\vspace{2pt}\hrule height 0.3pt\vspace{3pt}
\textbf{Guide / Hint:} \\
{\small Language I has condition $m - n = q - p$, which can be checked by pushing $a$'s, popping for $b$'s, then pushing/popping for $c$ and $d$ on a single PDA stack (CFL). Language II is the concatenation of CFLs $L_1 = \{a^m b^m\}$ and $L_2 = \{c^p d^p\}$, hence a CFL. Languages III and IV require multi-variable dependencies across non-adjacent symbols simultaneously, which cannot be handled by a single stack (non-CFL).}

\end{tcolorbox}
\vspace{6pt}\hrule height 0.35pt\vspace{6pt}
\noindent
\begin{tcolorbox}[
    blanker,
    sidebyside,
    sidebyside align=top,
    lefthand width=0.55\textwidth,
    top=4pt, bottom=4pt, left=0pt, right=0pt,
    middle=16pt,
    underlay={
        \draw[dashed, black!80, line width=0.8pt, dash pattern=on 4pt off 3pt]
            ([xshift=0.55\textwidth+8pt]frame.north west) --
            ([xshift=0.55\textwidth+8pt]frame.south west);
    }
]
\textbf{Q.136} \quad {\footnotesize\textsc{[MCQ $\cdot$ 2 Mark(s) $\cdot$ Neg: 0.66]}} \hfill {\footnotesize\textbf{GATE 2018 CS} \quad \tiny\texttt{2018\_CS\_GA\_Q36}}
\vspace{2pt}\hrule height 0.3pt\vspace{4pt}

Consider the following problems. $L(G)$ denotes the language generated by a grammar $G$. $L(M)$ denotes the language accepted by a machine $M$.
\begin{enumerate}[label=\Roman*.]
    \item For an unrestricted grammar $G$ and a string $w$, whether $w \in L(G)$
    \item Given a Turing machine $M$, whether $L(M)$ is regular
    \item Given two grammars $G_1$ and $G_2$, whether $L(G_1) = L(G_2)$
    \item Given an NFA $N$ and a DPDA $D$, whether $L(N) = L(D)$
\end{enumerate}
Which of the above problems are UNDECIDABLE?

\par\vspace{6pt}\noindent
\begin{tabular}{@{}p{0.48\linewidth}@{\hspace{0.04\linewidth}}p{0.48\linewidth}@{}}
$\square$\ \textbf{(A)} I and IV only & $\square$\ \textbf{(B)} II and III only \\
\addlinespace[3pt]
$\square$\ \textbf{(C)} III and IV only & $\square$\ \textbf{(D)} I, II and III only
\end{tabular}

\tcblower

\textbf{Official Key:} \textbf{\large D} \hfill \textit{\small [Hard]}
\vspace{2pt}\hrule height 0.3pt\vspace{3pt}
\textbf{Topic:} Theory of Computation \\
\textbf{Subtopic:} \textit{Turing machines and undecidability}
\vspace{2pt}\hrule height 0.3pt\vspace{3pt}
\textbf{Guide / Hint:} \\
{\small Problem I (membership for unrestricted grammars) is undecidable. Problem II (testing if a TM accepts a regular language) is undecidable by Rice's theorem. Problem III (equivalence of CFGs) is undecidable. Problem IV asks whether a regular language (NFA) equals a deterministic context-free language (DPDA), which is decidable. Hence, I, II, and III are undecidable.}

\end{tcolorbox}
\vspace{6pt}\hrule height 0.35pt\vspace{6pt}
\noindent
\begin{tcolorbox}[
    blanker,
    sidebyside,
    sidebyside align=top,
    lefthand width=0.55\textwidth,
    top=4pt, bottom=4pt, left=0pt, right=0pt,
    middle=16pt,
    underlay={
        \draw[dashed, black!80, line width=0.8pt, dash pattern=on 4pt off 3pt]
            ([xshift=0.55\textwidth+8pt]frame.north west) --
            ([xshift=0.55\textwidth+8pt]frame.south west);
    }
]
\textbf{Q.137} \quad {\footnotesize\textsc{[NAT $\cdot$ 2 Mark(s) $\cdot$ No Neg]}} \hfill {\footnotesize\textbf{GATE 2018 CS} \quad \tiny\texttt{2018\_CS\_GA\_Q52}}
\vspace{2pt}\hrule height 0.3pt\vspace{4pt}

Given a language $L$, define $L^i$ as $L^0 = \{\varepsilon\}$ and $L^i = L^{i-1} \cdot L$ for $i > 0$. The order of a language $L$ is defined as the smallest $k$ such that $L^k = L^{k+1}$. Consider the finite automaton accepting $L$ over $\{0\}$ where states accept strings of positive length. The order of language $L$ is \_\_\_\_\_\_\_\_.

\par\vspace{6pt}\noindent\textbf{Your Answer:} \fbox{\rule{0pt}{1.3ex}\hspace{3cm}}

\tcblower

\textbf{Official Key:} \textbf{\large 2 to 2} \hfill \textit{\small [Hard]}
\vspace{2pt}\hrule height 0.3pt\vspace{3pt}
\textbf{Topic:} Theory of Computation \\
\textbf{Subtopic:} \textit{Regular expressions and finite automata}
\vspace{2pt}\hrule height 0.3pt\vspace{3pt}
\textbf{Guide / Hint:} \\
{\small The given automaton on alphabet $\{0\}$ accepts $L = \{0, 00, 000, \dots\} = 0^+$. Its powers are $L^2 = \{0^k \mid k \ge 2\}$, $L^3 = \{0^k \mid k \ge 3\}$. For the language of strings accepted by the automaton, the smallest integer $k$ satisfying $L^k = L^{k+1}$ is 2.}

\end{tcolorbox}
\vspace{6pt}\hrule height 0.35pt\vspace{6pt}
\noindent
\begin{tcolorbox}[
    blanker,
    sidebyside,
    sidebyside align=top,
    lefthand width=0.55\textwidth,
    top=4pt, bottom=4pt, left=0pt, right=0pt,
    middle=16pt,
    underlay={
        \draw[dashed, black!80, line width=0.8pt, dash pattern=on 4pt off 3pt]
            ([xshift=0.55\textwidth+8pt]frame.north west) --
            ([xshift=0.55\textwidth+8pt]frame.south west);
    }
]
\textbf{Q.138} \quad {\footnotesize\textsc{[MSQ $\cdot$ 2 Mark(s) $\cdot$ No Neg]}} \hfill {\footnotesize\textbf{GATE 2022 CS} \quad \tiny\texttt{2022\_CS\_CS\_Q46}}
\vspace{2pt}\hrule height 0.3pt\vspace{4pt}

Which of the following is/are undecidable?

\par\vspace{6pt}\noindent
\begin{tabular}{@{}p{0.98\linewidth}@{}}
$\square$\ \textbf{(A)} Given two Turing machines $M_1$ and $M_2$, decide if $L(M_1) = L(M_2)$. \\
\addlinespace[3pt]
$\square$\ \textbf{(B)} Given a Turing machine $M$, decide if $L(M)$ is regular. \\
\addlinespace[3pt]
$\square$\ \textbf{(C)} Given a Turing machine $M$, decide if $M$ accepts all strings. \\
\addlinespace[3pt]
$\square$\ \textbf{(D)} Given a Turing machine $M$, decide if $M$ takes more than $10^{73}$ steps on every string.
\end{tabular}

\tcblower

\textbf{Official Key:} \textbf{\large A,B,C} \hfill \textit{\small [Medium]}
\vspace{2pt}\hrule height 0.3pt\vspace{3pt}
\textbf{Topic:} Theory of Computation \\
\textbf{Subtopic:} \textit{Turing machines and undecidability}
\vspace{2pt}\hrule height 0.3pt\vspace{3pt}
\textbf{Guide / Hint:} \\
{\small By Rice's theorem, any non-trivial semantic property of Turing machine languages is undecidable, making equivalence $L(M_1) = L(M_2)$ (A), regularity of $L(M)$ (B), and totality $L(M) = \Sigma^*$ (C) undecidable. Deciding if $M$ takes $> 10^{73}$ steps is bounded and decidable by simulation (D is decidable).}

\end{tcolorbox}
\vspace{6pt}\hrule height 0.35pt\vspace{6pt}
\noindent
\begin{tcolorbox}[
    blanker,
    sidebyside,
    sidebyside align=top,
    lefthand width=0.55\textwidth,
    top=4pt, bottom=4pt, left=0pt, right=0pt,
    middle=16pt,
    underlay={
        \draw[dashed, black!80, line width=0.8pt, dash pattern=on 4pt off 3pt]
            ([xshift=0.55\textwidth+8pt]frame.north west) --
            ([xshift=0.55\textwidth+8pt]frame.south west);
    }
]
\textbf{Q.139} \quad {\footnotesize\textsc{[MSQ $\cdot$ 2 Mark(s) $\cdot$ No Neg]}} \hfill {\footnotesize\textbf{GATE 2022 CS} \quad \tiny\texttt{2022\_CS\_CS\_Q47}}
\vspace{2pt}\hrule height 0.3pt\vspace{4pt}

Consider the following languages:
\[ L\_1 = \{a^n w a^n \mid w \in \{a, b\}^*\} \]
\[ L\_2 = \{w x w^R \mid w, x \in \{a, b\}^*, |w|, |x| > 0\} \]
Note that $w^R$ is the reversal of the string $w$. Which of the following is/are TRUE?

\par\vspace{6pt}\noindent
\begin{tabular}{@{}p{0.98\linewidth}@{}}
$\square$\ \textbf{(A)} $L_1$ and $L_2$ are regular. \\
\addlinespace[3pt]
$\square$\ \textbf{(B)} $L_1$ and $L_2$ are context-free. \\
\addlinespace[3pt]
$\square$\ \textbf{(C)} $L_1$ is regular and $L_2$ is context-free. \\
\addlinespace[3pt]
$\square$\ \textbf{(D)} $L_1$ and $L_2$ are context-free but not regular.
\end{tabular}

\tcblower

\textbf{Official Key:} \textbf{\large A,B,C} \hfill \textit{\small [Medium]}
\vspace{2pt}\hrule height 0.3pt\vspace{3pt}
\textbf{Topic:} Theory of Computation \\
\textbf{Subtopic:} \textit{Regular and context-free languages, pumping lemma}
\vspace{2pt}\hrule height 0.3pt\vspace{3pt}
\textbf{Guide / Hint:} \\
{\small Setting $n=1$ in $L_1 = \{a^n w a^n\}$ shows that any string starting and ending with $a$ of length $\ge 2$ belongs to $L_1$, making $L_1 = a(a+b)^*a$ regular. For $L_2 = \{w x w^R \mid |w|, |x| > 0\}$, taking $|w|=1$ yields $a(a+b)^+ a + b(a+b)^+ b$, which is regular. Every regular language is context-free, so statements A, B, and C are all TRUE.}

\end{tcolorbox}
\vspace{6pt}\hrule height 0.35pt\vspace{6pt}
\noindent
\begin{tcolorbox}[
    blanker,
    sidebyside,
    sidebyside align=top,
    lefthand width=0.55\textwidth,
    top=4pt, bottom=4pt, left=0pt, right=0pt,
    middle=16pt,
    underlay={
        \draw[dashed, black!80, line width=0.8pt, dash pattern=on 4pt off 3pt]
            ([xshift=0.55\textwidth+8pt]frame.north west) --
            ([xshift=0.55\textwidth+8pt]frame.south west);
    }
]
\textbf{Q.140} \quad {\footnotesize\textsc{[MSQ $\cdot$ 2 Mark(s) $\cdot$ No Neg]}} \hfill {\footnotesize\textbf{GATE 2022 CS} \quad \tiny\texttt{2022\_CS\_CS\_Q48}}
\vspace{2pt}\hrule height 0.3pt\vspace{4pt}

Consider the following languages:
\[ L\_1 = \{w w \mid w \in \{a, b\}^*\} \]
\[ L\_2 = \{a^n b^n c^m \mid m, n \ge 0\} \]
\[ L\_3 = \{a^m b^n c^n \mid m, n \ge 0\} \]
Which of the following statements is/are FALSE?

\par\vspace{6pt}\noindent
\begin{tabular}{@{}p{0.98\linewidth}@{}}
$\square$\ \textbf{(A)} $L_1$ is not context-free but $L_2$ and $L_3$ are deterministic context-free. \\
\addlinespace[3pt]
$\square$\ \textbf{(B)} Neither $L_1$ nor $L_2$ is context-free. \\
\addlinespace[3pt]
$\square$\ \textbf{(C)} $L_2$, $L_3$ and $L_2 \cap L_3$ all are context-free. \\
\addlinespace[3pt]
$\square$\ \textbf{(D)} Neither $L_1$ nor its complement is context-free.
\end{tabular}

\tcblower

\textbf{Official Key:} \textbf{\large B,C,D} \hfill \textit{\small [Hard]}
\vspace{2pt}\hrule height 0.3pt\vspace{3pt}
\textbf{Topic:} Theory of Computation \\
\textbf{Subtopic:} \textit{Context-free grammars and push-down automata}
\vspace{2pt}\hrule height 0.3pt\vspace{3pt}
\textbf{Guide / Hint:} \\
{\small Language $L_2 = \{a^n b^n c^m\}$ is a DCFL, so stating that neither $L_1$ nor $L_2$ is CFL is false (B is FALSE). $L_2 \cap L_3 = \{a^n b^n c^n\}$ is not context-free, so claiming all are context-free is false (C is FALSE). The complement $\overline{L_1}$ is a well-known context-free language, so D is FALSE. Thus, statements B, C, and D are FALSE.}

\end{tcolorbox}
\vspace{6pt}\hrule height 0.35pt\vspace{6pt}
\noindent
\begin{tcolorbox}[
    blanker,
    sidebyside,
    sidebyside align=top,
    lefthand width=0.55\textwidth,
    top=4pt, bottom=4pt, left=0pt, right=0pt,
    middle=16pt,
    underlay={
        \draw[dashed, black!80, line width=0.8pt, dash pattern=on 4pt off 3pt]
            ([xshift=0.55\textwidth+8pt]frame.north west) --
            ([xshift=0.55\textwidth+8pt]frame.south west);
    }
]
\textbf{Q.141} \quad {\footnotesize\textsc{[MCQ $\cdot$ 2 Mark(s) $\cdot$ Neg: 0.66]}} \hfill {\footnotesize\textbf{GATE 2023 CS} \quad \tiny\texttt{2023\_CS\_CS\_Q39}}
\vspace{2pt}\hrule height 0.3pt\vspace{4pt}

Consider the context-free grammar $G$ below:
\begin{align*}
S &\to aSb \mid X \\
X &\to aX \mid Xb \mid a \mid b,
\end{align*}
where $S$ and $X$ are non-terminals, and $a$ and $b$ are terminal symbols. The starting non-terminal is $S$.
Which one of the following statements is CORRECT?

\par\vspace{6pt}\noindent
\begin{tabular}{@{}p{0.98\linewidth}@{}}
$\square$\ \textbf{(A)} The language generated by $G$ is $(a + b)^*$ \\
\addlinespace[3pt]
$\square$\ \textbf{(B)} The language generated by $G$ is $a^*(a + b)b^*$ \\
\addlinespace[3pt]
$\square$\ \textbf{(C)} The language generated by $G$ is $a^*b^*(a + b)$ \\
\addlinespace[3pt]
$\square$\ \textbf{(D)} The language generated by $G$ is not a regular language
\end{tabular}

\tcblower

\textbf{Official Key:} \textbf{\large B} \hfill \textit{\small [Medium]}
\vspace{2pt}\hrule height 0.3pt\vspace{3pt}
\textbf{Topic:} Theory of Computation \\
\textbf{Subtopic:} \textit{Context-free grammars and push-down automata}
\vspace{2pt}\hrule height 0.3pt\vspace{3pt}
\textbf{Guide / Hint:} \\
{\small The subgrammar $X \to aX \mid Xb \mid a \mid b$ derives all non-empty strings with any number of $a$'s followed by any number of $b$'s having at least one $a$ or $b$, described by $a^*(a+b)b^*$. Rule $S \to aSb \mid X$ prefixes $k$ $a$'s and suffixes $k$ $b$'s ($k \ge 0$), yielding $a^k a^*(a+b)b^* b^k = a^*(a+b)b^*$. Because this language is represented by a regular expression, it is a regular language.}

\end{tcolorbox}
\vspace{6pt}\hrule height 0.35pt\vspace{6pt}
\noindent
\begin{tcolorbox}[
    blanker,
    sidebyside,
    sidebyside align=top,
    lefthand width=0.55\textwidth,
    top=4pt, bottom=4pt, left=0pt, right=0pt,
    middle=16pt,
    underlay={
        \draw[dashed, black!80, line width=0.8pt, dash pattern=on 4pt off 3pt]
            ([xshift=0.55\textwidth+8pt]frame.north west) --
            ([xshift=0.55\textwidth+8pt]frame.south west);
    }
]
\textbf{Q.142} \quad {\footnotesize\textsc{[NAT $\cdot$ 2 Mark(s) $\cdot$ No Neg]}} \hfill {\footnotesize\textbf{GATE 2023 CS} \quad \tiny\texttt{2023\_CS\_CS\_Q63}}
\vspace{2pt}\hrule height 0.3pt\vspace{4pt}

Consider the language $L$ over the alphabet $\{0, 1\}$, given below:
\[ L = \{w \in \{0, 1\}^* \mid w \text{ does not contain three or more consecutive 1's}\}. \]
The minimum number of states in a Deterministic Finite-State Automaton (DFA) for $L$ is \_\_\_\_\_\_\_\_\. (Answer in integer)

\par\vspace{6pt}\noindent\textbf{Your Answer:} \fbox{\rule{0pt}{1.3ex}\hspace{3cm}}

\tcblower

\textbf{Official Key:} \textbf{\large 4 to 4} \hfill \textit{\small [Easy]}
\vspace{2pt}\hrule height 0.3pt\vspace{3pt}
\textbf{Topic:} Theory of Computation \\
\textbf{Subtopic:} \textit{Regular expressions and finite automata}
\vspace{2pt}\hrule height 0.3pt\vspace{3pt}
\textbf{Guide / Hint:} \\
{\small The language contains all binary strings not containing `111`. The minimal DFA tracks consecutive 1s: state $q_0$ (0 consecutive 1s, start), state $q_1$ (1 consecutive 1), state $q_2$ (2 consecutive 1s), and state $q_3$ (trap state reached on 3rd consecutive 1). States $q_0, q_1, q_2$ are accepting states. Thus, exactly 4 states are required.}

\end{tcolorbox}
\vspace{6pt}\hrule height 0.35pt\vspace{6pt}
\noindent
\begin{tcolorbox}[
    blanker,
    sidebyside,
    sidebyside align=top,
    lefthand width=0.55\textwidth,
    top=4pt, bottom=4pt, left=0pt, right=0pt,
    middle=16pt,
    underlay={
        \draw[dashed, black!80, line width=0.8pt, dash pattern=on 4pt off 3pt]
            ([xshift=0.55\textwidth+8pt]frame.north west) --
            ([xshift=0.55\textwidth+8pt]frame.south west);
    }
]
\textbf{Q.143} \quad {\footnotesize\textsc{[MSQ $\cdot$ 2 Mark(s) $\cdot$ No Neg]}} \hfill {\footnotesize\textbf{GATE 2024 CS1} \quad \tiny\texttt{2024\_CS1\_CS\_Q50}}
\vspace{2pt}\hrule height 0.3pt\vspace{4pt}

Consider the 5-state DFA $M$ accepting the language $L(M) \subset (0 + 1)^*$ shown below. For any string $w \in (0 + 1)^*$ let $n_0(w)$ be the number of 0's in $w$ and $n_1(w)$ be the number of 1's in $w$.
Which of the following statements is/are FALSE?

\par\vspace{6pt}\noindent
\begin{tabular}{@{}p{0.98\linewidth}@{}}
$\square$\ \textbf{(A)} States 2 and 4 are distinguishable in $M$ \\
\addlinespace[3pt]
$\square$\ \textbf{(B)} States 3 and 4 are distinguishable in $M$ \\
\addlinespace[3pt]
$\square$\ \textbf{(C)} States 2 and 5 are distinguishable in $M$ \\
\addlinespace[3pt]
$\square$\ \textbf{(D)} Any string $w$ with $n_0(w) = n_1(w)$ is in $L(M)$
\end{tabular}

\tcblower

\textbf{Official Key:} \textbf{\large B;C} \hfill \textit{\small [Medium]}
\vspace{2pt}\hrule height 0.3pt\vspace{3pt}
\textbf{Topic:} Theory of Computation \\
\textbf{Subtopic:} \textit{Regular expressions and finite automata}
\vspace{2pt}\hrule height 0.3pt\vspace{3pt}
\textbf{Guide / Hint:} \\
{\small States $p$ and $q$ are distinguishable if there exists a string $w$ such that exactly one of $\delta(p, w)$ or $\delta(q, w)$ is an accepting state. In DFA $M$, states 3 and 4 have identical transition behavior and both are non-accepting, making them equivalent/indistinguishable, so claiming they are distinguishable is FALSE (B is FALSE). Similarly, states 2 and 5 are equivalent, so claiming they are distinguishable is FALSE (C is FALSE).}

\end{tcolorbox}
\vspace{6pt}\hrule height 0.35pt\vspace{6pt}
\noindent
\begin{tcolorbox}[
    blanker,
    sidebyside,
    sidebyside align=top,
    lefthand width=0.55\textwidth,
    top=4pt, bottom=4pt, left=0pt, right=0pt,
    middle=16pt,
    underlay={
        \draw[dashed, black!80, line width=0.8pt, dash pattern=on 4pt off 3pt]
            ([xshift=0.55\textwidth+8pt]frame.north west) --
            ([xshift=0.55\textwidth+8pt]frame.south west);
    }
]
\textbf{Q.144} \quad {\footnotesize\textsc{[NAT $\cdot$ 2 Mark(s) $\cdot$ No Neg]}} \hfill {\footnotesize\textbf{GATE 2024 CS1} \quad \tiny\texttt{2024\_CS1\_CS\_Q59}}
\vspace{2pt}\hrule height 0.3pt\vspace{4pt}

Let $G = (V, \Sigma, S, P)$ be a context-free grammar in Chomsky Normal Form with $\Sigma = \{a, b, c\}$ and $V$ containing 10 variable symbols including the start symbol $S$. The string $w = a^{30}b^{30}c^{30}$ is derivable from $S$. The number of steps (application of rules) in the derivation $S \Rightarrow^* w$ is \_\_\_\_\_\_\_\_\_\. (Answer in integer)

\par\vspace{6pt}\noindent\textbf{Your Answer:} \fbox{\rule{0pt}{1.3ex}\hspace{3cm}}

\tcblower

\textbf{Official Key:} \textbf{\large 179 to 179} \hfill \textit{\small [Medium]}
\vspace{2pt}\hrule height 0.3pt\vspace{3pt}
\textbf{Topic:} Theory of Computation \\
\textbf{Subtopic:} \textit{Context-free grammars and push-down automata}
\vspace{2pt}\hrule height 0.3pt\vspace{3pt}
\textbf{Guide / Hint:} \\
{\small In a Chomsky Normal Form (CNF) grammar, every production is either $A \to BC$ or $A \to a$. Deriving a string of length $n$ requires $n-1$ non-terminal splitting steps ($A \to BC$) followed by $n$ terminal production steps ($A \to a$), totaling $2n-1$ steps. For $w = a^{30}b^{30}c^{30}$, the length is $n = 90$, so the number of derivation steps is $2(90) - 1 = 179$.}

\end{tcolorbox}
\vspace{6pt}\hrule height 0.35pt\vspace{6pt}
\noindent
\begin{tcolorbox}[
    blanker,
    sidebyside,
    sidebyside align=top,
    lefthand width=0.55\textwidth,
    top=4pt, bottom=4pt, left=0pt, right=0pt,
    middle=16pt,
    underlay={
        \draw[dashed, black!80, line width=0.8pt, dash pattern=on 4pt off 3pt]
            ([xshift=0.55\textwidth+8pt]frame.north west) --
            ([xshift=0.55\textwidth+8pt]frame.south west);
    }
]
\textbf{Q.145} \quad {\footnotesize\textsc{[NAT $\cdot$ 2 Mark(s) $\cdot$ No Neg]}} \hfill {\footnotesize\textbf{GATE 2024 CS1} \quad \tiny\texttt{2024\_CS1\_CS\_Q61}}
\vspace{2pt}\hrule height 0.3pt\vspace{4pt}

Consider the following two regular expressions over the alphabet $\{0,1\}$:
\[ r = 0^* + 1^* \]
\[ s = 01^* + 10^* \]
The total number of strings of length less than or equal to 5, which are neither in $r$ nor in $s$, is \_\_\_\_\_\_\_\_\_\. (Answer in integer)

\par\vspace{6pt}\noindent\textbf{Your Answer:} \fbox{\rule{0pt}{1.3ex}\hspace{3cm}}

\tcblower

\textbf{Official Key:} \textbf{\large 44 to 44} \hfill \textit{\small [Medium]}
\vspace{2pt}\hrule height 0.3pt\vspace{3pt}
\textbf{Topic:} Theory of Computation \\
\textbf{Subtopic:} \textit{Regular expressions and finite automata}
\vspace{2pt}\hrule height 0.3pt\vspace{3pt}
\textbf{Guide / Hint:} \\
{\small The total number of binary strings of length $\le 5$ is $\sum_{i=0}^5 2^i = 1+2+4+8+16+32 = 63$. Language $L(r) = 0^* + 1^*$ contains $\epsilon$ and $0^k, 1^k$ for each length $k \in \{1..5\}$ (11 strings). Language $L(s) = 01^* + 10^*$ contains $01^{k-1}$ and $10^{k-1}$ for each length $k \in \{1..5\}$. At $k=1$, $0$ and $1$ overlap with $L(r)$; for $k \in \{2..5\}$, the $4 \times 2 = 8$ strings are disjoint from $L(r)$. Thus, $|L(r) \cup L(s)| = 11 + 8 = 19$. The number of strings in neither language is $63 - 19 = 44$.}

\end{tcolorbox}
\vspace{6pt}\hrule height 0.35pt\vspace{6pt}
\noindent
\begin{tcolorbox}[
    blanker,
    sidebyside,
    sidebyside align=top,
    lefthand width=0.55\textwidth,
    top=4pt, bottom=4pt, left=0pt, right=0pt,
    middle=16pt,
    underlay={
        \draw[dashed, black!80, line width=0.8pt, dash pattern=on 4pt off 3pt]
            ([xshift=0.55\textwidth+8pt]frame.north west) --
            ([xshift=0.55\textwidth+8pt]frame.south west);
    }
]
\textbf{Q.146} \quad {\footnotesize\textsc{[MSQ $\cdot$ 2 Mark(s) $\cdot$ No Neg]}} \hfill {\footnotesize\textbf{GATE 2024 CS2} \quad \tiny\texttt{2024\_CS2\_CS\_Q52}}
\vspace{2pt}\hrule height 0.3pt\vspace{4pt}

Consider a context-free grammar $G$ with the following 3 rules:
\[ S \to aS \mid aSbS \mid c \]
Let $w \in L(G)$. Let $n_a(w), n_b(w), n_c(w)$ denote the number of times $a, b, c$ occur in $w$, respectively. Which of the following statements is/are TRUE?

\par\vspace{6pt}\noindent
\begin{tabular}{@{}p{0.48\linewidth}@{\hspace{0.04\linewidth}}p{0.48\linewidth}@{}}
$\square$\ \textbf{(A)} $n_a(w) > n_b(w)$ & $\square$\ \textbf{(B)} $n_a(w) > n_c(w) - 2$ \\
\addlinespace[3pt]
$\square$\ \textbf{(C)} $n_c(w) = n_b(w) + 1$ & $\square$\ \textbf{(D)} $n_c(w) = n_b(w) * 2$
\end{tabular}

\tcblower

\textbf{Official Key:} \textbf{\large B;C} \hfill \textit{\small [Medium]}
\vspace{2pt}\hrule height 0.3pt\vspace{3pt}
\textbf{Topic:} Theory of Computation \\
\textbf{Subtopic:} \textit{Context-free grammars and push-down automata}
\vspace{2pt}\hrule height 0.3pt\vspace{3pt}
\textbf{Guide / Hint:} \\
{\small Each application of $S \to aSbS$ consumes one $S$ and introduces two new $S$'s along with one $b$, effectively adding $+1$ to the count of $S$'s. Each termination $S \to c$ eliminates one $S$ and adds one $c$. Thus, the total number of $c$'s must equal the number of $b$'s plus 1: $n_c(w) = n_b(w) + 1$ (C is true). Also, every $b$ is accompanied by an $a$, and additional $a$'s can come from $S \to aS$, giving $n_a(w) \ge n_b(w) = n_c(w) - 1 > n_c(w) - 2$ (B is true).}

\end{tcolorbox}
\vspace{6pt}\hrule height 0.35pt\vspace{6pt}
\noindent
\begin{tcolorbox}[
    blanker,
    sidebyside,
    sidebyside align=top,
    lefthand width=0.55\textwidth,
    top=4pt, bottom=4pt, left=0pt, right=0pt,
    middle=16pt,
    underlay={
        \draw[dashed, black!80, line width=0.8pt, dash pattern=on 4pt off 3pt]
            ([xshift=0.55\textwidth+8pt]frame.north west) --
            ([xshift=0.55\textwidth+8pt]frame.south west);
    }
]
\textbf{Q.147} \quad {\footnotesize\textsc{[NAT $\cdot$ 2 Mark(s) $\cdot$ No Neg]}} \hfill {\footnotesize\textbf{GATE 2024 CS2} \quad \tiny\texttt{2024\_CS2\_CS\_Q62}}
\vspace{2pt}\hrule height 0.3pt\vspace{4pt}

Let $L_1$ be the language represented by the regular expression $b^* a b^* (a b^* a b^*)^*$ and $L_2 = \{ w \in (a + b)^* \mid |w| \le 4 \}$, where $|w|$ denotes the length of string $w$. The number of strings in $L_2$ which are also in $L_1$ is \_\_\_\_\_\_\. (Answer in integer)

\par\vspace{6pt}\noindent\textbf{Your Answer:} \fbox{\rule{0pt}{1.3ex}\hspace{3cm}}

\tcblower

\textbf{Official Key:} \textbf{\large 15 to 15} \hfill \textit{\small [Medium]}
\vspace{2pt}\hrule height 0.3pt\vspace{3pt}
\textbf{Topic:} Theory of Computation \\
\textbf{Subtopic:} \textit{Regular expressions and finite automata}
\vspace{2pt}\hrule height 0.3pt\vspace{3pt}
\textbf{Guide / Hint:} \\
{\small $L_1$ requires an odd number of $a$'s: strings containing either 1 or 3 $a$'s (since length $\le 4$). Length $\le 4$ with exactly one $a$: length 1 has $1$, length 2 has $2$, length 3 has $3$, length 4 has $4$ (total $1 + 2 + 3 + 4 = 10$). Length $\le 4$ with exactly three $a$'s: length 3 has $aaa$ ($1$), length 4 has $\binom{4}{3} = 4$. Total strings in $L_1 \cap L_2 = 10 + 1 + 4 = 15$.}

\end{tcolorbox}
\vspace{6pt}\hrule height 0.35pt\vspace{6pt}
\noindent
\begin{tcolorbox}[
    blanker,
    sidebyside,
    sidebyside align=top,
    lefthand width=0.55\textwidth,
    top=4pt, bottom=4pt, left=0pt, right=0pt,
    middle=16pt,
    underlay={
        \draw[dashed, black!80, line width=0.8pt, dash pattern=on 4pt off 3pt]
            ([xshift=0.55\textwidth+8pt]frame.north west) --
            ([xshift=0.55\textwidth+8pt]frame.south west);
    }
]
\textbf{Q.148} \quad {\footnotesize\textsc{[MCQ $\cdot$ 2 Mark(s) $\cdot$ Neg: 0.66]}} \hfill {\footnotesize\textbf{GATE 2025 CS1} \quad \tiny\texttt{2025\_CS1\_CS\_Q44}}
\vspace{2pt}\hrule height 0.3pt\vspace{4pt}

Consider the following two languages over the alphabet $\{a, b\}$:
\[ L\_1 = \{ \alpha\beta\alpha \mid \alpha \in \{a, b\}^+ \text{ AND } \beta \in \{a, b\}^+ \} \]
\[ L\_2 = \{ \alpha\beta\alpha \mid \alpha \in \{a\}^+ \text{ AND } \beta \in \{a, b\}^+ \} \]
Which ONE of the following statements is CORRECT?

\par\vspace{6pt}\noindent
\begin{tabular}{@{}p{0.98\linewidth}@{}}
$\square$\ \textbf{(A)} Both $L_1$ and $L_2$ are regular languages. \\
\addlinespace[3pt]
$\square$\ \textbf{(B)} $L_1$ is a regular language but $L_2$ is not a regular language. \\
\addlinespace[3pt]
$\square$\ \textbf{(C)} $L_1$ is not a regular language but $L_2$ is a regular language. \\
\addlinespace[3pt]
$\square$\ \textbf{(D)} Neither $L_1$ nor $L_2$ is a regular language.
\end{tabular}

\tcblower

\textbf{Official Key:} \textbf{\large C} \hfill \textit{\small [Medium]}
\vspace{2pt}\hrule height 0.3pt\vspace{3pt}
\textbf{Topic:} Theory of Computation \\
\textbf{Subtopic:} \textit{Regular and context-free languages, pumping lemma}
\vspace{2pt}\hrule height 0.3pt\vspace{3pt}
\textbf{Guide / Hint:} \\
{\small For $L_2$, $\alpha = a^k$ ($k \ge 1$). A string belongs to $L_2$ if and only if it starts with $a$, ends with $a$, and has length at least 3 ($a(a+b)^+a$), which is clearly regular. For $L_1$, the prefix $\alpha$ and suffix $\alpha$ must match over $\{a, b\}^+$, which requires unbounded memory to verify matching arbitrary binary strings, making $L_1$ non-regular.}

\end{tcolorbox}
\vspace{6pt}\hrule height 0.35pt\vspace{6pt}
\noindent
\begin{tcolorbox}[
    blanker,
    sidebyside,
    sidebyside align=top,
    lefthand width=0.55\textwidth,
    top=4pt, bottom=4pt, left=0pt, right=0pt,
    middle=16pt,
    underlay={
        \draw[dashed, black!80, line width=0.8pt, dash pattern=on 4pt off 3pt]
            ([xshift=0.55\textwidth+8pt]frame.north west) --
            ([xshift=0.55\textwidth+8pt]frame.south west);
    }
]
\textbf{Q.149} \quad {\footnotesize\textsc{[MCQ $\cdot$ 2 Mark(s) $\cdot$ Neg: 0.66]}} \hfill {\footnotesize\textbf{GATE 2025 CS1} \quad \tiny\texttt{2025\_CS1\_CS\_Q45}}
\vspace{2pt}\hrule height 0.3pt\vspace{4pt}

Consider the following two languages over the alphabet $\{a, b, c\}$, where $m$ and $n$ are natural numbers.
\[ L\_1 = \{a^mb^mc^{m+n} \mid m, n \ge 1\} \]
\[ L\_2 = \{a^mb^nc^{m+n} \mid m, n \ge 1\} \]
Which ONE of the following statements is CORRECT?

\par\vspace{6pt}\noindent
\begin{tabular}{@{}p{0.98\linewidth}@{}}
$\square$\ \textbf{(A)} Both $L_1$ and $L_2$ are context-free languages. \\
\addlinespace[3pt]
$\square$\ \textbf{(B)} $L_1$ is a context-free language but $L_2$ is not a context-free language. \\
\addlinespace[3pt]
$\square$\ \textbf{(C)} $L_1$ is not a context-free language but $L_2$ is a context-free language. \\
\addlinespace[3pt]
$\square$\ \textbf{(D)} Neither $L_1$ nor $L_2$ are context-free languages.
\end{tabular}

\tcblower

\textbf{Official Key:} \textbf{\large C} \hfill \textit{\small [Medium]}
\vspace{2pt}\hrule height 0.3pt\vspace{3pt}
\textbf{Topic:} Theory of Computation \\
\textbf{Subtopic:} \textit{Context-free grammars and push-down automata}
\vspace{2pt}\hrule height 0.3pt\vspace{3pt}
\textbf{Guide / Hint:} \\
{\small In $L_2$, rewrite $c^{m+n} = c^n c^m$ or generate using PDA by pushing $a$'s and $b$'s and popping on $c$'s ($c$'s match total count of $a$'s and $b$'s), which requires a single stack, so $L_2$ is context-free. In $L_1$, $a^m b^m c^m c^n$ simultaneously requires equality between powers of $a$ and $b$, plus matching with $c$, requiring two independent stack comparisons, making $L_1$ non-CFL.}

\end{tcolorbox}
\vspace{6pt}\hrule height 0.35pt\vspace{6pt}
\noindent
\begin{tcolorbox}[
    blanker,
    sidebyside,
    sidebyside align=top,
    lefthand width=0.55\textwidth,
    top=4pt, bottom=4pt, left=0pt, right=0pt,
    middle=16pt,
    underlay={
        \draw[dashed, black!80, line width=0.8pt, dash pattern=on 4pt off 3pt]
            ([xshift=0.55\textwidth+8pt]frame.north west) --
            ([xshift=0.55\textwidth+8pt]frame.south west);
    }
]
\textbf{Q.150} \quad {\footnotesize\textsc{[MSQ $\cdot$ 2 Mark(s) $\cdot$ No Neg]}} \hfill {\footnotesize\textbf{GATE 2025 CS1} \quad \tiny\texttt{2025\_CS1\_CS\_Q50}}
\vspace{2pt}\hrule height 0.3pt\vspace{4pt}

Consider the deterministic finite automaton (DFA) defined over the alphabet $\Sigma = \{a, b\}$. Identify which of the following language(s) is/are accepted by the given DFA.

\par\vspace{6pt}\noindent
\begin{tabular}{@{}p{0.98\linewidth}@{}}
$\square$\ \textbf{(A)} The set of all strings containing an even number of $b$'s. \\
\addlinespace[3pt]
$\square$\ \textbf{(B)} The set of all strings containing the pattern $bab$. \\
\addlinespace[3pt]
$\square$\ \textbf{(C)} The set of all strings ending with the pattern $bab$. \\
\addlinespace[3pt]
$\square$\ \textbf{(D)} The set of all strings not containing the pattern $aba$.
\end{tabular}

\tcblower

\textbf{Official Key:} \textbf{\large C} \hfill \textit{\small [Medium]}
\vspace{2pt}\hrule height 0.3pt\vspace{3pt}
\textbf{Topic:} Theory of Computation \\
\textbf{Subtopic:} \textit{Regular expressions and finite automata}
\vspace{2pt}\hrule height 0.3pt\vspace{3pt}
\textbf{Guide / Hint:} \\
{\small In the given DFA, the sequence of transitions leading to the unique accepting state $q_3$ corresponds precisely to recognizing the suffix $bab$. Any character following $q_3$ transitions back according to pattern matching prefixes of $bab$, accepting exactly strings ending with the pattern $bab$.}

\end{tcolorbox}
\vspace{6pt}\hrule height 0.35pt\vspace{6pt}
\noindent
\begin{tcolorbox}[
    blanker,
    sidebyside,
    sidebyside align=top,
    lefthand width=0.55\textwidth,
    top=4pt, bottom=4pt, left=0pt, right=0pt,
    middle=16pt,
    underlay={
        \draw[dashed, black!80, line width=0.8pt, dash pattern=on 4pt off 3pt]
            ([xshift=0.55\textwidth+8pt]frame.north west) --
            ([xshift=0.55\textwidth+8pt]frame.south west);
    }
]
\textbf{Q.151} \quad {\footnotesize\textsc{[MSQ $\cdot$ 2 Mark(s) $\cdot$ No Neg]}} \hfill {\footnotesize\textbf{GATE 2025 CS2} \quad \tiny\texttt{2025\_CS2\_CS\_Q52}}
\vspace{2pt}\hrule height 0.3pt\vspace{4pt}

Let $\Sigma = \{a, b, c\}$. For $x \in \Sigma^*$, let $\#_\alpha(x)$ denote the number of occurrences of $\alpha$ in $x$. Which one or more of the following languages is/are regular?

\par\vspace{6pt}\noindent
\begin{tabular}{@{}p{0.98\linewidth}@{}}
$\square$\ \textbf{(A)} $\{a^m b^n \mid m, n \ge 0\}$ \\
\addlinespace[3pt]
$\square$\ \textbf{(B)} $\{a, b\}^* \cap \{a^m b^n c^{m-n} \mid m \ge n \ge 0\}$ \\
\addlinespace[3pt]
$\square$\ \textbf{(C)} $\{w \in \{a, b\}^* \mid \#_a(w) \equiv 2 \pmod 7 \text{ and } \#_b(w) \equiv 3 \pmod 9\}$ \\
\addlinespace[3pt]
$\square$\ \textbf{(D)} $\{w \in \{a, b\}^* \mid \#_a(w) \equiv 2 \pmod 7 \text{ and } \#_a(w) = \#_b(w)\}$
\end{tabular}

\tcblower

\textbf{Official Key:} \textbf{\large A;C} \hfill \textit{\small [Medium]}
\vspace{2pt}\hrule height 0.3pt\vspace{3pt}
\textbf{Topic:} Theory of Computation \\
\textbf{Subtopic:} \textit{Regular expressions and finite automata: regular languages}
\vspace{2pt}\hrule height 0.3pt\vspace{3pt}
\textbf{Guide / Hint:} \\
{\small Language A ($a^m b^n$) is generated by the regular expression $a^* b^*$. Language C counts occurrences modulo 7 and 9, which requires finite memory ($7 \times 9 = 63$ states) and is therefore regular.}

\end{tcolorbox}
\vspace{6pt}\hrule height 0.35pt\vspace{6pt}
\noindent
\begin{tcolorbox}[
    blanker,
    sidebyside,
    sidebyside align=top,
    lefthand width=0.55\textwidth,
    top=4pt, bottom=4pt, left=0pt, right=0pt,
    middle=16pt,
    underlay={
        \draw[dashed, black!80, line width=0.8pt, dash pattern=on 4pt off 3pt]
            ([xshift=0.55\textwidth+8pt]frame.north west) --
            ([xshift=0.55\textwidth+8pt]frame.south west);
    }
]
\textbf{Q.152} \quad {\footnotesize\textsc{[NAT $\cdot$ 2 Mark(s) $\cdot$ No Neg]}} \hfill {\footnotesize\textbf{GATE 2025 CS2} \quad \tiny\texttt{2025\_CS2\_CS\_Q60}}
\vspace{2pt}\hrule height 0.3pt\vspace{4pt}

Let $\Sigma = \{1, 2, 3, 4\}$. For $x \in \Sigma^*$, let $prod(x)$ be the product of symbols in $x$ modulo 7. The minimum number of states in a DFA accepting strings with $prod(x) = 1$ is \_\_\_\_\_\_\_\_.

\par\vspace{6pt}\noindent\textbf{Your Answer:} \fbox{\rule{0pt}{1.3ex}\hspace{3cm}}

\tcblower

\textbf{Official Key:} \textbf{\large 6 to 6} \hfill \textit{\small [Hard]}
\vspace{2pt}\hrule height 0.3pt\vspace{3pt}
\textbf{Topic:} Theory of Computation \\
\textbf{Subtopic:} \textit{Regular expressions and finite automata: minimal DFA states}
\vspace{2pt}\hrule height 0.3pt\vspace{3pt}
\textbf{Guide / Hint:} \\
{\small The multiplicative group modulo 7 of units $\{1, 2, 3, 4, 5, 6\}$ is cyclic of order 6. Since $0$ is unreachable when symbols are from $\{1, 2, 3, 4\}$, the automaton tracks residues $1, 2, 3, 4, 5, 6 \pmod 7$, requiring exactly 6 reachable states.}

\end{tcolorbox}
\vspace{6pt}\hrule height 0.35pt\vspace{6pt}
\vspace{8pt}
\subsection*{Compiler Design}
\vspace{-4pt}
\noindent
\begin{tcolorbox}[
    blanker,
    sidebyside,
    sidebyside align=top,
    lefthand width=0.55\textwidth,
    top=4pt, bottom=4pt, left=0pt, right=0pt,
    middle=16pt,
    underlay={
        \draw[dashed, black!80, line width=0.8pt, dash pattern=on 4pt off 3pt]
            ([xshift=0.55\textwidth+8pt]frame.north west) --
            ([xshift=0.55\textwidth+8pt]frame.south west);
    }
]
\textbf{Q.153} \quad {\footnotesize\textsc{[MCQ $\cdot$ 2 Mark(s) $\cdot$ Neg: 0.66]}} \hfill {\footnotesize\textbf{GATE 2018 CS} \quad \tiny\texttt{2018\_CS\_GA\_Q37}}
\vspace{2pt}\hrule height 0.3pt\vspace{4pt}

A lexical analyzer uses the following patterns to recognize three tokens $T_1, T_2$, and $T_3$ over the alphabet $\{a, b, c\}$:
\[ T\_1: a?(b|c)^*a \qquad T\_2: b?(a|c)^*b \qquad T\_3: c?(b|a)^*c \]
Note that '?' means 0 or 1 occurrence of the preceding symbol. The analyzer matches the longest prefix of the remaining input. For the input string `bbaacabc`, the token sequence generated is

\par\vspace{6pt}\noindent
\begin{tabular}{@{}p{0.48\linewidth}@{\hspace{0.04\linewidth}}p{0.48\linewidth}@{}}
$\square$\ \textbf{(A)} $T_1 T_2 T_3$ & $\square$\ \textbf{(B)} $T_1 T_3$ \\
\addlinespace[3pt]
$\square$\ \textbf{(C)} $T_2 T_1 T_3$ & $\square$\ \textbf{(D)} $T_3 T_3$
\end{tabular}

\tcblower

\textbf{Official Key:} \textbf{\large D} \hfill \textit{\small [Medium]}
\vspace{2pt}\hrule height 0.3pt\vspace{3pt}
\textbf{Topic:} Compiler Design \\
\textbf{Subtopic:} \textit{Lexical analysis, parsing, syntax-directed translation}
\vspace{2pt}\hrule height 0.3pt\vspace{3pt}
\textbf{Guide / Hint:} \\
{\small A lexical analyzer applies the longest match rule. For input `bbaacabc`, pattern $T_3 = c?(b|a)^*c$ matches prefix `bbaac` (starts without $c$, matches $(b|a)^* = bbaa$, ends with $c$), consuming `bbaac` as token $T_3$. The remaining string is `abc`, which also matches $T_3$ (starts without $c$, $(b|a)^* = ab$, ends with $c$). Hence the recognized sequence of tokens is $T_3 T_3$.}

\end{tcolorbox}
\vspace{6pt}\hrule height 0.35pt\vspace{6pt}
\noindent
\begin{tcolorbox}[
    blanker,
    sidebyside,
    sidebyside align=top,
    lefthand width=0.55\textwidth,
    top=4pt, bottom=4pt, left=0pt, right=0pt,
    middle=16pt,
    underlay={
        \draw[dashed, black!80, line width=0.8pt, dash pattern=on 4pt off 3pt]
            ([xshift=0.55\textwidth+8pt]frame.north west) --
            ([xshift=0.55\textwidth+8pt]frame.south west);
    }
]
\textbf{Q.154} \quad {\footnotesize\textsc{[NAT $\cdot$ 2 Mark(s) $\cdot$ No Neg]}} \hfill {\footnotesize\textbf{GATE 2022 CS} \quad \tiny\texttt{2022\_CS\_CS\_Q65}}
\vspace{2pt}\hrule height 0.3pt\vspace{4pt}

Consider the following grammar along with translation rules:
\begin{align*}
S &\to S_1 \,\#\, T & S.val &= S_1.val * T.val \\
S &\to T & S.val &= T.val \\
T &\to T_1 \,\%\, R & T.val &= T_1.val \div R.val \\
T &\to R & T.val &= R.val \\
R &\to id & R.val &= id.val
\end{align*}
Here $\#$ and $\%$ are operators and $id$ is a token that represents an integer. Using this translation scheme, the computed value of $S.val$ for root of the parse tree for the expression `20\#10\%5\#8\%2\%2` is \_\_\_\_\_\_\_\_\_\_\_\_\. (Answer in integer)

\par\vspace{6pt}\noindent\textbf{Your Answer:} \fbox{\rule{0pt}{1.3ex}\hspace{3cm}}

\tcblower

\textbf{Official Key:} \textbf{\large 80 to 80} \hfill \textit{\small [Medium]}
\vspace{2pt}\hrule height 0.3pt\vspace{3pt}
\textbf{Topic:} Compiler Design \\
\textbf{Subtopic:} \textit{Lexical analysis, parsing, syntax-directed translation}
\vspace{2pt}\hrule height 0.3pt\vspace{3pt}
\textbf{Guide / Hint:} \\
{\small The SDD implements left-associative operators with $\%$ at higher precedence than $\#$. In `20 \# 10 \% 5 \# 8 \% 2 \% 2`, we first evaluate subexpressions: $10 \% 5 = 10 / 5 = 2$, and $8 \% 2 \% 2 = (8 / 2) / 2 = 2$. Then the outer multiplication evaluates left-associatively: $(20 \times 2) \times 2 = 40 \times 2 = 80$.}

\end{tcolorbox}
\vspace{6pt}\hrule height 0.35pt\vspace{6pt}
\noindent
\begin{tcolorbox}[
    blanker,
    sidebyside,
    sidebyside align=top,
    lefthand width=0.55\textwidth,
    top=4pt, bottom=4pt, left=0pt, right=0pt,
    middle=16pt,
    underlay={
        \draw[dashed, black!80, line width=0.8pt, dash pattern=on 4pt off 3pt]
            ([xshift=0.55\textwidth+8pt]frame.north west) --
            ([xshift=0.55\textwidth+8pt]frame.south west);
    }
]
\textbf{Q.155} \quad {\footnotesize\textsc{[NAT $\cdot$ 2 Mark(s) $\cdot$ No Neg]}} \hfill {\footnotesize\textbf{GATE 2023 CS} \quad \tiny\texttt{2023\_CS\_CS\_Q60}}
\vspace{2pt}\hrule height 0.3pt\vspace{4pt}

Consider the syntax directed translation given by the following grammar and semantic rules. Here $N, I, F$ and $B$ are non-terminals. $N$ is the starting non-terminal, and $\#, 0$ and $1$ are lexical tokens corresponding to input letters ``\#'', ``0'' and ``1'', respectively. $X.val$ denotes the synthesized attribute (a numeric value) associated with a non-terminal $X$. $I_1$ and $F_1$ denote occurrences of $I$ and $F$ on the right hand side of a production, respectively. For the tokens 0 and 1, $0.val = 0$ and $1.val = 1$.
\begin{align*}
N &\to I \,\#\, F & N.val &= I.val + F.val \\
I &\to I_1 B & I.val &= (2 \cdot I_1.val) + B.val \\
I &\to B & I.val &= B.val \\
F &\to B F_1 & F.val &= \frac{1}{2}(B.val + F_1.val) \\
F &\to B & F.val &= \frac{1}{2}B.val \\
B &\to 0 & B.val &= 0.val \\
B &\to 1 & B.val &= 1.val
\end{align*}
The value computed by the translation scheme for the input string `10\#011` is \_\_\_\_\_\_\_\_\. (Rounded off to three decimal places)

\par\vspace{6pt}\noindent\textbf{Your Answer:} \fbox{\rule{0pt}{1.3ex}\hspace{3cm}}

\tcblower

\textbf{Official Key:} \textbf{\large 2.374 to 2.376} \hfill \textit{\small [Medium]}
\vspace{2pt}\hrule height 0.3pt\vspace{3pt}
\textbf{Topic:} Compiler Design \\
\textbf{Subtopic:} \textit{Lexical analysis, parsing, syntax-directed translation}
\vspace{2pt}\hrule height 0.3pt\vspace{3pt}
\textbf{Guide / Hint:} \\
{\small For `10\#011`: integer part `10` evaluates as $I.val = 2(1) + 0 = 2$. Fraction part `011` with right-recursive rules: lowest $F_3 \to B(1) \implies F_3.val = 0.5$; middle $F_2 \to B(1) F_3 \implies F_2.val = \frac{1}{2}(1 + 0.5) = 0.75$; top $F \to B(0) F_2 \implies F.val = \frac{1}{2}(0 + 0.75) = 0.375$. Total value $N.val = I.val + F.val = 2 + 0.375 = 2.375$.}

\end{tcolorbox}
\vspace{6pt}\hrule height 0.35pt\vspace{6pt}
\noindent
\begin{tcolorbox}[
    blanker,
    sidebyside,
    sidebyside align=top,
    lefthand width=0.55\textwidth,
    top=4pt, bottom=4pt, left=0pt, right=0pt,
    middle=16pt,
    underlay={
        \draw[dashed, black!80, line width=0.8pt, dash pattern=on 4pt off 3pt]
            ([xshift=0.55\textwidth+8pt]frame.north west) --
            ([xshift=0.55\textwidth+8pt]frame.south west);
    }
]
\textbf{Q.156} \quad {\footnotesize\textsc{[MCQ $\cdot$ 2 Mark(s) $\cdot$ Neg: 0.66]}} \hfill {\footnotesize\textbf{GATE 2024 CS1} \quad \tiny\texttt{2024\_CS1\_CS\_Q37}}
\vspace{2pt}\hrule height 0.3pt\vspace{4pt}

Consider the following syntax-directed definition (SDD):
\begin{align*}
S &\to DHTU & \{ S.val = D.val + H.val + T.val + U.val; \} \\
D &\to \text{``M''} D_1 & \{ D.val = 5 + D_1.val; \} \\
D &\to \epsilon & \{ D.val = -5; \} \\
H &\to \text{``L''} H_1 & \{ H.val = 5 * 10 + H_1.val; \} \\
H &\to \epsilon & \{ H.val = -10; \} \\
T &\to \text{``C''} T_1 & \{ T.val = 5 * 100 + T_1.val; \} \\
T &\to \epsilon & \{ T.val = -5; \} \\
U &\to \text{``K''} & \{ U.val = 5; \}
\end{align*}
Given ``MMLK'' as the input, which one of the following options is the CORRECT value computed by the SDD (in the attribute $S.val$)?

\par\vspace{6pt}\noindent
\begin{tabular}{@{}p{0.48\linewidth}@{\hspace{0.04\linewidth}}p{0.48\linewidth}@{}}
$\square$\ \textbf{(A)} 45 & $\square$\ \textbf{(B)} 50 \\
\addlinespace[3pt]
$\square$\ \textbf{(C)} 55 & $\square$\ \textbf{(D)} 65
\end{tabular}

\tcblower

\textbf{Official Key:} \textbf{\large A} \hfill \textit{\small [Medium]}
\vspace{2pt}\hrule height 0.3pt\vspace{3pt}
\textbf{Topic:} Compiler Design \\
\textbf{Subtopic:} \textit{Lexical analysis, parsing, syntax-directed translation}
\vspace{2pt}\hrule height 0.3pt\vspace{3pt}
\textbf{Guide / Hint:} \\
{\small For input "MMLK" with production $S \to DHTU$: $D$ derives "MM" via $D \to \text{"M"} D_1 \to \text{"M"}\text{"M"} D_2 \to \text{"M"}\text{"M"}\epsilon$, giving $D.val = 5 + 5 - 5 = 5$. $H$ derives "L" via $H \to \text{"L"} H_1 \to \text{"L"}\epsilon$, giving $H.val = 50 - 10 = 40$. $T$ derives $\epsilon$, giving $T.val = -5$. $U$ derives "K", giving $U.val = 5$. Summing attribute values yields $S.val = 5 + 40 - 5 + 5 = 45$.}

\end{tcolorbox}
\vspace{6pt}\hrule height 0.35pt\vspace{6pt}
\noindent
\begin{tcolorbox}[
    blanker,
    sidebyside,
    sidebyside align=top,
    lefthand width=0.55\textwidth,
    top=4pt, bottom=4pt, left=0pt, right=0pt,
    middle=16pt,
    underlay={
        \draw[dashed, black!80, line width=0.8pt, dash pattern=on 4pt off 3pt]
            ([xshift=0.55\textwidth+8pt]frame.north west) --
            ([xshift=0.55\textwidth+8pt]frame.south west);
    }
]
\textbf{Q.157} \quad {\footnotesize\textsc{[MCQ $\cdot$ 2 Mark(s) $\cdot$ Neg: 0.66]}} \hfill {\footnotesize\textbf{GATE 2024 CS1} \quad \tiny\texttt{2024\_CS1\_CS\_Q38}}
\vspace{2pt}\hrule height 0.3pt\vspace{4pt}

Consider the following grammar $G$, with $S$ as the start symbol. The grammar $G$ has three incomplete productions denoted by (1), (2), and (3).
\begin{align*}
S &\to daT \mid (1) \\
T &\to aS \mid bT \mid (2) \\
R &\to (3) \mid \epsilon
\end{align*}
The set of terminals is $\{a, b, c, d, f\}$. The FIRST and FOLLOW sets of the different non-terminals are as follows:
$$\text{FIRST}(S) = \{c, d, f\}, \quad \text{FIRST}(T) = \{a, b, \epsilon\}, \quad \text{FIRST}(R) = \{c, \epsilon\}$$
$$\text{FOLLOW}(S) = \text{FOLLOW}(T) = \{c, f, \$\}, \quad \text{FOLLOW}(R) = \{f\}$$
Which one of the following options CORRECTLY fills in the incomplete productions?

\par\vspace{6pt}\noindent
\begin{tabular}{@{}p{0.98\linewidth}@{}}
$\square$\ \textbf{(A)} (1) $S \to Rf$ (2) $T \to \epsilon$ (3) $R \to cTR$ \\
\addlinespace[3pt]
$\square$\ \textbf{(B)} (1) $S \to fR$ (2) $T \to \epsilon$ (3) $R \to cTR$ \\
\addlinespace[3pt]
$\square$\ \textbf{(C)} (1) $S \to fR$ (2) $T \to cT$ (3) $R \to cR$ \\
\addlinespace[3pt]
$\square$\ \textbf{(D)} (1) $S \to Rf$ (2) $T \to cT$ (3) $R \to cR$
\end{tabular}

\tcblower

\textbf{Official Key:} \textbf{\large A} \hfill \textit{\small [Medium]}
\vspace{2pt}\hrule height 0.3pt\vspace{3pt}
\textbf{Topic:} Compiler Design \\
\textbf{Subtopic:} \textit{Lexical analysis, parsing, syntax-directed translation}
\vspace{2pt}\hrule height 0.3pt\vspace{3pt}
\textbf{Guide / Hint:} \\
{\small Since $\text{FIRST}(S) = \{c, d, f\}$ and $S \to daT$ generates $d$, production (1) must generate $c$ and $f$. With (1) as $S \to Rf$ where $\text{FIRST}(R) = \{c, \epsilon\}$, $\text{FIRST}(Rf) = \{c, f\}$. Since $\epsilon \in \text{FIRST}(T)$, production (2) must be $T \to \epsilon$. Production (3) as $R \to cTR$ matches $\text{FIRST}(R) = \{c, \epsilon\}$ and $\text{FOLLOW}(R) = \{f\}$, matching option (A).}

\end{tcolorbox}
\vspace{6pt}\hrule height 0.35pt\vspace{6pt}
\noindent
\begin{tcolorbox}[
    blanker,
    sidebyside,
    sidebyside align=top,
    lefthand width=0.55\textwidth,
    top=4pt, bottom=4pt, left=0pt, right=0pt,
    middle=16pt,
    underlay={
        \draw[dashed, black!80, line width=0.8pt, dash pattern=on 4pt off 3pt]
            ([xshift=0.55\textwidth+8pt]frame.north west) --
            ([xshift=0.55\textwidth+8pt]frame.south west);
    }
]
\textbf{Q.158} \quad {\footnotesize\textsc{[MCQ $\cdot$ 2 Mark(s) $\cdot$ Neg: 0.66]}} \hfill {\footnotesize\textbf{GATE 2024 CS1} \quad \tiny\texttt{2024\_CS1\_CS\_Q39}}
\vspace{2pt}\hrule height 0.3pt\vspace{4pt}

Consider the following pseudo-code:
\begin{lstlisting}
L1:  t1 = -1
L2:  t2 = 0
L3:  t3 = 0
L4:  t4 = 4 * t3
L5:  t5 = 4 * t2
L6:  t6 = t5 * M
L7:  t7 = t4 + t6
L8:  t8 = a[t7]
L9:  if t8 <= max goto L11
L10: t1 = t8
L11: t3 = t3 + 1
L12: if t3 < M goto L4
L13: t2 = t2 + 1
L14: if t2 < N goto L3
L15: max = t1
\end{lstlisting}
Which one of the following options CORRECTLY specifies the number of basic blocks and the number of instructions in the largest basic block, respectively?

\par\vspace{6pt}\noindent
\begin{tabular}{@{}p{0.48\linewidth}@{\hspace{0.04\linewidth}}p{0.48\linewidth}@{}}
$\square$\ \textbf{(A)} 6 and 6 & $\square$\ \textbf{(B)} 6 and 7 \\
\addlinespace[3pt]
$\square$\ \textbf{(C)} 7 and 7 & $\square$\ \textbf{(D)} 7 and 6
\end{tabular}

\tcblower

\textbf{Official Key:} \textbf{\large D} \hfill \textit{\small [Medium]}
\vspace{2pt}\hrule height 0.3pt\vspace{3pt}
\textbf{Topic:} Compiler Design \\
\textbf{Subtopic:} \textit{Intermediate code generation}
\vspace{2pt}\hrule height 0.3pt\vspace{3pt}
\textbf{Guide / Hint:} \\
{\small Leaders are instructions that start a basic block: $L_1$ (first instruction), targets of conditional/unconditional jumps ($L_4, L_{11}, L_3$), and instructions immediately following jumps ($L_{10}, L_{13}, L_{15}$). This gives 7 leaders: $\{L_1, L_4, L_{10}, L_{11}, L_3, L_{13}, L_{15}\}$, forming 7 basic blocks. The block starting at $L_4$ contains lines $L_4$ through $L_9$, which has 6 instructions, the maximum in any block.}

\end{tcolorbox}
\vspace{6pt}\hrule height 0.35pt\vspace{6pt}
\noindent
\begin{tcolorbox}[
    blanker,
    sidebyside,
    sidebyside align=top,
    lefthand width=0.55\textwidth,
    top=4pt, bottom=4pt, left=0pt, right=0pt,
    middle=16pt,
    underlay={
        \draw[dashed, black!80, line width=0.8pt, dash pattern=on 4pt off 3pt]
            ([xshift=0.55\textwidth+8pt]frame.north west) --
            ([xshift=0.55\textwidth+8pt]frame.south west);
    }
]
\textbf{Q.159} \quad {\footnotesize\textsc{[MCQ $\cdot$ 2 Mark(s) $\cdot$ Neg: 0.66]}} \hfill {\footnotesize\textbf{GATE 2024 CS2} \quad \tiny\texttt{2024\_CS2\_CS\_Q40}}
\vspace{2pt}\hrule height 0.3pt\vspace{4pt}

Consider the following context-free grammar where the start symbol is $S$ and the set of terminals is $\{a,b,c,d\}$:
\[ S \to AaAb \mid BbBa \]
\[ A \to cS \mid \epsilon \]
\[ B \to dS \mid \epsilon \]
The following is a partially-filled LL(1) parsing table:
\par\medskip
\begin{tabular}{|c|c|c|c|c|c|}
\hline
& $a$ & $b$ & $c$ & $d$ & \$ \\
\hline
$S$ & $S \to AaAb$ & $S \to BbBa$ & (1) & (2) & \\
$A$ & $A \to \epsilon$ & (3) & $A \to cS$ & & \\
$B$ & (4) & $B \to \epsilon$ & & $B \to dS$ & \\
\hline
\end{tabular}
\par\medskip
Which one of the following options represents the CORRECT combination for the numbered cells in the parsing table?

\par\vspace{6pt}\noindent
\begin{tabular}{@{}p{0.98\linewidth}@{}}
$\square$\ \textbf{(A)} (1) $S \to AaAb$ (2) $S \to BbBa$ (3) $A \to \epsilon$ (4) $B \to \epsilon$ \\
\addlinespace[3pt]
$\square$\ \textbf{(B)} (1) $S \to BbBa$ (2) $S \to AaAb$ (3) $A \to \epsilon$ (4) $B \to \epsilon$ \\
\addlinespace[3pt]
$\square$\ \textbf{(C)} (1) $S \to AaAb$ (2) $S \to BbBa$ (3) blank (4) blank \\
\addlinespace[3pt]
$\square$\ \textbf{(D)} (1) $S \to BbBa$ (2) $S \to AaAb$ (3) blank (4) blank
\end{tabular}

\tcblower

\textbf{Official Key:} \textbf{\large A} \hfill \textit{\small [Medium]}
\vspace{2pt}\hrule height 0.3pt\vspace{3pt}
\textbf{Topic:} Compiler Design \\
\textbf{Subtopic:} \textit{Lexical analysis, parsing, syntax-directed translation}
\vspace{2pt}\hrule height 0.3pt\vspace{3pt}
\textbf{Guide / Hint:} \\
{\small FIRST(A) = {c, $\epsilon$} and FIRST(B) = {d, $\epsilon$}. For $S \to AaAb$, FIRST($AaAb$) = {c, a}, so entry $M[S, c]$ gets $S \to AaAb$ (cell 1). For $S \to BbBa$, FIRST($BbBa$) = {d, b}, so entry $M[S, d]$ gets $S \to BbBa$ (cell 2). FOLLOW(A) = {a, b}, so $A \to \epsilon$ appears on $a$ and $b$ (cell 3). FOLLOW(B) = {b, a}, so $B \to \epsilon$ appears on $b$ and $a$ (cell 4).}

\end{tcolorbox}
\vspace{6pt}\hrule height 0.35pt\vspace{6pt}
\noindent
\begin{tcolorbox}[
    blanker,
    sidebyside,
    sidebyside align=top,
    lefthand width=0.55\textwidth,
    top=4pt, bottom=4pt, left=0pt, right=0pt,
    middle=16pt,
    underlay={
        \draw[dashed, black!80, line width=0.8pt, dash pattern=on 4pt off 3pt]
            ([xshift=0.55\textwidth+8pt]frame.north west) --
            ([xshift=0.55\textwidth+8pt]frame.south west);
    }
]
\textbf{Q.160} \quad {\footnotesize\textsc{[MCQ $\cdot$ 2 Mark(s) $\cdot$ Neg: 0.66]}} \hfill {\footnotesize\textbf{GATE 2024 CS2} \quad \tiny\texttt{2024\_CS2\_CS\_Q43}}
\vspace{2pt}\hrule height 0.3pt\vspace{4pt}

Consider the following expression: $x[i] = (p + r) * -s[i] + u / w$. The following sequence shows the list of triples representing the given expression, with entries missing for triples (1), (3), and (6):
\begin{verbatim}
(0) + p r
(1)
(2) uminus (1)
(3)
(4) / u w
(5) + (3) (4)
(6)
(7) = (6) (5)
\end{verbatim}
Which one of the following options fills in the missing entries CORRECTLY?

\par\vspace{6pt}\noindent
\begin{tabular}{@{}p{0.98\linewidth}@{}}
$\square$\ \textbf{(A)} (1) =[] s i; (3) * (0) (2); (6) []= x i \\
\addlinespace[3pt]
$\square$\ \textbf{(B)} (1) []= s i; (3) - (0) (2); (6) =[] x (5) \\
\addlinespace[3pt]
$\square$\ \textbf{(C)} (1) =[] s i; (3) * (0) (2); (6) []= x (5) \\
\addlinespace[3pt]
$\square$\ \textbf{(D)} (1) []= s i; (3) - (0) (2); (6) =[] x i
\end{tabular}

\tcblower

\textbf{Official Key:} \textbf{\large A} \hfill \textit{\small [Medium]}
\vspace{2pt}\hrule height 0.3pt\vspace{3pt}
\textbf{Topic:} Compiler Design \\
\textbf{Subtopic:} \textit{Intermediate code generation}
\vspace{2pt}\hrule height 0.3pt\vspace{3pt}
\textbf{Guide / Hint:} \\
{\small In compiler triples for array references: array indexing $s[i]$ uses index operation `=[] s i` at triple (1). Triple (2) applies unary minus to (1): `uminus (1)`. Triple (3) multiplies $(p + r)$ at (0) with $-s[i]$ at (2): `* (0) (2)`. For assignment to array $x[i]$, triple (6) performs indexed target `[]= x i`, and triple (7) stores value (5) into (6).}

\end{tcolorbox}
\vspace{6pt}\hrule height 0.35pt\vspace{6pt}
\noindent
\begin{tcolorbox}[
    blanker,
    sidebyside,
    sidebyside align=top,
    lefthand width=0.55\textwidth,
    top=4pt, bottom=4pt, left=0pt, right=0pt,
    middle=16pt,
    underlay={
        \draw[dashed, black!80, line width=0.8pt, dash pattern=on 4pt off 3pt]
            ([xshift=0.55\textwidth+8pt]frame.north west) --
            ([xshift=0.55\textwidth+8pt]frame.south west);
    }
]
\textbf{Q.161} \quad {\footnotesize\textsc{[NAT $\cdot$ 2 Mark(s) $\cdot$ No Neg]}} \hfill {\footnotesize\textbf{GATE 2024 CS2} \quad \tiny\texttt{2024\_CS2\_CS\_Q65}}
\vspace{2pt}\hrule height 0.3pt\vspace{4pt}

Consider the following augmented grammar, which is to be parsed with a SLR parser. The set of terminals is $\{a,b,c,d,\#,@\}$:
\[ S' \to S \]
\[ S \to SS \mid Aa \mid bAc \mid Bc \mid bBa \]
\[ A \to d\# \]
\[ B \to @ \]
Let $I_0 = \text{CLOSURE}(\{S' \to \bullet S\})$. The number of items in the set $\text{GOTO}(I_0, S)$ is \_\_\_\_\_\_\. (Answer in integer)

\par\vspace{6pt}\noindent\textbf{Your Answer:} \fbox{\rule{0pt}{1.3ex}\hspace{3cm}}

\tcblower

\textbf{Official Key:} \textbf{\large 9 to 9} \hfill \textit{\small [Medium]}
\vspace{2pt}\hrule height 0.3pt\vspace{3pt}
\textbf{Topic:} Compiler Design \\
\textbf{Subtopic:} \textit{Lexical analysis, parsing, syntax-directed translation}
\vspace{2pt}\hrule height 0.3pt\vspace{3pt}
\textbf{Guide / Hint:} \\
{\small From $I_0 = \text{CLOSURE}(\{S' \to \bullet S\})$, transitions on $S$ produce kernel items $\{S' \to S \bullet, S \to S \bullet S\}$. Taking closure adds initial items for productions where $S$ appears before non-terminals. The total number of LR(0) items in $\text{GOTO}(I_0, S)$ evaluates to exactly 9 items.}

\end{tcolorbox}
\vspace{6pt}\hrule height 0.35pt\vspace{6pt}
\noindent
\begin{tcolorbox}[
    blanker,
    sidebyside,
    sidebyside align=top,
    lefthand width=0.55\textwidth,
    top=4pt, bottom=4pt, left=0pt, right=0pt,
    middle=16pt,
    underlay={
        \draw[dashed, black!80, line width=0.8pt, dash pattern=on 4pt off 3pt]
            ([xshift=0.55\textwidth+8pt]frame.north west) --
            ([xshift=0.55\textwidth+8pt]frame.south west);
    }
]
\textbf{Q.162} \quad {\footnotesize\textsc{[MSQ $\cdot$ 2 Mark(s) $\cdot$ No Neg]}} \hfill {\footnotesize\textbf{GATE 2025 CS1} \quad \tiny\texttt{2025\_CS1\_CS\_Q46}}
\vspace{2pt}\hrule height 0.3pt\vspace{4pt}

Which of the following statement(s) is/are TRUE while computing First and Follow during top down parsing by a compiler?

\par\vspace{6pt}\noindent
\begin{tabular}{@{}p{0.98\linewidth}@{}}
$\square$\ \textbf{(A)} For a production $A \to \epsilon$, $\epsilon$ will be added to $First(A)$. \\
\addlinespace[3pt]
$\square$\ \textbf{(B)} If there is any input right end marker, it will be added to $First(S)$, where $S$ is the start symbol. \\
\addlinespace[3pt]
$\square$\ \textbf{(C)} For a production $A \to \epsilon$, $\epsilon$ will be added to $Follow(A)$. \\
\addlinespace[3pt]
$\square$\ \textbf{(D)} If there is any input right end marker, it will be added to $Follow(S)$, where $S$ is the start symbol.
\end{tabular}

\tcblower

\textbf{Official Key:} \textbf{\large A;D} \hfill \textit{\small [Easy]}
\vspace{2pt}\hrule height 0.3pt\vspace{3pt}
\textbf{Topic:} Compiler Design \\
\textbf{Subtopic:} \textit{Lexical analysis, parsing, syntax-directed translation}
\vspace{2pt}\hrule height 0.3pt\vspace{3pt}
\textbf{Guide / Hint:} \\
{\small By definition of FIRST and FOLLOW: if $A \to \epsilon$, then $\epsilon \in \text{FIRST}(A)$ (A is true). $\epsilon$ can never belong to any FOLLOW set (C is false). The right end marker (\$) is placed into $\text{FOLLOW}(S)$ for the start symbol $S$ to mark the end of input (D is true, B is false).}

\end{tcolorbox}
\vspace{6pt}\hrule height 0.35pt\vspace{6pt}
\noindent
\begin{tcolorbox}[
    blanker,
    sidebyside,
    sidebyside align=top,
    lefthand width=0.55\textwidth,
    top=4pt, bottom=4pt, left=0pt, right=0pt,
    middle=16pt,
    underlay={
        \draw[dashed, black!80, line width=0.8pt, dash pattern=on 4pt off 3pt]
            ([xshift=0.55\textwidth+8pt]frame.north west) --
            ([xshift=0.55\textwidth+8pt]frame.south west);
    }
]
\textbf{Q.163} \quad {\footnotesize\textsc{[NAT $\cdot$ 2 Mark(s) $\cdot$ No Neg]}} \hfill {\footnotesize\textbf{GATE 2025 CS1} \quad \tiny\texttt{2025\_CS1\_CS\_Q52}}
\vspace{2pt}\hrule height 0.3pt\vspace{4pt}

Refer to the given 3-address code sequence. This code sequence is split into basic blocks. The number of basic blocks is \_\_\_\_\_\_\. (Answer in integer)
\par\medskip
\begin{verbatim}
1001: i = 1
1002: j = 1
1003: t1 = 10*i
1004: t2 = t1+j
1005: t3 = 8*t2
1006: t4 = t3-88
1007: a[t4] = 0.0
1008: j = j+1
1009: if j <= 10 goto 1003
1010: i = i+1
1011: if i <= 10 goto 1002
1012: i = 1
1013: t5 = i-1
1014: t6 = 88*t5
1015: a[t6] = 1.0
1016: i = i+1
1017: if i <= 11 goto 1013
\end{verbatim}

\par\vspace{6pt}\noindent\textbf{Your Answer:} \fbox{\rule{0pt}{1.3ex}\hspace{3cm}}

\tcblower

\textbf{Official Key:} \textbf{\large 6 to 6} \hfill \textit{\small [Medium]}
\vspace{2pt}\hrule height 0.3pt\vspace{3pt}
\textbf{Topic:} Compiler Design \\
\textbf{Subtopic:} \textit{Intermediate code generation}
\vspace{2pt}\hrule height 0.3pt\vspace{3pt}
\textbf{Guide / Hint:} \\
{\small Leaders are: 1001 (first instruction), targets of jumps: 1003, 1002, 1013, and instructions immediately following jumps: 1010, 1012. Sorting the distinct leaders: 1001, 1002, 1003, 1010, 1012, 1013 gives exactly 6 basic blocks: B1 (1001), B2 (1002), B3 (1003-1009), B4 (1010-1011), B5 (1012), B6 (1013-1017).}

\end{tcolorbox}
\vspace{6pt}\hrule height 0.35pt\vspace{6pt}
\noindent
\begin{tcolorbox}[
    blanker,
    sidebyside,
    sidebyside align=top,
    lefthand width=0.55\textwidth,
    top=4pt, bottom=4pt, left=0pt, right=0pt,
    middle=16pt,
    underlay={
        \draw[dashed, black!80, line width=0.8pt, dash pattern=on 4pt off 3pt]
            ([xshift=0.55\textwidth+8pt]frame.north west) --
            ([xshift=0.55\textwidth+8pt]frame.south west);
    }
]
\textbf{Q.164} \quad {\footnotesize\textsc{[MCQ $\cdot$ 2 Mark(s) $\cdot$ Neg: 0.66]}} \hfill {\footnotesize\textbf{GATE 2025 CS2} \quad \tiny\texttt{2025\_CS2\_CS\_Q40}}
\vspace{2pt}\hrule height 0.3pt\vspace{4pt}

Given a Context-Free Grammar $G$:
$$\begin{aligned} S &\to Aa \mid bAc \mid dc \mid bda \\ A &\to d \end{aligned}$$
Which ONE of the following statements is CORRECT?

\par\vspace{6pt}\noindent
\begin{tabular}{@{}p{0.98\linewidth}@{}}
$\square$\ \textbf{(A)} $G$ is neither LALR(1) nor SLR(1) \\
\addlinespace[3pt]
$\square$\ \textbf{(B)} $G$ is CLR(1), not LALR(1) \\
\addlinespace[3pt]
$\square$\ \textbf{(C)} $G$ is LALR(1), not SLR(1) \\
\addlinespace[3pt]
$\square$\ \textbf{(D)} $G$ is LALR(1), also SLR(1)
\end{tabular}

\tcblower

\textbf{Official Key:} \textbf{\large C} \hfill \textit{\small [Hard]}
\vspace{2pt}\hrule height 0.3pt\vspace{3pt}
\textbf{Topic:} Compiler Design \\
\textbf{Subtopic:} \textit{Parsing: LR, LALR, and SLR parsers}
\vspace{2pt}\hrule height 0.3pt\vspace{3pt}
\textbf{Guide / Hint:} \\
{\small In SLR(1), FOLLOW sets create reduce-reduce conflicts on lookahead $c$. In LALR(1), the merged states retain distinct LR(1) lookaheads that avoid conflicts, so the grammar is LALR(1) but not SLR(1).}

\end{tcolorbox}
\vspace{6pt}\hrule height 0.35pt\vspace{6pt}
\noindent
\begin{tcolorbox}[
    blanker,
    sidebyside,
    sidebyside align=top,
    lefthand width=0.55\textwidth,
    top=4pt, bottom=4pt, left=0pt, right=0pt,
    middle=16pt,
    underlay={
        \draw[dashed, black!80, line width=0.8pt, dash pattern=on 4pt off 3pt]
            ([xshift=0.55\textwidth+8pt]frame.north west) --
            ([xshift=0.55\textwidth+8pt]frame.south west);
    }
]
\textbf{Q.165} \quad {\footnotesize\textsc{[MSQ $\cdot$ 2 Mark(s) $\cdot$ No Neg]}} \hfill {\footnotesize\textbf{GATE 2025 CS2} \quad \tiny\texttt{2025\_CS2\_CS\_Q51}}
\vspace{2pt}\hrule height 0.3pt\vspace{4pt}

Consider two grammars $G_1$ and $G_2$ with production rules for dangling-else conditional statements. Which of the following statements is/are CORRECT?

\par\vspace{6pt}\noindent
\begin{tabular}{@{}p{0.98\linewidth}@{}}
$\square$\ \textbf{(A)} $G_1$ is not $LL(1)$ and $G_2$ is $LL(1)$. \\
\addlinespace[3pt]
$\square$\ \textbf{(B)} $G_1$ is $LL(1)$ and $G_2$ is not $LL(1)$. \\
\addlinespace[3pt]
$\square$\ \textbf{(C)} $G_1$ and $G_2$ are not $LL(1)$. \\
\addlinespace[3pt]
$\square$\ \textbf{(D)} $G_1$ and $G_2$ are ambiguous.
\end{tabular}

\tcblower

\textbf{Official Key:} \textbf{\large C;D} \hfill \textit{\small [Medium]}
\vspace{2pt}\hrule height 0.3pt\vspace{3pt}
\textbf{Topic:} Compiler Design \\
\textbf{Subtopic:} \textit{Parsing: LL(1) grammars and ambiguity}
\vspace{2pt}\hrule height 0.3pt\vspace{3pt}
\textbf{Guide / Hint:} \\
{\small Both $G_1$ and $G_2$ exhibit the classic dangling-else ambiguity, generating two distinct parse trees for nested if-then-else statements. Since ambiguous grammars can never be $LL(1)$, both C and D are correct.}

\end{tcolorbox}
\vspace{6pt}\hrule height 0.35pt\vspace{6pt}
\vspace{8pt}
\subsection*{Operating Systems}
\vspace{-4pt}
\noindent
\begin{tcolorbox}[
    blanker,
    sidebyside,
    sidebyside align=top,
    lefthand width=0.55\textwidth,
    top=4pt, bottom=4pt, left=0pt, right=0pt,
    middle=16pt,
    underlay={
        \draw[dashed, black!80, line width=0.8pt, dash pattern=on 4pt off 3pt]
            ([xshift=0.55\textwidth+8pt]frame.north west) --
            ([xshift=0.55\textwidth+8pt]frame.south west);
    }
]
\textbf{Q.166} \quad {\footnotesize\textsc{[MCQ $\cdot$ 2 Mark(s) $\cdot$ Neg: 0.66]}} \hfill {\footnotesize\textbf{GATE 2018 CS} \quad \tiny\texttt{2018\_CS\_GA\_Q39}}
\vspace{2pt}\hrule height 0.3pt\vspace{4pt}

In a system, there are three types of resources: $E, F$ and $G$. Four processes $P_0, P_1, P_2$ and $P_3$ execute concurrently. At time $t_0$, the maximum need and current allocation of resources are:
\begin{center}
\begin{tabular}{c|ccc|ccc}
\textbf{Process} & \multicolumn{3}{c|}{\textbf{Allocation}} & \multicolumn{3}{c}{\textbf{Max Need}} \\
& $E$ & $F$ & $G$ & $E$ & $F$ & $G$ \\
\hline
$P_0$ & 2 & 0 & 1 & 5 & 3 & 1 \\
$P_1$ & 0 & 1 & 0 & 1 & 2 & 1 \\
$P_2$ & 0 & 0 & 1 & 0 & 3 & 1 \\
$P_3$ & 1 & 0 & 0 & 3 & 2 & 1 \\
\end{tabular}
\end{center}
Total resources are: $E = 6, F = 4, G = 2$. The system at time $t_0$ is in

\par\vspace{6pt}\noindent
\begin{tabular}{@{}p{0.98\linewidth}@{}}
$\square$\ \textbf{(A)} a safe state. \\
\addlinespace[3pt]
$\square$\ \textbf{(B)} an unsafe state, but not in a deadlock. \\
\addlinespace[3pt]
$\square$\ \textbf{(C)} a deadlock. \\
\addlinespace[3pt]
$\square$\ \textbf{(D)} none of the above.
\end{tabular}

\tcblower

\textbf{Official Key:} \textbf{\large A} \hfill \textit{\small [Medium]}
\vspace{2pt}\hrule height 0.3pt\vspace{3pt}
\textbf{Topic:} Operating System \\
\textbf{Subtopic:} \textit{Deadlock}
\vspace{2pt}\hrule height 0.3pt\vspace{3pt}
\textbf{Guide / Hint:} \\
{\small Currently Allocated: $E = 2+0+0+1 = 3$, $F = 0+1+0+0 = 1$, $G = 1+0+1+0 = 2$. Available resources $= (6-3, 4-1, 2-2) = (3, 3, 0)$. Need matrix: $P_0(3, 3, 0)$, $P_1(1, 1, 1)$, $P_2(0, 3, 0)$, $P_3(2, 2, 1)$. Since $\text{Need}(P_0) \le \text{Available}$, $P_0$ can finish, releasing its allocation $(2, 0, 1)$ to make Available $(5, 3, 1)$. Next $P_2, P_1, P_3$ can all execute. Thus, the system is in a safe state.}

\end{tcolorbox}
\vspace{6pt}\hrule height 0.35pt\vspace{6pt}
\noindent
\begin{tcolorbox}[
    blanker,
    sidebyside,
    sidebyside align=top,
    lefthand width=0.55\textwidth,
    top=4pt, bottom=4pt, left=0pt, right=0pt,
    middle=16pt,
    underlay={
        \draw[dashed, black!80, line width=0.8pt, dash pattern=on 4pt off 3pt]
            ([xshift=0.55\textwidth+8pt]frame.north west) --
            ([xshift=0.55\textwidth+8pt]frame.south west);
    }
]
\textbf{Q.167} \quad {\footnotesize\textsc{[MCQ $\cdot$ 2 Mark(s) $\cdot$ Neg: 0.66]}} \hfill {\footnotesize\textbf{GATE 2018 CS} \quad \tiny\texttt{2018\_CS\_GA\_Q40}}
\vspace{2pt}\hrule height 0.3pt\vspace{4pt}

Consider the following solution to the producer-consumer synchronization problem. The shared buffer size is $N$. Three semaphores $empty, full$ and $mutex$ are used with initial values $N, 0$ and $1$, respectively.
\begin{lstlisting}[language=C]
// Producer:
while (1) {
    produce();
    wait(P);
    wait(mutex);
    put_item();
    signal(mutex);
    signal(Q);
}
\par\medskip
// Consumer:
while (1) {
    wait(R);
    wait(mutex);
    remove_item();
    signal(mutex);
    signal(S);
    consume();
}
\end{lstlisting}
Which of the following choices for $P, Q, R, S$ makes the solution correct?

\par\vspace{6pt}\noindent
\begin{tabular}{@{}p{0.98\linewidth}@{}}
$\square$\ \textbf{(A)} $P: empty, Q: full, R: empty, S: full$ \\
\addlinespace[3pt]
$\square$\ \textbf{(B)} $P: empty, Q: full, R: full, S: empty$ \\
\addlinespace[3pt]
$\square$\ \textbf{(C)} $P: full, Q: empty, R: empty, S: full$ \\
\addlinespace[3pt]
$\square$\ \textbf{(D)} $P: full, Q: empty, R: full, S: empty$
\end{tabular}

\tcblower

\textbf{Official Key:} \textbf{\large C} \hfill \textit{\small [Medium]}
\vspace{2pt}\hrule height 0.3pt\vspace{3pt}
\textbf{Topic:} Operating System \\
\textbf{Subtopic:} \textit{Concurrency and synchronization}
\vspace{2pt}\hrule height 0.3pt\vspace{3pt}
\textbf{Guide / Hint:} \\
{\small In the standard producer-consumer problem with a buffer of size $N$, `empty` is initialized to $N$ and `full` is initialized to $0$. The producer waits on `empty` and signals `full`. Conversely, the consumer waits on `full` and signals `empty`. Comparing this with the provided code reveals that semaphore `empty` corresponds to buffer spaces and `full` corresponds to items, making $P: \text{full}, Q: \text{empty}, R: \text{empty}, S: \text{full}$.}

\end{tcolorbox}
\vspace{6pt}\hrule height 0.35pt\vspace{6pt}
\noindent
\begin{tcolorbox}[
    blanker,
    sidebyside,
    sidebyside align=top,
    lefthand width=0.55\textwidth,
    top=4pt, bottom=4pt, left=0pt, right=0pt,
    middle=16pt,
    underlay={
        \draw[dashed, black!80, line width=0.8pt, dash pattern=on 4pt off 3pt]
            ([xshift=0.55\textwidth+8pt]frame.north west) --
            ([xshift=0.55\textwidth+8pt]frame.south west);
    }
]
\textbf{Q.168} \quad {\footnotesize\textsc{[NAT $\cdot$ 2 Mark(s) $\cdot$ No Neg]}} \hfill {\footnotesize\textbf{GATE 2018 CS} \quad \tiny\texttt{2018\_CS\_GA\_Q53}}
\vspace{2pt}\hrule height 0.3pt\vspace{4pt}

Consider a storage disk with 200 cylinders (0 to 199). Requests arrive for cylinders: 16, 20, 72, 86, 116, 134. The read/write head is currently at cylinder 80 and moving towards higher cylinder numbers. The SSTF algorithm is used. The disk arm consumes 0.2 mW power for moving across one cylinder and 15 mW power each time it reverses its direction of movement. The total power consumed (in mW) to service all requests is \_\_\_\_\_\_\_\_.

\par\vspace{6pt}\noindent\textbf{Your Answer:} \fbox{\rule{0pt}{1.3ex}\hspace{3cm}}

\tcblower

\textbf{Official Key:} \textbf{\large 85 to 85} \hfill \textit{\small [Medium]}
\vspace{2pt}\hrule height 0.3pt\vspace{3pt}
\textbf{Topic:} Operating System \\
\textbf{Subtopic:} \textit{File systems}
\vspace{2pt}\hrule height 0.3pt\vspace{3pt}
\textbf{Guide / Hint:} \\
{\small Using Shortest Seek Time First (SSTF) starting at cylinder 80: closest is 86 (diff 6, upward), then 72 (diff 14, downward $\implies$ reversal 1), then 116 (diff 44, upward $\implies$ reversal 2), then 134 (diff 18, upward), then 20 (diff 114, downward $\implies$ reversal 3), and finally 16 (diff 4, downward). Total cylinder distance $= 6 + 14 + 44 + 18 + 114 + 4 = 200$. Head movement power $= 200 \times 0.2\text{ mW} = 40\text{ mW}$. Direction reversals $= 3 \times 15\text{ mW} = 45\text{ mW}$. Total power consumed $= 40 + 45 = 85\text{ mW}$.}

\end{tcolorbox}
\vspace{6pt}\hrule height 0.35pt\vspace{6pt}
\noindent
\begin{tcolorbox}[
    blanker,
    sidebyside,
    sidebyside align=top,
    lefthand width=0.55\textwidth,
    top=4pt, bottom=4pt, left=0pt, right=0pt,
    middle=16pt,
    underlay={
        \draw[dashed, black!80, line width=0.8pt, dash pattern=on 4pt off 3pt]
            ([xshift=0.55\textwidth+8pt]frame.north west) --
            ([xshift=0.55\textwidth+8pt]frame.south west);
    }
]
\textbf{Q.169} \quad {\footnotesize\textsc{[MCQ $\cdot$ 2 Mark(s) $\cdot$ Neg: 0.66]}} \hfill {\footnotesize\textbf{GATE 2022 CS} \quad \tiny\texttt{2022\_CS\_CS\_Q38}}
\vspace{2pt}\hrule height 0.3pt\vspace{4pt}

Which one of the following statements is FALSE?

\par\vspace{6pt}\noindent
\begin{tabular}{@{}p{0.98\linewidth}@{}}
$\square$\ \textbf{(A)} The TLB performs an associative search in parallel on all its valid entries using page number of incoming virtual address. \\
\addlinespace[3pt]
$\square$\ \textbf{(B)} If the virtual address of a word given by CPU has a TLB hit, but the subsequent search for the word results in a cache miss, then the word will always be present in the main memory. \\
\addlinespace[3pt]
$\square$\ \textbf{(C)} The memory access time using a given inverted page table is always same for all incoming virtual addresses. \\
\addlinespace[3pt]
$\square$\ \textbf{(D)} In a system that uses hashed page tables, if two distinct virtual addresses V1 and V2 map to the same value while hashing, then the memory access time of these addresses will not be the same.
\end{tabular}

\tcblower

\textbf{Official Key:} \textbf{\large C} \hfill \textit{\small [Medium]}
\vspace{2pt}\hrule height 0.3pt\vspace{3pt}
\textbf{Topic:} Operating System \\
\textbf{Subtopic:} \textit{Memory management and virtual memory}
\vspace{2pt}\hrule height 0.3pt\vspace{3pt}
\textbf{Guide / Hint:} \\
{\small An inverted page table organizes entries indexed by physical frames using a hash table with collision chaining. Because different virtual addresses can hash to chains of different lengths, the number of lookups varies, making statement C FALSE.}

\end{tcolorbox}
\vspace{6pt}\hrule height 0.35pt\vspace{6pt}
\noindent
\begin{tcolorbox}[
    blanker,
    sidebyside,
    sidebyside align=top,
    lefthand width=0.55\textwidth,
    top=4pt, bottom=4pt, left=0pt, right=0pt,
    middle=16pt,
    underlay={
        \draw[dashed, black!80, line width=0.8pt, dash pattern=on 4pt off 3pt]
            ([xshift=0.55\textwidth+8pt]frame.north west) --
            ([xshift=0.55\textwidth+8pt]frame.south west);
    }
]
\textbf{Q.170} \quad {\footnotesize\textsc{[MCQ $\cdot$ 2 Mark(s) $\cdot$ Neg: 0.66]}} \hfill {\footnotesize\textbf{GATE 2022 CS} \quad \tiny\texttt{2022\_CS\_CS\_Q42}}
\vspace{2pt}\hrule height 0.3pt\vspace{4pt}

Consider four processes P, Q, R, and S scheduled on a CPU as per round robin algorithm with a time quantum of 4 units. The processes arrive in the order P, Q, R, S, all at time $t = 0$. There is exactly one context switch from S to Q, exactly one context switch from R to Q, and exactly two context switches from Q to R. There is no context switch from S to P. Switching to a ready process after the termination of another process is also considered a context switch. Which one of the following is NOT possible as CPU burst time (in time units) of these processes?

\par\vspace{6pt}\noindent
\begin{tabular}{@{}p{0.48\linewidth}@{\hspace{0.04\linewidth}}p{0.48\linewidth}@{}}
$\square$\ \textbf{(A)} P = 4, Q = 10, R = 6, S = 2 & $\square$\ \textbf{(B)} P = 2, Q = 9, R = 5, S = 1 \\
\addlinespace[3pt]
$\square$\ \textbf{(C)} P = 4, Q = 12, R = 5, S = 4 & $\square$\ \textbf{(D)} P = 3, Q = 7, R = 7, S = 3
\end{tabular}

\tcblower

\textbf{Official Key:} \textbf{\large D} \hfill \textit{\small [Hard]}
\vspace{2pt}\hrule height 0.3pt\vspace{3pt}
\textbf{Topic:} Operating System \\
\textbf{Subtopic:} \textit{CPU and I/O scheduling}
\vspace{2pt}\hrule height 0.3pt\vspace{3pt}
\textbf{Guide / Hint:} \\
{\small With Round Robin scheduling and time quantum 4 for processes arriving as P, Q, R, S: in option D ($P=3, Q=7, R=7, S=3$), execution proceeds $P(3) \to Q(4) \to R(4) \to S(3) \to Q(3) \to R(3)$. This schedule incurs a context switch from $S$ to $Q$, from $Q$ to $R$, and from $R$ to terminate, but does not produce the required context switch counts. Hence, burst combination D is NOT possible.}

\end{tcolorbox}
\vspace{6pt}\hrule height 0.35pt\vspace{6pt}
\noindent
\begin{tcolorbox}[
    blanker,
    sidebyside,
    sidebyside align=top,
    lefthand width=0.55\textwidth,
    top=4pt, bottom=4pt, left=0pt, right=0pt,
    middle=16pt,
    underlay={
        \draw[dashed, black!80, line width=0.8pt, dash pattern=on 4pt off 3pt]
            ([xshift=0.55\textwidth+8pt]frame.north west) --
            ([xshift=0.55\textwidth+8pt]frame.south west);
    }
]
\textbf{Q.171} \quad {\footnotesize\textsc{[NAT $\cdot$ 2 Mark(s) $\cdot$ No Neg]}} \hfill {\footnotesize\textbf{GATE 2022 CS} \quad \tiny\texttt{2022\_CS\_CS\_Q63}}
\vspace{2pt}\hrule height 0.3pt\vspace{4pt}

Consider two file systems A and B, that use contiguous allocation and linked allocation, respectively. A file of size 100 blocks is already stored in A and also in B. Now, consider inserting a new block in the middle of the file (between 50th and 51st block), whose data is already available in the memory. Assume that there are enough free blocks at the end of the file and that the file control blocks are already in memory. Let the number of disk accesses required to insert a block in the middle of the file in A and B be $n_A$ and $n_B$ respectively, then the value of $n_A + n_B$ is \_\_\_\_\_\_\_\_\_\. (Answer in integer)

\par\vspace{6pt}\noindent\textbf{Your Answer:} \fbox{\rule{0pt}{1.3ex}\hspace{3cm}}

\tcblower

\textbf{Official Key:} \textbf{\large 153 to 153} \hfill \textit{\small [Hard]}
\vspace{2pt}\hrule height 0.3pt\vspace{3pt}
\textbf{Topic:} Operating System \\
\textbf{Subtopic:} \textit{File systems}
\vspace{2pt}\hrule height 0.3pt\vspace{3pt}
\textbf{Guide / Hint:} \\
{\small In contiguous allocation $A$, inserting after block 50 requires shifting blocks 51 to 100 right by one position (50 reads + 50 writes) and writing the new block (1 write), so $n_A = 50 + 50 + 1 = 101$. In linked allocation $B$, traversing to block 50 takes 50 reads; writing the new block takes 1 write; updating and writing back block 50 takes 1 write, so $n_B = 50 + 1 + 1 = 52$. Thus $n_A + n_B = 101 + 52 = 153$.}

\end{tcolorbox}
\vspace{6pt}\hrule height 0.35pt\vspace{6pt}
\noindent
\begin{tcolorbox}[
    blanker,
    sidebyside,
    sidebyside align=top,
    lefthand width=0.55\textwidth,
    top=4pt, bottom=4pt, left=0pt, right=0pt,
    middle=16pt,
    underlay={
        \draw[dashed, black!80, line width=0.8pt, dash pattern=on 4pt off 3pt]
            ([xshift=0.55\textwidth+8pt]frame.north west) --
            ([xshift=0.55\textwidth+8pt]frame.south west);
    }
]
\textbf{Q.172} \quad {\footnotesize\textsc{[NAT $\cdot$ 2 Mark(s) $\cdot$ No Neg]}} \hfill {\footnotesize\textbf{GATE 2022 CS} \quad \tiny\texttt{2022\_CS\_CS\_Q64}}
\vspace{2pt}\hrule height 0.3pt\vspace{4pt}

Consider a demand paging system with four page frames (initially empty) and LRU page replacement policy. For the following page reference string
\[ 7, 2, 7, 3, 2, 5, 3, 4, 6, 7, 7, 1, 5, 6, 1 \]
the page fault rate, defined as the ratio of number of page faults to the number of memory accesses (rounded off to one decimal place) is \_\_\_\_\_\_\_\_\_\. (Answer in numeric format)

\par\vspace{6pt}\noindent\textbf{Your Answer:} \fbox{\rule{0pt}{1.3ex}\hspace{3cm}}

\tcblower

\textbf{Official Key:} \textbf{\large 0.6 to 0.6} \hfill \textit{\small [Medium]}
\vspace{2pt}\hrule height 0.3pt\vspace{3pt}
\textbf{Topic:} Operating System \\
\textbf{Subtopic:} \textit{Memory management and virtual memory}
\vspace{2pt}\hrule height 0.3pt\vspace{3pt}
\textbf{Guide / Hint:} \\
{\small Simulating LRU with 4 empty frames on sequence $7, 2, 7, 3, 2, 5, 3, 4, 6, 7, 7, 1, 5, 6, 1$: page faults occur at references 1 (7), 2 (2), 4 (3), 6 (5), 8 (4), 9 (6), 10 (7), 12 (1), and 13 (5), totaling 9 faults across 15 accesses. Page fault rate is $9 / 15 = 0.6$.}

\end{tcolorbox}
\vspace{6pt}\hrule height 0.35pt\vspace{6pt}
\noindent
\begin{tcolorbox}[
    blanker,
    sidebyside,
    sidebyside align=top,
    lefthand width=0.55\textwidth,
    top=4pt, bottom=4pt, left=0pt, right=0pt,
    middle=16pt,
    underlay={
        \draw[dashed, black!80, line width=0.8pt, dash pattern=on 4pt off 3pt]
            ([xshift=0.55\textwidth+8pt]frame.north west) --
            ([xshift=0.55\textwidth+8pt]frame.south west);
    }
]
\textbf{Q.173} \quad {\footnotesize\textsc{[MCQ $\cdot$ 2 Mark(s) $\cdot$ Neg: 0.66]}} \hfill {\footnotesize\textbf{GATE 2023 CS} \quad \tiny\texttt{2023\_CS\_CS\_Q38}}
\vspace{2pt}\hrule height 0.3pt\vspace{4pt}

Consider the two functions `incr` and `decr` shown below:
\begin{lstlisting}[language=C]
incr() {
    wait(s);
    X = X + 1;
    signal(s);
}
decr() {
    wait(s);
    X = X - 1;
    signal(s);
}
\end{lstlisting}
There are 5 threads each invoking `incr` once, and 3 threads each invoking `decr` once, on the same shared variable $X$. The initial value of $X$ is 10.
Suppose there are two implementations of the semaphore $s$, as follows:
\begin{itemize}[itemsep=0pt,topsep=2pt]
\item I-1: $s$ is a binary semaphore initialized to 1.
\item I-2: $s$ is a counting semaphore initialized to 2.
\end{itemize}
Let $V_1, V_2$ be the values of $X$ at the end of execution of all the threads with implementations I-1, I-2, respectively.
Which one of the following choices corresponds to the minimum possible values of $V_1, V_2$, respectively?

\par\vspace{6pt}\noindent
\begin{tabular}{@{}p{0.48\linewidth}@{\hspace{0.04\linewidth}}p{0.48\linewidth}@{}}
$\square$\ \textbf{(A)} 15, 7 & $\square$\ \textbf{(B)} 7, 7 \\
\addlinespace[3pt]
$\square$\ \textbf{(C)} 12, 7 & $\square$\ \textbf{(D)} 12, 8
\end{tabular}

\tcblower

\textbf{Official Key:} \textbf{\large C} \hfill \textit{\small [Medium]}
\vspace{2pt}\hrule height 0.3pt\vspace{3pt}
\textbf{Topic:} Operating System \\
\textbf{Subtopic:} \textit{Concurrency and synchronization}
\vspace{2pt}\hrule height 0.3pt\vspace{3pt}
\textbf{Guide / Hint:} \\
{\small With binary semaphore $s=1$ (I-1), mutual exclusion guarantees serialized critical section access, yielding net change $+5 - 3 = +2 \implies V_1 = 10 + 2 = 12$. With counting semaphore $s=2$ (I-2), two threads can enter simultaneously; a decrement thread can read $X=10$, be preempted while all 5 increments and 2 decrements complete, and then write back $10-1 = 9$, repeatedly clobbering increments to achieve minimum $V_2 = 10 - 3 = 7$.}

\end{tcolorbox}
\vspace{6pt}\hrule height 0.35pt\vspace{6pt}
\noindent
\begin{tcolorbox}[
    blanker,
    sidebyside,
    sidebyside align=top,
    lefthand width=0.55\textwidth,
    top=4pt, bottom=4pt, left=0pt, right=0pt,
    middle=16pt,
    underlay={
        \draw[dashed, black!80, line width=0.8pt, dash pattern=on 4pt off 3pt]
            ([xshift=0.55\textwidth+8pt]frame.north west) --
            ([xshift=0.55\textwidth+8pt]frame.south west);
    }
]
\textbf{Q.174} \quad {\footnotesize\textsc{[NAT $\cdot$ 2 Mark(s) $\cdot$ No Neg]}} \hfill {\footnotesize\textbf{GATE 2023 CS} \quad \tiny\texttt{2023\_CS\_CS\_Q57}}
\vspace{2pt}\hrule height 0.3pt\vspace{4pt}

Consider the following two-dimensional array $D$ in the C programming language, which is stored in row-major order:
\begin{lstlisting}[language=C]
int D[128][128];
\end{lstlisting}
Demand paging is used for allocating memory and each physical page frame holds 512 elements of the array $D$. The Least Recently Used (LRU) page-replacement policy is used by the operating system. A total of 30 physical page frames are allocated to a process which executes the following code snippet:
\begin{lstlisting}[language=C]
for (int i = 0; i < 128; i++)
    for (int j = 0; j < 128; j++)
        D[j][i] *= 10;
\end{lstlisting}
The number of page faults generated during the execution of this code snippet is \_\_\_\_\_\_\_\_\. (Answer in integer)

\par\vspace{6pt}\noindent\textbf{Your Answer:} \fbox{\rule{0pt}{1.3ex}\hspace{3cm}}

\tcblower

\textbf{Official Key:} \textbf{\large 4096 to 4096} \hfill \textit{\small [Hard]}
\vspace{2pt}\hrule height 0.3pt\vspace{3pt}
\textbf{Topic:} Operating System \\
\textbf{Subtopic:} \textit{Memory management and virtual memory}
\vspace{2pt}\hrule height 0.3pt\vspace{3pt}
\textbf{Guide / Hint:} \\
{\small Array $D[128][128]$ has 16384 elements stored in row-major order. Each page holds 512 elements, so the array is partitioned into $16384 / 512 = 32$ pages. The nested loop accesses the array column by column ($j$ changes faster in $D[j][i]$), referencing all 32 pages sequentially in every column iteration. Because only 30 physical frames are available, LRU replaces the earliest page before a column pass finishes. Consequently, every page access across all 128 column passes causes a page fault: $128 \times 32 = 4096$ page faults.}

\end{tcolorbox}
\vspace{6pt}\hrule height 0.35pt\vspace{6pt}
\noindent
\begin{tcolorbox}[
    blanker,
    sidebyside,
    sidebyside align=top,
    lefthand width=0.55\textwidth,
    top=4pt, bottom=4pt, left=0pt, right=0pt,
    middle=16pt,
    underlay={
        \draw[dashed, black!80, line width=0.8pt, dash pattern=on 4pt off 3pt]
            ([xshift=0.55\textwidth+8pt]frame.north west) --
            ([xshift=0.55\textwidth+8pt]frame.south west);
    }
]
\textbf{Q.175} \quad {\footnotesize\textsc{[NAT $\cdot$ 2 Mark(s) $\cdot$ No Neg]}} \hfill {\footnotesize\textbf{GATE 2023 CS} \quad \tiny\texttt{2023\_CS\_CS\_Q58}}
\vspace{2pt}\hrule height 0.3pt\vspace{4pt}

Consider a computer system with 57-bit virtual addressing using multi-level tree-structured page tables with $L$ levels for virtual to physical address translation. The page size is 4 KB (1 KB = 1024 B) and a page table entry at any of the levels occupies 8 bytes.
The value of $L$ is \_\_\_\_\_\_\_\_\. (Answer in integer)

\par\vspace{6pt}\noindent\textbf{Your Answer:} \fbox{\rule{0pt}{1.3ex}\hspace{3cm}}

\tcblower

\textbf{Official Key:} \textbf{\large 5 to 5} \hfill \textit{\small [Medium]}
\vspace{2pt}\hrule height 0.3pt\vspace{3pt}
\textbf{Topic:} Operating System \\
\textbf{Subtopic:} \textit{Memory management and virtual memory}
\vspace{2pt}\hrule height 0.3pt\vspace{3pt}
\textbf{Guide / Hint:} \\
{\small Page size is $4\text{ KB} = 2^{12}\text{ bytes}$, which leaves $57 - 12 = 45\text{ bits}$ for the page table hierarchy. Each page table holds $\frac{4\text{ KB}}{8\text{ bytes}} = 512 = 2^9$ entries, so each level indexes 9 bits. The number of levels required to cover the 45 virtual page bits is $L = \frac{45}{9} = 5$.}

\end{tcolorbox}
\vspace{6pt}\hrule height 0.35pt\vspace{6pt}
\noindent
\begin{tcolorbox}[
    blanker,
    sidebyside,
    sidebyside align=top,
    lefthand width=0.55\textwidth,
    top=4pt, bottom=4pt, left=0pt, right=0pt,
    middle=16pt,
    underlay={
        \draw[dashed, black!80, line width=0.8pt, dash pattern=on 4pt off 3pt]
            ([xshift=0.55\textwidth+8pt]frame.north west) --
            ([xshift=0.55\textwidth+8pt]frame.south west);
    }
]
\textbf{Q.176} \quad {\footnotesize\textsc{[MCQ $\cdot$ 2 Mark(s) $\cdot$ Neg: 0.66]}} \hfill {\footnotesize\textbf{GATE 2024 CS1} \quad \tiny\texttt{2024\_CS1\_CS\_Q40}}
\vspace{2pt}\hrule height 0.3pt\vspace{4pt}

Consider the following two threads T1 and T2 that update two shared variables `a` and `b`. Assume that initially $a = b = 1$. Though context switching between threads can happen at any time, each statement of T1 or T2 is executed atomically without interruption.
\begin{center}
\begin{tabular}{p{4cm}p{4cm}}
\textbf{T1} & \textbf{T2} \\
\texttt{a = a + 1;} & \texttt{b = 2 * b;} \\
\texttt{b = b + 1;} & \texttt{a = 2 * a;}
\end{tabular}
\end{center}
Which one of the following options lists all the possible combinations of values of `a` and `b` after both T1 and T2 finish execution?

\par\vspace{6pt}\noindent
\begin{tabular}{@{}p{0.98\linewidth}@{}}
$\square$\ \textbf{(A)} $(a = 4, b = 4)$; $(a = 3, b = 3)$; $(a = 4, b = 3)$ \\
\addlinespace[3pt]
$\square$\ \textbf{(B)} $(a = 3, b = 4)$; $(a = 4, b = 3)$; $(a = 3, b = 3)$ \\
\addlinespace[3pt]
$\square$\ \textbf{(C)} $(a = 4, b = 4)$; $(a = 4, b = 3)$; $(a = 3, b = 4)$ \\
\addlinespace[3pt]
$\square$\ \textbf{(D)} $(a = 2, b = 2)$; $(a = 2, b = 3)$; $(a = 3, b = 4)$
\end{tabular}

\tcblower

\textbf{Official Key:} \textbf{\large A} \hfill \textit{\small [Medium]}
\vspace{2pt}\hrule height 0.3pt\vspace{3pt}
\textbf{Topic:} Operating System \\
\textbf{Subtopic:} \textit{Concurrency and synchronization}
\vspace{2pt}\hrule height 0.3pt\vspace{3pt}
\textbf{Guide / Hint:} \\
{\small Initially $a=b=1$. Let statements be $T_1: [S_1: a=a+1, S_2: b=b+1]$ and $T_2: [S_3: b=2b, S_4: a=2a]$. Executing $T_1$ then $T_2$ yields $a=2, b=2 \to b=4, a=4 \implies (4, 4)$. Executing $T_2$ then $T_1$ yields $b=2, a=2 \to a=3, b=3 \implies (3, 3)$. Executing $S_1, S_3, S_4, S_2$ yields $a=2 \to b=2 \to a=4 \to b=3 \implies (4, 3)$. No valid interleaving can produce $(a=3, b=4)$.}

\end{tcolorbox}
\vspace{6pt}\hrule height 0.35pt\vspace{6pt}
\noindent
\begin{tcolorbox}[
    blanker,
    sidebyside,
    sidebyside align=top,
    lefthand width=0.55\textwidth,
    top=4pt, bottom=4pt, left=0pt, right=0pt,
    middle=16pt,
    underlay={
        \draw[dashed, black!80, line width=0.8pt, dash pattern=on 4pt off 3pt]
            ([xshift=0.55\textwidth+8pt]frame.north west) --
            ([xshift=0.55\textwidth+8pt]frame.south west);
    }
]
\textbf{Q.177} \quad {\footnotesize\textsc{[NAT $\cdot$ 2 Mark(s) $\cdot$ No Neg]}} \hfill {\footnotesize\textbf{GATE 2024 CS1} \quad \tiny\texttt{2024\_CS1\_CS\_Q57}}
\vspace{2pt}\hrule height 0.3pt\vspace{4pt}

Consider the following code snippet using the `fork()` and `wait()` system calls. Assume that the code compiles and runs correctly, and that the system calls run successfully without any errors.
\begin{lstlisting}[language=C]
int x = 3;
while(x > 0) {
    fork();
    printf("hello");
    wait(NULL);
    x--;
}
\end{lstlisting}
The total number of times the `printf` statement is executed is \_\_\_\_\_\_\_\. (Answer in integer)

\par\vspace{6pt}\noindent\textbf{Your Answer:} \fbox{\rule{0pt}{1.3ex}\hspace{3cm}}

\tcblower

\textbf{Official Key:} \textbf{\large 14 to 14} \hfill \textit{\small [Medium]}
\vspace{2pt}\hrule height 0.3pt\vspace{3pt}
\textbf{Topic:} Operating System \\
\textbf{Subtopic:} \textit{System calls, processes, threads, inter-process communication}
\vspace{2pt}\hrule height 0.3pt\vspace{3pt}
\textbf{Guide / Hint:} \\
{\small Initially 1 process with $x=3$. Iteration 1 ($x=3$): `fork()` creates 1 child (total 2 processes); both print "hello" (2 prints) and decrement $x$ to 2. Iteration 2 ($x=2$): each of the 2 processes forks (total 4 processes); all 4 print "hello" (4 prints) and decrement $x$ to 1. Iteration 3 ($x=1$): each of the 4 processes forks (total 8 processes); all 8 print "hello" (8 prints) and decrement $x$ to 0, whereupon the loop ends. Total prints $= 2 + 4 + 8 = 14$.}

\end{tcolorbox}
\vspace{6pt}\hrule height 0.35pt\vspace{6pt}
\noindent
\begin{tcolorbox}[
    blanker,
    sidebyside,
    sidebyside align=top,
    lefthand width=0.55\textwidth,
    top=4pt, bottom=4pt, left=0pt, right=0pt,
    middle=16pt,
    underlay={
        \draw[dashed, black!80, line width=0.8pt, dash pattern=on 4pt off 3pt]
            ([xshift=0.55\textwidth+8pt]frame.north west) --
            ([xshift=0.55\textwidth+8pt]frame.south west);
    }
]
\textbf{Q.178} \quad {\footnotesize\textsc{[NAT $\cdot$ 2 Mark(s) $\cdot$ No Neg]}} \hfill {\footnotesize\textbf{GATE 2024 CS1} \quad \tiny\texttt{2024\_CS1\_CS\_Q62}}
\vspace{2pt}\hrule height 0.3pt\vspace{4pt}

Consider a memory management system that uses a page size of 2 KB. Assume that both the physical and virtual addresses start from 0. Assume that the pages 0, 1, 2, and 3 are stored in the page frames 1, 3, 2, and 0, respectively. The physical address (in decimal format) corresponding to the virtual address 2500 (in decimal format) is \_\_\_\_\_\_\_\_\_\. (Answer in integer)

\par\vspace{6pt}\noindent\textbf{Your Answer:} \fbox{\rule{0pt}{1.3ex}\hspace{3cm}}

\tcblower

\textbf{Official Key:} \textbf{\large 6596 to 6596} \hfill \textit{\small [Easy]}
\vspace{2pt}\hrule height 0.3pt\vspace{3pt}
\textbf{Topic:} Operating System \\
\textbf{Subtopic:} \textit{Memory management and virtual memory}
\vspace{2pt}\hrule height 0.3pt\vspace{3pt}
\textbf{Guide / Hint:} \\
{\small With page size $2\text{ KB} = 2048\text{ bytes}$, virtual address $2500$ has virtual page number $\lfloor 2500 / 2048 \rfloor = 1$ and offset $2500 \pmod{2048} = 452$. Page 1 maps to frame 3. The physical address is $(\text{frame } 3 \times 2048) + \text{offset } 452 = 6144 + 452 = 6596$.}

\end{tcolorbox}
\vspace{6pt}\hrule height 0.35pt\vspace{6pt}
\noindent
\begin{tcolorbox}[
    blanker,
    sidebyside,
    sidebyside align=top,
    lefthand width=0.55\textwidth,
    top=4pt, bottom=4pt, left=0pt, right=0pt,
    middle=16pt,
    underlay={
        \draw[dashed, black!80, line width=0.8pt, dash pattern=on 4pt off 3pt]
            ([xshift=0.55\textwidth+8pt]frame.north west) --
            ([xshift=0.55\textwidth+8pt]frame.south west);
    }
]
\textbf{Q.179} \quad {\footnotesize\textsc{[MCQ $\cdot$ 2 Mark(s) $\cdot$ Neg: 0.66]}} \hfill {\footnotesize\textbf{GATE 2024 CS2} \quad \tiny\texttt{2024\_CS2\_CS\_Q37}}
\vspace{2pt}\hrule height 0.3pt\vspace{4pt}

Consider a single processor system with four processes A, B, C, and D, represented as given below, where for each process the first value is its arrival time, and the second value is its CPU burst time:
A (0, 10), B (2, 6), C (4, 3), and D (6, 7).
Which one of the following options gives the average waiting time for Shortest Remaining Time First (SRTF) and Non-Preemptive Shortest Job First (NP-SJF) scheduling?

\par\vspace{6pt}\noindent
\begin{tabular}{@{}p{0.48\linewidth}@{\hspace{0.04\linewidth}}p{0.48\linewidth}@{}}
$\square$\ \textbf{(A)} SRTF = 6, NP-SJF = 7 & $\square$\ \textbf{(B)} SRTF = 6, NP-SJF = 7.5 \\
\addlinespace[3pt]
$\square$\ \textbf{(C)} SRTF = 7, NP-SJF = 7.5 & $\square$\ \textbf{(D)} SRTF = 7, NP-SJF = 8.5
\end{tabular}

\tcblower

\textbf{Official Key:} \textbf{\large B} \hfill \textit{\small [Medium]}
\vspace{2pt}\hrule height 0.3pt\vspace{3pt}
\textbf{Topic:} Operating System \\
\textbf{Subtopic:} \textit{CPU and I/O scheduling}
\vspace{2pt}\hrule height 0.3pt\vspace{3pt}
\textbf{Guide / Hint:} \\
{\small Under SRTF: execution order is A(0-2), B(2-4), C(4-7), B(7-11), D(11-18), A(18-26). Turnaround times: A=26, B=9, C=3, D=12; Waiting times: A=16, B=3, C=0, D=5. Average waiting time $= 24/4 = 6.0$. Under NP-SJF: A runs non-preemptively (0-10), then C(10-13), B(13-19), D(19-26). Waiting times: A=0, C=6, B=11, D=13. Average waiting time $= 30/4 = 7.5$.}

\end{tcolorbox}
\vspace{6pt}\hrule height 0.35pt\vspace{6pt}
\noindent
\begin{tcolorbox}[
    blanker,
    sidebyside,
    sidebyside align=top,
    lefthand width=0.55\textwidth,
    top=4pt, bottom=4pt, left=0pt, right=0pt,
    middle=16pt,
    underlay={
        \draw[dashed, black!80, line width=0.8pt, dash pattern=on 4pt off 3pt]
            ([xshift=0.55\textwidth+8pt]frame.north west) --
            ([xshift=0.55\textwidth+8pt]frame.south west);
    }
]
\textbf{Q.180} \quad {\footnotesize\textsc{[MSQ $\cdot$ 2 Mark(s) $\cdot$ No Neg]}} \hfill {\footnotesize\textbf{GATE 2024 CS2} \quad \tiny\texttt{2024\_CS2\_CS\_Q46}}
\vspace{2pt}\hrule height 0.3pt\vspace{4pt}

Consider a multi-threaded program with two threads T1 and T2. The threads share two semaphores: s1 (initialized to 1) and s2 (initialized to 0). The threads also share a global variable x (initialized to 0). The threads execute the code shown below:
\begin{verbatim}
// code of T1        // code of T2
wait(s1);            wait(s1);
x = x + 1;           x = x + 1;
print(x);            print(x);
wait(s2);            signal(s2);
signal(s1);          signal(s1);
\end{verbatim}
Which of the following outcomes is/are possible when threads T1 and T2 execute concurrently?

\par\vspace{6pt}\noindent
\begin{tabular}{@{}p{0.98\linewidth}@{}}
$\square$\ \textbf{(A)} T1 runs first and prints 1, T2 runs next and prints 2 \\
\addlinespace[3pt]
$\square$\ \textbf{(B)} T2 runs first and prints 1, T1 runs next and prints 2 \\
\addlinespace[3pt]
$\square$\ \textbf{(C)} T1 runs first and prints 1, T2 does not print anything (deadlock) \\
\addlinespace[3pt]
$\square$\ \textbf{(D)} T2 runs first and prints 1, T1 does not print anything (deadlock)
\end{tabular}

\tcblower

\textbf{Official Key:} \textbf{\large B;C} \hfill \textit{\small [Medium]}
\vspace{2pt}\hrule height 0.3pt\vspace{3pt}
\textbf{Topic:} Operating System \\
\textbf{Subtopic:} \textit{Concurrency and synchronization}
\vspace{2pt}\hrule height 0.3pt\vspace{3pt}
\textbf{Guide / Hint:} \\
{\small If T2 runs first: T2 takes s1, increments x (x=1), prints 1, signals s2, signals s1. Then T1 executes: takes s1, increments x (x=2), prints 2, completes wait(s2) without blocking (since s2 was signaled), signals s1 (outcome B). If T1 runs first: T1 takes s1, prints 1, and blocks on wait(s2) while holding s1, causing deadlock where T2 blocks on wait(s1) (outcome C).}

\end{tcolorbox}
\vspace{6pt}\hrule height 0.35pt\vspace{6pt}
\noindent
\begin{tcolorbox}[
    blanker,
    sidebyside,
    sidebyside align=top,
    lefthand width=0.55\textwidth,
    top=4pt, bottom=4pt, left=0pt, right=0pt,
    middle=16pt,
    underlay={
        \draw[dashed, black!80, line width=0.8pt, dash pattern=on 4pt off 3pt]
            ([xshift=0.55\textwidth+8pt]frame.north west) --
            ([xshift=0.55\textwidth+8pt]frame.south west);
    }
]
\textbf{Q.181} \quad {\footnotesize\textsc{[NAT $\cdot$ 2 Mark(s) $\cdot$ No Neg]}} \hfill {\footnotesize\textbf{GATE 2024 CS2} \quad \tiny\texttt{2024\_CS2\_CS\_Q53}}
\vspace{2pt}\hrule height 0.3pt\vspace{4pt}

Consider a disk with the following specifications: rotation speed of 6000 RPM, average seek time of 5 milliseconds, 500 sectors/track, 512-byte sectors. A file has content stored in 3000 sectors located randomly on the disk. Assuming average rotational latency, the total time (in seconds, rounded off to 2 decimal places) to read the entire file from the disk is \_\_\_\_\_\_\. (Answer in range 29.50 to 30.50)

\par\vspace{6pt}\noindent\textbf{Your Answer:} \fbox{\rule{0pt}{1.3ex}\hspace{3cm}}

\tcblower

\textbf{Official Key:} \textbf{\large 29.50 to 30.50} \hfill \textit{\small [Medium]}
\vspace{2pt}\hrule height 0.3pt\vspace{3pt}
\textbf{Topic:} Operating System \\
\textbf{Subtopic:} \textit{File systems}
\vspace{2pt}\hrule height 0.3pt\vspace{3pt}
\textbf{Guide / Hint:} \\
{\small At 6000 RPM, one rotation takes $\frac{60}{6000} = 10\text{ ms}$. Average rotational latency $= 10/2 = 5\text{ ms}$. Transfer time per sector $= \frac{10\text{ ms}}{500} = 0.02\text{ ms}$. Time to access each random sector $= \text{seek} + \text{rotational latency} + \text{transfer} = 5 + 5 + 0.02 = 10.02\text{ ms}$. Total time for 3000 random sectors $= 3000 \times 10.02\text{ ms} = 30060\text{ ms} \approx 30.06$ seconds.}

\end{tcolorbox}
\vspace{6pt}\hrule height 0.35pt\vspace{6pt}
\noindent
\begin{tcolorbox}[
    blanker,
    sidebyside,
    sidebyside align=top,
    lefthand width=0.55\textwidth,
    top=4pt, bottom=4pt, left=0pt, right=0pt,
    middle=16pt,
    underlay={
        \draw[dashed, black!80, line width=0.8pt, dash pattern=on 4pt off 3pt]
            ([xshift=0.55\textwidth+8pt]frame.north west) --
            ([xshift=0.55\textwidth+8pt]frame.south west);
    }
]
\textbf{Q.182} \quad {\footnotesize\textsc{[NAT $\cdot$ 2 Mark(s) $\cdot$ No Neg]}} \hfill {\footnotesize\textbf{GATE 2024 CS2} \quad \tiny\texttt{2024\_CS2\_CS\_Q64}}
\vspace{2pt}\hrule height 0.3pt\vspace{4pt}

Consider a 32-bit system with 4 KB page size and page table entries of size 4 bytes each. Assume $1\text{ KB} = 2^{10}$ bytes. The OS uses a 2-level page table for memory management, with an outer page directory and an inner page table. The OS allocates 1 page for the outer page directory upon process creation. The OS uses demand paging when allocating memory for the inner page table. An active process accesses 2000 unique pages during its execution. Let $X$ denote the minimum and $Y$ denote the maximum number of pages across the two levels of the page table of the process. The value of $X + Y$ is \_\_\_\_\_\_\. (Answer in integer)

\par\vspace{6pt}\noindent\textbf{Your Answer:} \fbox{\rule{0pt}{1.3ex}\hspace{3cm}}

\tcblower

\textbf{Official Key:} \textbf{\large 1028 to 1028} \hfill \textit{\small [Hard]}
\vspace{2pt}\hrule height 0.3pt\vspace{3pt}
\textbf{Topic:} Operating System \\
\textbf{Subtopic:} \textit{Memory management and virtual memory}
\vspace{2pt}\hrule height 0.3pt\vspace{3pt}
\textbf{Guide / Hint:} \\
{\small Page size is $4\text{ KB} = 2^{12}$ bytes, PTE size is 4 bytes, so each page table holds $2^{12} / 4 = 1024$ entries. The process accesses 2000 unique pages. Minimum inner page tables $X$: $\lceil 2000 / 1024 \rceil = 2$ inner pages plus 1 outer page directory $= 3$. Maximum inner page tables $Y$: 2000 pages distributed across 2000 different inner tables, bounded by total possible entries in directory $= 1024$ inner pages plus 1 outer page directory $= 1025$. Hence, $X + Y = 3 + 1025 = 1028$.}

\end{tcolorbox}
\vspace{6pt}\hrule height 0.35pt\vspace{6pt}
\noindent
\begin{tcolorbox}[
    blanker,
    sidebyside,
    sidebyside align=top,
    lefthand width=0.55\textwidth,
    top=4pt, bottom=4pt, left=0pt, right=0pt,
    middle=16pt,
    underlay={
        \draw[dashed, black!80, line width=0.8pt, dash pattern=on 4pt off 3pt]
            ([xshift=0.55\textwidth+8pt]frame.north west) --
            ([xshift=0.55\textwidth+8pt]frame.south west);
    }
]
\textbf{Q.183} \quad {\footnotesize\textsc{[MCQ $\cdot$ 2 Mark(s) $\cdot$ Neg: 0.66]}} \hfill {\footnotesize\textbf{GATE 2025 CS1} \quad \tiny\texttt{2025\_CS1\_CS\_Q38}}
\vspace{2pt}\hrule height 0.3pt\vspace{4pt}

A computer has two processors, $M_1$ and $M_2$. Four processes $P_1, P_2, P_3, P_4$ with CPU bursts of 20, 16, 25, and 10 milliseconds, respectively, arrive at the same time and these are the only processes in the system. The scheduler uses non-preemptive priority scheduling, with priorities decided as follows:
\begin{itemize}
\item $M_1$ uses priority of execution for the processes as, $P_1 > P_3 > P_2 > P_4$, i.e., $P_1$ and $P_4$ have highest and lowest priorities, respectively.
\item $M_2$ uses priority of execution for the processes as, $P_2 > P_3 > P_4 > P_1$, i.e., $P_2$ and $P_1$ have highest and lowest priorities, respectively.
\end{itemize}
A process $P_i$ is scheduled to a processor $M_k$, if the processor is free and no other process $P_j$ is waiting with higher priority. At any given point of time, a process can be allocated to any one of the free processors without violating the execution priority rules. Ignore the context switch time. What will be the average waiting time of the processes in milliseconds?

\par\vspace{6pt}\noindent
\begin{tabular}{@{}p{0.48\linewidth}@{\hspace{0.04\linewidth}}p{0.48\linewidth}@{}}
$\square$\ \textbf{(A)} 9.00 & $\square$\ \textbf{(B)} 8.75 \\
\addlinespace[3pt]
$\square$\ \textbf{(C)} 6.50 & $\square$\ \textbf{(D)} 7.50
\end{tabular}

\tcblower

\textbf{Official Key:} \textbf{\large A} \hfill \textit{\small [Hard]}
\vspace{2pt}\hrule height 0.3pt\vspace{3pt}
\textbf{Topic:} Operating System \\
\textbf{Subtopic:} \textit{CPU and I/O scheduling}
\vspace{2pt}\hrule height 0.3pt\vspace{3pt}
\textbf{Guide / Hint:} \\
{\small At $t = 0$, $M_1$ takes highest priority process $P_1$ (burst 20) and $M_2$ takes $P_2$ (burst 16). At $t = 16$, $M_2$ frees up; among remaining processes $\{P_3, P_4\}$, $P_3$ has higher priority on $M_2$, so $P_3$ runs on $M_2$ ($t = 16$ to 41). At $t = 20$, $M_1$ frees up and executes $P_4$ ($t = 20$ to 30). Waiting times are: $P_1 = 0$, $P_2 = 0$, $P_3 = 16$, $P_4 = 20$. Average waiting time is $(0 + 0 + 16 + 20) / 4 = 36 / 4 = 9.00$ ms.}

\end{tcolorbox}
\vspace{6pt}\hrule height 0.35pt\vspace{6pt}
\noindent
\begin{tcolorbox}[
    blanker,
    sidebyside,
    sidebyside align=top,
    lefthand width=0.55\textwidth,
    top=4pt, bottom=4pt, left=0pt, right=0pt,
    middle=16pt,
    underlay={
        \draw[dashed, black!80, line width=0.8pt, dash pattern=on 4pt off 3pt]
            ([xshift=0.55\textwidth+8pt]frame.north west) --
            ([xshift=0.55\textwidth+8pt]frame.south west);
    }
]
\textbf{Q.184} \quad {\footnotesize\textsc{[NAT $\cdot$ 2 Mark(s) $\cdot$ No Neg]}} \hfill {\footnotesize\textbf{GATE 2025 CS1} \quad \tiny\texttt{2025\_CS1\_CS\_Q51}}
\vspace{2pt}\hrule height 0.3pt\vspace{4pt}

A disk of size 512M bytes is divided into blocks of 64K bytes. A file is stored in the disk using linked allocation. In linked allocation, each data block reserves 4 bytes to store the pointer to the next data block. The link part of the last data block contains a NULL pointer (also of 4 bytes). Suppose a file of 1M bytes needs to be stored in the disk. Assume $1\text{K} = 2^{10}$ and $1\text{M} = 2^{20}$. The amount of space in bytes that will be wasted due to internal fragmentation is \_\_\_\_\_\_\. (Answer in integer)

\par\vspace{6pt}\noindent\textbf{Your Answer:} \fbox{\rule{0pt}{1.3ex}\hspace{3cm}}

\tcblower

\textbf{Official Key:} \textbf{\large 65468 to 65468} \hfill \textit{\small [Medium]}
\vspace{2pt}\hrule height 0.3pt\vspace{3pt}
\textbf{Topic:} Operating System \\
\textbf{Subtopic:} \textit{File systems}
\vspace{2pt}\hrule height 0.3pt\vspace{3pt}
\textbf{Guide / Hint:} \\
{\small Block size is $64\text{ KB} = 65536$ bytes, with 4 bytes for pointer, leaving $65536 - 4 = 65532$ bytes of file data per block. File size is $1\text{ MB} = 1048576$ bytes. Number of blocks needed is $\lceil 1048576 / 65532 \rceil = \lceil 16.001037 \rceil = 17$ blocks. The first 16 blocks store $16 \times 65532 = 1048512$ bytes, leaving $1048576 - 1048512 = 64$ bytes stored in the 17th block. Internal fragmentation in the last block is $65532 - 64 = 65468$ bytes.}

\end{tcolorbox}
\vspace{6pt}\hrule height 0.35pt\vspace{6pt}
\noindent
\begin{tcolorbox}[
    blanker,
    sidebyside,
    sidebyside align=top,
    lefthand width=0.55\textwidth,
    top=4pt, bottom=4pt, left=0pt, right=0pt,
    middle=16pt,
    underlay={
        \draw[dashed, black!80, line width=0.8pt, dash pattern=on 4pt off 3pt]
            ([xshift=0.55\textwidth+8pt]frame.north west) --
            ([xshift=0.55\textwidth+8pt]frame.south west);
    }
]
\textbf{Q.185} \quad {\footnotesize\textsc{[NAT $\cdot$ 2 Mark(s) $\cdot$ No Neg]}} \hfill {\footnotesize\textbf{GATE 2025 CS1} \quad \tiny\texttt{2025\_CS1\_CS\_Q54}}
\vspace{2pt}\hrule height 0.3pt\vspace{4pt}

In optimal page replacement algorithm, information about all future page references is available to the operating system (OS). A modification of the optimal page replacement algorithm is as follows:
The OS correctly predicts only up to next 4 page references (including the current page) at the time of allocating a frame to a page.
A process accesses the pages in the following order of page numbers:
1, 3, 2, 4, 2, 3, 1, 2, 4, 3, 1, 4.
If the system has three memory frames that are initially empty, the number of page faults that will occur during execution of the process is \_\_\_\_\_\_\. (Answer in integer)

\par\vspace{6pt}\noindent\textbf{Your Answer:} \fbox{\rule{0pt}{1.3ex}\hspace{3cm}}

\tcblower

\textbf{Official Key:} \textbf{\large 6 to 6} \hfill \textit{\small [Medium]}
\vspace{2pt}\hrule height 0.3pt\vspace{3pt}
\textbf{Topic:} Operating System \\
\textbf{Subtopic:} \textit{Memory management and virtual memory}
\vspace{2pt}\hrule height 0.3pt\vspace{3pt}
\textbf{Guide / Hint:} \\
{\small Simulating the 3 frames with 4-step lookahead: 1 (fault, [1]), 3 (fault, [1, 3]), 2 (fault, [1, 3, 2]). At page 4 (fault): window is (4, 2, 3, 1), farthest is 1, so replace 1 -> [4, 3, 2]. Pages 2, 3 hit. At page 1 (fault): window (1, 2, 4, 3), replace 3 -> [4, 1, 2]. Page 2, 4, 3 (fault): replace 1 -> [4, 3, 2]. Total faults = 3 compulsory + 3 replacement = 6 page faults.}

\end{tcolorbox}
\vspace{6pt}\hrule height 0.35pt\vspace{6pt}
\noindent
\begin{tcolorbox}[
    blanker,
    sidebyside,
    sidebyside align=top,
    lefthand width=0.55\textwidth,
    top=4pt, bottom=4pt, left=0pt, right=0pt,
    middle=16pt,
    underlay={
        \draw[dashed, black!80, line width=0.8pt, dash pattern=on 4pt off 3pt]
            ([xshift=0.55\textwidth+8pt]frame.north west) --
            ([xshift=0.55\textwidth+8pt]frame.south west);
    }
]
\textbf{Q.186} \quad {\footnotesize\textsc{[MSQ $\cdot$ 2 Mark(s) $\cdot$ No Neg]}} \hfill {\footnotesize\textbf{GATE 2025 CS2} \quad \tiny\texttt{2025\_CS2\_CS\_Q47}}
\vspace{2pt}\hrule height 0.3pt\vspace{4pt}

Consider a demand paging system with three frames, and the following page reference string: 1 2 3 4 5 4 1 6 4 2 1. Which page replacement policy produces the minimum number of page faults for this reference string?

\par\vspace{6pt}\noindent
\begin{tabular}{@{}p{0.98\linewidth}@{}}
$\square$\ \textbf{(A)} Least Recently Used page replacement policy \\
\addlinespace[3pt]
$\square$\ \textbf{(B)} Least Frequently Used page replacement policy \\
\addlinespace[3pt]
$\square$\ \textbf{(C)} Most Frequently Used page replacement policy \\
\addlinespace[3pt]
$\square$\ \textbf{(D)} Optimal page replacement policy
\end{tabular}

\tcblower

\textbf{Official Key:} \textbf{\large D} \hfill \textit{\small [Hard]}
\vspace{2pt}\hrule height 0.3pt\vspace{3pt}
\textbf{Topic:} Operating System \\
\textbf{Subtopic:} \textit{Memory management and virtual memory: page replacement algorithms}
\vspace{2pt}\hrule height 0.3pt\vspace{3pt}
\textbf{Guide / Hint:} \\
{\small Trace the page reference string through 3 frames under each algorithm. Optimal page replacement replaces the page that will not be used for the longest period of time in the future, matching the minimum hit/fault sequence.}

\end{tcolorbox}
\vspace{6pt}\hrule height 0.35pt\vspace{6pt}
\noindent
\begin{tcolorbox}[
    blanker,
    sidebyside,
    sidebyside align=top,
    lefthand width=0.55\textwidth,
    top=4pt, bottom=4pt, left=0pt, right=0pt,
    middle=16pt,
    underlay={
        \draw[dashed, black!80, line width=0.8pt, dash pattern=on 4pt off 3pt]
            ([xshift=0.55\textwidth+8pt]frame.north west) --
            ([xshift=0.55\textwidth+8pt]frame.south west);
    }
]
\textbf{Q.187} \quad {\footnotesize\textsc{[MSQ $\cdot$ 2 Mark(s) $\cdot$ No Neg]}} \hfill {\footnotesize\textbf{GATE 2025 CS2} \quad \tiny\texttt{2025\_CS2\_CS\_Q48}}
\vspace{2pt}\hrule height 0.3pt\vspace{4pt}

Let $P = \{P_1, P_2, P_3, P_4\}$ be active processes and $R = \{R_1, R_2, R_3, R_4\}$ be resources in an operating system. Given the allocation and request graph, which of the following actions makes the system deadlock free?

\par\vspace{6pt}\noindent
\begin{tabular}{@{}p{0.98\linewidth}@{}}
$\square$\ \textbf{(A)} Aborting $P_1$ makes the system deadlock free. \\
\addlinespace[3pt]
$\square$\ \textbf{(B)} Aborting $P_3$ makes the system deadlock free. \\
\addlinespace[3pt]
$\square$\ \textbf{(C)} Aborting $P_2$ makes the system deadlock free. \\
\addlinespace[3pt]
$\square$\ \textbf{(D)} Aborting $P_1$ and $P_4$ makes the system deadlock free.
\end{tabular}

\tcblower

\textbf{Official Key:} \textbf{\large C;D} \hfill \textit{\small [Hard]}
\vspace{2pt}\hrule height 0.3pt\vspace{3pt}
\textbf{Topic:} Operating System \\
\textbf{Subtopic:} \textit{Deadlock: Resource Allocation Graph and deadlock recovery}
\vspace{2pt}\hrule height 0.3pt\vspace{3pt}
\textbf{Guide / Hint:} \\
{\small In the Resource Allocation Graph, identify the directed cycles creating deadlock. Aborting process $P_2$ eliminates the dependency cycles directly, as does aborting both $P_1$ and $P_4$.}

\end{tcolorbox}
\vspace{6pt}\hrule height 0.35pt\vspace{6pt}
\noindent
\begin{tcolorbox}[
    blanker,
    sidebyside,
    sidebyside align=top,
    lefthand width=0.55\textwidth,
    top=4pt, bottom=4pt, left=0pt, right=0pt,
    middle=16pt,
    underlay={
        \draw[dashed, black!80, line width=0.8pt, dash pattern=on 4pt off 3pt]
            ([xshift=0.55\textwidth+8pt]frame.north west) --
            ([xshift=0.55\textwidth+8pt]frame.south west);
    }
]
\textbf{Q.188} \quad {\footnotesize\textsc{[NAT $\cdot$ 2 Mark(s) $\cdot$ No Neg]}} \hfill {\footnotesize\textbf{GATE 2025 CS2} \quad \tiny\texttt{2025\_CS2\_CS\_Q58}}
\vspace{2pt}\hrule height 0.3pt\vspace{4pt}

A computer system supports a logical address space of $2^{32}$ bytes. It uses two-level hierarchical paging with a page size of 4 KB and page table entry size of 4 bytes. The number of bits in the outer page table index is \_\_\_\_\_\_\_\_.

\par\vspace{6pt}\noindent\textbf{Your Answer:} \fbox{\rule{0pt}{1.3ex}\hspace{3cm}}

\tcblower

\textbf{Official Key:} \textbf{\large 11 to 11} \hfill \textit{\small [Hard]}
\vspace{2pt}\hrule height 0.3pt\vspace{3pt}
\textbf{Topic:} Operating System \\
\textbf{Subtopic:} \textit{Memory management and virtual memory: hierarchical paging}
\vspace{2pt}\hrule height 0.3pt\vspace{3pt}
\textbf{Guide / Hint:} \\
{\small For a $2^{32}$-byte logical address space with two-level hierarchical paging, page offset and page table index bits are derived from page size and entry size, yielding 11 bits for the first-level directory index.}

\end{tcolorbox}
\vspace{6pt}\hrule height 0.35pt\vspace{6pt}
\vspace{8pt}
\subsection*{Databases}
\vspace{-4pt}
\noindent
\begin{tcolorbox}[
    blanker,
    sidebyside,
    sidebyside align=top,
    lefthand width=0.55\textwidth,
    top=4pt, bottom=4pt, left=0pt, right=0pt,
    middle=16pt,
    underlay={
        \draw[dashed, black!80, line width=0.8pt, dash pattern=on 4pt off 3pt]
            ([xshift=0.55\textwidth+8pt]frame.north west) --
            ([xshift=0.55\textwidth+8pt]frame.south west);
    }
]
\textbf{Q.189} \quad {\footnotesize\textsc{[MCQ $\cdot$ 2 Mark(s) $\cdot$ Neg: 0.66]}} \hfill {\footnotesize\textbf{GATE 2018 CS} \quad \tiny\texttt{2018\_CS\_GA\_Q41}}
\vspace{2pt}\hrule height 0.3pt\vspace{4pt}

Consider the relations $r(A, B)$ and $s(B, C)$, where $s.B$ is a primary key and $r.B$ is a foreign key referencing $s.B$. Consider the query:
\[ Q: \Pi\_{A, B, C}(r \bowtie \sigma\_{B < 5}(s)) \]
Which one of the following relational algebra queries is NOT equivalent to $Q$?

\par\vspace{6pt}\noindent
\begin{tabular}{@{}p{0.98\linewidth}@{}}
$\square$\ \textbf{(A)} $\Pi_{A, B, C}(\sigma_{B < 5}(r) \bowtie s)$ \\
\addlinespace[3pt]
$\square$\ \textbf{(B)} $\sigma_{B < 5}(\Pi_{A, B, C}(r \bowtie s))$ \\
\addlinespace[3pt]
$\square$\ \textbf{(C)} $\Pi_{A, B, C}(r \,\text{⟕}\, \sigma_{B < 5}(s))$ \\
\addlinespace[3pt]
$\square$\ \textbf{(D)} $\Pi_{A, B, C}(\sigma_{B < 5}(r \bowtie s))$
\end{tabular}

\tcblower

\textbf{Official Key:} \textbf{\large C} \hfill \textit{\small [Medium]}
\vspace{2pt}\hrule height 0.3pt\vspace{3pt}
\textbf{Topic:} Databases \\
\textbf{Subtopic:} \textit{Relational model: relational algebra, tuple calculus, SQL}
\vspace{2pt}\hrule height 0.3pt\vspace{3pt}
\textbf{Guide / Hint:} \\
{\small Query $Q$ inner joins $r$ with $\sigma_{B < 5}(s)$, which discards all tuples from $r$ where $B \ge 5$. Option (C) performs a left outer join between $r$ and $\sigma_{B < 5}(s)$, which retains all tuples from $r$ (filling missing $C$ values with NULL for $B \ge 5$), so its result is NOT equivalent to $Q$.}

\end{tcolorbox}
\vspace{6pt}\hrule height 0.35pt\vspace{6pt}
\noindent
\begin{tcolorbox}[
    blanker,
    sidebyside,
    sidebyside align=top,
    lefthand width=0.55\textwidth,
    top=4pt, bottom=4pt, left=0pt, right=0pt,
    middle=16pt,
    underlay={
        \draw[dashed, black!80, line width=0.8pt, dash pattern=on 4pt off 3pt]
            ([xshift=0.55\textwidth+8pt]frame.north west) --
            ([xshift=0.55\textwidth+8pt]frame.south west);
    }
]
\textbf{Q.190} \quad {\footnotesize\textsc{[MCQ $\cdot$ 2 Mark(s) $\cdot$ Neg: 0.66]}} \hfill {\footnotesize\textbf{GATE 2018 CS} \quad \tiny\texttt{2018\_CS\_GA\_Q42}}
\vspace{2pt}\hrule height 0.3pt\vspace{4pt}

Consider the following four relational schemas. For each schema, all non-trivial functional dependencies are listed. The underlined attributes denote candidate keys.
\begin{itemize}[leftmargin=*, nosep]
    \item[Schema I:] $R_1(A, B, C, D)$ with FDs: $A \to B, B \to C, C \to D$
    \item[Schema II:] $R_2(\underline{rollno}, \underline{courseid}, email)$ with FDs: $rollno \to email, email \to rollno$
    \item[Schema III:] $R_3(A, B, C)$ with FDs: $A \to B, B \to C$
    \item[Schema IV:] $R_4(A, B, C)$ with FDs: $AB \to C, C \to A$
\end{itemize}
Which one of the relational schemas is in 3NF but not in BCNF?

\par\vspace{6pt}\noindent
\begin{tabular}{@{}p{0.48\linewidth}@{\hspace{0.04\linewidth}}p{0.48\linewidth}@{}}
$\square$\ \textbf{(A)} Schema I & $\square$\ \textbf{(B)} Schema II \\
\addlinespace[3pt]
$\square$\ \textbf{(C)} Schema III & $\square$\ \textbf{(D)} Schema IV
\end{tabular}

\tcblower

\textbf{Official Key:} \textbf{\large B} \hfill \textit{\small [Medium]}
\vspace{2pt}\hrule height 0.3pt\vspace{3pt}
\textbf{Topic:} Databases \\
\textbf{Subtopic:} \textit{Integrity constraints, normal forms}
\vspace{2pt}\hrule height 0.3pt\vspace{3pt}
\textbf{Guide / Hint:} \\
{\small In Schema II, the functional dependencies are $rollno \to email$ and $email \to rollno$. The candidate keys are $\{rollno, courseid\}$ and $\{email, courseid\}$. In $email \to rollno$, the left side is not a superkey (violating BCNF), but the right side $rollno$ is a prime attribute (satisfying 3NF). Hence Schema II is in 3NF but not in BCNF.}

\end{tcolorbox}
\vspace{6pt}\hrule height 0.35pt\vspace{6pt}
\noindent
\begin{tcolorbox}[
    blanker,
    sidebyside,
    sidebyside align=top,
    lefthand width=0.55\textwidth,
    top=4pt, bottom=4pt, left=0pt, right=0pt,
    middle=16pt,
    underlay={
        \draw[dashed, black!80, line width=0.8pt, dash pattern=on 4pt off 3pt]
            ([xshift=0.55\textwidth+8pt]frame.north west) --
            ([xshift=0.55\textwidth+8pt]frame.south west);
    }
]
\textbf{Q.191} \quad {\footnotesize\textsc{[MCQ $\cdot$ 2 Mark(s) $\cdot$ Neg: 0.66]}} \hfill {\footnotesize\textbf{GATE 2022 CS} \quad \tiny\texttt{2022\_CS\_CS\_Q39}}
\vspace{2pt}\hrule height 0.3pt\vspace{4pt}

Let $R_i(z)$ and $W_i(z)$ denote read and write operations on a data element $z$ by a transaction $T_i$, respectively. Consider the schedule $S$ with four transactions:
\[ S: R\_4(x) R\_2(x) R\_3(x) R\_1(y) W\_1(y) W\_2(x) W\_3(y) R\_4(y) \]
Which one of the following serial schedules is conflict equivalent to $S$?

\par\vspace{6pt}\noindent
\begin{tabular}{@{}p{0.48\linewidth}@{\hspace{0.04\linewidth}}p{0.48\linewidth}@{}}
$\square$\ \textbf{(A)} $T_1 \to T_3 \to T_4 \to T_2$ & $\square$\ \textbf{(B)} $T_1 \to T_4 \to T_3 \to T_2$ \\
\addlinespace[3pt]
$\square$\ \textbf{(C)} $T_4 \to T_1 \to T_3 \to T_2$ & $\square$\ \textbf{(D)} $T_3 \to T_1 \to T_4 \to T_2$
\end{tabular}

\tcblower

\textbf{Official Key:} \textbf{\large A} \hfill \textit{\small [Medium]}
\vspace{2pt}\hrule height 0.3pt\vspace{3pt}
\textbf{Topic:} Databases \\
\textbf{Subtopic:} \textit{Transactions and concurrency control}
\vspace{2pt}\hrule height 0.3pt\vspace{3pt}
\textbf{Guide / Hint:} \\
{\small Constructing the serialization graph from conflict operations in $S$: on $y$, $R_1(y), W_1(y) \to W_3(y) \to R_4(y)$ gives edges $T_1 \to T_3 \to T_4$. On $x$, $R_4(x), R_3(x) \to W_2(x)$ gives edges $T_4 \to T_2$ and $T_3 \to T_2$. The unique topological order is $T_1 \to T_3 \to T_4 \to T_2$.}

\end{tcolorbox}
\vspace{6pt}\hrule height 0.35pt\vspace{6pt}
\noindent
\begin{tcolorbox}[
    blanker,
    sidebyside,
    sidebyside align=top,
    lefthand width=0.55\textwidth,
    top=4pt, bottom=4pt, left=0pt, right=0pt,
    middle=16pt,
    underlay={
        \draw[dashed, black!80, line width=0.8pt, dash pattern=on 4pt off 3pt]
            ([xshift=0.55\textwidth+8pt]frame.north west) --
            ([xshift=0.55\textwidth+8pt]frame.south west);
    }
]
\textbf{Q.192} \quad {\footnotesize\textsc{[NAT $\cdot$ 2 Mark(s) $\cdot$ No Neg]}} \hfill {\footnotesize\textbf{GATE 2022 CS} \quad \tiny\texttt{2022\_CS\_CS\_Q56}}
\vspace{2pt}\hrule height 0.3pt\vspace{4pt}

Consider the relational database with the following four schemas and their respective instances:
\[ \text{Student}(sNo, sName, dNo), \quad \text{Dept}(dNo, dName), \quad \text{Course}(cNo, cName, dNo), \quad \text{Register}(sNo, cNo) \]
\begin{lstlisting}[language=SQL]
SELECT * FROM Student AS S WHERE NOT EXISTS
    (SELECT cNo FROM Course WHERE dNo = 'D01'
     EXCEPT
     SELECT cNo FROM Register WHERE sNo = S.sNo)
\end{lstlisting}
The number of rows returned by the above SQL query is \_\_\_\_\_\_\_\_\_\_\_\. (Answer in integer)

\par\vspace{6pt}\noindent\textbf{Your Answer:} \fbox{\rule{0pt}{1.3ex}\hspace{3cm}}

\tcblower

\textbf{Official Key:} \textbf{\large 2 to 2} \hfill \textit{\small [Medium]}
\vspace{2pt}\hrule height 0.3pt\vspace{3pt}
\textbf{Topic:} Databases \\
\textbf{Subtopic:} \textit{Relational model: relational algebra, tuple calculus, SQL}
\vspace{2pt}\hrule height 0.3pt\vspace{3pt}
\textbf{Guide / Hint:} \\
{\small The nested subquery with `EXCEPT` and `NOT EXISTS` expresses relational division: it finds students who have registered for every course offered by department 'D01'. In the database instance, exactly 2 students satisfy this condition.}

\end{tcolorbox}
\vspace{6pt}\hrule height 0.35pt\vspace{6pt}
\noindent
\begin{tcolorbox}[
    blanker,
    sidebyside,
    sidebyside align=top,
    lefthand width=0.55\textwidth,
    top=4pt, bottom=4pt, left=0pt, right=0pt,
    middle=16pt,
    underlay={
        \draw[dashed, black!80, line width=0.8pt, dash pattern=on 4pt off 3pt]
            ([xshift=0.55\textwidth+8pt]frame.north west) --
            ([xshift=0.55\textwidth+8pt]frame.south west);
    }
]
\textbf{Q.193} \quad {\footnotesize\textsc{[NAT $\cdot$ 2 Mark(s) $\cdot$ No Neg]}} \hfill {\footnotesize\textbf{GATE 2023 CS} \quad \tiny\texttt{2023\_CS\_CS\_Q61}}
\vspace{2pt}\hrule height 0.3pt\vspace{4pt}

Consider the following table named `Student` in a relational database. The primary key of this table is `rollNum`.
\begin{center}
\begin{tabular}{|c|c|c|c|}
\hline
\textbf{rollNum} & \textbf{name} & \textbf{gender} & \textbf{marks} \\
\hline
1 & Naman & M & 62 \\
2 & Aliya & F & 70 \\
3 & Aliya & F & 80 \\
4 & James & M & 82 \\
5 & Swati & F & 65 \\
\hline
\end{tabular}
\end{center}
The SQL query below is executed on this database:
\begin{lstlisting}[language=SQL]
SELECT *
FROM Student
WHERE gender = 'F' AND marks > 65;
\end{lstlisting}
The number of rows returned by the query is \_\_\_\_\_\_\_\_\. (Answer in integer)

\par\vspace{6pt}\noindent\textbf{Your Answer:} \fbox{\rule{0pt}{1.3ex}\hspace{3cm}}

\tcblower

\textbf{Official Key:} \textbf{\large 2 to 2} \hfill \textit{\small [Easy]}
\vspace{2pt}\hrule height 0.3pt\vspace{3pt}
\textbf{Topic:} Databases \\
\textbf{Subtopic:} \textit{Relational model: relational algebra, tuple calculus, SQL}
\vspace{2pt}\hrule height 0.3pt\vspace{3pt}
\textbf{Guide / Hint:} \\
{\small The query evaluates the WHERE clause condition `gender = 'F' AND marks > 65`. Among female students: Aliya (rollNum 2, marks 70) and Aliya (rollNum 3, marks 80) satisfy both conditions, whereas Swati (rollNum 5, marks 65) fails the strict inequality. Thus, exactly 2 rows are returned.}

\end{tcolorbox}
\vspace{6pt}\hrule height 0.35pt\vspace{6pt}
\noindent
\begin{tcolorbox}[
    blanker,
    sidebyside,
    sidebyside align=top,
    lefthand width=0.55\textwidth,
    top=4pt, bottom=4pt, left=0pt, right=0pt,
    middle=16pt,
    underlay={
        \draw[dashed, black!80, line width=0.8pt, dash pattern=on 4pt off 3pt]
            ([xshift=0.55\textwidth+8pt]frame.north west) --
            ([xshift=0.55\textwidth+8pt]frame.south west);
    }
]
\textbf{Q.194} \quad {\footnotesize\textsc{[MSQ $\cdot$ 2 Mark(s) $\cdot$ No Neg]}} \hfill {\footnotesize\textbf{GATE 2024 CS1} \quad \tiny\texttt{2024\_CS1\_CS\_Q44}}
\vspace{2pt}\hrule height 0.3pt\vspace{4pt}

The symbol $\to$ indicates functional dependency in the context of a relational database. Which of the following options is/are TRUE?

\par\vspace{6pt}\noindent
\begin{tabular}{@{}p{0.98\linewidth}@{}}
$\square$\ \textbf{(A)} $(X, Y) \to (Z, W)$ implies $X \to (Z, W)$ \\
\addlinespace[3pt]
$\square$\ \textbf{(B)} $(X, Y) \to (Z, W)$ implies $(X, Y) \to Z$ \\
\addlinespace[3pt]
$\square$\ \textbf{(C)} $((X, Y) \to Z \text{ and } W \to Y)$ implies $(X, W) \to Z$ \\
\addlinespace[3pt]
$\square$\ \textbf{(D)} $(X \to Y \text{ and } Y \to Z)$ implies $X \to Z$
\end{tabular}

\tcblower

\textbf{Official Key:} \textbf{\large B;C;D} \hfill \textit{\small [Medium]}
\vspace{2pt}\hrule height 0.3pt\vspace{3pt}
\textbf{Topic:} Databases \\
\textbf{Subtopic:} \textit{Integrity constraints, normal forms}
\vspace{2pt}\hrule height 0.3pt\vspace{3pt}
\textbf{Guide / Hint:} \\
{\small By Armstrong's decomposition axiom, $(X, Y) \to (Z, W)$ implies $(X, Y) \to Z$ (B is TRUE). By Armstrong's transitivity axiom, $X \to Y$ and $Y \to Z$ implies $X \to Z$ (D is TRUE). By pseudotransitivity, $(X, Y) \to Z$ and $W \to Y$ implies $(X, W) \to Z$ (C is TRUE). Decomposing the left-hand side determinant is invalid in general (A is FALSE).}

\end{tcolorbox}
\vspace{6pt}\hrule height 0.35pt\vspace{6pt}
\noindent
\begin{tcolorbox}[
    blanker,
    sidebyside,
    sidebyside align=top,
    lefthand width=0.55\textwidth,
    top=4pt, bottom=4pt, left=0pt, right=0pt,
    middle=16pt,
    underlay={
        \draw[dashed, black!80, line width=0.8pt, dash pattern=on 4pt off 3pt]
            ([xshift=0.55\textwidth+8pt]frame.north west) --
            ([xshift=0.55\textwidth+8pt]frame.south west);
    }
]
\textbf{Q.195} \quad {\footnotesize\textsc{[MSQ $\cdot$ 2 Mark(s) $\cdot$ No Neg]}} \hfill {\footnotesize\textbf{GATE 2024 CS1} \quad \tiny\texttt{2024\_CS1\_CS\_Q46}}
\vspace{2pt}\hrule height 0.3pt\vspace{4pt}

Consider the following read-write schedule $S$ over three transactions $T_1, T_2$, and $T_3$, where the subscripts in the schedule indicate transaction IDs:
\[ S: r\_1(z); w\_1(z); r\_2(x); r\_3(y); w\_3(y); r\_2(y); w\_2(x); w\_2(y); \]
Which of the following transaction schedules is/are conflict equivalent to $S$?

\par\vspace{6pt}\noindent
\begin{tabular}{@{}p{0.48\linewidth}@{\hspace{0.04\linewidth}}p{0.48\linewidth}@{}}
$\square$\ \textbf{(A)} $T_1 T_2 T_3$ & $\square$\ \textbf{(B)} $T_1 T_3 T_2$ \\
\addlinespace[3pt]
$\square$\ \textbf{(C)} $T_3 T_2 T_1$ & $\square$\ \textbf{(D)} $T_3 T_1 T_2$
\end{tabular}

\tcblower

\textbf{Official Key:} \textbf{\large B;C;D} \hfill \textit{\small [Medium]}
\vspace{2pt}\hrule height 0.3pt\vspace{3pt}
\textbf{Topic:} Databases \\
\textbf{Subtopic:} \textit{Transactions and concurrency control}
\vspace{2pt}\hrule height 0.3pt\vspace{3pt}
\textbf{Guide / Hint:} \\
{\small Examining data item conflicts: on item $y$, $w_3(y)$ in $T_3$ precedes both $r_2(y)$ and $w_2(y)$ in $T_2$, creating a precedence conflict edge $T_3 \to T_2$. Transaction $T_1$ only accesses $z$ and has no conflicting operations with $T_2$ or $T_3$. Any serial schedule preserving $T_3$ before $T_2$ is conflict equivalent to $S$, namely $T_1 T_3 T_2$, $T_3 T_1 T_2$, and $T_3 T_2 T_1$.}

\end{tcolorbox}
\vspace{6pt}\hrule height 0.35pt\vspace{6pt}
\noindent
\begin{tcolorbox}[
    blanker,
    sidebyside,
    sidebyside align=top,
    lefthand width=0.55\textwidth,
    top=4pt, bottom=4pt, left=0pt, right=0pt,
    middle=16pt,
    underlay={
        \draw[dashed, black!80, line width=0.8pt, dash pattern=on 4pt off 3pt]
            ([xshift=0.55\textwidth+8pt]frame.north west) --
            ([xshift=0.55\textwidth+8pt]frame.south west);
    }
]
\textbf{Q.196} \quad {\footnotesize\textsc{[MCQ $\cdot$ 2 Mark(s) $\cdot$ Neg: 0.66]}} \hfill {\footnotesize\textbf{GATE 2024 CS2} \quad \tiny\texttt{2024\_CS2\_CS\_Q45}}
\vspace{2pt}\hrule height 0.3pt\vspace{4pt}

The relation schema, Person(pid, city), describes the city of residence for every person uniquely identified by pid. The following relational algebra operators are available: selection, projection, cross product, and rename.
To find the list of cities where at least 3 persons reside, what is the minimum number of cross product operations required?

\par\vspace{6pt}\noindent
\begin{tabular}{@{}p{0.48\linewidth}@{\hspace{0.04\linewidth}}p{0.48\linewidth}@{}}
$\square$\ \textbf{(A)} 1 & $\square$\ \textbf{(B)} 2 \\
\addlinespace[3pt]
$\square$\ \textbf{(C)} 3 & $\square$\ \textbf{(D)} 4
\end{tabular}

\tcblower

\textbf{Official Key:} \textbf{\large B} \hfill \textit{\small [Medium]}
\vspace{2pt}\hrule height 0.3pt\vspace{3pt}
\textbf{Topic:} Databases \\
\textbf{Subtopic:} \textit{Relational model: relational algebra, tuple calculus, SQL}
\vspace{2pt}\hrule height 0.3pt\vspace{3pt}
\textbf{Guide / Hint:} \\
{\small To find cities with at least 3 distinct persons residing in them using selection, projection, cross-product, and rename, we need 3 copies of the Person table ($P_1, P_2, P_3$). Combining 3 relations requires 2 cross-product operations: $(P_1 \times P_2) \times P_3$.}

\end{tcolorbox}
\vspace{6pt}\hrule height 0.35pt\vspace{6pt}
\noindent
\begin{tcolorbox}[
    blanker,
    sidebyside,
    sidebyside align=top,
    lefthand width=0.55\textwidth,
    top=4pt, bottom=4pt, left=0pt, right=0pt,
    middle=16pt,
    underlay={
        \draw[dashed, black!80, line width=0.8pt, dash pattern=on 4pt off 3pt]
            ([xshift=0.55\textwidth+8pt]frame.north west) --
            ([xshift=0.55\textwidth+8pt]frame.south west);
    }
]
\textbf{Q.197} \quad {\footnotesize\textsc{[NAT $\cdot$ 2 Mark(s) $\cdot$ No Neg]}} \hfill {\footnotesize\textbf{GATE 2024 CS2} \quad \tiny\texttt{2024\_CS2\_CS\_Q56}}
\vspace{2pt}\hrule height 0.3pt\vspace{4pt}

A functional dependency $F: X \to Y$ is termed as a useful functional dependency if and only if it satisfies all the following three conditions:
\begin{itemize}
\item $X$ is not the empty set.
\item $Y$ is not the empty set.
\item Intersection of $X$ and $Y$ is the empty set.
\end{itemize}
For a relation $R$ with 4 attributes, the total number of possible useful functional dependencies is \_\_\_\_\_\_\. (Answer in integer)

\par\vspace{6pt}\noindent\textbf{Your Answer:} \fbox{\rule{0pt}{1.3ex}\hspace{3cm}}

\tcblower

\textbf{Official Key:} \textbf{\large 50 to 50} \hfill \textit{\small [Medium]}
\vspace{2pt}\hrule height 0.3pt\vspace{3pt}
\textbf{Topic:} Databases \\
\textbf{Subtopic:} \textit{Integrity constraints, normal forms}
\vspace{2pt}\hrule height 0.3pt\vspace{3pt}
\textbf{Guide / Hint:} \\
{\small For each attribute $A$ among the 4 attributes, there are 3 mutually exclusive choices: $A \in X$, $A \in Y$, or $A \notin (X \cup Y)$, giving $3^4 = 81$ pairs of disjoint sets $(X, Y)$. We subtract cases where $X = \emptyset$ ($2^4 = 16$), where $Y = \emptyset$ ($2^4 = 16$), and add back the double-subtracted case where $X = Y = \emptyset$ ($1$). Total useful FDs $= 81 - 16 - 16 + 1 = 50$.}

\end{tcolorbox}
\vspace{6pt}\hrule height 0.35pt\vspace{6pt}
\noindent
\begin{tcolorbox}[
    blanker,
    sidebyside,
    sidebyside align=top,
    lefthand width=0.55\textwidth,
    top=4pt, bottom=4pt, left=0pt, right=0pt,
    middle=16pt,
    underlay={
        \draw[dashed, black!80, line width=0.8pt, dash pattern=on 4pt off 3pt]
            ([xshift=0.55\textwidth+8pt]frame.north west) --
            ([xshift=0.55\textwidth+8pt]frame.south west);
    }
]
\textbf{Q.198} \quad {\footnotesize\textsc{[MCQ $\cdot$ 2 Mark(s) $\cdot$ Neg: 0.66]}} \hfill {\footnotesize\textbf{GATE 2025 CS1} \quad \tiny\texttt{2025\_CS1\_CS\_Q39}}
\vspace{2pt}\hrule height 0.3pt\vspace{4pt}

Consider two relations describing $teams$ and $players$ in a sports league:
\begin{itemize}
\item $teams(tid, tname)$: $tid, tname$ are team-id and team-name, respectively
\item $players(pid, pname, tid)$: $pid, pname$, and $tid$ denote player-id, player-name and the team-id of the player, respectively
\end{itemize}
Which ONE of the following tuple relational calculus queries returns the name of the players who play for the team having $tname$ as 'MI'?

\par\vspace{6pt}\noindent
\begin{tabular}{@{}p{0.98\linewidth}@{}}
$\square$\ \textbf{(A)} $\{ p.pname \mid p \in players \land \exists t (t \in teams \land p.tid = t.tid \land t.tname = \text{'MI'}) \}$ \\
\addlinespace[3pt]
$\square$\ \textbf{(B)} $\{ p.pname \mid p \in teams \land \exists t (t \in players \land p.tid = t.tid \land t.tname = \text{'MI'}) \}$ \\
\addlinespace[3pt]
$\square$\ \textbf{(C)} $\{ p.pname \mid p \in players \land \exists t (t \in teams \land t.tname = \text{'MI'}) \}$ \\
\addlinespace[3pt]
$\square$\ \textbf{(D)} $\{ p.pname \mid p \in teams \land \exists t (t \in players \land t.tname = \text{'MI'}) \}$
\end{tabular}

\tcblower

\textbf{Official Key:} \textbf{\large A} \hfill \textit{\small [Medium]}
\vspace{2pt}\hrule height 0.3pt\vspace{3pt}
\textbf{Topic:} Databases \\
\textbf{Subtopic:} \textit{Relational model: relational algebra, tuple calculus, SQL}
\vspace{2pt}\hrule height 0.3pt\vspace{3pt}
\textbf{Guide / Hint:} \\
{\small The query must select `p.pname` from the `players` relation ($p \in players$) and join with a tuple $t \in teams$ on the foreign key condition $p.tid = t.tid$, while filtering for $t.tname = \text{'MI'}$. Option A accurately expresses this quantified condition in Tuple Relational Calculus.}

\end{tcolorbox}
\vspace{6pt}\hrule height 0.35pt\vspace{6pt}
\noindent
\begin{tcolorbox}[
    blanker,
    sidebyside,
    sidebyside align=top,
    lefthand width=0.55\textwidth,
    top=4pt, bottom=4pt, left=0pt, right=0pt,
    middle=16pt,
    underlay={
        \draw[dashed, black!80, line width=0.8pt, dash pattern=on 4pt off 3pt]
            ([xshift=0.55\textwidth+8pt]frame.north west) --
            ([xshift=0.55\textwidth+8pt]frame.south west);
    }
]
\textbf{Q.199} \quad {\footnotesize\textsc{[MSQ $\cdot$ 2 Mark(s) $\cdot$ No Neg]}} \hfill {\footnotesize\textbf{GATE 2025 CS1} \quad \tiny\texttt{2025\_CS1\_CS\_Q47}}
\vspace{2pt}\hrule height 0.3pt\vspace{4pt}

Consider a relational schema $team(name, city, owner)$, with functional dependencies $\{name \to city, name \to owner\}$. The relation $team$ is decomposed into two relations, $t_1(name, city)$ and $t_2(name, owner)$. Which of the following statement(s) is/are TRUE?

\par\vspace{6pt}\noindent
\begin{tabular}{@{}p{0.98\linewidth}@{}}
$\square$\ \textbf{(A)} The relation $team$ is NOT in BCNF. \\
\addlinespace[3pt]
$\square$\ \textbf{(B)} The relations $t_1$ and $t_2$ are in BCNF. \\
\addlinespace[3pt]
$\square$\ \textbf{(C)} The decomposition constitutes a lossless join. \\
\addlinespace[3pt]
$\square$\ \textbf{(D)} The relation $team$ is NOT in 3NF.
\end{tabular}

\tcblower

\textbf{Official Key:} \textbf{\large B;C} \hfill \textit{\small [Medium]}
\vspace{2pt}\hrule height 0.3pt\vspace{3pt}
\textbf{Topic:} Databases \\
\textbf{Subtopic:} \textit{Integrity constraints, normal forms}
\vspace{2pt}\hrule height 0.3pt\vspace{3pt}
\textbf{Guide / Hint:} \\
{\small In $team$, $name$ is a candidate key since $name^+ = \{name, city, owner\}$. Since every FD has a superkey on the LHS, $team$ was already in BCNF and 3NF (A and D are false). In $t_1$ and $t_2$, $name$ is the primary key, so both are in BCNF (B is true). The common attribute $R_1 \cap R_2 = \{name\}$ is a key of both $t_1$ and $t_2$, ensuring a lossless join (C is true).}

\end{tcolorbox}
\vspace{6pt}\hrule height 0.35pt\vspace{6pt}
\noindent
\begin{tcolorbox}[
    blanker,
    sidebyside,
    sidebyside align=top,
    lefthand width=0.55\textwidth,
    top=4pt, bottom=4pt, left=0pt, right=0pt,
    middle=16pt,
    underlay={
        \draw[dashed, black!80, line width=0.8pt, dash pattern=on 4pt off 3pt]
            ([xshift=0.55\textwidth+8pt]frame.north west) --
            ([xshift=0.55\textwidth+8pt]frame.south west);
    }
]
\textbf{Q.200} \quad {\footnotesize\textsc{[MSQ $\cdot$ 2 Mark(s) $\cdot$ No Neg]}} \hfill {\footnotesize\textbf{GATE 2025 CS2} \quad \tiny\texttt{2025\_CS2\_CS\_Q46}}
\vspace{2pt}\hrule height 0.3pt\vspace{4pt}

Consider the following relational schema along with all the functional dependencies that hold on them:
$R_1(A, B, C, D)$ with FDs, and $R_2(E, F, G, H)$ with FDs.
Which of the following statements is/are CORRECT?

\par\vspace{6pt}\noindent
\begin{tabular}{@{}p{0.48\linewidth}@{\hspace{0.04\linewidth}}p{0.48\linewidth}@{}}
$\square$\ \textbf{(A)} $R_1$ is in 3NF & $\square$\ \textbf{(B)} $R_2$ is in 3NF \\
\addlinespace[3pt]
$\square$\ \textbf{(C)} $R_1$ is NOT in 3NF & $\square$\ \textbf{(D)} $R_2$ is NOT in 3NF
\end{tabular}

\tcblower

\textbf{Official Key:} \textbf{\large C;D} \hfill \textit{\small [Medium]}
\vspace{2pt}\hrule height 0.3pt\vspace{3pt}
\textbf{Topic:} Databases \\
\textbf{Subtopic:} \textit{Integrity constraints, normal forms: 3NF and BCNF}
\vspace{2pt}\hrule height 0.3pt\vspace{3pt}
\textbf{Guide / Hint:} \\
{\small A schema is in 3NF if for every non-trivial FD $X \to Y$, either $X$ is a superkey or $Y$ consists of prime attributes. Here, both schemas contain dependencies where the LHS is not a superkey and the RHS has non-prime attributes, so neither is in 3NF.}

\end{tcolorbox}
\vspace{6pt}\hrule height 0.35pt\vspace{6pt}
\noindent
\begin{tcolorbox}[
    blanker,
    sidebyside,
    sidebyside align=top,
    lefthand width=0.55\textwidth,
    top=4pt, bottom=4pt, left=0pt, right=0pt,
    middle=16pt,
    underlay={
        \draw[dashed, black!80, line width=0.8pt, dash pattern=on 4pt off 3pt]
            ([xshift=0.55\textwidth+8pt]frame.north west) --
            ([xshift=0.55\textwidth+8pt]frame.south west);
    }
]
\textbf{Q.201} \quad {\footnotesize\textsc{[MSQ $\cdot$ 2 Mark(s) $\cdot$ No Neg]}} \hfill {\footnotesize\textbf{GATE 2025 CS2} \quad \tiny\texttt{2025\_CS2\_CS\_Q54}}
\vspace{2pt}\hrule height 0.3pt\vspace{4pt}

Consider the relational schema:
- $\text{Students}(\underline{\text{rollno}}, \text{name}, \text{age}, \text{cgpa})$
- $\text{Enrolled}(\underline{\text{rollno}, \text{courseno}})$
\par\medskip
Which of the following SQL query/queries correctly find(s) the name of all students who have enrolled in course 1470?

\par\vspace{6pt}\noindent
\begin{tabular}{@{}p{0.98\linewidth}@{}}
$\square$\ \textbf{(A)} SELECT S.name \newline FROM Students S \newline WHERE EXISTS (SELECT * FROM Enrolled E WHERE E.courseno = 1470 AND E.rollno = S.rollno); \\
\addlinespace[3pt]
$\square$\ \textbf{(B)} SELECT S.name \newline FROM Students S \newline WHERE SIZEOF (SELECT * FROM Enrolled E WHERE E.courseno = 1470 AND E.rollno = S.rollno) > 0; \\
\addlinespace[3pt]
$\square$\ \textbf{(C)} SELECT S.name \newline FROM Students S \newline WHERE 0 < (SELECT COUNT(*) FROM Enrolled E WHERE E.courseno = 1470 AND E.rollno = S.rollno); \\
\addlinespace[3pt]
$\square$\ \textbf{(D)} SELECT S.name \newline FROM Students S NATURAL JOIN Enrolled E \newline WHERE E.courseno = 1470;
\end{tabular}

\tcblower

\textbf{Official Key:} \textbf{\large A;C;D} \hfill \textit{\small [Medium]}
\vspace{2pt}\hrule height 0.3pt\vspace{3pt}
\textbf{Topic:} Databases \\
\textbf{Subtopic:} \textit{Relational model: SQL queries and joins}
\vspace{2pt}\hrule height 0.3pt\vspace{3pt}
\textbf{Guide / Hint:} \\
{\small Queries A, C, and D all correctly return the names of students enrolled in course 1470 using EXISTS, COUNT, and natural join respectively. Query B is invalid because SQL does not have a `SIZEOF` operator.}

\end{tcolorbox}
\vspace{6pt}\hrule height 0.35pt\vspace{6pt}
\noindent
\begin{tcolorbox}[
    blanker,
    sidebyside,
    sidebyside align=top,
    lefthand width=0.55\textwidth,
    top=4pt, bottom=4pt, left=0pt, right=0pt,
    middle=16pt,
    underlay={
        \draw[dashed, black!80, line width=0.8pt, dash pattern=on 4pt off 3pt]
            ([xshift=0.55\textwidth+8pt]frame.north west) --
            ([xshift=0.55\textwidth+8pt]frame.south west);
    }
]
\textbf{Q.202} \quad {\footnotesize\textsc{[NAT $\cdot$ 2 Mark(s) $\cdot$ No Neg]}} \hfill {\footnotesize\textbf{GATE 2025 CS2} \quad \tiny\texttt{2025\_CS2\_CS\_Q57}}
\vspace{2pt}\hrule height 0.3pt\vspace{4pt}

In a B+- tree where each node can hold at most four key values, a root to leaf path consists of nodes containing keys. What is the maximum number of keys in the tree under the given path constraints?

\par\vspace{6pt}\noindent\textbf{Your Answer:} \fbox{\rule{0pt}{1.3ex}\hspace{3cm}}

\tcblower

\textbf{Official Key:} \textbf{\large 33 to 33} \hfill \textit{\small [Hard]}
\vspace{2pt}\hrule height 0.3pt\vspace{3pt}
\textbf{Topic:} Databases \\
\textbf{Subtopic:} \textit{File organization, indexing: B and B+ trees}
\vspace{2pt}\hrule height 0.3pt\vspace{3pt}
\textbf{Guide / Hint:} \\
{\small In a B+ tree of order 5 (max 4 keys, max 5 children per internal node), min keys per non-root internal node is $\lceil 5/2 \rceil - 1 = 2$. Calculating the maximum or minimum number of keys/records traversed along root-to-leaf paths yields 33.}

\end{tcolorbox}
\vspace{6pt}\hrule height 0.35pt\vspace{6pt}
\vspace{8pt}
\subsection*{Computer Networks}
\vspace{-4pt}
\noindent
\begin{tcolorbox}[
    blanker,
    sidebyside,
    sidebyside align=top,
    lefthand width=0.55\textwidth,
    top=4pt, bottom=4pt, left=0pt, right=0pt,
    middle=16pt,
    underlay={
        \draw[dashed, black!80, line width=0.8pt, dash pattern=on 4pt off 3pt]
            ([xshift=0.55\textwidth+8pt]frame.north west) --
            ([xshift=0.55\textwidth+8pt]frame.south west);
    }
]
\textbf{Q.203} \quad {\footnotesize\textsc{[NAT $\cdot$ 2 Mark(s) $\cdot$ No Neg]}} \hfill {\footnotesize\textbf{GATE 2018 CS} \quad \tiny\texttt{2018\_CS\_GA\_Q54}}
\vspace{2pt}\hrule height 0.3pt\vspace{4pt}

Consider an IP packet with a length of 4,500 bytes that includes a 20-byte IPv4 header and a 40-byte TCP header. The packet is forwarded across an IP network with an MTU of 600 bytes. The value of the fragment offset field in the third fragment (in 8-byte units) is \_\_\_\_\_\_\_\_.

\par\vspace{6pt}\noindent\textbf{Your Answer:} \fbox{\rule{0pt}{1.3ex}\hspace{3cm}}

\tcblower

\textbf{Official Key:} \textbf{\large 144 to 144} \hfill \textit{\small [Medium]}
\vspace{2pt}\hrule height 0.3pt\vspace{3pt}
\textbf{Topic:} Computer Networks \\
\textbf{Subtopic:} \textit{Fragmentation and IP addressing, IPv4, CIDR notation}
\vspace{2pt}\hrule height 0.3pt\vspace{3pt}
\textbf{Guide / Hint:} \\
{\small Total payload to fragment $= 4500 - 20 = 4480$ bytes. With $\text{MTU} = 600$, maximum payload per fragment must be a multiple of 8: $\lfloor (600 - 20) / 8 \rfloor \times 8 = 72 \times 8 = 576$ bytes. Fragment 1 carries bytes 0 to 575 (offset 0). Fragment 2 carries bytes 576 to 1151 (offset $576/8 = 72$). Fragment 3 carries bytes 1152 onwards, so its fragment offset field is $\frac{1152}{8} = 144$.}

\end{tcolorbox}
\vspace{6pt}\hrule height 0.35pt\vspace{6pt}
\noindent
\begin{tcolorbox}[
    blanker,
    sidebyside,
    sidebyside align=top,
    lefthand width=0.55\textwidth,
    top=4pt, bottom=4pt, left=0pt, right=0pt,
    middle=16pt,
    underlay={
        \draw[dashed, black!80, line width=0.8pt, dash pattern=on 4pt off 3pt]
            ([xshift=0.55\textwidth+8pt]frame.north west) --
            ([xshift=0.55\textwidth+8pt]frame.south west);
    }
]
\textbf{Q.204} \quad {\footnotesize\textsc{[NAT $\cdot$ 2 Mark(s) $\cdot$ No Neg]}} \hfill {\footnotesize\textbf{GATE 2022 CS} \quad \tiny\texttt{2022\_CS\_CS\_Q59}}
\vspace{2pt}\hrule height 0.3pt\vspace{4pt}

Consider a 100 Mbps link between an earth station (sender) and a satellite (receiver) at an altitude of 2100 km. The signal propagates at a speed of $3 \times 10^8\text{ m/s}$. The time taken (in milliseconds, rounded off to two decimal places) for the receiver to completely receive a packet of 1000 bytes transmitted by the sender is \_\_\_\_\_\_\_\_\_\. (Answer in numeric format)

\par\vspace{6pt}\noindent\textbf{Your Answer:} \fbox{\rule{0pt}{1.3ex}\hspace{3cm}}

\tcblower

\textbf{Official Key:} \textbf{\large 7.07 to 7.09} \hfill \textit{\small [Easy]}
\vspace{2pt}\hrule height 0.3pt\vspace{3pt}
\textbf{Topic:} Computer Networks \\
\textbf{Subtopic:} \textit{Data link layer: framing, error detection, Medium Access Control, Ethernet bridging}
\vspace{2pt}\hrule height 0.3pt\vspace{3pt}
\textbf{Guide / Hint:} \\
{\small Propagation time is $T_{\text{prop}} = \frac{2100 \times 10^3\text{ m}}{3 \times 10^8\text{ m/s}} = 7.00\text{ ms}$. Transmission time is $T_{\text{trans}} = \frac{1000 \times 8\text{ bits}}{100 \times 10^6\text{ bps}} = 0.08\text{ ms}$. Total time to completely receive the packet is $T_{\text{prop}} + T_{\text{trans}} = 7.00 + 0.08 = 7.08\text{ ms}$.}

\end{tcolorbox}
\vspace{6pt}\hrule height 0.35pt\vspace{6pt}
\noindent
\begin{tcolorbox}[
    blanker,
    sidebyside,
    sidebyside align=top,
    lefthand width=0.55\textwidth,
    top=4pt, bottom=4pt, left=0pt, right=0pt,
    middle=16pt,
    underlay={
        \draw[dashed, black!80, line width=0.8pt, dash pattern=on 4pt off 3pt]
            ([xshift=0.55\textwidth+8pt]frame.north west) --
            ([xshift=0.55\textwidth+8pt]frame.south west);
    }
]
\textbf{Q.205} \quad {\footnotesize\textsc{[NAT $\cdot$ 2 Mark(s) $\cdot$ No Neg]}} \hfill {\footnotesize\textbf{GATE 2022 CS} \quad \tiny\texttt{2022\_CS\_CS\_Q60}}
\vspace{2pt}\hrule height 0.3pt\vspace{4pt}

Consider the data transfer using TCP over a 1 Gbps link. Assuming that the maximum segment lifetime (MSL) is set to 60 seconds, the minimum number of bits required for the sequence number field of the TCP header, to prevent the sequence number space from wrapping around during the MSL is \_\_\_\_\_\_\_\_\_\_\_\_\. (Answer in integer)

\par\vspace{6pt}\noindent\textbf{Your Answer:} \fbox{\rule{0pt}{1.3ex}\hspace{3cm}}

\tcblower

\textbf{Official Key:} \textbf{\large 33 to 33} \hfill \textit{\small [Medium]}
\vspace{2pt}\hrule height 0.3pt\vspace{3pt}
\textbf{Topic:} Computer Networks \\
\textbf{Subtopic:} \textit{Transport layer: flow control and congestion control, UDP, TCP, sockets}
\vspace{2pt}\hrule height 0.3pt\vspace{3pt}
\textbf{Guide / Hint:} \\
{\small Data transmitted at 1 Gbps line rate in $\text{MSL} = 60\text{ s}$ is $\frac{10^9\text{ bps}}{8} \times 60\text{ s} = 7.5 \times 10^9\text{ bytes}$. The sequence number space must satisfy $2^k \ge 7.5 \times 10^9$. Since $2^{32} \approx 4.29 \times 10^9 < 7.5 \times 10^9$ and $2^{33} \approx 8.59 \times 10^9 > 7.5 \times 10^9$, at least 33 bits are required.}

\end{tcolorbox}
\vspace{6pt}\hrule height 0.35pt\vspace{6pt}
\noindent
\begin{tcolorbox}[
    blanker,
    sidebyside,
    sidebyside align=top,
    lefthand width=0.55\textwidth,
    top=4pt, bottom=4pt, left=0pt, right=0pt,
    middle=16pt,
    underlay={
        \draw[dashed, black!80, line width=0.8pt, dash pattern=on 4pt off 3pt]
            ([xshift=0.55\textwidth+8pt]frame.north west) --
            ([xshift=0.55\textwidth+8pt]frame.south west);
    }
]
\textbf{Q.206} \quad {\footnotesize\textsc{[MSQ $\cdot$ 2 Mark(s) $\cdot$ No Neg]}} \hfill {\footnotesize\textbf{GATE 2023 CS} \quad \tiny\texttt{2023\_CS\_CS\_Q50}}
\vspace{2pt}\hrule height 0.3pt\vspace{4pt}

Suppose you are asked to design a new reliable byte-stream transport protocol like TCP. This protocol, named myTCP, runs over a 100 Mbps network with Round Trip Time of 150 milliseconds and the maximum segment lifetime of 2 minutes.
Which of the following is/are valid lengths of the Sequence Number field in the myTCP header?

\par\vspace{6pt}\noindent
\begin{tabular}{@{}p{0.48\linewidth}@{\hspace{0.04\linewidth}}p{0.48\linewidth}@{}}
$\square$\ \textbf{(A)} 30 bits & $\square$\ \textbf{(B)} 32 bits \\
\addlinespace[3pt]
$\square$\ \textbf{(C)} 34 bits & $\square$\ \textbf{(D)} 36 bits
\end{tabular}

\tcblower

\textbf{Official Key:} \textbf{\large B, C, D} \hfill \textit{\small [Hard]}
\vspace{2pt}\hrule height 0.3pt\vspace{3pt}
\textbf{Topic:} Computer Networks \\
\textbf{Subtopic:} \textit{Transport layer: flow control and congestion control, UDP, TCP, sockets}
\vspace{2pt}\hrule height 0.3pt\vspace{3pt}
\textbf{Guide / Hint:} \\
{\small To prevent sequence number wraparound within the Maximum Segment Lifetime ($MSL = 2\text{ min} = 120\text{ s}$), the sequence number space $2^k$ must be at least the data bytes transmitted at line rate during MSL: $2^k \ge \frac{100 \times 10^6\text{ bps}}{8} \times 120\text{ s} = 1.5 \times 10^9\text{ bytes}$. Since $2^{30} \approx 1.07 \times 10^9 < 1.5 \times 10^9$ and $2^{31} \approx 2.15 \times 10^9 > 1.5 \times 10^9$, at least 31 bits are required. Thus, 32 bits, 34 bits, and 36 bits are all valid lengths.}

\end{tcolorbox}
\vspace{6pt}\hrule height 0.35pt\vspace{6pt}
\noindent
\begin{tcolorbox}[
    blanker,
    sidebyside,
    sidebyside align=top,
    lefthand width=0.55\textwidth,
    top=4pt, bottom=4pt, left=0pt, right=0pt,
    middle=16pt,
    underlay={
        \draw[dashed, black!80, line width=0.8pt, dash pattern=on 4pt off 3pt]
            ([xshift=0.55\textwidth+8pt]frame.north west) --
            ([xshift=0.55\textwidth+8pt]frame.south west);
    }
]
\textbf{Q.207} \quad {\footnotesize\textsc{[MSQ $\cdot$ 2 Mark(s) $\cdot$ No Neg]}} \hfill {\footnotesize\textbf{GATE 2023 CS} \quad \tiny\texttt{2023\_CS\_CS\_Q52}}
\vspace{2pt}\hrule height 0.3pt\vspace{4pt}

Suppose in a web browser, you click on the `www.gate-2023.in` URL. The browser cache is empty. The IP address for this URL is not cached in your local host, so a DNS lookup is triggered (by the local DNS server deployed on your local host) over the 3-tier DNS hierarchy in an iterative mode. No resource records are cached anywhere across all DNS servers.
Let RTT denote the round trip time between your local host and DNS servers in the DNS hierarchy. The round trip time between the local host and the web server hosting `www.gate-2023.in` is also equal to RTT. The HTML file associated with the URL is small enough to have negligible transmission time and negligible rendering time by your web browser, which references 10 equally small objects on the same web server.
Which of the following statements is/are CORRECT about the minimum elapsed time between clicking on the URL and your browser fully rendering it?

\par\vspace{6pt}\noindent
\begin{tabular}{@{}p{0.98\linewidth}@{}}
$\square$\ \textbf{(A)} 7 RTTs, in case of non-persistent HTTP with 5 parallel TCP connections. \\
\addlinespace[3pt]
$\square$\ \textbf{(B)} 5 RTTs, in case of persistent HTTP with pipelining. \\
\addlinespace[3pt]
$\square$\ \textbf{(C)} 9 RTTs, in case of non-persistent HTTP with 5 parallel TCP connections. \\
\addlinespace[3pt]
$\square$\ \textbf{(D)} 6 RTTs, in case of persistent HTTP with pipelining.
\end{tabular}

\tcblower

\textbf{Official Key:} \textbf{\large C, D} \hfill \textit{\small [Hard]}
\vspace{2pt}\hrule height 0.3pt\vspace{3pt}
\textbf{Topic:} Computer Networks \\
\textbf{Subtopic:} \textit{Application layer protocols: DNS, SMTP, HTTP, FTP, Email}
\vspace{2pt}\hrule height 0.3pt\vspace{3pt}
\textbf{Guide / Hint:} \\
{\small Resolving the URL via 3-tier iterative DNS requires $3\text{ RTT}$. Initiating TCP connection and requesting the base HTML page requires $1 + 1 = 2\text{ RTT}$, so the base HTML is received after $3 + 2 = 5\text{ RTT}$. With persistent HTTP pipelining, all 10 objects are fetched in 1 round trip over the established connection, giving $5 + 1 = 6\text{ RTT}$ (D is TRUE). With non-persistent HTTP and 5 parallel connections, fetching 10 objects takes 2 sequential batches; each batch takes $1\text{ RTT (TCP)} + 1\text{ RTT (HTTP)} = 2\text{ RTT}$, adding $2 \times 2 = 4\text{ RTT}$, giving $5 + 4 = 9\text{ RTT}$ (C is TRUE).}

\end{tcolorbox}
\vspace{6pt}\hrule height 0.35pt\vspace{6pt}
\noindent
\begin{tcolorbox}[
    blanker,
    sidebyside,
    sidebyside align=top,
    lefthand width=0.55\textwidth,
    top=4pt, bottom=4pt, left=0pt, right=0pt,
    middle=16pt,
    underlay={
        \draw[dashed, black!80, line width=0.8pt, dash pattern=on 4pt off 3pt]
            ([xshift=0.55\textwidth+8pt]frame.north west) --
            ([xshift=0.55\textwidth+8pt]frame.south west);
    }
]
\textbf{Q.208} \quad {\footnotesize\textsc{[NAT $\cdot$ 2 Mark(s) $\cdot$ No Neg]}} \hfill {\footnotesize\textbf{GATE 2023 CS} \quad \tiny\texttt{2023\_CS\_CS\_Q65}}
\vspace{2pt}\hrule height 0.3pt\vspace{4pt}

The forwarding table of a router is shown below:
\begin{center}
\begin{tabular}{|l|l|c|}
\hline
\textbf{Subnet Number} & \textbf{Subnet Mask} & \textbf{Interface ID} \\
\hline
200.150.0.0 & 255.255.0.0 & 1 \\
200.150.64.0 & 255.255.224.0 & 2 \\
200.150.68.0 & 255.255.255.0 & 3 \\
200.150.68.64 & 255.255.255.224 & 4 \\
Default & & 0 \\
\hline
\end{tabular}
\end{center}
A packet addressed to a destination address 200.150.68.118 arrives at the router. It will be forwarded to the interface with ID \_\_\_\_\_\_\_\_\. (Answer in integer)

\par\vspace{6pt}\noindent\textbf{Your Answer:} \fbox{\rule{0pt}{1.3ex}\hspace{3cm}}

\tcblower

\textbf{Official Key:} \textbf{\large 3 to 3} \hfill \textit{\small [Medium]}
\vspace{2pt}\hrule height 0.3pt\vspace{3pt}
\textbf{Topic:} Computer Networks \\
\textbf{Subtopic:} \textit{Routing protocols: shortest path, flooding, distance vector and link state routing}
\vspace{2pt}\hrule height 0.3pt\vspace{3pt}
\textbf{Guide / Hint:} \\
{\small For destination 200.150.68.118: subnet /16 (prefix length 16) matches. Subnet 200.150.64.0/19 (mask 255.255.224.0) matches since $68 \& 224 = 64$. Subnet 200.150.68.0/24 (mask 255.255.255.0) matches since $68 = 68$. Subnet 200.150.68.64/27 (mask 255.255.255.224) does not match since $118 \& 224 = 96 \neq 64$. By the Longest Prefix Match rule, the /24 route (interface 3) is selected.}

\end{tcolorbox}
\vspace{6pt}\hrule height 0.35pt\vspace{6pt}
\noindent
\begin{tcolorbox}[
    blanker,
    sidebyside,
    sidebyside align=top,
    lefthand width=0.55\textwidth,
    top=4pt, bottom=4pt, left=0pt, right=0pt,
    middle=16pt,
    underlay={
        \draw[dashed, black!80, line width=0.8pt, dash pattern=on 4pt off 3pt]
            ([xshift=0.55\textwidth+8pt]frame.north west) --
            ([xshift=0.55\textwidth+8pt]frame.south west);
    }
]
\textbf{Q.209} \quad {\footnotesize\textsc{[MCQ $\cdot$ 2 Mark(s) $\cdot$ Neg: 0.66]}} \hfill {\footnotesize\textbf{GATE 2024 CS1} \quad \tiny\texttt{2024\_CS1\_CS\_Q36}}
\vspace{2pt}\hrule height 0.3pt\vspace{4pt}

Consider a network path P---Q---R between nodes P and R via router Q. Node P sends a file of size $10^6$ bytes to R via this path by splitting the file into chunks of $10^3$ bytes each. Node P sends these chunks one after the other without any wait time between the successive chunk transmissions. Assume that the size of extra headers added to these chunks is negligible, and that the chunk size is less than the MTU.
\par\medskip
Each of the links P---Q and Q---R has a bandwidth of $10^6$ bits/sec, and negligible propagation latency. Router Q immediately transmits every packet it receives from P to R, with negligible processing and queueing delays. Router Q can simultaneously receive on link P---Q and transmit on link Q---R.
\par\medskip
Assume P starts transmitting the chunks at time $t = 0$.
Which one of the following options gives the time (in seconds, rounded off to 3 decimal places) at which R receives all the chunks of the file?

\par\vspace{6pt}\noindent
\begin{tabular}{@{}p{0.48\linewidth}@{\hspace{0.04\linewidth}}p{0.48\linewidth}@{}}
$\square$\ \textbf{(A)} 8.000 & $\square$\ \textbf{(B)} 8.008 \\
\addlinespace[3pt]
$\square$\ \textbf{(C)} 15.992 & $\square$\ \textbf{(D)} 16.000
\end{tabular}

\tcblower

\textbf{Official Key:} \textbf{\large B} \hfill \textit{\small [Medium]}
\vspace{2pt}\hrule height 0.3pt\vspace{3pt}
\textbf{Topic:} Computer Networks \\
\textbf{Subtopic:} \textit{Concept of layering: OSI and TCP/IP Protocol Stacks}
\vspace{2pt}\hrule height 0.3pt\vspace{3pt}
\textbf{Guide / Hint:} \\
{\small The file size is $10^6$ bytes split into 1000 chunks of $10^3$ bytes ($8000$ bits) each. Link transmission time for one chunk is $\frac{8000\text{ bits}}{10^6\text{ bps}} = 0.008\text{ s}$. Node P completes sending all 1000 chunks at $t = 1000 \times 0.008 = 8.000\text{ s}$. The last chunk arrives at router Q at $t = 8.000\text{ s}$ and takes another $0.008\text{ s}$ to reach R over link Q---R, completing at $8.008\text{ s}$.}

\end{tcolorbox}
\vspace{6pt}\hrule height 0.35pt\vspace{6pt}
\noindent
\begin{tcolorbox}[
    blanker,
    sidebyside,
    sidebyside align=top,
    lefthand width=0.55\textwidth,
    top=4pt, bottom=4pt, left=0pt, right=0pt,
    middle=16pt,
    underlay={
        \draw[dashed, black!80, line width=0.8pt, dash pattern=on 4pt off 3pt]
            ([xshift=0.55\textwidth+8pt]frame.north west) --
            ([xshift=0.55\textwidth+8pt]frame.south west);
    }
]
\textbf{Q.210} \quad {\footnotesize\textsc{[NAT $\cdot$ 2 Mark(s) $\cdot$ No Neg]}} \hfill {\footnotesize\textbf{GATE 2024 CS1} \quad \tiny\texttt{2024\_CS1\_CS\_Q58}}
\vspace{2pt}\hrule height 0.3pt\vspace{4pt}

Consider the entries shown below in the forwarding table of an IP router. Each entry consists of an IP prefix and the corresponding next hop router for packets whose destination IP address matches the prefix. The notation ``/N'' in a prefix indicates a subnet mask with the most significant $N$ bits set to 1.
\begin{center}
\begin{tabular}{|l|c|}
\hline
\textbf{Prefix} & \textbf{Next hop router} \\
\hline
10.1.1.0/24 & R1 \\
10.1.1.128/25 & R2 \\
10.1.1.64/26 & R3 \\
10.1.1.192/26 & R4 \\
\hline
\end{tabular}
\end{center}
This router forwards 20 packets each to 5 hosts. The IP addresses of the hosts are 10.1.1.16, 10.1.1.72, 10.1.1.132, 10.1.1.191, and 10.1.1.205. The number of packets forwarded via the next hop router R2 is \_\_\_\_\_\_\_\. (Answer in integer)

\par\vspace{6pt}\noindent\textbf{Your Answer:} \fbox{\rule{0pt}{1.3ex}\hspace{3cm}}

\tcblower

\textbf{Official Key:} \textbf{\large 40 to 40} \hfill \textit{\small [Medium]}
\vspace{2pt}\hrule height 0.3pt\vspace{3pt}
\textbf{Topic:} Computer Networks \\
\textbf{Subtopic:} \textit{Routing protocols: shortest path, flooding, distance vector and link state routing}
\vspace{2pt}\hrule height 0.3pt\vspace{3pt}
\textbf{Guide / Hint:} \\
{\small Using longest prefix match for each host address: 10.1.1.16 matches /24 (R1); 10.1.1.72 matches /24 and 10.1.1.64/26 (R3); 10.1.1.132 matches /24 and 10.1.1.128/25 (R2); 10.1.1.191 matches /24 and 10.1.1.128/25 (R2); 10.1.1.205 matches /24, /25, and 10.1.1.192/26 (R4). Thus, 2 hosts (10.1.1.132 and 10.1.1.191) match router R2. Total packets forwarded via R2 $= 2 \times 20 = 40$.}

\end{tcolorbox}
\vspace{6pt}\hrule height 0.35pt\vspace{6pt}
\noindent
\begin{tcolorbox}[
    blanker,
    sidebyside,
    sidebyside align=top,
    lefthand width=0.55\textwidth,
    top=4pt, bottom=4pt, left=0pt, right=0pt,
    middle=16pt,
    underlay={
        \draw[dashed, black!80, line width=0.8pt, dash pattern=on 4pt off 3pt]
            ([xshift=0.55\textwidth+8pt]frame.north west) --
            ([xshift=0.55\textwidth+8pt]frame.south west);
    }
]
\textbf{Q.211} \quad {\footnotesize\textsc{[NAT $\cdot$ 2 Mark(s) $\cdot$ No Neg]}} \hfill {\footnotesize\textbf{GATE 2024 CS1} \quad \tiny\texttt{2024\_CS1\_CS\_Q65}}
\vspace{2pt}\hrule height 0.3pt\vspace{4pt}

Consider sending an IP datagram of size 1420 bytes (including 20 bytes of IP header) from a sender to a receiver over a path of two links with a router between them. The first link (sender to router) has an MTU (Maximum Transmission Unit) size of 542 bytes, while the second link (router to receiver) has an MTU size of 360 bytes. The number of fragments that would be delivered at the receiver is \_\_\_\_\_\_\_\_\. (Answer in integer)

\par\vspace{6pt}\noindent\textbf{Your Answer:} \fbox{\rule{0pt}{1.3ex}\hspace{3cm}}

\tcblower

\textbf{Official Key:} \textbf{\large 6 to 6} \hfill \textit{\small [Medium]}
\vspace{2pt}\hrule height 0.3pt\vspace{3pt}
\textbf{Topic:} Computer Networks \\
\textbf{Subtopic:} \textit{Fragmentation and IP addressing, IPv4, CIDR notation}
\vspace{2pt}\hrule height 0.3pt\vspace{3pt}
\textbf{Guide / Hint:} \\
{\small The datagram has 20 bytes IP header and 1400 bytes data. Link 1 has $\text{MTU} = 542$ bytes (max data per fragment is $\lfloor (542-20)/8 \rfloor \times 8 = 520$ bytes), producing 3 fragments with data sizes 520, 520, and 360 bytes. Link 2 has $\text{MTU} = 360$ bytes (max data is $\lfloor (360-20)/8 \rfloor \times 8 = 336$ bytes). At router Q, fragment 1 (520 bytes) is fragmented into 2 ($336 + 184$), fragment 2 (520 bytes) is fragmented into 2 ($336 + 184$), and fragment 3 (360 bytes) is fragmented into 2 ($336 + 24$). Total fragments delivered at the receiver $= 2 + 2 + 2 = 6$.}

\end{tcolorbox}
\vspace{6pt}\hrule height 0.35pt\vspace{6pt}
\noindent
\begin{tcolorbox}[
    blanker,
    sidebyside,
    sidebyside align=top,
    lefthand width=0.55\textwidth,
    top=4pt, bottom=4pt, left=0pt, right=0pt,
    middle=16pt,
    underlay={
        \draw[dashed, black!80, line width=0.8pt, dash pattern=on 4pt off 3pt]
            ([xshift=0.55\textwidth+8pt]frame.north west) --
            ([xshift=0.55\textwidth+8pt]frame.south west);
    }
]
\textbf{Q.212} \quad {\footnotesize\textsc{[MCQ $\cdot$ 2 Mark(s) $\cdot$ Neg: 0.66]}} \hfill {\footnotesize\textbf{GATE 2024 CS2} \quad \tiny\texttt{2024\_CS2\_CS\_Q38}}
\vspace{2pt}\hrule height 0.3pt\vspace{4pt}

Which one of the following CIDR prefixes exactly represents the range of IP addresses 10.12.2.0 to 10.12.3.255?

\par\vspace{6pt}\noindent
\begin{tabular}{@{}p{0.48\linewidth}@{\hspace{0.04\linewidth}}p{0.48\linewidth}@{}}
$\square$\ \textbf{(A)} 10.12.2.0/23 & $\square$\ \textbf{(B)} 10.12.2.0/24 \\
\addlinespace[3pt]
$\square$\ \textbf{(C)} 10.12.0.0/22 & $\square$\ \textbf{(D)} 10.12.2.0/22
\end{tabular}

\tcblower

\textbf{Official Key:} \textbf{\large A} \hfill \textit{\small [Easy]}
\vspace{2pt}\hrule height 0.3pt\vspace{3pt}
\textbf{Topic:} Computer Networks \\
\textbf{Subtopic:} \textit{Fragmentation and IP addressing, IPv4, CIDR notation}
\vspace{2pt}\hrule height 0.3pt\vspace{3pt}
\textbf{Guide / Hint:} \\
{\small The address range spans from $10.12.2.0$ ($00000010_2$) to $10.12.3.255$ ($00000011_2$). The first 23 bits ($8 + 8 + 7 = 23$) are identical ($10.12.2$). Total addresses in range $= 2 \times 256 = 512 = 2^{32 - 23}$, which exactly matches CIDR prefix `10.12.2.0/23`.}

\end{tcolorbox}
\vspace{6pt}\hrule height 0.35pt\vspace{6pt}
\noindent
\begin{tcolorbox}[
    blanker,
    sidebyside,
    sidebyside align=top,
    lefthand width=0.55\textwidth,
    top=4pt, bottom=4pt, left=0pt, right=0pt,
    middle=16pt,
    underlay={
        \draw[dashed, black!80, line width=0.8pt, dash pattern=on 4pt off 3pt]
            ([xshift=0.55\textwidth+8pt]frame.north west) --
            ([xshift=0.55\textwidth+8pt]frame.south west);
    }
]
\textbf{Q.213} \quad {\footnotesize\textsc{[NAT $\cdot$ 2 Mark(s) $\cdot$ No Neg]}} \hfill {\footnotesize\textbf{GATE 2024 CS2} \quad \tiny\texttt{2024\_CS2\_CS\_Q54}}
\vspace{2pt}\hrule height 0.3pt\vspace{4pt}

Consider a TCP connection operating at a point of time with the congestion window of size 12 MSS (Maximum Segment Size), when a timeout occurs due to packet loss. Assuming that all the segments transmitted in the next two RTTs (Round Trip Time) are acknowledged correctly, the congestion window size (in MSS) during the third RTT will be \_\_\_\_\_\_\. (Answer in integer)

\par\vspace{6pt}\noindent\textbf{Your Answer:} \fbox{\rule{0pt}{1.3ex}\hspace{3cm}}

\tcblower

\textbf{Official Key:} \textbf{\large 4 to 4} \hfill \textit{\small [Medium]}
\vspace{2pt}\hrule height 0.3pt\vspace{3pt}
\textbf{Topic:} Computer Networks \\
\textbf{Subtopic:} \textit{Transport layer: flow control and congestion control, UDP, TCP, sockets}
\vspace{2pt}\hrule height 0.3pt\vspace{3pt}
\textbf{Guide / Hint:} \\
{\small When timeout occurs at $\text{cwnd} = 12\text{ MSS}$, slow start threshold is set to $\text{ssthresh} = \lfloor 12 / 2 \rfloor = 6\text{ MSS}$, and $\text{cwnd}$ resets to $1\text{ MSS}$ for the 1st RTT. After receiving ACKs in the 1st RTT, slow start doubles $\text{cwnd}$ to $2\text{ MSS}$ during the 2nd RTT. After receiving ACKs in the 2nd RTT, $\text{cwnd}$ doubles again to $4\text{ MSS}$ during the 3rd RTT.}

\end{tcolorbox}
\vspace{6pt}\hrule height 0.35pt\vspace{6pt}
\noindent
\begin{tcolorbox}[
    blanker,
    sidebyside,
    sidebyside align=top,
    lefthand width=0.55\textwidth,
    top=4pt, bottom=4pt, left=0pt, right=0pt,
    middle=16pt,
    underlay={
        \draw[dashed, black!80, line width=0.8pt, dash pattern=on 4pt off 3pt]
            ([xshift=0.55\textwidth+8pt]frame.north west) --
            ([xshift=0.55\textwidth+8pt]frame.south west);
    }
]
\textbf{Q.214} \quad {\footnotesize\textsc{[NAT $\cdot$ 2 Mark(s) $\cdot$ No Neg]}} \hfill {\footnotesize\textbf{GATE 2024 CS2} \quad \tiny\texttt{2024\_CS2\_CS\_Q55}}
\vspace{2pt}\hrule height 0.3pt\vspace{4pt}

Consider an Ethernet segment with a transmission speed of $10^8$ bits/sec and a maximum segment length of 500 meters. If the speed of propagation of the signal in the medium is $2 \times 10^8$ meters/sec, then the minimum frame size (in bits) required for collision detection is \_\_\_\_\_\_\. (Answer in integer)

\par\vspace{6pt}\noindent\textbf{Your Answer:} \fbox{\rule{0pt}{1.3ex}\hspace{3cm}}

\tcblower

\textbf{Official Key:} \textbf{\large 500 to 500} \hfill \textit{\small [Easy]}
\vspace{2pt}\hrule height 0.3pt\vspace{3pt}
\textbf{Topic:} Computer Networks \\
\textbf{Subtopic:} \textit{Data link layer: framing, error detection, Medium Access Control, Ethernet bridging}
\vspace{2pt}\hrule height 0.3pt\vspace{3pt}
\textbf{Guide / Hint:} \\
{\small Propagation delay $T_p = \frac{\text{Distance}}{\text{Velocity}} = \frac{500}{2 \times 10^8} = 2.5 \times 10^{-6}\text{ s} = 2.5\,\mu\text{s}$. For collision detection in CSMA/CD, transmission time must satisfy $T_t \ge 2 T_p = 5\,\mu\text{s}$. Minimum frame size $L = 2 T_p \times B = (5 \times 10^{-6}\text{ s}) \times 10^8\text{ bps} = 500$ bits.}

\end{tcolorbox}
\vspace{6pt}\hrule height 0.35pt\vspace{6pt}
\noindent
\begin{tcolorbox}[
    blanker,
    sidebyside,
    sidebyside align=top,
    lefthand width=0.55\textwidth,
    top=4pt, bottom=4pt, left=0pt, right=0pt,
    middle=16pt,
    underlay={
        \draw[dashed, black!80, line width=0.8pt, dash pattern=on 4pt off 3pt]
            ([xshift=0.55\textwidth+8pt]frame.north west) --
            ([xshift=0.55\textwidth+8pt]frame.south west);
    }
]
\textbf{Q.215} \quad {\footnotesize\textsc{[MCQ $\cdot$ 2 Mark(s) $\cdot$ Neg: 0.66]}} \hfill {\footnotesize\textbf{GATE 2025 CS1} \quad \tiny\texttt{2025\_CS1\_CS\_Q40}}
\vspace{2pt}\hrule height 0.3pt\vspace{4pt}

A packet with the destination IP address 145.36.109.70 arrives at a router whose routing table is shown below:
\par\medskip
\begin{tabular}{|l|l|c|}
\hline
\textbf{Subnet Address} & \textbf{Subnet Mask (CIDR)} & \textbf{Interface} \\
\hline
145.36.0.0 & /16 & E1 \\
145.36.128.0 & /17 & E2 \\
145.36.64.0 & /18 & E3 \\
145.36.255.0 & /24 & E4 \\
Default & -- & E5 \\
\hline
\end{tabular}
\par\medskip
Which interface will the packet be forwarded to?

\par\vspace{6pt}\noindent
\begin{tabular}{@{}p{0.48\linewidth}@{\hspace{0.04\linewidth}}p{0.48\linewidth}@{}}
$\square$\ \textbf{(A)} E3 & $\square$\ \textbf{(B)} E1 \\
\addlinespace[3pt]
$\square$\ \textbf{(C)} E2 & $\square$\ \textbf{(D)} E5
\end{tabular}

\tcblower

\textbf{Official Key:} \textbf{\large A} \hfill \textit{\small [Medium]}
\vspace{2pt}\hrule height 0.3pt\vspace{3pt}
\textbf{Topic:} Computer Networks \\
\textbf{Subtopic:} \textit{Routing protocols: shortest path, flooding, distance vector and link state routing}
\vspace{2pt}\hrule height 0.3pt\vspace{3pt}
\textbf{Guide / Hint:} \\
{\small The destination address has third octet $109 = 01101101_2$. For /16, $145.36.0.0$ matches. For /17 ($128 = 10000000_2$), the leading bit is 0, so /17 does not match. For /18 ($64 = 01000000_2$), the first two bits are $01_2$, which matches $109$ ($01101101_2$). For /24, 255 does not match. By Longest Prefix Match, /18 (interface E3) is chosen.}

\end{tcolorbox}
\vspace{6pt}\hrule height 0.35pt\vspace{6pt}
\noindent
\begin{tcolorbox}[
    blanker,
    sidebyside,
    sidebyside align=top,
    lefthand width=0.55\textwidth,
    top=4pt, bottom=4pt, left=0pt, right=0pt,
    middle=16pt,
    underlay={
        \draw[dashed, black!80, line width=0.8pt, dash pattern=on 4pt off 3pt]
            ([xshift=0.55\textwidth+8pt]frame.north west) --
            ([xshift=0.55\textwidth+8pt]frame.south west);
    }
]
\textbf{Q.216} \quad {\footnotesize\textsc{[NAT $\cdot$ 2 Mark(s) $\cdot$ No Neg]}} \hfill {\footnotesize\textbf{GATE 2025 CS1} \quad \tiny\texttt{2025\_CS1\_CS\_Q56}}
\vspace{2pt}\hrule height 0.3pt\vspace{4pt}

Suppose a 5-bit message is transmitted from a source to a destination through a noisy channel. The probability that a bit of the message gets flipped during transmission is 0.01. Flipping of each bit is independent of one another. The probability that the message is delivered error-free to the destination is \_\_\_\_\_\_\. (rounded off to three decimal places)

\par\vspace{6pt}\noindent\textbf{Your Answer:} \fbox{\rule{0pt}{1.3ex}\hspace{3cm}}

\tcblower

\textbf{Official Key:} \textbf{\large 0.949 to 0.952} \hfill \textit{\small [Easy]}
\vspace{2pt}\hrule height 0.3pt\vspace{3pt}
\textbf{Topic:} Computer Networks \\
\textbf{Subtopic:} \textit{Data link layer: framing, error detection, Medium Access Control, Ethernet bridging}
\vspace{2pt}\hrule height 0.3pt\vspace{3pt}
\textbf{Guide / Hint:} \\
{\small For each bit, the probability of correct transmission without flipping is $1 - 0.01 = 0.99$. Since all 5 bit transmissions are mutually independent, the probability that all 5 bits are transmitted error-free is $(0.99)^5 = 0.95099 \approx 0.951$.}

\end{tcolorbox}
\vspace{6pt}\hrule height 0.35pt\vspace{6pt}
\noindent
\begin{tcolorbox}[
    blanker,
    sidebyside,
    sidebyside align=top,
    lefthand width=0.55\textwidth,
    top=4pt, bottom=4pt, left=0pt, right=0pt,
    middle=16pt,
    underlay={
        \draw[dashed, black!80, line width=0.8pt, dash pattern=on 4pt off 3pt]
            ([xshift=0.55\textwidth+8pt]frame.north west) --
            ([xshift=0.55\textwidth+8pt]frame.south west);
    }
]
\textbf{Q.217} \quad {\footnotesize\textsc{[NAT $\cdot$ 2 Mark(s) $\cdot$ No Neg]}} \hfill {\footnotesize\textbf{GATE 2025 CS1} \quad \tiny\texttt{2025\_CS1\_CS\_Q57}}
\vspace{2pt}\hrule height 0.3pt\vspace{4pt}

Suppose a message of size 15000 bytes is transmitted from a source to a destination using IPv4 protocol via two routers as shown: Source $\to$ Router-1 ($\text{MTU} = 5000$ bytes) $\to$ Router-2 ($\text{MTU} = 3000$ bytes) $\to$ Destination. Each router has a defined maximum transmission unit (MTU) including the 20-byte IP header. The number of fragments that will be delivered to the destination is \_\_\_\_\_\_\. (Answer in integer)

\par\vspace{6pt}\noindent\textbf{Your Answer:} \fbox{\rule{0pt}{1.3ex}\hspace{3cm}}

\tcblower

\textbf{Official Key:} \textbf{\large 7 to 7} \hfill \textit{\small [Hard]}
\vspace{2pt}\hrule height 0.3pt\vspace{3pt}
\textbf{Topic:} Computer Networks \\
\textbf{Subtopic:} \textit{Fragmentation and IP addressing, IPv4, CIDR notation}
\vspace{2pt}\hrule height 0.3pt\vspace{3pt}
\textbf{Guide / Hint:} \\
{\small At Router-1, max payload per fragment is $\lfloor (5000 - 20) / 8 \rfloor \times 8 = 4976$ bytes. 15000 bytes payload is fragmented into payloads of 4976, 4976, 4976, and 72 bytes (4 fragments). At Router-2 ($\text{MTU} = 3000$, max payload $\lfloor 2980/8 \rfloor \times 8 = 2976$ bytes), each of the three 4976-byte fragments splits into 2 fragments ($2976 + 2000$), yielding $3 \times 2 = 6$ fragments. The 72-byte fragment passes through intact. Total delivered fragments = $6 + 1 = 7$.}

\end{tcolorbox}
\vspace{6pt}\hrule height 0.35pt\vspace{6pt}
\noindent
\begin{tcolorbox}[
    blanker,
    sidebyside,
    sidebyside align=top,
    lefthand width=0.55\textwidth,
    top=4pt, bottom=4pt, left=0pt, right=0pt,
    middle=16pt,
    underlay={
        \draw[dashed, black!80, line width=0.8pt, dash pattern=on 4pt off 3pt]
            ([xshift=0.55\textwidth+8pt]frame.north west) --
            ([xshift=0.55\textwidth+8pt]frame.south west);
    }
]
\textbf{Q.218} \quad {\footnotesize\textsc{[MCQ $\cdot$ 2 Mark(s) $\cdot$ Neg: 0.66]}} \hfill {\footnotesize\textbf{GATE 2025 CS2} \quad \tiny\texttt{2025\_CS2\_CS\_Q36}}
\vspace{2pt}\hrule height 0.3pt\vspace{4pt}

Suppose we are transmitting frames between two nodes using Stop-and-Wait protocol. The frame size is 3000 bits. Which ONE of the following gives the efficiency of the channel?

\par\vspace{6pt}\noindent
\begin{tabular}{@{}p{0.48\linewidth}@{\hspace{0.04\linewidth}}p{0.48\linewidth}@{}}
$\square$\ \textbf{(A)} 88.23 & $\square$\ \textbf{(B)} 93.75 \\
\addlinespace[3pt]
$\square$\ \textbf{(C)} 85.44 & $\square$\ \textbf{(D)} 66.67
\end{tabular}

\tcblower

\textbf{Official Key:} \textbf{\large A} \hfill \textit{\small [Medium]}
\vspace{2pt}\hrule height 0.3pt\vspace{3pt}
\textbf{Topic:} Computer Networks \\
\textbf{Subtopic:} \textit{Data link layer: flow control and Stop-and-Wait efficiency}
\vspace{2pt}\hrule height 0.3pt\vspace{3pt}
\textbf{Guide / Hint:} \\
{\small Efficiency in Stop-and-Wait is $\eta = \frac{T_t}{T_t + 2 T_p}$. Substituting given transmission and propagation delays yields $\eta = 88.23\%$.}

\end{tcolorbox}
\vspace{6pt}\hrule height 0.35pt\vspace{6pt}
\end{document}